% Generated mechanically from frozen input inputs/p10/ja/spaces-simplicial.tex.
% Do not edit this generated file; edit the bound locale source instead.

% OK, start here.
%

\setcounter{chapter}{84}
\chapter{単体的空間}



\phantomsection
\label{spaces-simplicial-section-phantom}


\section{序論}
\label{spaces-simplicial-section-introduction}

\noindent
本章では、単体的位相空間、単体的環付き空間、単体的スキーム、および
単体的代数空間に関する理論を展開する。
単体的空間の理論を用いると、別の方法では証明が難しく見える
局所大域原理を証明できる場合がある。
応用例は論文
\cite{faltings_finiteness}、\cite{Kiehl}、および \cite{HodgeIII} に見られる。

\medskip\noindent
以下では、読者が『単体的方法』章の基本概念と結果に
習熟しているものと仮定する。『単体的方法』節
\ref{simplicial-section-introduction} を参照されたい。
とくに、今後も単体的対象を $X_\bullet$ ではなく $X$ と書く。









\section{単体的位相空間}
\label{spaces-simplicial-section-simplicial-top}

\noindent
{\it 単体的空間}とは位相空間の圏における単体的対象であり、
この圏の射は位相空間の連続写像である。
（代数空間の圏における単体的対象を「単体的代数空間」と呼ぶ。）
単体的空間 $X$ は次のように図示できる。
$$
\xymatrix{
X_2
\ar@<2ex>[r]
\ar@<0ex>[r]
\ar@<-2ex>[r]
&
X_1
\ar@<1ex>[r]
\ar@<-1ex>[r]
\ar@<1ex>[l]
\ar@<-1ex>[l]
&
X_0
\ar@<0ex>[l]
}
$$
ここには二つの射 $d^1_0, d^1_1 : X_1 \to X_0$ と
一つの射 $s^0_0 : X_0 \to X_1$ があり、以下同様である。
$d^n_i : X_n \to X_{n - 1}$ は「第 $i$ 座標を忘れる射影」と、
$s^n_j : X_n \to X_{n + 1}$ は第 $j$ 座標を繰り返す対角写像と
考えるべきであることに注意したい。

\medskip\noindent
$X$ を単体的空間とする。次のようにして $X$ にサイト
$X_{Zar}$\footnote{この記法は『サイト』例
\ref{sites-example-site-topological} および
『トポロジー』定義 
\ref{topologies-definition-big-small-Zariski} の記法と類似している。}
を対応させる。
\begin{enumerate}
\item $X_{Zar}$ の対象とは、ある $n$ に対する $X_n$ の開集合 $U$ である。
\item $X_{Zar}$ の射 $U \to V$ は、$U \subset X_n$、$V \subset X_m$
となる $n, m$ と、$\varphi$ が $X(\varphi)(U) \subset V$ を満たす
$\varphi : [m] \to [n]$ により与えられる。
\item $X_{Zar}$ における被覆 $\{U_i \to U\}$ とは、
$U, U_i \subset X_n$ が開集合で、写像 $U_i \to U$ が
$\text{id} : [n] \to [n]$ により与えられ、かつ
$U = \bigcup U_i$ となるものをいう。
\end{enumerate}
とくに、$U \to V$ が $\varphi$ により与えられる $X_{Zar}$ の射ならば、
$X(\varphi) : X_n \to X_m$ は実際に位相空間の連続写像
$U \to V$ を誘導する。

\noindent
上の構成は、任意の位相空間の図式にサイトを対応させる構成の
特別な場合であることは明らかである。そこで標準的な引理を述べる。

\begin{lemma}
\label{spaces-simplicial-lemma-simplicial-site}
$X$ を単体的空間とする。このとき上で定義した $X_{Zar}$ はサイトである。
\end{lemma}

\begin{proof}
省略する。
\end{proof}

\noindent
$X$ を単体的空間とし、$\mathcal{F}$ を $X_{Zar}$ 上の層とする。
被覆の定義から、$\mathcal{F}$ を $X_n$ の開集合に制限すると、
位相空間 $X_n$ 上の層 $\mathcal{F}_n$ が定まることは明らかである。
各 $\varphi : [m] \to [n]$ に対し、$U \subset X_n$、$V \subset X_m$、
$X(\varphi)(U) \subset V$ を満たす組についての $\mathcal{F}$ の制限写像は、
$X(\varphi)$-写像
$\mathcal{F}(\varphi) : \mathcal{F}_m \to \mathcal{F}_n$ を定める。
『層』定義 
\ref{sheaves-definition-f-map} を参照されたい。
さらに、$\varphi : [m] \to [n]$ と $\psi : [l] \to [m]$ が与えられると
$$
\mathcal{F}(\varphi) \circ \mathcal{F}(\psi) =
\mathcal{F}(\varphi \circ \psi)
$$
である（左辺では $f$-写像の合成を用いる。『層』定義
\ref{sheaves-definition-composition-f-maps} を参照されたい）。
逆も明らかに成り立つ。すなわち、上のような系\linebreak
$(\{\mathcal{F}_n\}_{n \geq 0},
\{\mathcal{F}(\varphi)\}_{\varphi \in \text{射}(\Delta)})$
が表示された等式を満たせば、$X_{Zar}$ 上の層が得られる。

\begin{lemma}
\label{spaces-simplicial-lemma-describe-sheaves-simplicial-site}
$X$ を単体的空間とする。次の二つの圏は同値である。
\begin{enumerate}
\item $\Sh(X_{Zar})$。
\item 上で述べた系 $(\mathcal{F}_n, \mathcal{F}(\varphi))$ の圏。
\end{enumerate}
\end{lemma}

\begin{proof}
上の議論を参照されたい。
\end{proof}

\begin{lemma}
\label{spaces-simplicial-lemma-simplicial-space-site-functorial}
$f : Y \to X$ を単体的空間の射とする。開集合 $U \subset X_n$ に
開集合 $f_n^{-1}(U) \subset Y_n$ を対応させる関手
$u : X_{Zar} \to Y_{Zar}$ は、サイトの射
$f_{Zar} : Y_{Zar} \to X_{Zar}$ を定める。
\end{lemma}

\begin{proof}
$u$ が連続関手であることは明らかである。したがって関手
$f_{Zar, *} = u^s$ および $f_{Zar}^{-1} = u_s$ を得る。
『サイト』節 
\ref{sites-section-morphism-sites} を参照されたい。
これがサイトの射を与えることを見るには、$u_s$ が完全であることを
示さなければならない。これには『サイト』補題
\ref{sites-lemma-directed-morphism} を用いる。
$V \subset Y_n$ を開部分集合とする。圏
$\mathcal{I}_V^u$（『サイト』節
\ref{sites-section-functoriality-PSh} を参照）は、次のような組
$(U, \varphi)$ からなる。すなわち、$\varphi : [m] \to [n]$ であり、
$U \subset X_m$ は $Y(\varphi)(V) \subset f_m^{-1}(U)$ を満たす開集合である。
また、射 $(U, \varphi) \to (U', \varphi')$ は、
$X(\psi)(U) \subset U'$ および $\varphi \circ \psi = \varphi'$ を満たす
$\psi : [m'] \to [m]$ により与えられる。
示すべきことは、$\mathcal{I}_V^u$ が余フィルター的であることである。

\medskip\noindent
『圏論』定義 
\ref{categories-definition-codirected} の条件を確認する。
$(X_n, \text{id}_{[n]})$ は対象なので、条件 (1) は成り立つ。
$(U, \varphi)$ を対象とする。条件
$Y(\varphi)(V) \subset f_m^{-1}(U)$ は
$V \subset f_n^{-1}(X(\varphi)^{-1}(U))$ と同値である。
したがって、$\psi = \varphi$ とおくことにより与えられる射
$(X(\varphi)^{-1}(U), \text{id}_{[n]}) \to (U, \varphi)$ を得る。
さらに、$(U, \text{id}_{[n]})$ と $(U', \text{id}_{[n]})$ という形の
二つの対象が与えられると、両方へ写る対象
$(U \cap U', \text{id}_{[n]})$ が存在する。よって条件 (2) が成り立つ。
条件 (3) を確認するため、$\psi, \psi' : [m'] \to [m]$ により与えられる
二つの射 $a, a': (U, \varphi) \to (U', \varphi')$ をとる。
$\varphi$ により与えられる射
$(X(\varphi)^{-1}(U), \text{id}_{[n]}) \to (U, \varphi)$ を前合成すると、
$\varphi \circ \psi = \varphi' = \varphi \circ \psi'$ なので
$a, a'$ は等化される。これで証明が完了する。
\end{proof}

\begin{lemma}
\label{spaces-simplicial-lemma-describe-functoriality}
$f : Y \to X$ を単体的空間の射とする。補題
\ref{spaces-simplicial-lemma-describe-sheaves-simplicial-site} における層の記述を用いると、
補題 \ref{spaces-simplicial-lemma-simplicial-space-site-functorial} の射 $f_{Zar}$ は
次のように記述できる。
\begin{enumerate}
\item $\mathcal{G}$ が $Y$ 上の層ならば、
$(f_{Zar, *}\mathcal{G})_n = f_{n, *}\mathcal{G}_n$ である。
\item $\mathcal{F}$ が $X$ 上の層ならば、
$(f_{Zar}^{-1}\mathcal{F})_n = f_n^{-1}\mathcal{F}_n$ である。
\end{enumerate}
\end{lemma}

\begin{proof}
第一の主張は定義から直ちに従う。第二の主張については、補題
\ref{spaces-simplicial-lemma-simplicial-space-site-functorial} の証明において、
開集合 $V \subset Y_n$ に対する圏 $(\mathcal{I}_V^u)^{opp}$ が、
$f_n^{-1}(U) \supset V$ を満たす開集合 $U \subset X_n$ と包含写像を
射とする圏を余終部分圏として含むことを示した点に注意する。
したがって、$u_p\mathcal{F}$ を $Y_n$ の開集合に制限したものは、
『層』補題 \ref{sheaves-lemma-pullback-presheaves} で定義された前層
$f_{n, p}\mathcal{F}_n$ である。$f_{Zar}^{-1}\mathcal{F} = u_s\mathcal{F}$ は
$u_p\mathcal{F}$ の層化であり、層化には被覆のみが用いられ、かつ
$Y_{Zar}$ における被覆には同じ $Y_n$ 上の開集合間の包含のみが用いられる。
一方、$f_n^{-1}\mathcal{F}_n$ は対応する意味で
$f_{n, p}\mathcal{F}_n$ の層化であるから、主張が従う。『層』節
\ref{sheaves-section-presheaves-functorial} を参照されたい。
\end{proof}

\noindent
$X$ を位相空間とする。『サイト』例
\ref{sites-example-site-topological} では、$X$ の開集合を対象、包含を射、
開被覆を被覆とするサイトを $X_{Zar}$ と表した。トポス
$\Sh(X_{Zar})$ を $X$ 上の層の圏と同一視する。

\begin{lemma}
\label{spaces-simplicial-lemma-restriction-to-components}
$X$ を単体的空間とする。関手
$X_{n, Zar} \to X_{Zar}$、$U \mapsto U$ は連続かつ余連続である。
付随するトポスの射 $g_n : \Sh(X_n) \to \Sh(X_{Zar})$ は次を満たす。
\begin{enumerate}
\item $g_n^{-1}$ は $X$ 上の層 $\mathcal{F}$ に $X_n$ 上の層
$\mathcal{F}_n$ を対応させる。
\item $g_n^{-1} : \Sh(X_{Zar}) \to \Sh(X_n)$ は左随伴
$g^{Sh}_{n!}$ をもつ。
\item $g^{Sh}_{n!}$ は有限連結極限と可換である。
\item $g_n^{-1} : \textit{Ab}(X_{Zar}) \to \textit{Ab}(X_n)$ は
左随伴 $g_{n!}$ をもつ。
\item $g_{n!}$ は完全である。
\end{enumerate}
\end{lemma}

\begin{proof}
主張に挙げた関手の性質に加えて、圏 $X_{n, Zar}$ はファイバー積と
等化子をもち、この関手はそれらと可換である（$X_{Zar}$ はすべての
ファイバー積をもつとは限らないことに注意せよ）。したがって、補題は
『サイト』節 \ref{sites-section-cocontinuous-functors} および
\ref{sites-section-cocontinuous-morphism-topoi} と、
『サイト上の加群』節 \ref{sites-modules-section-exactness-lower-shriek}
の議論から従う。より正確には、『サイト』補題
\ref{sites-lemma-cocontinuous-morphism-topoi}、
\ref{sites-lemma-when-shriek}、\ref{sites-lemma-preserve-equalizers}、および
『サイト上の加群』補題
\ref{sites-modules-lemma-g-shriek-adjoint}、
\ref{sites-modules-lemma-exactness-lower-shriek} を用いる。
\end{proof}

\begin{lemma}
\label{spaces-simplicial-lemma-restriction-injective-to-component}
$X$ を単体的空間とする。$\mathcal{I}$ が $X_{Zar}$ 上の入射的アーベル層
ならば、$\mathcal{I}_n$ は $X_n$ 上の入射的アーベル層である。
\end{lemma}

\begin{proof}
これは『ホモロジー』補題
\ref{homology-lemma-adjoint-preserve-injectives} と補題
\ref{spaces-simplicial-lemma-restriction-to-components} から従う。
\end{proof}

\begin{lemma}
\label{spaces-simplicial-lemma-restriction-to-components-functorial}
$f : Y \to X$ を単体的空間の射とする。このとき
$$
\xymatrix{
\Sh(Y_n) \ar[d] \ar[r]_{f_n} & \Sh(X_n) \ar[d] \\
\Sh(Y_{Zar}) \ar[r]^{f_{Zar}} & \Sh(X_{Zar})
}
$$
はトポスの可換図式である。
\end{lemma}

\begin{proof}
補題 \ref{spaces-simplicial-lemma-describe-functoriality} および
\ref{spaces-simplicial-lemma-restriction-to-components} における引戻し関手の記述から
直ちに従う。
\end{proof}

\begin{lemma}
\label{spaces-simplicial-lemma-augmentation}
$Y$ を単体的空間とし、$a : Y \to X$ を増大
（『単体』定義 \ref{simplicial-definition-augmentation}）とする。
$a_n : Y_n \to X$ を対応する位相空間の射とする。次の性質をもつ
標準的なトポスの射
$$
a : \Sh(Y_{Zar}) \to \Sh(X)
$$
が存在する。
\begin{enumerate}
\item $a^{-1}\mathcal{F}$ は、$Y_n$ への制限が
$a_n^{-1}\mathcal{F}$ となる層である。
\item すべての $\varphi : [m] \to [n]$ に対して
$a_m \circ Y(\varphi) = a_n$ である。
\item 補題 \ref{spaces-simplicial-lemma-restriction-to-components} の $g_n$ に対し、
トポスの射として $a \circ g_n = a_n$ である。
\item $\mathcal{G} \in \Sh(Y_{Zar})$ に対する $a_*\mathcal{G}$ は、
二つの写像 $a_{0, *}\mathcal{G}_0 \to a_{1, *}\mathcal{G}_1$ の
等化子である。
\end{enumerate}
\end{lemma}

\begin{proof}
主張 (2) は任意の圏における単体的対象の増大について成り立つ。
したがって $Y(\varphi)^{-1} a_m^{-1} \mathcal{F} = a_n^{-1}\mathcal{F}$
であり、これは $a_m^{-1}\mathcal{F}$ から $a_n^{-1}\mathcal{F}$ への
$Y(\varphi)$-写像を定める。そこで (1) を $a^{-1}\mathcal{F}$ の定義
（補題 \ref{spaces-simplicial-lemma-describe-sheaves-simplicial-site} を用いる）とし、
(4) を $a_*$ の定義とすることができる。これがトポスの射を定めるならば、
構成により $g_n^{-1} \circ a^{-1} = a_n^{-1}$ となるので (3) が従う。
$a$ がトポスの射であることを確認するには、$a^{-1}$ が $a_*$ の
左随伴であること、および $a^{-1}$ が完全であることを示せばよい。
後者は関手 $a_n^{-1}$ の完全性から直ちに従う。

\medskip\noindent
$\mathcal{F}$ を $\Sh(X)$ の対象、$\mathcal{G}$ を
$\Sh(Y_{Zar})$ の対象とする。
$\beta : a^{-1}\mathcal{F} \to \mathcal{G}$ が与えられたとき、その成分
$\beta_n : a_n^{-1}\mathcal{F} \to \mathcal{G}_n$ を考えることができる。
これらの写像に随伴する写像を
$\beta_n : \mathcal{F} \to a_{n, *}\mathcal{G}_n$ と書く。
単体構造との両立性から、$\beta_0$ は $a_*\mathcal{G}$ を経由する。
逆に、写像 $\alpha : \mathcal{F} \to a_*\mathcal{G}$ が与えられたとする。
任意の $n$ に対して $\varphi : [0] \to [n]$ を一つ選ぶ。このとき合成
$$
\mathcal{F} \xrightarrow{\alpha} a_*\mathcal{G}
\to a_{0, *}\mathcal{G}_0 \xrightarrow{\mathcal{G}(\varphi)}
a_{n, *}\mathcal{G}_n
$$
を考えることができる。これらに随伴する写像
$a_n^{-1}\mathcal{F} \to \mathcal{G}_n$ は、層の射
$a^{-1}\mathcal{F} \to \mathcal{G}$ を定める。上で与えた構成が互いに逆な
全単射
$$
\Mor_{\Sh(Y_{Zar})}(a^{-1}\mathcal{F}, \mathcal{G}) =
\Mor_{\Sh(X)}(\mathcal{F}, a_*\mathcal{G})
$$
を定めることの証明は省略する。これで証明が完了する。ここで興味深いのは、
このトポスの射が、それらのトポスを定めるサイト間の明白な幾何学的関手には
対応しないという点である。
\end{proof}

\begin{lemma}
\label{spaces-simplicial-lemma-simplicial-resolution-Z}
$X$ を単体的位相空間とする。境界写像を $\sum (-1)^i d^n_i$ とする
$X_{Zar}$ 上のアーベル前層の複体
$$
\ldots \to \mathbf{Z}_{X_2} \to \mathbf{Z}_{X_1} \to \mathbf{Z}_{X_0}
$$
は定数前層 $\mathbf{Z}$ の分解である。
\end{lemma}

\begin{proof}
$U \subset X_m$ を $X_{Zar}$ の対象とする。上の複体の $U$ における値は
アーベル群の複体
$$
\ldots \to
\mathbf{Z}[\Mor_\Delta([2], [m])] \to
\mathbf{Z}[\Mor_\Delta([1], [m])] \to
\mathbf{Z}[\Mor_\Delta([0], [m])]
$$
である。言い換えると、これは単体集合 $\Delta[m]$ 上の自由アーベル群に
付随する複体である。『単体』例
\ref{simplicial-example-simplex-simplicial-set} を参照されたい。
$\Delta[m]$ は $\Delta[0]$ とホモトピー同値であり（『単体』例
\ref{simplicial-example-simplex-contractible}）、また「自由アーベル群をとる」
操作は関手なので、上の複体は $\Delta[0]$ 上の自由アーベル群と
ホモトピー同値である（『単体』注意
\ref{simplicial-remark-homotopy-better} および補題
\ref{simplicial-lemma-homotopy-equivalence-s-N}）。この複体は正の次数で
非輪状であり、次数 $0$ では $\mathbf{Z}$ に等しい。
\end{proof}

\begin{lemma}
\label{spaces-simplicial-lemma-simplicial-sheaf-cohomology}
$X$ を単体的位相空間とする。$\mathcal{F}$ を $X$ 上のアーベル層とする。
スペクトル系列 $(E_r, d_r)_{r \geq 0}$ で
$$
E_1^{p, q} = H^q(X_p, \mathcal{F}_p)
$$
をもち、$H^{p + q}(X_{Zar}, \mathcal{F})$ に収束するものが存在する。
このスペクトル系列は $\mathcal{F}$ に関して関手的である。
\end{lemma}

\begin{proof}
$\mathcal{F} \to \mathcal{I}^\bullet$ を入射分解とする。項が
$$
A^{p, q} = \mathcal{I}^q(X_p)
$$
であり、第一の微分が写像 $d^{p + 1}_i$ に沿う写像
$\mathcal{I}_p^q \to \mathcal{I}_{p + 1}^q$ の交代和で与えられる
二重複体を考える。補題
\ref{spaces-simplicial-lemma-describe-sheaves-simplicial-site} を参照されたい。ここで
$$
A^{p, q} = \Gamma(X_p, \mathcal{I}_p^q) =
\Mor_{\textit{PSh}}(h_{X_p}, \mathcal{I}^q) =
\Mor_{\textit{PAb}}(\mathbf{Z}_{X_p}, \mathcal{I}^q)
$$
である。したがって、補題 \ref{spaces-simplicial-lemma-simplicial-resolution-Z} および
『サイト上のコホモロジー』補題
\ref{sites-cohomology-lemma-injective-abelian-sheaf-injective-presheaf}
から、二重複体の各行は正の次数で完全であり、次数 $0$ では
$\Gamma(X_{Zar}, \mathcal{I}^q)$ になる。一方、制限関手は完全なので
（補題 \ref{spaces-simplicial-lemma-restriction-to-components}）、写像
$$
\mathcal{F}_p \to \mathcal{I}_p^\bullet
$$
は分解である。層 $\mathcal{I}_p^q$ は $X_p$ 上の入射的アーベル層である
（補題 \ref{spaces-simplicial-lemma-restriction-injective-to-component}）。したがって各列の
コホモロジーは群 $H^q(X_p, \mathcal{F}_p)$ を計算する。
『ホモロジー』補題 \ref{homology-lemma-first-quadrant-ss} および
\ref{homology-lemma-double-complex-gives-resolution} を適用して結論を得る。
\end{proof}







\begin{lemma}
\label{spaces-simplicial-lemma-augmentation-pushforward-higher-direct-image}
$X$ を単体的空間とし、$a : X \to Y$ を増大とする。
$\mathcal{F}$ を $X_{Zar}$ 上のアーベル層とする。このとき
$R^na_*\mathcal{F}$ は前層
$$
V \longmapsto H^n((X \times_Y V)_{Zar}, \mathcal{F}|_{(X \times_Y V)_{Zar}})
$$
に伴う層である。
\end{lemma}

\begin{proof}
これは『コホモロジー』補題
\ref{cohomology-lemma-describe-higher-direct-images} または
『サイト上のコホモロジー』補題
\ref{sites-cohomology-lemma-higher-direct-images} の類似であり、読者には
この証明を飛ばすことを強く勧める。$X_{Zar}$ 上で $\mathcal{F}$ の
入射分解を一つ選び、定義を用いると、次を示せば十分であることが分かる。
(1) $X_{Zar}$ 上の入射的アーベル層を $(X \times_Y V)_{Zar}$ に制限したものは
入射的アーベル層である。(2) $a_*\mathcal{F}$ は規則
$$
V \longmapsto H^0((X \times_Y V)_{Zar}, \mathcal{F}|_{(X \times_Y V)_{Zar}})
$$
に等しい。(2) は次の事実から従う。
\begin{enumerate}
\item[(2a)] 補題 \ref{spaces-simplicial-lemma-augmentation} により、$a_*\mathcal{F}$ は
二つの写像 $a_{0, *}\mathcal{F}_0 \to a_{1, *}\mathcal{F}_1$ の等化子である。
\item[(2b)] $a_{0, *}\mathcal{F}_0(V) =
H^0(a_0^{-1}(V), \mathcal{F}_0)$ および
$a_{1, *}\mathcal{F}_1(V) = H^0(a_1^{-1}(V), \mathcal{F}_1)$ である。
\item[(2c)] $X_0 \times_Y V = a_0^{-1}(V)$ および
$X_1 \times_Y V = a_1^{-1}(V)$ である。
\item[(2d)] たとえば補題 \ref{spaces-simplicial-lemma-simplicial-sheaf-cohomology} により、
$H^0((X \times_Y V)_{Zar}, \mathcal{F}|_{(X \times_Y V)_{Zar}})$ は
二つの写像
$H^0(X_0 \times_Y V, \mathcal{F}_0) \to H^0(X_1 \times_Y V, \mathcal{F}_1)$
の等化子である。
\end{enumerate}
(1) は、制限関手
$\textit{Ab}(X_{Zar}) \to \textit{Ab}((X \times_Y V)_{Zar})$ の
完全な左随伴
$j_! : \textit{Ab}((X \times_Y V)_{Zar}) \to \textit{Ab}(X_{Zar})$
（零による延長）を定義し、『ホモロジー』補題
\ref{homology-lemma-adjoint-preserve-injectives} を用いれば従う。
\end{proof}

\noindent
$X$ を位相空間とする。値が $X$ である定数単体的位相空間を
$X_\bullet$ と表す。補題 \ref{spaces-simplicial-lemma-describe-sheaves-simplicial-site} により、
$X_{\bullet, Zar}$ 上の層は、$X$ 上の層の圏における余単体的対象と同じものである。

\begin{lemma}
\label{spaces-simplicial-lemma-constant-simplicial-space}
$X$ を位相空間とし、$X_\bullet$ を値が $X$ である定数単体的位相空間とする。
関手
$$
X_{\bullet, Zar} \longrightarrow X_{Zar},\quad
U \longmapsto U
$$
は連続かつ余連続であり、トポスの射
$g : \Sh(X_{\bullet, Zar}) \to \Sh(X)$ と、$g^{-1}$ の左随伴 $g_!$ を定める。
さらに次が成り立つ。
\begin{enumerate}
\item $g^{-1}$ は $X$ 上の層に $X$ 上の定数余単体的層を対応させる。
\item $g_!$ は $X_{\bullet, Zar}$ 上の層 $\mathcal{F}$ に層
$\mathcal{F}_0$ を対応させる。
\item $g_*$ は $X_{\bullet, Zar}$ 上の層 $\mathcal{F}$ に、二つの写像
$\mathcal{F}_0 \to \mathcal{F}_1$ の等化子を対応させる。
\end{enumerate}
\end{lemma}

\begin{proof}
関手についての主張は容易に確認できる。$g$ と $g_!$ の存在は
『サイト』補題 \ref{sites-lemma-cocontinuous-morphism-topoi} および
\ref{sites-lemma-when-shriek} から従う。$g^{-1}$ の記述は
『サイト』補題 \ref{sites-lemma-when-shriek} から直ちに従う。
$g_*$ と $g_!$ の記述は、与えられた関手がそれぞれ $g^{-1}$ の
右随伴と左随伴であることから従う。
\end{proof}

\section{単体的サイトとトポス}
\label{spaces-simplicial-section-simplicial-sites}

\noindent
{\it 単体的サイト}を、対象がサイトで射がサイトの射である（大きな）圏の
単体的対象として定義するのは自然に思われる。『サイト』定義
\ref{sites-definition-site}、\ref{sites-definition-morphism-sites}、および
『サイト』補題 \ref{sites-lemma-composition-morphisms-sites} における
射の合成を参照されたい。しかし、考慮に値する変種がいくつかある。
(a) 代わりにサイト間の余連続関手（『サイト』節
\ref{sites-section-cocontinuous-functors} および
\ref{sites-section-cocontinuous-morphism-topoi}）を用いることもできる。
(b) サイトの射の間に $2$-射の概念を導入した、適切なサイトの $2$-圏で
作業することもできる。(c) 余連続関手から構成される $2$-圏で作業することも
できる。これらの変種の一つを定義として選ぶ代わりに、必要に応じて
理論を展開することにする。

\medskip\noindent
{\it 単体的トポス}は、おそらく $\Delta^{opp}$ からトポスの $2$-圏への
擬関手として定義すべきであろう。『圏論』定義
\ref{categories-definition-functor-into-2-category}、ならびに『サイト』節
\ref{sites-section-topoi} および \ref{sites-section-2-category} を
参照されたい。可能な限り、このような代物を直接扱うことは避ける。

\medskip\noindent
{\bf 場合 A.}
$\mathcal{C}$ を、対象がサイトで射がサイトの射である圏の単体的対象とする。
これは、$\Delta$ の各射 $\varphi : [m] \to [n]$ に対してサイトの射
$f_\varphi : \mathcal{C}_n \to \mathcal{C}_m$ が存在することを意味する。
この射は逆向きの連続関手により与えられ、それを
$u_\varphi : \mathcal{C}_m \to \mathcal{C}_n$ と表す。

\begin{lemma}
\label{spaces-simplicial-lemma-simplicial-site-site}
$\mathcal{C}$ をサイトの圏における単体的対象とする。上の記法のもとで、
サイト $\mathcal{C}_{total}$ を次のように構成する。
\begin{enumerate}
\item $\mathcal{C}_{total}$ の対象とは、ある $n$ に対する
$\mathcal{C}_n$ の対象 $U$ である。
\item $\mathcal{C}_{total}$ の射 $(\varphi, f) : U \to V$ は、
$U \in \Ob(\mathcal{C}_n)$、$V \in \Ob(\mathcal{C}_m)$ を満たす写像
$\varphi : [m] \to [n]$ と、$\mathcal{C}_n$ の射
$f : U \to u_\varphi(V)$ により与えられる。
\item $\mathcal{C}_{total}$ における被覆
$\{(\text{id}, f_i) :  U_i \to U\}$ は、ある $n$ と
$\mathcal{C}_n$ の被覆 $\{f_i : U_i \to U\}$ により与えられる。
\end{enumerate}
\end{lemma}

\begin{proof}
$(\varphi, f) : U \to V$ と $(\psi, g) : V \to W$ の合成は
$(\varphi \circ \psi, u_\varphi(g) \circ f)$ で与えられる。
ここでは $u_\varphi \circ u_\psi = u_{\varphi \circ \psi}$ を用いた。

\medskip\noindent
(3) のような被覆 $\{(\text{id}, f_i) :  U_i \to U\}$ と、
$W \in \Ob(\mathcal{C}_m)$ を満たす射 $(\varphi, g) : W \to U$ をとる。
圏 $\mathcal{C}_{total}$ において
$$
W \times_{(\varphi, g), U, (\text{id}, f_i)} U_i =
W \times_{g, u_\varphi(U), u_\varphi(f_i)} u_\varphi(U_i)
$$
であると主張する。これは意味をもつ。実際、サイトの射の定義により、
$u_\varphi$ は被覆を被覆に移すので、$\mathcal{C}_m$ における必要な
ファイバー積が存在する。同じ議論から上の主張が従う（詳細は省略する）。
したがって、被覆の集まりは基底変換で安定である。サイトの他の公理は
直ちに従う。
\end{proof}

\noindent
{\bf 場合 B.}
$\mathcal{C}$ を、対象がサイトで射が余連続関手である圏の単体的対象とする。
これは、$\Delta$ の各射 $\varphi : [m] \to [n]$ に対して、
$u_\varphi : \mathcal{C}_n \to \mathcal{C}_m$ と表す余連続関手が
存在することを意味する。付随するトポスの射を
$f_\varphi : \Sh(\mathcal{C}_n) \to \Sh(\mathcal{C}_m)$ と表す。

\begin{lemma}
\label{spaces-simplicial-lemma-simplicial-cocontinuous-site}
$\mathcal{C}$ を、対象がサイトで射が余連続関手である圏の単体的対象とする。
上の記法のもとで、関手
$u_\varphi : \mathcal{C}_n \to \mathcal{C}_m$ が『サイト』注意
\ref{sites-remark-cartesian-cocontinuous} の性質 $P$ をもつと仮定する。
このときサイト $\mathcal{C}_{total}$ を次のように構成できる。
\begin{enumerate}
\item $\mathcal{C}_{total}$ の対象とは、ある $n$ に対する
$\mathcal{C}_n$ の対象 $U$ である。
\item $\mathcal{C}_{total}$ の射 $(\varphi, f) : U \to V$ は、
$U \in \Ob(\mathcal{C}_n)$、$V \in \Ob(\mathcal{C}_m)$ を満たす写像
$\varphi : [m] \to [n]$ と、$\mathcal{C}_m$ の射
$f : u_\varphi(U) \to V$ により与えられる。
\item $\mathcal{C}_{total}$ における被覆
$\{(\text{id}, f_i) :  U_i \to U\}$ は、ある $n$ と
$\mathcal{C}_n$ の被覆 $\{f_i : U_i \to U\}$ により与えられる。
\end{enumerate}
\end{lemma}

\begin{proof}
$(\varphi, f) : U \to V$ と $(\psi, g) : V \to W$ の合成は
$(\varphi \circ \psi, g \circ u_\psi(f))$ で与えられる。
ここでは $u_\psi \circ u_\varphi = u_{\varphi \circ \psi}$ を用いた。

\medskip\noindent
(3) のような被覆 $\{(\text{id}, f_i) :  U_i \to U\}$ と、
$W \in \Ob(\mathcal{C}_m)$ を満たす射 $(\varphi, g) : W \to U$ をとる。
圏 $\mathcal{C}_{total}$ において
$$
W \times_{(\varphi, g), U, (\text{id}, f_i)} U_i =
W \times_{g, U, f_i} U_i
$$
であると主張する。右辺は『サイト』注意
\ref{sites-remark-cartesian-cocontinuous} で定義された
$\mathcal{C}_m$ の対象であり、性質 $P$ により存在する。
この種のファイバー積と関手の合成との両立性から主張が従う
（詳細は省略する）。性質 $P$ により族
$\{W \times_{g, U, f_i} U_i \to W\}$ は $\mathcal{C}_m$ の被覆なので、
被覆の集まりは基底変換で安定である。サイトの他の公理は直ちに従う。
\end{proof}

\begin{situation}
\label{spaces-simplicial-situation-simplicial-site}
ここでは次の二つの場合のいずれかを扱う。
\begin{enumerate}
\item[(A)] $\mathcal{C}$ は、対象がサイトで射がサイトの射である圏の
単体的対象である。$\Delta$ の各射 $\varphi : [m] \to [n]$ に対して、
連続関手 $u_\varphi : \mathcal{C}_m \to \mathcal{C}_n$ により与えられる
サイトの射 $f_\varphi : \mathcal{C}_n \to \mathcal{C}_m$ が存在する。
\item[(B)] $\mathcal{C}$ は、対象がサイトで射が『サイト』注意
\ref{sites-remark-cartesian-cocontinuous} の性質 $P$ をもつ余連続関手である
圏の単体的対象である。$\Delta$ の各射 $\varphi : [m] \to [n]$ に対して、
トポスの射 $f_\varphi : \Sh(\mathcal{C}_n) \to \Sh(\mathcal{C}_m)$ を
誘導する余連続関手 $u_\varphi : \mathcal{C}_n \to \mathcal{C}_m$ が存在する。
\end{enumerate}
通常どおり、引戻しと順像をそれぞれ $f_\varphi^{-1}$、
$f_{\varphi, *}$ と表す。補題 \ref{spaces-simplicial-lemma-simplicial-site-site}（場合 A）
または補題 \ref{spaces-simplicial-lemma-simplicial-cocontinuous-site}（場合 B）で定義された
サイトを $\mathcal{C}_{total}$ と表す。
\end{situation}

\noindent
$\mathcal{C}$ を状況 \ref{spaces-simplicial-situation-simplicial-site} のとおりとし、
$\mathcal{F}$ を $\mathcal{C}_{total}$ 上の層とする。被覆の定義から、
$\mathcal{F}$ を $\mathcal{C}_n$ の対象に制限すると、サイト
$\mathcal{C}_n$ 上の層 $\mathcal{F}_n$ が定まることは明らかである。
各 $\varphi : [m] \to [n]$ に対し、$U \in \Ob(\mathcal{C}_n)$、
$V \in \Ob(\mathcal{C}_m)$ を満たす射 $(\varphi, f) : U \to V$ に沿う
$\mathcal{F}$ の制限写像は、次の集合の元 $\mathcal{F}(\varphi)$ を定める。
$$
\Mor_{\Sh(\mathcal{C}_m)}(\mathcal{F}_m, f_{\varphi, *}\mathcal{F}_n) =
\Mor_{\Sh(\mathcal{C}_n)}(f_\varphi^{-1}\mathcal{F}_m, \mathcal{F}_n)
$$
さらに、$\varphi : [m] \to [n]$ と $\psi : [l] \to [m]$ が与えられると、
図式
$$
\vcenter{
\xymatrix@C=1.4em{
\mathcal{F}_l \ar[rr]_{\mathcal{F}(\varphi \circ \psi)}
\ar[rd]_{\mathcal{F}(\psi)}
& & f_{\varphi \circ \psi, *} \mathcal{F}_n \\
& f_{\psi, *}\mathcal{F}_m \ar[ur]_{f_{\psi, *}\mathcal{F}(\varphi)}
}
}
\quad\text{および}\quad
\vcenter{
\xymatrix@C=1.4em{
f_{\varphi \circ \psi}^{-1}\mathcal{F}_l
\ar[rr]_{\mathcal{F}(\varphi \circ \psi)}
\ar[rd]_{f_\varphi^{-1}\mathcal{F}(\psi)}
& & \mathcal{F}_n \\
& f_\varphi^{-1}\mathcal{F}_m \ar[ur]_{\mathcal{F}(\varphi)}
}
}
$$
は可換である。逆も明らかに成り立つ。すなわち、上の可換性条件を満たす系
$(\{\mathcal{F}_n\}_{n \geq 0},
\{\mathcal{F}(\varphi)\}_{\varphi \in \text{射}(\Delta)})$
が与えられれば、$\mathcal{C}_{total}$ 上の層が得られる。

\begin{lemma}
\label{spaces-simplicial-lemma-describe-sheaves-simplicial-site-site}
状況 \ref{spaces-simplicial-situation-simplicial-site} において、次の二つの圏は同値である。
\begin{enumerate}
\item $\Sh(\mathcal{C}_{total})$。
\item 上で述べた系 $(\mathcal{F}_n, \mathcal{F}(\varphi))$ の圏。
\end{enumerate}
とくに、トポス $\Sh(\mathcal{C}_{total})$ はトポス
$\Sh(\mathcal{C}_n)$ とトポスの射 $f_\varphi$ のみに依存する。
\end{lemma}

\begin{proof}
上の議論を参照されたい。
\end{proof}

\begin{lemma}
\label{spaces-simplicial-lemma-restriction-to-components-site}
状況 \ref{spaces-simplicial-situation-simplicial-site} において、関手
$\mathcal{C}_n \to \mathcal{C}_{total}$、$U \mapsto U$ は連続かつ
余連続である。付随するトポスの射
$g_n : \Sh(\mathcal{C}_n) \to \Sh(\mathcal{C}_{total})$ は次を満たす。
\begin{enumerate}
\item $g_n^{-1}$ は $\mathcal{C}_{total}$ 上の層 $\mathcal{F}$ に
$\mathcal{C}_n$ 上の層 $\mathcal{F}_n$ を対応させる。
\item $g_n^{-1} : \Sh(\mathcal{C}_{total}) \to \Sh(\mathcal{C}_n)$ は
左随伴 $g^{Sh}_{n!}$ をもつ。
\item $\Sh(\mathcal{C}_n)$ の $\mathcal{G}$ に対し、
$g_{n!}^{Sh}\mathcal{G}$ の $\mathcal{C}_m$ への制限は
$\coprod\nolimits_{\varphi : [n] \to [m]} f_\varphi^{-1}\mathcal{G}$ である。
\item $g_{n!}^{Sh}$ は有限連結極限と可換である。
\item $g_n^{-1} : \textit{Ab}(\mathcal{C}_{total}) \to
\textit{Ab}(\mathcal{C}_n)$ は左随伴 $g_{n!}$ をもつ。
\item $\textit{Ab}(\mathcal{C}_n)$ の $\mathcal{G}$ に対し、
$g_{n!}\mathcal{G}$ の $\mathcal{C}_m$ への制限は
$\bigoplus\nolimits_{\varphi : [n] \to [m]} f_\varphi^{-1}\mathcal{G}$ である。
\item $g_{n!}$ は完全である。
\end{enumerate}
\end{lemma}

\begin{proof}
場合 A。$\{U_i \to U\}_{i \in I}$ が $\mathcal{C}_n$ の被覆ならば、
定義（補題 \ref{spaces-simplicial-lemma-simplicial-site-site}）により、その像
$\{U_i \to U\}_{i \in I}$ は $\mathcal{C}_{total}$ の被覆である。
$\mathcal{C}_n$ の射 $V \to U$ に対して、$\mathcal{C}_n$ における
ファイバー積 $V \times_U U_i$ は $\mathcal{C}_{total}$ における
ファイバー積でもある（補題 \ref{spaces-simplicial-lemma-simplicial-site-site} の証明中の
主張による）。したがって、この関手は連続である。一方、この関手は
$\mathcal{C}_n$ における $U$ の被覆と $\mathcal{C}_{total}$ における
$U$ の被覆との間に全単射を定める。したがって、この関手は余連続でもある。

\medskip\noindent
場合 B。$\{U_i \to U\}_{i \in I}$ が $\mathcal{C}_n$ の被覆ならば、
定義（補題 \ref{spaces-simplicial-lemma-simplicial-cocontinuous-site}）により、その像
$\{U_i \to U\}_{i \in I}$ は $\mathcal{C}_{total}$ の被覆である。
$\mathcal{C}_n$ の射 $V \to U$ に対して、$\mathcal{C}_n$ における
ファイバー積 $V \times_U U_i$ は $\mathcal{C}_{total}$ における
ファイバー積でもある（補題 \ref{spaces-simplicial-lemma-simplicial-cocontinuous-site} の
証明中の主張による）。したがって、この関手は連続である。一方、この関手は
$\mathcal{C}_n$ における $U$ の被覆と $\mathcal{C}_{total}$ における
$U$ の被覆との間に全単射を定める。したがって、この関手は余連続でもある。

\medskip\noindent
ここまでで、場合 A と B のいずれについても、(1) と
$g^{Sh}_{n!}$、$g_{n!}$ の存在は『サイト』補題
\ref{sites-lemma-cocontinuous-morphism-topoi}、\ref{sites-lemma-when-shriek}
および『サイト上の加群』補題
\ref{sites-modules-lemma-g-shriek-adjoint} から従う。

\medskip\noindent
(3) の証明。$\mathcal{G}$ を $\mathcal{C}_n$ 上の層とする。
$\mathcal{C}_{total}$ 上の層 $\mathcal{H}$ で、その次数 $m$ 部分が
補題の主張 (3) にある層
$$
\mathcal{H}_m = \coprod\nolimits_{\varphi : [n] \to [m]}
f_\varphi^{-1}\mathcal{G}
$$
であるものを考える。写像 $\psi : [m] \to [m']$ が与えられると、写像
$\mathcal{H}(\psi) : f_\psi^{-1}\mathcal{H}_m \to \mathcal{H}_{m'}$ は、
各成分上で同一視
$$
f_\psi^{-1} f_\varphi^{-1} \mathcal{G} \to
f_{\psi \circ \varphi}^{-1}\mathcal{G}
$$
により与えられる。$\mathcal{C}_{total}$ 上の層の写像
$\alpha : \mathcal{H} \to \mathcal{F}$ が与えられると、
$\alpha_n$ を $\mathcal{H}_n$ の成分 $\mathcal{G}$ に制限することに
対応する写像 $\mathcal{G} \to \mathcal{F}_n$ が得られる。
逆に、$\mathcal{C}_n$ 上の層の写像
$\beta : \mathcal{G} \to \mathcal{F}_n$ が与えられると、
$\alpha : \mathcal{H} \to \mathcal{F}$ を次のように定義できる。
$\alpha_m$ の各成分
$$
f_\varphi^{-1}\mathcal{G} \to \mathcal{F}_m
$$
を、$\mathcal{F}(\varphi) \circ f_\varphi^{-1}\beta$ に随伴する写像とする。
これら二つの構成が互いに逆な写像
$$
\Mor_{\Sh(\mathcal{C}_n)}(\mathcal{G}, \mathcal{F}_n) =
\Mor_{\Sh(\mathcal{C}_{total})}(\mathcal{H}, \mathcal{F})
$$
を与えることの確認は省略する。したがって所望のとおり
$\mathcal{H} = g^{Sh}_{n!}\mathcal{G}$ である。

\medskip\noindent
% P10-E0606 P10-E0607 P10-E0608
(4) の証明。$\mathcal{G}$ が $\mathcal{C}_n$ 上のアーベル層ならば、
上とまったく同じ方法で進める。ただし $\mathcal{H}$ を
$\mathcal{C}_{total}$ 上のアーベル層で、その次数 $m$ 部分が
$$
\bigoplus\nolimits_{\varphi : [n] \to [m]} f_\varphi^{-1}\mathcal{G}
$$
であるものとして定義し、遷移写像を上とまったく同様に定義する。全単射
$$
\Mor_{\textit{Ab}(\mathcal{C}_n)}(\mathcal{G}, \mathcal{F}_n) =
\Mor_{\textit{Ab}(\mathcal{C}_{total})}(\mathcal{H}, \mathcal{F})
$$
も上とまったく同様に証明される。したがって所望のとおり
$\mathcal{H} = g_{n!}\mathcal{G}$ である。

\medskip\noindent
$g^{Sh}_{n!}$ と $g_{n!}$ の完全性に関する性質は、これらの関手について
与えた公式から従う。
\end{proof}

\begin{lemma}
\label{spaces-simplicial-lemma-restriction-injective-to-component-site}
\begin{slogan}
単体的サイト上の入射的アーベル層は各成分上で入射的である
\end{slogan}
% P10-E0609
状況 \ref{spaces-simplicial-situation-simplicial-site} において、
$\mathcal{I}$ が $\textit{Ab}(\mathcal{C}_{total})$ において入射的ならば、
$\mathcal{I}_n$ は $\textit{Ab}(\mathcal{C}_n)$ において入射的である。
$\mathcal{I}^\bullet$ が $\textit{Ab}(\mathcal{C}_{total})$ における
K-入射的複体ならば、$\mathcal{I}_n^\bullet$ は
$\textit{Ab}(\mathcal{C}_n)$ における K-入射的複体である。
\end{lemma}

\begin{proof}
第一の主張は『ホモロジー』補題
\ref{homology-lemma-adjoint-preserve-injectives} と補題
\ref{spaces-simplicial-lemma-restriction-to-components-site} から従う。第二の主張は
『導来圏』補題 \ref{derived-lemma-adjoint-preserve-K-injectives} と補題
\ref{spaces-simplicial-lemma-restriction-to-components-site} から従う。
\end{proof}

\section{単体的サイトの増大}
\label{spaces-simplicial-section-augmentation-simplicial-sites}

\noindent
節 \ref{spaces-simplicial-section-simplicial-sites} で述べた方法を続け、同節で扱った
場合 A と B における増大の意味を明らかにする。

\begin{remark}
\label{spaces-simplicial-remark-augmentation-site}
状況 \ref{spaces-simplicial-situation-simplicial-site} における
{\it サイト $\mathcal{D}$ への増大 $a_0$}とは、次を意味する。
\begin{enumerate}
\item[(A)] $a_0 : \mathcal{C}_0 \to \mathcal{D}$ は連続関手
$u_0 : \mathcal{D} \to \mathcal{C}_0$ により与えられるサイトの射であり、
すべての $\varphi, \psi : [0] \to [n]$ に対して
$u_\varphi \circ u_0 = u_\psi \circ u_0$ である。
\item[(B)] $a_0 : \Sh(\mathcal{C}_0) \to \Sh(\mathcal{D})$ は余連続関手
$u_0 : \mathcal{C}_0 \to \mathcal{D}$ により与えられるトポスの射であり、
すべての $\varphi, \psi : [0] \to [n]$ に対して
$u_0 \circ u_\varphi = u_0 \circ u_\psi$ である。
\end{enumerate}
\end{remark}

\begin{lemma}
\label{spaces-simplicial-lemma-augmentation-site}
状況 \ref{spaces-simplicial-situation-simplicial-site} において、$a_0$ を注意
\ref{spaces-simplicial-remark-augmentation-site} のようなサイト $\mathcal{D}$ への増大とする。
このとき $a_0$ は次を誘導する。
\begin{enumerate}
\item すべての $n \geq 0$ に対するトポスの射
$a_n : \Sh(\mathcal{C}_n) \to \Sh(\mathcal{D})$。
\item トポスの射 $a : \Sh(\mathcal{C}_{total}) \to \Sh(\mathcal{D})$。
\end{enumerate}
これらは次を満たす。
\begin{enumerate}
\item すべての $\varphi : [m] \to [n]$ に対して
$a_m \circ f_\varphi = a_n$ である。
\item $g_n : \Sh(\mathcal{C}_n) \to \Sh(\mathcal{C}_{total})$ が補題
\ref{spaces-simplicial-lemma-restriction-to-components-site} のとおりならば、
$a \circ g_n = a_n$ である。
\item $\mathcal{F} \in \Sh(\mathcal{C}_{total})$ に対する
$a_*\mathcal{F}$ は、二つの写像
$a_{0, *}\mathcal{F}_0 \to a_{1, *}\mathcal{F}_1$ の等化子である。
\end{enumerate}
\end{lemma}

\begin{proof}
場合 A。$u_n : \mathcal{D} \to \mathcal{C}_n$ を、
$\varphi : [0] \to [n]$ に対する関手 $u_\varphi \circ u_0$ の共通の値とする。
このとき $u_n$ はサイトの射 $a_n : \mathcal{C}_n \to \mathcal{D}$ に
対応する。『サイト』補題 \ref{sites-lemma-composition-morphisms-sites} を
参照されたい。同じ補題から、すべての $\varphi : [m] \to [n]$ に対して
$a_m \circ f_\varphi = a_n$ であることが分かる。

\medskip\noindent
場合 B。$u_n : \mathcal{C}_n \to \mathcal{D}$ を、
$\varphi : [0] \to [n]$ に対する関手 $u_0 \circ u_\varphi$ の共通の値とする。
このとき $u_n$ は余連続であり、したがってトポスの射
% P10-E0172
$a_n : \Sh(\mathcal{C}_n) \to \Sh(\mathcal{D)}$ を定める。
『サイト』補題 \ref{sites-lemma-composition-cocontinuous} を参照されたい。
同じ補題から、すべての $\varphi : [m] \to [n]$ に対して
$a_m \circ f_\varphi = a_n$ であることが分かる。

\medskip\noindent
集合の層 $\mathcal{G}$ に、成分が $a_n^{-1}\mathcal{G}$ であり、
$\varphi : [m] \to [n]$ に対する遷移写像 $\mathcal{F}(\varphi)$ が同一視
$$
f_\varphi^{-1}\mathcal{F}_m =
f_\varphi^{-1} a_m^{-1}\mathcal{G} =
a_n^{-1}\mathcal{G} =
\mathcal{F}_n
$$
である層 $\mathcal{F} = a^{-1}\mathcal{G}$ を対応させる関手
$a^{-1} : \Sh(\mathcal{D}) \to \Sh(\mathcal{C}_{total})$ を考える。
補題 \ref{spaces-simplicial-lemma-describe-sheaves-simplicial-site-site} における
$\Sh(\mathcal{C}_{total})$ の記述を参照されたい。関手 $a_n^{-1}$ は
完全なので、$a^{-1}$ は完全関手である。最後に、
$a_* : \Sh(\mathcal{C}_{total}) \to \Sh(\mathcal{D})$ として、
% P10-E0173
$\Sh(\mathcal{D})$ 上の層 $\mathcal{F}$ に
$$
\xymatrix{
a_*\mathcal{F} \ar@{=}[r] &
\text{等化子}(a_{0, *}\mathcal{F}_0
\ar@<1ex>[r] \ar@<-1ex>[r] &
a_{1, *}\mathcal{F}_1)
}
$$
を対応させる関手をとる。ここで二つの写像は、二つの写像
$\varphi : [0] \to [1]$ から
$$
a_{0, *}\mathcal{F}_0 \to
a_{0, *}f_{\varphi, *} f_\varphi^{-1}\mathcal{F}_0
\xrightarrow{\mathcal{F}(\varphi)}
a_{0, *}f_{\varphi, *} \mathcal{F}_1 = a_{1, *}\mathcal{F}_1
$$
を介して得られる。最初の矢印は
$1 \to f_{\varphi, *} f_\varphi^{-1}$ から来る。
% P10-E0174
$\mathcal{G}_\bullet$ を値が $\mathcal{G}$ である定数余単体的層と表し、
$a_{\bullet, *}\mathcal{F}$ を次数 $n$ の成分が
$a_{n, *}\mathcal{F}_n$ である余単体的層と表す。
% P10-E0175
トポスの射 $a_n$ に対する通常の随伴により、写像
$a^{-1}\mathcal{G} \to \mathcal{F}$ を与えることは、余単体的層の写像
$$
\mathcal{G}_\bullet \longrightarrow a_{\bullet, *}\mathcal{F}
$$
を与えることと同じである。『単体』補題
\ref{simplicial-lemma-augmentation-howto} の双対により、これは写像
$\mathcal{G} \to a_*\mathcal{F}$ を与えることと同じである。
したがって $a^{-1}$ と $a_*$ は随伴関手であり、トポスの射
$a$\footnote{場合 B では、射 $a$ は $\mathcal{C}_n$ の $U$ を
$u_n(U)$ に送る余連続関手
$\mathcal{C}_{total} \to \mathcal{D}$ に対応する。}を得る。
% P10-E0176
等式 $a \circ g_n = f_n$ は定義から直ちに従う。
\end{proof}

\section{単体的サイトの射}
\label{spaces-simplicial-section-morphism-simplicial-sites}

\noindent
節 \ref{spaces-simplicial-section-simplicial-sites} で述べた方法を続け、同節で扱った
場合 A と B における単体的サイトの射の意味を明らかにする。

\begin{remark}
\label{spaces-simplicial-remark-morphism-simplicial-sites}
$\mathcal{C}_n, f_\varphi, u_\varphi$ および
$\mathcal{C}'_n, f'_\varphi, u'_\varphi$ を状況
\ref{spaces-simplicial-situation-simplicial-site} のとおりとする。
{\it 単体的サイト間の射 $h$}とは、次を意味する。
\begin{enumerate}
\item[(A)] サイトの射 $h_n : \mathcal{C}_n \to \mathcal{C}'_n$ であって、
すべての $\varphi : [m] \to [n]$ に対してサイトの射として
$f'_\varphi \circ h_n = h_m \circ f_\varphi$ を満たすもの。
\item[(B)] トポスの射
$h_n : \Sh(\mathcal{C}_n) \to \Sh(\mathcal{C}'_n)$ を誘導する余連続関手
$v_n : \mathcal{C}_n \to \mathcal{C}'_n$ であって、すべての
$\varphi : [m] \to [n]$ に対して関手として
$u'_\varphi \circ v_n = v_m \circ u_\varphi$ を満たすもの。
\end{enumerate}
いずれの場合も、トポスの射として
$f'_\varphi \circ h_n = h_m \circ f_\varphi$ である。場合 B については
『サイト』補題 \ref{sites-lemma-composition-cocontinuous}、場合 A については
『サイト』定義 \ref{sites-definition-composition-morphisms-sites} を
参照されたい。
\end{remark}

\begin{lemma}
\label{spaces-simplicial-lemma-morphism-simplicial-sites}
$\mathcal{C}_n, f_\varphi, u_\varphi$ および
$\mathcal{C}'_n, f'_\varphi, u'_\varphi$ を状況
\ref{spaces-simplicial-situation-simplicial-site} のとおりとする。$h$ を注意
\ref{spaces-simplicial-remark-morphism-simplicial-sites} のような単体的サイト間の射とする。
このときトポスの射
$$
h_{total} : \Sh(\mathcal{C}_{total}) \to \Sh(\mathcal{C}'_{total})
$$
と可換図式
$$
\xymatrix{
\Sh(\mathcal{C}_n) \ar[d]_{g_n} \ar[r]_{h_n} &
\Sh(\mathcal{C}'_n) \ar[d]^{g'_n} \\
\Sh(\mathcal{C}_{total}) \ar[r]^{h_{total}} &
\Sh(\mathcal{C}'_{total})
}
$$
を得る。さらに
$(g'_n)^{-1} \circ h_{total, *} = h_{n, *} \circ g_n^{-1}$ である。
\end{lemma}

\begin{proof}
場合 A。$h_n$ が連続関手
$v_n : \mathcal{C}'_n \to \mathcal{C}_n$ に対応するとする。
次数 $n$ で $v_n$ を用いることにより、関手
$v_{total} : \mathcal{C}'_{total} \to \mathcal{C}_{total}$ を定義できる。
これは明らかに連続関手である（補題 \ref{spaces-simplicial-lemma-simplicial-site-site} の
被覆の定義を参照）。『サイト』節
\ref{sites-section-continuous-functors} および
\ref{sites-section-morphism-sites} で構成され研究された随伴関手の対を
$h_{total}^{-1} = v_{total, s} :
\Sh(\mathcal{C}'_{total}) \to \Sh(\mathcal{C}_{total})$ および
$h_{total, *} = v_{total}^s = v_{total}^p :
\Sh(\mathcal{C}_{total}) \to \Sh(\mathcal{C}'_{total})$ とする。
$h_{total}$ がトポスの射であることを見るには、なお
$h_{total}^{-1}$ が完全であることを確認しなければならない。まず
$(g'_n)^{-1} \circ h_{total, *} = h_{n, *} \circ g_n^{-1}$
であることに注意する。これは $\mathcal{C}'_n$ の対象 $U$ 上の切断を
計算すれば直ちに分かる。したがって、補題
\ref{spaces-simplicial-lemma-describe-sheaves-simplicial-site-site} のように
$\mathcal{C}_{total}$ 上の層 $\mathcal{F}$ を系
$(\mathcal{F}_n, \mathcal{F}(\varphi))$ とみなすと、
$h_{total, *}\mathcal{F}$ は系
$(h_{n, *}\mathcal{F}_n, h_{n, *}\mathcal{F}(\varphi))$ に対応する。
明らかに、関手
$(\mathcal{F}'_n, \mathcal{F}'(\varphi)) \to
(h_n^{-1}\mathcal{F}'_n, h_n^{-1}\mathcal{F}'(\varphi))$
はその左随伴である。随伴の一意性により、系に対する
$h_{total}^{-1}$ はこの規則で与えられる。とくに、同補題による
$\mathcal{C}_{total}$ 上の層の記述と関手 $h_n^{-1}$ の完全性から、
$h_{total}^{-1}$ は完全であり、所望のトポスの射を得る。最後に
$g_n^{-1} \circ h_{total}^{-1} =
h_n^{-1} \circ (g'_n)^{-1}$ も得られ、これにより補題の表示図式が
可換であることが示される。

\medskip\noindent
場合 B。次数 $n$ で $v_n$ を用いることにより、関手
$v_{total} : \mathcal{C}_{total} \to \mathcal{C}'_{total}$ を得る。
これは明らかに余連続関手である（補題
\ref{spaces-simplicial-lemma-simplicial-cocontinuous-site} の被覆の定義を参照）。
$h_{total}$ を $v_{total}$ に付随するトポスの射とする。補題の表示図式の
可換性は『サイト』補題 \ref{sites-lemma-composition-cocontinuous} から
直ちに従う。左随伴をとると、補題の最後の等式は
$$
h_{total}^{-1} \circ (g'_n)^{Sh}_! = g^{Sh}_{n!} \circ h_n^{-1}
$$
となる。これは、補題 \ref{spaces-simplicial-lemma-restriction-to-components-site} における
関手 $(g'_n)^{Sh}_!$ と $g^{Sh}_{n!}$ の明示的記述、
$\varphi : [m] \to [n]$ に対する等式
$h_n^{-1} \circ (f'_\varphi)^{-1} =
f_\varphi^{-1} \circ h_m^{-1}$、および $h_{total}^{-1}$ が単体的サイトの
次数 $n$ 部分への制限と可換であることが既に分かっているという事実から
直ちに従う。
\end{proof}

\begin{lemma}
\label{spaces-simplicial-lemma-direct-image-morphism-simplicial-sites}
% P10-E0610
補題 \ref{spaces-simplicial-lemma-morphism-simplicial-sites} の記法と仮定のもとで、
\linebreak
$K \in D(\mathcal{C}_{total})$ に対し
$(g'_n)^{-1}Rh_{total, *}K = Rh_{n, *}g_n^{-1}K$ である。
\end{lemma}

\begin{proof}
$\mathcal{I}^\bullet$ を $K$ を表す $\mathcal{C}_{total}$ 上の
K-入射的複体とする。このとき $g_n^{-1}K$ は
$g_n^{-1}\mathcal{I}^\bullet = \mathcal{I}_n^\bullet$ により表され、
これは補題 \ref{spaces-simplicial-lemma-restriction-injective-to-component-site} により
K-入射的である。
% P10-E0177
補題 \ref{spaces-simplicial-lemma-morphism-simplicial-sites} により
$(g'_n)^{-1}h_{total, *}\mathcal{I}^\bullet =
h_{n, *}g_n^{-1}\mathcal{I}_n^\bullet$ であり、所望の等式を得る。
\end{proof}

\begin{remark}
\label{spaces-simplicial-remark-morphism-augmentation-simplicial-sites}
$\mathcal{C}_n, f_\varphi, u_\varphi$ および
$\mathcal{C}'_n, f'_\varphi, u'_\varphi$ を状況
\ref{spaces-simplicial-situation-simplicial-site} のとおりとする。$a_0$ と $a'_0$ をそれぞれ
注意 \ref{spaces-simplicial-remark-augmentation-site} のようなサイト $\mathcal{D}$ と
$\mathcal{D}'$ への増大とする。$h$ を注意
\ref{spaces-simplicial-remark-morphism-simplicial-sites} のような単体的サイト間の射とする。
トポスの射 $h_{-1} : \Sh(\mathcal{D}) \to \Sh(\mathcal{D}')$ が
次を満たすとき、これは{\it $h$、$a_0$、$a'_0$ と両立する}という。
\begin{enumerate}
\item[(A)] $h_{-1}$ はサイトの射
$h_{-1} : \mathcal{D} \to \mathcal{D}'$ から来ており、サイトの射として
$a'_0 \circ h_0 = h_{-1} \circ a_0$ である。
\item[(B)] $h_{-1}$ は余連続関手
$v_{-1} : \mathcal{D} \to \mathcal{D}'$ から来ており、関手として
$u'_0 \circ v_0 = v_{-1} \circ u_0$ である。
\end{enumerate}
いずれの場合も、トポスの射として
$a'_0 \circ h_0 = h_{-1} \circ a_0$ である。場合 B については
『サイト』補題 \ref{sites-lemma-composition-cocontinuous}、場合 A については
『サイト』定義 \ref{sites-definition-composition-morphisms-sites} を
参照されたい。
\end{remark}

\begin{lemma}
\label{spaces-simplicial-lemma-morphism-augmentation-simplicial-sites}
$\mathcal{C}_n, f_\varphi, u_\varphi, \mathcal{D}, a_0$、
$\mathcal{C}'_n, f'_\varphi, u'_\varphi, \mathcal{D}', a'_0$、および
$n \geq -1$ に対する $h_n$ を注意
\ref{spaces-simplicial-remark-morphism-augmentation-simplicial-sites} のとおりとする。
このとき可換図式
$$
\xymatrix{
\Sh(\mathcal{C}_{total}) \ar[d]_a \ar[r]_{h_{total}} &
\Sh(\mathcal{C}'_{total}) \ar[d]^{a'} \\
\Sh(\mathcal{D}) \ar[r]^{h_{-1}} &
\Sh(\mathcal{D}')
}
$$
を得る。
\end{lemma}

\begin{proof}
% P10-E0611
射 $h$ は補題 \ref{spaces-simplicial-lemma-morphism-simplicial-sites} で定義され、
射 $a$ と $a'$ は補題 \ref{spaces-simplicial-lemma-augmentation-site} で定義される。
したがって、示すべきことは図式の可換性だけである。そのために
$a^{-1} \circ h_{-1}^{-1} = h_{total}^{-1} \circ (a')^{-1}$ を示す。
補題 \ref{spaces-simplicial-lemma-morphism-simplicial-sites} の可換図式と、補題
\ref{spaces-simplicial-lemma-describe-sheaves-simplicial-site-site} による成分を用いた
$\Sh(\mathcal{C}_{total})$ と $\Sh(\mathcal{C}'_{total})$ の記述から、
すべての $n$ に対して図式
$$
\xymatrix{
\Sh(\mathcal{C}_n) \ar[d]_{a_n} \ar[r]_{h_n} &
\Sh(\mathcal{C}'_n) \ar[d]^{a'_n} \\
\Sh(\mathcal{D}) \ar[r]^{h_{-1}} &
\Sh(\mathcal{D}')
}
$$
が可換であることを示せば十分である。$n = 0$ の場合は注意
\ref{spaces-simplicial-remark-morphism-augmentation-simplicial-sites} における仮定である。
$n > 0$ のときは $\varphi : [0] \to [n]$ を一つ選ぶ。すると所望の可換性は
$n = 0$ の場合、関係 $a_n = a_0 \circ f_\varphi$ および
$a'_n = a'_0 \circ f'_\varphi$、ならびに可換関係
$f'_\varphi \circ h_n = h_0 \circ f_\varphi$ から従う。
\end{proof}

\section{環付き単体的サイト}
\label{spaces-simplicial-section-simplicial-sites-modules}

\noindent
単体的トポスに環の層を備えさせる。

\begin{lemma}
\label{spaces-simplicial-lemma-restriction-module-to-components-site}
% P10-E0612
状況 \ref{spaces-simplicial-situation-simplicial-site} において、$\mathcal{O}$ を
$\mathcal{C}_{total}$ 上の環の層とする。標準的な環付きトポスの射
$g_n : (\Sh(\mathcal{C}_n), \mathcal{O}_n) \to
(\Sh(\mathcal{C}_{total}), \mathcal{O})$
で、その底のトポス上では補題
\ref{spaces-simplicial-lemma-restriction-to-components-site} の射 $g_n$ と一致するものが
存在する。関手
$g_n^* : \textit{Mod}(\mathcal{O}) \to \textit{Mod}(\mathcal{O}_n)$
は左随伴 $g_{n!}$ をもつ。
% P10-E0613
$\textit{Mod}(\mathcal{O}_n)$-加群の $\mathcal{G}$ に対し、
$g_{n!}\mathcal{G}$ の $\mathcal{C}_m$ への制限は
$$
\bigoplus\nolimits_{\varphi : [n] \to [m]} f_\varphi^*\mathcal{G}
$$
である。ここで
$f_\varphi : (\Sh(\mathcal{C}_m), \mathcal{O}_m) \to
(\Sh(\mathcal{C}_n), \mathcal{O}_n)$ は環付きトポスの射であり、
底のトポス上では先に定義した $f_\varphi$ と一致し、環の層上では写像
$\mathcal{O}(\varphi) : f_\varphi^{-1}\mathcal{O}_n \to \mathcal{O}_m$
を用いる。
\end{lemma}

\begin{proof}
補題 \ref{spaces-simplicial-lemma-restriction-to-components-site} により
$g_n^{-1}\mathcal{O} = \mathcal{O}_n$ なので、上の環付きトポスの射を得る。
『サイト上の加群』補題 \ref{sites-modules-lemma-lower-shriek-modules} により
随伴 $g_{n!}$ を得る。$g_{n!}$ の公式を証明するため、まず
$\mathcal{C}_{total}$ 上の $\mathcal{O}$-加群の層 $\mathcal{H}$ で、
その次数 $m$ の成分が $\mathcal{O}_m$-加群
$$
\mathcal{H}_m =
\bigoplus\nolimits_{\varphi : [n] \to [m]} f_\varphi^*\mathcal{G}
$$
であるものを定義する。写像 $\psi : [m] \to [m']$ が与えられると、
写像 $\mathcal{H}(\psi) : f_\psi^{-1}\mathcal{H}_m \to \mathcal{H}_{m'}$
は各成分上で
$$
f_\psi^{-1} f_\varphi^*\mathcal{G} \to
f_\psi^* f_\varphi^*\mathcal{G} \to
f_{\psi \circ \varphi}^*\mathcal{G}
$$
により与えられる。この写像
$f_\psi^{-1}\mathcal{H}_m \to \mathcal{H}_{m'}$ は
$\mathcal{O}(\psi) : f_\psi^{-1}\mathcal{O}_m \to \mathcal{O}_{m'}$
に関して半線形なので、実際に $\mathcal{O}$-加群を定める
（補題 \ref{spaces-simplicial-lemma-describe-sheaves-simplicial-site-site} を用いる）。
次に、補題 \ref{spaces-simplicial-lemma-restriction-to-components-site} の証明と同様に進めて、
$$
\Mor_{\mathcal{O}_n}(\mathcal{G}, \mathcal{F}_n) =
\Mor_{\mathcal{O}}(\mathcal{H}, \mathcal{F})
$$
を直接証明する。したがって所望のとおり
$\mathcal{H} = g_{n!}\mathcal{G}$ である。
\end{proof}

\begin{lemma}
\label{spaces-simplicial-lemma-restriction-injective-to-component-limp}
% P10-E0614
状況 \ref{spaces-simplicial-situation-simplicial-site} において、$\mathcal{O}$ を
$\mathcal{C}_{total}$ 上の環の層とする。$\mathcal{I}$ が
$\textit{Mod}(\mathcal{O})$ において入射的ならば、$\mathcal{I}_n$ は
$\mathcal{C}_n$ 上の完全非輪状層である。
\end{lemma}

\begin{proof}
これは、『サイト上のコホモロジー』補題
\ref{sites-cohomology-lemma-pullback-injective-limp} を、包含関手
$\mathcal{C}_n \to \mathcal{C}_{total}$ と、補題
\ref{spaces-simplicial-lemma-restriction-to-components-site} で証明されたその性質に
適用すれば従う。
\end{proof}

\begin{lemma}
\label{spaces-simplicial-lemma-exactness-g-shriek-modules}
補題 \ref{spaces-simplicial-lemma-restriction-module-to-components-site} と同じ仮定のもとで、
すべての $\varphi : [n] \to [m]$ に対して写像
$f_\varphi^{-1}\mathcal{O}_n \to \mathcal{O}_m$ が平坦ならば、関手
$g_{n!} : \textit{Mod}(\mathcal{O}_n) \to \textit{Mod}(\mathcal{O})$
は完全である。
\end{lemma}

\begin{proof}
$g_{n!}\mathcal{G}$ は $\mathcal{O}$-加群であり、その次数 $m$ 部分は
$\mathcal{O}_m$-加群
$$
\bigoplus\nolimits_{\varphi : [n] \to [m]} f_\varphi^*\mathcal{G}
$$
であることを思い出そう。ここで環付きトポスの射
$f_\varphi : (\Sh(\mathcal{C}_m), \mathcal{O}_m) \to
(\Sh(\mathcal{C}_n), \mathcal{O}_n)$ は、補題の主張にある写像
$f_\varphi^{-1}\mathcal{O}_n \to \mathcal{O}_m$ を用いる。
これらの写像が平坦ならば、$f_\varphi^*$ は完全である
（『サイト上の加群』補題
\ref{sites-modules-lemma-flat-pullback-exact}）。サイト
$\mathcal{C}_{total}$ の定義から、これらの関手が所望の完全性を
もつことが分かり、結論を得る。
\end{proof}

\begin{lemma}
\label{spaces-simplicial-lemma-restriction-injective-to-component-site-module}
% P10-E0615
状況 \ref{spaces-simplicial-situation-simplicial-site} において、$\mathcal{O}$ を
$\mathcal{C}_{total}$ 上の環の層とし、すべての
$\varphi : [n] \to [m]$ に対して
$f_\varphi^{-1}\mathcal{O}_n \to \mathcal{O}_m$ が平坦であるとする。
$\mathcal{I}$ が $\textit{Mod}(\mathcal{O})$ において入射的ならば、
$\mathcal{I}_n$ は $\textit{Mod}(\mathcal{O}_n)$ において入射的である。
\end{lemma}

\begin{proof}
これは『ホモロジー』補題
\ref{homology-lemma-adjoint-preserve-injectives} と補題
\ref{spaces-simplicial-lemma-exactness-g-shriek-modules} から従う。
\end{proof}

\section{環付き単体的サイトの射}
\label{spaces-simplicial-section-morphism-simplicial-sites-modules}

\noindent
節 \ref{spaces-simplicial-section-morphism-simplicial-sites} の議論を続ける。

\begin{remark}
\label{spaces-simplicial-remark-morphism-simplicial-sites-modules}
$\mathcal{C}_n, f_\varphi, u_\varphi$ および
$\mathcal{C}'_n, f'_\varphi, u'_\varphi$ を状況
\ref{spaces-simplicial-situation-simplicial-site} のとおりとする。
% P10-E0616
$\mathcal{O}$ と $\mathcal{O}'$ をそれぞれ $\mathcal{C}_{total}$ と
$\mathcal{C}'_{total}$ 上の環の層とする。$h$ が注意
\ref{spaces-simplicial-remark-morphism-simplicial-sites} のような単体的サイト間の射であり、
$h^\sharp : h_{total}^{-1}\mathcal{O}' \to \mathcal{O}$、または同値な形で
$h^\sharp : \mathcal{O}' \to h_{total, *}\mathcal{O}$ が環の層の準同型で
あるとき、$(h, h^\sharp)$ を{\it 環付き単体的サイト間の射}という。
\end{remark}

\begin{lemma}
\label{spaces-simplicial-lemma-morphism-simplicial-sites-modules}
$\mathcal{C}_n, f_\varphi, u_\varphi$ および
$\mathcal{C}'_n, f'_\varphi, u'_\varphi$ を状況
\ref{spaces-simplicial-situation-simplicial-site} のとおりとする。
% P10-E0617
$\mathcal{O}$ と $\mathcal{O}'$ をそれぞれ $\mathcal{C}_{total}$ と
$\mathcal{C}'_{total}$ 上の環の層とする。
% P10-E0618
$(h, h^\sharp)$ を注意 \ref{spaces-simplicial-remark-morphism-simplicial-sites-modules} の
ような単体的サイト間の射とする。このとき環付きトポスの射
$$
h_{total} :
(\Sh(\mathcal{C}_{total}), \mathcal{O})
\to
(\Sh(\mathcal{C}'_{total}), \mathcal{O}')
$$
と、環付きトポスの可換図式
$$
\xymatrix{
(\Sh(\mathcal{C}_n), \mathcal{O}_n) \ar[d]_{g_n} \ar[r]_{h_n} &
(\Sh(\mathcal{C}'_n), \mathcal{O}'_n) \ar[d]^{g'_n} \\
(\Sh(\mathcal{C}_{total}), \mathcal{O}) \ar[r]^{h_{total}} &
(\Sh(\mathcal{C}'_{total}), \mathcal{O}')
}
$$
を得る。ここで $g_n$ と $g'_n$ は補題
\ref{spaces-simplicial-lemma-restriction-module-to-components-site} のとおりである。
さらに
% P10-E0619
$(g'_n)^* \circ h_{total, *} = h_{n, *} \circ g_n^*$
が関手 $\textit{Mod}(\mathcal{O}) \to \textit{Mod}(\mathcal{O}'_n)$ として
成り立つ。
\end{lemma}

\begin{proof}
環の層を追跡すれば、補題 \ref{spaces-simplicial-lemma-morphism-simplicial-sites} および
\ref{spaces-simplicial-lemma-restriction-module-to-components-site} から従う。
小さな注意点は、$h_n$ を環付きトポスの射として定義するために
$h_n^\sharp = g_n^{-1}h^\sharp :
g_n^{-1}h_{total}^{-1}\mathcal{O}' \to g_n^{-1}\mathcal{O}$
とおくことである。これは\linebreak
$g_n^{-1}h_{total}^{-1}\mathcal{O}' =
\allowbreak h_n^{-1}(g'_n)^{-1}\mathcal{O}' =
\allowbreak
h_n^{-1}\mathcal{O}'_n$ および
$g_n^{-1}\mathcal{O} = \mathcal{O}_n$ なので意味をもつ。
$\mathcal{F}$ が $\mathcal{O}$-加群の層ならば
$g_n^*\mathcal{F} = g_n^{-1}\mathcal{F}$ であり、$g'_n$ についても
同様であることに注意する。これは、
$(g'_n)^* \circ h_{total, *} = h_{n, *} \circ g_n^*$ が補題
\ref{spaces-simplicial-lemma-morphism-simplicial-sites} の対応する主張から従う理由を
説明している。
\end{proof}

\begin{lemma}
\label{spaces-simplicial-lemma-direct-image-morphism-simplicial-sites-modules}
% P10-E0620
補題 \ref{spaces-simplicial-lemma-morphism-simplicial-sites-modules} の記法と仮定のもとで、
\linebreak
$K \in D(\mathcal{O})$ に対し
$(g'_n)^*Rh_{total, *}K = Rh_{n, *}g_n^*K$ である。
\end{lemma}

\begin{proof}
補題 \ref{spaces-simplicial-lemma-restriction-module-to-components-site} の構成により
$g_n^{-1}\mathcal{O} = \mathcal{O}_n$ なので、$g_n^* = g_n^{-1}$ であることを
思い出そう。とくに $g_n^*$ は完全であり、$Lg_n^*$ は加群の任意の代表複体に
$g_n^*$ を適用することで与えられる。$g'_n$ についても同様である。
標準的な基底変換写像
$(g'_n)^*Rh_{total, *}K \to Rh_{n, *}g_n^*K$ が存在する。『サイト上の
コホモロジー』注意 \ref{sites-cohomology-remark-base-change} を参照されたい。
『サイト上のコホモロジー』補題
\ref{sites-cohomology-lemma-modules-abelian-unbounded} により、これの
$D(\mathcal{C}'_n)$ における像は写像
$(g'_n)^{-1}Rh_{total, *}K_{ab} \to Rh_{n, *}g_n^{-1}K_{ab}$ である。
ここで $K_{ab}$ は $D(\mathcal{C}_{total})$ における $K$ の像である。
これが同型であることは補題
\ref{spaces-simplicial-lemma-direct-image-morphism-simplicial-sites} で証明したので、結論が従う。
\end{proof}






\section{単体的サイト上のコホモロジー}
\label{spaces-simplicial-section-cohomology-simplicial-sites}

\noindent
$\mathcal{C}$ を状況 \ref{spaces-simplicial-situation-simplicial-site} のとおりとする。
% P10-E0621
以下の補題の主張では、$g_n : \Sh(\mathcal{C}_n) \to
\Sh(\mathcal{C}_{total})$ を補題
\ref{spaces-simplicial-lemma-restriction-to-components-site} のトポスの射とする。
$\varphi : [m] \to [n]$ が $\Delta$ の射ならば、トポスの図式
$$
\xymatrix{
\Sh(\mathcal{C}_n) \ar[rd]_{g_n} \ar[rr]_{f_\varphi} & &
\Sh(\mathcal{C}_m) \ar[ld]^{g_m} \\
& \Sh(\mathcal{C}_{total})
}
$$
は可換ではないが、写像
$\mathcal{F}(\varphi) : f_\varphi^{-1}\mathcal{F}_m \to \mathcal{F}_n$
から生じる $2$-射 $g_n \to g_m \circ f_\varphi$ が存在する。
『サイト』節 \ref{sites-section-2-category} を参照されたい。

\begin{lemma}
\label{spaces-simplicial-lemma-simplicial-resolution-Z-site}
状況 \ref{spaces-simplicial-situation-simplicial-site} において上の記法を用いると、
$\mathcal{C}_{total}$ 上のアーベル層の複体
$$
\ldots \to g_{2!}\mathbf{Z} \to g_{1!}\mathbf{Z} \to g_{0!}\mathbf{Z}
$$
で、$\mathcal{C}_{total}$ 上の値 $\mathbf{Z}$ の定数層の分解をなすものが存在する。
\end{lemma}

\begin{proof}
補題 \ref{spaces-simplicial-lemma-restriction-to-components-site} における関手 $g_{n!}$ の記述を、
以下では断りなく用いる。複体の写像として $\sum (-1)^i d^n_i$ をとる。
ここで $d^n_i : g_{n!}\mathbf{Z} \to g_{n - 1!}\mathbf{Z}$ は、
$\delta^n_i : [n - 1] \to [n]$ とラベル付けされた因子に対応する写像
$\mathbf{Z} \to
\bigoplus_{[n - 1] \to [n]} \mathbf{Z} = g_n^{-1}g_{n - 1!}\mathbf{Z}$
に随伴する写像である。この複体に $g_m^{-1}$ を適用すると、
$\mathcal{C}_m$ 上の複体
$$
\ldots \to
\bigoplus\nolimits_{\alpha \in \Mor_\Delta([2], [m])]} \mathbf{Z} \to
\bigoplus\nolimits_{\alpha \in \Mor_\Delta([1], [m])]} \mathbf{Z} \to
\bigoplus\nolimits_{\alpha \in \Mor_\Delta([0], [m])]} \mathbf{Z}
$$
が得られる。言い換えると、これは単体集合 $\Delta[m]$ 上の自由アーベル層に
付随する複体である。『単体』例
\ref{simplicial-example-simplex-simplicial-set} を参照されたい。
$\Delta[m]$ は $\Delta[0]$ とホモトピー同値であり（『単体』例
\ref{simplicial-example-simplex-contractible}）、また「自由アーベル層をとる」操作は
関手なので、上の複体は $\Delta[0]$ 上の自由アーベル層とホモトピー同値である
（『単体』注意 \ref{simplicial-remark-homotopy-better} および補題
\ref{simplicial-lemma-homotopy-equivalence-s-N}）。この複体は正の次数で
非輪状であり、次数 $0$ では $\mathbf{Z}$ に等しい。
\end{proof}

\begin{lemma}
\label{spaces-simplicial-lemma-cech-complex}
% P10-E0622
状況 \ref{spaces-simplicial-situation-simplicial-site} において、$\mathcal{F}$ を
$\mathcal{C}_{total}$ 上のアーベル層とする。このとき標準的な複体
$$
0 \to \Gamma(\mathcal{C}_{total}, \mathcal{F}) \to
\Gamma(\mathcal{C}_0, \mathcal{F}_0) \to
\Gamma(\mathcal{C}_1, \mathcal{F}_1) \to
\Gamma(\mathcal{C}_2, \mathcal{F}_2) \to \ldots
$$
が存在し、これは次数 $-1, 0$ で完全である。また $\mathcal{F}$ が入射的ならば
すべての次数で完全である。
\end{lemma}

\begin{proof}
$\Hom(\mathbf{Z}, \mathcal{F}) = \Gamma(\mathcal{C}_{total}, \mathcal{F})$
および
$\Hom(g_{n!}\mathbf{Z}, \mathcal{F}) = \Gamma(\mathcal{C}_n, \mathcal{F}_n)$
であることに注意する。したがって、この補題は補題
\ref{spaces-simplicial-lemma-simplicial-resolution-Z-site} と、$\mathcal{F}$ が入射的ならば
$\Hom(-, \mathcal{F})$ が完全であるという事実から直ちに従う。
\end{proof}

\begin{lemma}
\label{spaces-simplicial-lemma-simplicial-sheaf-cohomology-site}
% P10-E0623
状況 \ref{spaces-simplicial-situation-simplicial-site} において、$K$ を
$D^+(\mathcal{C}_{total})$ の対象とする。このときスペクトル系列
$(E_r, d_r)_{r \geq 0}$ で
$$
E_1^{p, q} = H^q(\mathcal{C}_p, K_p),\quad
d_1^{p, q} : E_1^{p, q} \to E_1^{p + 1, q}
$$
をもち、$H^{p + q}(\mathcal{C}_{total}, K)$ に収束するものが存在する。
このスペクトル系列は $K$ に関して関手的である。
\end{lemma}

\begin{proof}
$K$ を表す入射的対象からなる下に有界な複体を $\mathcal{I}^\bullet$ とする。
項が
$$
A^{p, q} = \Gamma(\mathcal{C}_p, \mathcal{I}^q_p)
$$
である二重複体を考える。ここで水平方向の射は補題
\ref{spaces-simplicial-lemma-cech-complex} から、垂直方向の射は複体
$\mathcal{I}^\bullet$ の微分から来る。二重複体の各行は正の次数で完全であり、
次数 $0$ では $\Gamma(\mathcal{C}_{total}, \mathcal{I}^q)$ になる。
一方、$\mathcal{C}_p$ への制限は完全なので（補題
\ref{spaces-simplicial-lemma-restriction-to-components-site}）、複体 $\mathcal{I}_p^\bullet$ は
$D(\mathcal{C}_p)$ において $K_p$ を表す。層 $\mathcal{I}_p^q$ は
$\mathcal{C}_p$ 上の入射的アーベル層である（補題
\ref{spaces-simplicial-lemma-restriction-injective-to-component-site}）。したがって各列の
コホモロジーは群 $H^q(\mathcal{C}_p, K_p)$ を計算する。
『ホモロジー』補題 \ref{homology-lemma-first-quadrant-ss} および
\ref{homology-lemma-double-complex-gives-resolution} を適用して結論を得る。
\end{proof}

\begin{remark}
\label{spaces-simplicial-remark-simplicial-cohomology-site}
% P10-E0624
補題 \ref{spaces-simplicial-lemma-simplicial-sheaf-cohomology-site} と同じ仮定と記法を用いるが、
$K$ が $D(\mathcal{C}_{total})$ において下に有界であるとは仮定しない。この場合にも
自然なスペクトル系列が存在すると主張する。実際、$\mathcal{I}^\bullet$ を、
$K$ を表す、各項が入射的な $\mathcal{C}_{total}$ 上の層の K-入射的複体とする。
すると
\begin{align*}
R\Gamma(\mathcal{C}_{total}, K)
& =
R\Hom(\mathbf{Z}, K) \\
& =
R\Hom(
\ldots \to g_{2!}\mathbf{Z} \to g_{1!}\mathbf{Z} \to g_{0!}\mathbf{Z}, K) \\
& =
\Gamma(\mathcal{C}_{total},
\SheafHom^\bullet(
\ldots \to g_{2!}\mathbf{Z} \to g_{1!}\mathbf{Z} \to g_{0!}\mathbf{Z},
\mathcal{I}^\bullet)) \\
& =
\text{Tot}_\pi(A^{\bullet, \bullet})
\end{align*}
である。ここで $A^{\bullet, \bullet}$ は、項が
$A^{p, q} = \Gamma(\mathcal{C}_p, \mathcal{I}^q_p)$ である二重複体であり、
$\text{Tot}_\pi$ はこの二重複体の積全化を表す。第一の等式は任意のサイトで
成り立つ。第二の等式は補題 \ref{spaces-simplicial-lemma-simplicial-resolution-Z-site} により
成り立つ。第三の等式は $\mathcal{I}^\bullet$ が K-入射的であることから
成り立つ。『サイト上のコホモロジー』節
\ref{sites-cohomology-section-hom-complexes} および
\ref{sites-cohomology-section-internal-hom} を参照されたい。最後の等式は
$\SheafHom^\bullet$ の構成と
$\Hom(g_{p!}\mathbf{Z}, \mathcal{I}^q) =
\Gamma(\mathcal{C}_p, \mathcal{I}^q_p)$ であることから成り立つ。
次に、$\text{Tot}_\pi(A^{\bullet, \bullet})$ を
$F^i\text{Tot}^n_\pi(A^{\bullet, \bullet}) =
\prod_{p + q = n,\ p \geq i} A^{p, q}$ をもつフィルター付き複体とみなすことで、
所望のスペクトル系列を得る。得られるスペクトル系列は、『ホモロジー』節
\ref{homology-section-double-complex} のスペクトル系列
$({}'E_r, {}'d_r)_{r \geq 0}$（そこでは直和全化の場合を扱っている）と同様に
振る舞うが、正則性、有界性、収束性および収束先に関する問題は別である。
とくに、補題 \ref{spaces-simplicial-lemma-simplicial-sheaf-cohomology-site} と同様に
$E_1^{p, q} = H^q(\mathcal{C}_p, K_p)$ を得る。
\end{remark}

\begin{lemma}
\label{spaces-simplicial-lemma-simplicial-sheaf-cohomology-site-zero}
% P10-E0625
状況 \ref{spaces-simplicial-situation-simplicial-site} において、$K$ を
$D(\mathcal{C}_{total})$ の対象とする。
\begin{enumerate}
\item すべての $p \geq 0$ に対して
$H^{-p}(\mathcal{C}_p, K_p) = 0$ ならば、
$H^0(\mathcal{C}_{total}, K) = 0$ である。
\item すべての $p \geq 0$ に対して $R\Gamma(\mathcal{C}_p, K_p) = 0$
ならば、$R\Gamma(\mathcal{C}_{total}, K) = 0$ である。
\end{enumerate}
\end{lemma}

\begin{proof}
注意 \ref{spaces-simplicial-remark-simplicial-cohomology-site} の記法を用いると、
$R\Gamma(\mathcal{C}_{total}, K)$ は
$\text{Tot}_\pi(A^{\bullet, \bullet})$ によって表される。
(1) の仮定は $H^{-p}(A^{p, \bullet}) = 0$ を意味する。
したがって (1) の消滅は『代数詳論』補題
\ref{more-algebra-lemma-vanishing-coh-prod-totalization} から従う。
(2) は (1) とシフトをとることから従う。
\end{proof}

\begin{lemma}
\label{spaces-simplicial-lemma-sanity-check}
$\mathcal{C}$ を状況 \ref{spaces-simplicial-situation-simplicial-site} のとおりとする。
$U \in \Ob(\mathcal{C}_n)$ および
$\mathcal{F} \in \textit{Ab}(\mathcal{C}_{total})$ とする。このとき
$H^p(U, \mathcal{F}) = H^p(U, g_n^{-1}\mathcal{F})$ である。ここで左辺では
$U$ を $\mathcal{C}_{total}$ の対象とみなしている。
\end{lemma}

\begin{proof}
「$U$ を $\mathcal{C}_{total}$ の対象とみなす」ことは、場合 (A) では補題
\ref{spaces-simplicial-lemma-simplicial-site-site}、場合 (B) では補題
\ref{spaces-simplicial-lemma-simplicial-cocontinuous-site} における $\mathcal{C}_{total}$ の構成により
説明される。すると等式は補題
\ref{spaces-simplicial-lemma-restriction-injective-to-component-site} とコホモロジーの定義から従う。
\end{proof}






\section{単体的サイトのコホモロジーと増大}
\label{spaces-simplicial-section-cohomology-augmentation-simplicial-sites}

\noindent
状況 \ref{spaces-simplicial-situation-simplicial-site} のような単体的サイト
$\mathcal{C}$ を考える。$a_0$ を注意 \ref{spaces-simplicial-remark-augmentation-site} のような
サイト $\mathcal{D}$ への増大とする。補題 \ref{spaces-simplicial-lemma-augmentation-site} により、
トポスの射
$$
a : \Sh(\mathcal{C}_{total}) \longrightarrow \Sh(\mathcal{D})
$$
および、補題 \ref{spaces-simplicial-lemma-restriction-to-components-site} のようなトポスの射
$g_n : \Sh(\mathcal{C}_n) \to \Sh(\mathcal{C}_{total})$
を得る。合成 $a \circ g_n$ を
$a_n : \Sh(\mathcal{C}_n) \to \Sh(\mathcal{D})$ と表す。
さらに、単体的構造はトポスの射
$f_\varphi : \Sh(\mathcal{C}_n) \to \Sh(\mathcal{C}_m)$ を与え、
% P10-E0629
すべての $\varphi : [m] \to [n]$ に対して
$a_n \circ f_\varphi = a_m$ が成り立つ。

\begin{lemma}
\label{spaces-simplicial-lemma-simplicial-resolution-augmentation}
状況 \ref{spaces-simplicial-situation-simplicial-site} において、$a_0$ を注意
\ref{spaces-simplicial-remark-augmentation-site} のようなサイト $\mathcal{D}$ への増大とする。
$\mathcal{D}$ 上の任意のアーベル層 $\mathcal{G}$ に対し、
$\mathcal{C}_{total}$ 上のアーベル層の完全複体
$$
\ldots \to
g_{2!}(a_2^{-1}\mathcal{G}) \to
g_{1!}(a_1^{-1}\mathcal{G}) \to
g_{0!}(a_0^{-1}\mathcal{G}) \to
a^{-1}\mathcal{G} \to 0
$$
が存在する。
\end{lemma}

\begin{proof}
読者には、まず補題 \ref{spaces-simplicial-lemma-simplicial-resolution-Z-site} の証明を読むことを
勧める。以下では、補題 \ref{spaces-simplicial-lemma-augmentation-site} と、補題
\ref{spaces-simplicial-lemma-restriction-to-components-site} における関手 $g_{n!}$ の記述を
断りなく用いる。とくに $g_{n!}(a_n^{-1}\mathcal{G})$ は、その
$\mathcal{C}_m$ への制限が層
$$
\bigoplus\nolimits_{\varphi : [n] \to [m]} f_\varphi^{-1}a_n^{-1}\mathcal{G} =
\bigoplus\nolimits_{\varphi : [n] \to [m]} a_m^{-1}\mathcal{G}
$$
であるような $\mathcal{C}_{total}$ 上の層である。
複体の写像として $\sum (-1)^i d^n_i$ をとる。ここで
$d^n_i : g_{n!}(a_n^{-1}\mathcal{G}) \to
g_{n - 1!}(a_{n - 1}^{-1}\mathcal{G})$ は、
$\delta^n_i : [n - 1] \to [n]$ とラベル付けされた因子に対応する写像
$a_n^{-1}\mathcal{G} \to \bigoplus_{[n - 1] \to [n]} a_n^{-1}\mathcal{G} =
g_n^{-1}g_{n - 1!}(a_{n - 1}^{-1}\mathcal{G})$
に随伴する写像である。写像
$g_{0!}(a_0^{-1}\mathcal{G}) \to a^{-1}\mathcal{G}$ は
$a_0^{-1}\mathcal{G}$ の恒等写像に随伴する。
すると、次数 $\ldots, 2, 1, 0$ の鎖複体に $g_m^{-1}$ を適用して得られるのは、
$\mathcal{C}_m$ 上の複体
$$
\begin{aligned}
&\ldots \to
\bigoplus\nolimits_{\alpha \in \Mor_\Delta([2], [m])]} a_m^{-1}\mathcal{G}
\\
&\to
\bigoplus\nolimits_{\alpha \in \Mor_\Delta([1], [m])]} a_m^{-1}\mathcal{G}
\\
&\to
\bigoplus\nolimits_{\alpha \in \Mor_\Delta([0], [m])]} a_m^{-1}\mathcal{G}
\end{aligned}
$$
である。これは $a_m^{-1}\mathcal{G}$ と複体
$$
\ldots \to
\bigoplus\nolimits_{\alpha \in \Mor_\Delta([2], [m])]} \mathbf{Z} \to
\bigoplus\nolimits_{\alpha \in \Mor_\Delta([1], [m])]} \mathbf{Z} \to
\bigoplus\nolimits_{\alpha \in \Mor_\Delta([0], [m])]} \mathbf{Z}
$$
との定数層 $\mathbf{Z}$ 上のテンソル積に等しい。後者は補題
\ref{spaces-simplicial-lemma-simplicial-resolution-Z-site} の証明で扱った複体である。
そこで、この複体は次数 $0$ に置かれた $\mathbf{Z}$ とホモトピー同値であることを
見たので、証明が完了する。
\end{proof}

\begin{lemma}
\label{spaces-simplicial-lemma-augmentation-cech-complex}
状況 \ref{spaces-simplicial-situation-simplicial-site} において、$a_0$ を注意
\ref{spaces-simplicial-remark-augmentation-site} のようなサイト $\mathcal{D}$ への増大とする。
$\mathcal{C}_{total}$ 上のアーベル層 $\mathcal{F}$ に対し、
$\mathcal{D}$ 上の標準的な複体
$$
0 \to a_*\mathcal{F} \to a_{0, *}\mathcal{F}_0 \to a_{1, *}\mathcal{F}_1 \to
a_{2, *}\mathcal{F}_2 \to \ldots
$$
が存在し、これは次数 $-1, 0$ で完全である。また $\mathcal{F}$ が入射的ならば
すべての次数で完全である。
\end{lemma}

\begin{proof}
複体を構成するには、米田の補題により、$\mathcal{D}$ 上の任意のアーベル層
$\mathcal{G}$ に対して、$\mathcal{G}$ に関して関手的な複体
$$
0 \to \Hom(\mathcal{G}, a_*\mathcal{F}) \to
\Hom(\mathcal{G}, a_{0, *}\mathcal{F}_0) \to
\Hom(\mathcal{G}, a_{1, *}\mathcal{F}_1) \to \ldots
$$
を構成すれば十分である。そのためには補題
\ref{spaces-simplicial-lemma-simplicial-resolution-augmentation} の完全複体に
$\Hom(-, \mathcal{F})$ を適用し、引戻しと順像の随伴性を用いる。
次数 $-1, 0$ における完全性は、$\Hom(-, \mathcal{F})$ が左完全なので
この構成から従う。$\mathcal{F}$ が入射的アーベル層ならば、
$\Hom(-, \mathcal{F})$ が完全なので複体は完全である。
\end{proof}

\begin{lemma}
\label{spaces-simplicial-lemma-augmentation-spectral-sequence}
状況 \ref{spaces-simplicial-situation-simplicial-site} において、$a_0$ を注意
\ref{spaces-simplicial-remark-augmentation-site} のようなサイト $\mathcal{D}$ への増大とする。
任意の $K$ を $D^+(\mathcal{C}_{total})$ の対象とすると、スペクトル系列
$(E_r, d_r)_{r \geq 0}$ で
$$
E_1^{p, q} = R^qa_{p, *} K_p,\quad
d_1^{p, q} : E_1^{p, q} \to E_1^{p + 1, q}
$$
をもち、$R^{p + q}a_*K$ に収束するものが存在する。このスペクトル系列は
$K$ に関して関手的である。
\end{lemma}

\begin{proof}
$K$ を表す入射的対象からなる下に有界な複体を $\mathcal{I}^\bullet$ とする。
項が
$$
A^{p, q} = a_{p, *}\mathcal{I}^q_p
$$
である二重複体を考える。ここで水平方向の射は補題
\ref{spaces-simplicial-lemma-augmentation-cech-complex} から、垂直方向の射は複体
$\mathcal{I}^\bullet$ の微分から来る。二重複体の各行は正の次数で完全であり、
次数 $0$ では $a_*\mathcal{I}^q$ になる。一方、$\mathcal{C}_p$ への制限は
完全なので（補題 \ref{spaces-simplicial-lemma-restriction-to-components-site}）、複体
$\mathcal{I}_p^\bullet$ は $D(\mathcal{C}_p)$ において $K_p$ を表す。
層 $\mathcal{I}_p^q$ は $\mathcal{C}_p$ 上の入射的アーベル層である
（補題 \ref{spaces-simplicial-lemma-restriction-injective-to-component-site}）。したがって各列の
コホモロジーは $R^qa_{p, *}K_p$ を計算する。『ホモロジー』補題
\ref{homology-lemma-first-quadrant-ss} および
\ref{homology-lemma-double-complex-gives-resolution} を適用して結論を得る。
\end{proof}






\section{環付き単体的サイト上のコホモロジー}
\label{spaces-simplicial-section-cohomology-simplicial-sites-modules}

\noindent
本節は、加群層に対する節 \ref{spaces-simplicial-section-cohomology-simplicial-sites} の類似である。

\medskip\noindent
状況 \ref{spaces-simplicial-situation-simplicial-site} において、$\mathcal{O}$ を
$\mathcal{C}_{total}$ 上の環の層とする。
% P10-E0626
以下の補題の主張では、
$g_n : (\Sh(\mathcal{C}_n), \mathcal{O}_n) \to
(\Sh(\mathcal{C}_{total}), \mathcal{O})$ を補題
\ref{spaces-simplicial-lemma-restriction-module-to-components-site} の環付きトポスの射とする。
$\varphi : [m] \to [n]$ が $\Delta$ の射ならば、環付きトポスの図式
$$
\xymatrix{
(\Sh(\mathcal{C}_n), \mathcal{O}_n) \ar[rd]_{g_n} \ar[rr]_{f_\varphi} & &
(\Sh(\mathcal{C}_m), \mathcal{O}_m) \ar[ld]^{g_m} \\
& (\Sh(\mathcal{C}_{total}), \mathcal{O})
}
$$
は可換ではないが、写像
$\mathcal{F}(\varphi) : f_\varphi^{-1}\mathcal{F}_m \to \mathcal{F}_n$
から生じる $2$-射 $g_n \to g_m \circ f_\varphi$ が存在する。
『サイト』節 \ref{sites-section-2-category} を参照されたい。

\begin{lemma}
\label{spaces-simplicial-lemma-simplicial-resolution-ringed}
状況 \ref{spaces-simplicial-situation-simplicial-site} において、$\mathcal{O}$ を
$\mathcal{C}_{total}$ 上の環の層とする。このとき $\mathcal{O}$-加群の複体
$$
\ldots \to g_{2!}\mathcal{O}_2 \to g_{1!}\mathcal{O}_1 \to g_{0!}\mathcal{O}_0
$$
で、$\mathcal{O}$ の分解をなすものが存在する。ここで $g_{n!}$ は補題
\ref{spaces-simplicial-lemma-restriction-module-to-components-site} のとおりである。
\end{lemma}

\begin{proof}
% P10-E0627
補題 \ref{spaces-simplicial-lemma-restriction-to-components-site} で与えられた $g_{n!}$ の記述を
用いる。複体の写像として $\sum (-1)^i d^n_i$ をとる。ここで
$d^n_i : g_{n!}\mathcal{O}_n \to g_{n - 1!}\mathcal{O}_{n - 1}$ は、
$\delta^n_i : [n - 1] \to [n]$ とラベル付けされた因子に対応する写像
$\mathcal{O}_n \to \bigoplus_{[n - 1] \to [n]} \mathcal{O}_n =
g_n^*g_{n - 1!}\mathcal{O}_{n - 1}$ に随伴する写像である。
この複体に $g_m^{-1}$ を適用すると、$\mathcal{C}_m$ 上の複体
$$
\begin{aligned}
&\ldots \to
\bigoplus\nolimits_{\alpha \in \Mor_\Delta([2], [m])]} \mathcal{O}_m
\\
&\to
\bigoplus\nolimits_{\alpha \in \Mor_\Delta([1], [m])]} \mathcal{O}_m
\\
&\to
\bigoplus\nolimits_{\alpha \in \Mor_\Delta([0], [m])]} \mathcal{O}_m
\end{aligned}
$$
が得られる。言い換えると、これは単体集合 $\Delta[m]$ 上の自由
$\mathcal{O}_m$-加群に付随する複体である。『単体』例
\ref{simplicial-example-simplex-simplicial-set} を参照されたい。
$\Delta[m]$ は $\Delta[0]$ とホモトピー同値であり（『単体』例
\ref{simplicial-example-simplex-contractible}）、また
% P10-E0628
「自由アーベル層をとる」操作は関手なので、上の複体は $\Delta[0]$ 上の
自由アーベル層とホモトピー同値である（『単体』注意
\ref{simplicial-remark-homotopy-better} および補題
\ref{simplicial-lemma-homotopy-equivalence-s-N}）。この複体は正の次数で
非輪状であり、次数 $0$ では $\mathcal{O}_m$ に等しい。
\end{proof}

\begin{lemma}
\label{spaces-simplicial-lemma-cech-complex-modules}
状況 \ref{spaces-simplicial-situation-simplicial-site} において、$\mathcal{O}$ を環の層とする。
$\mathcal{F}$ を $\mathcal{O}$-加群の層とする。このとき標準的な複体
$$
0 \to \Gamma(\mathcal{C}_{total}, \mathcal{F}) \to
\Gamma(\mathcal{C}_0, \mathcal{F}_0) \to
\Gamma(\mathcal{C}_1, \mathcal{F}_1) \to
\Gamma(\mathcal{C}_2, \mathcal{F}_2) \to \ldots
$$
が存在し、これは次数 $-1, 0$ で完全である。また $\mathcal{F}$ が入射的
$\mathcal{O}$-加群ならば、すべての次数で完全である。
\end{lemma}

\begin{proof}
$\Hom(\mathcal{O}, \mathcal{F}) = \Gamma(\mathcal{C}_{total}, \mathcal{F})$
および
$\Hom(g_{n!}\mathcal{O}_n, \mathcal{F}) = \Gamma(\mathcal{C}_n, \mathcal{F}_n)$
であることに注意する。したがって、この補題は補題
\ref{spaces-simplicial-lemma-simplicial-resolution-ringed} と、$\mathcal{F}$ が入射的ならば
$\Hom(-, \mathcal{F})$ が完全であるという事実から直ちに従う。
\end{proof}

\begin{lemma}
\label{spaces-simplicial-lemma-simplicial-module-cohomology-site}
状況 \ref{spaces-simplicial-situation-simplicial-site} において、$\mathcal{O}$ を環の層とする。
$K$ を $D^+(\mathcal{O})$ の対象とすると、スペクトル系列
$(E_r, d_r)_{r \geq 0}$ で
$$
E_1^{p, q} = H^q(\mathcal{C}_p, K_p),\quad
d_1^{p, q} : E_1^{p, q} \to E_1^{p + 1, q}
$$
をもち、$H^{p + q}(\mathcal{C}_{total}, K)$ に収束するものが存在する。
このスペクトル系列は $K$ に関して関手的である。
\end{lemma}

\begin{proof}
$K$ を表す入射的 $\mathcal{O}$-加群からなる下に有界な複体を
$\mathcal{I}^\bullet$ とする。項が
$$
A^{p, q} = \Gamma(\mathcal{C}_p, \mathcal{I}^q_p)
$$
である二重複体を考える。ここで水平方向の射は補題
\ref{spaces-simplicial-lemma-cech-complex-modules} から、垂直方向の射は複体
$\mathcal{I}^\bullet$ の微分から来る。$\textit{Mod}(\mathcal{O}_\mathcal{D})$
上で
$\Gamma(\mathcal{D}, -) =
\Hom_{\mathcal{O}_\mathcal{D}}(\mathcal{O}_\mathcal{D}, -)$
であることに注意する。したがって補題は、二重複体の各行が正の次数で完全で
あり、次数 $0$ では $\Gamma(\mathcal{C}_{total}, \mathcal{I}^q)$ になることを
述べている。ゆえに、この二重複体に付随する全複体は、『ホモロジー』補題
\ref{homology-lemma-double-complex-gives-resolution} により
$R\Gamma(\mathcal{C}_{total}, K)$ を計算する。一方、$\mathcal{C}_p$ への制限は
完全なので（補題 \ref{spaces-simplicial-lemma-restriction-to-components-site}）、複体
$\mathcal{I}_p^\bullet$ は $D(\mathcal{C}_p)$ において $K_p$ を表す。
層 $\mathcal{I}_p^q$ は $\mathcal{C}_p$ 上で完全非輪状である（補題
\ref{spaces-simplicial-lemma-restriction-injective-to-component-limp}）。したがって Leray の
非輪状性補題（『導来圏』補題 \ref{derived-lemma-leray-acyclicity}）および
『サイト上のコホモロジー』補題
\ref{sites-cohomology-lemma-limp-acyclic} により、各列のコホモロジーは群
$H^q(\mathcal{C}_p, K_p)$ を計算する。『ホモロジー』補題
\ref{homology-lemma-first-quadrant-ss} を適用して結論を得る。
\end{proof}

\begin{lemma}
\label{spaces-simplicial-lemma-sanity-check-modules}
状況 \ref{spaces-simplicial-situation-simplicial-site} において、$\mathcal{O}$ を環の層とする。
$U \in \Ob(\mathcal{C}_n)$ および
$\mathcal{F} \in \textit{Mod}(\mathcal{O})$ とする。このとき
$H^p(U, \mathcal{F}) = H^p(U, g_n^*\mathcal{F})$ である。ここで左辺では
$U$ を $\mathcal{C}_{total}$ の対象とみなしている。
\end{lemma}

\begin{proof}
「$U$ を $\mathcal{C}_{total}$ の対象とみなす」ことは、場合 (A) では補題
\ref{spaces-simplicial-lemma-simplicial-site-site}、場合 (B) では補題
\ref{spaces-simplicial-lemma-simplicial-cocontinuous-site} における $\mathcal{C}_{total}$ の構成により
説明される。どちらの場合にも関手 $\mathcal{C}_n \to \mathcal{C}$ は連続かつ
余連続であり（補題 \ref{spaces-simplicial-lemma-restriction-to-components-site}）、また定義により
$g_n^{-1}\mathcal{O} = \mathcal{O}_n$ である。したがって結果は
『サイト上のコホモロジー』補題
\ref{sites-cohomology-lemma-pullback-same-cohomology} の特別な場合である。
\end{proof}






\section{環付き単体的サイトのコホモロジーと増大}
\label{spaces-simplicial-section-cohomology-augmentation-ringed-simplicial-sites}

\noindent
本節は、加群層に対する節
\ref{spaces-simplicial-section-cohomology-augmentation-simplicial-sites} の類似である。

\medskip\noindent
状況 \ref{spaces-simplicial-situation-simplicial-site} のような単体的サイト
$\mathcal{C}$ を考える。$a_0$ を注意 \ref{spaces-simplicial-remark-augmentation-site} のような
サイト $\mathcal{D}$ への増大とする。$\mathcal{O}$ を
$\mathcal{C}_{total}$ 上の環の層とし、$\mathcal{O}_\mathcal{D}$ を
$\mathcal{D}$ 上の環の層とする。射
$$
a^\sharp : \mathcal{O}_\mathcal{D} \longrightarrow a_*\mathcal{O}
$$
が与えられていると仮定する。ここで $a$ は補題
\ref{spaces-simplicial-lemma-augmentation-site} のとおりである。したがって環付きトポスの射
$$
a :
(\Sh(\mathcal{C}_{total}), \mathcal{O})
\longrightarrow
(\Sh(\mathcal{D}), \mathcal{O}_\mathcal{D})
$$
を得る。$g_n : (\Sh(\mathcal{C}_n), \mathcal{O}_n) \to
(\Sh(\mathcal{C}_{total}), \mathcal{O})$ を補題
\ref{spaces-simplicial-lemma-restriction-module-to-components-site} のような環付きトポスの射と
みなす。環付きトポスの射として合成 $a_n = a \circ g_n$ をとると
（補題 \ref{spaces-simplicial-lemma-augmentation-site}）、
$$
a_n :
(\Sh(\mathcal{C}_n), \mathcal{O}_n)
\longrightarrow
(\Sh(\mathcal{D}), \mathcal{O}_\mathcal{D})
$$
を得る。遷移写像 $f_\varphi^{-1}\mathcal{O}_m \to \mathcal{O}_n$ を用いて、
環付きトポスの射
$$
f_\varphi : (\Sh(\mathcal{C}_n), \mathcal{O}_n) \to
(\Sh(\mathcal{C}_m), \mathcal{O}_m)
$$
を得る。これは、すべての $\varphi : [m] \to [n]$ に対して
環付きトポスの射として $a_m \circ f_\varphi = a_n$ を満たす。

\begin{lemma}
\label{spaces-simplicial-lemma-flat-augmentation-modules}
% P10-E0630
上の記法を用いる。射
$a : (\Sh(\mathcal{C}_{total}), \mathcal{O}) \to
(\Sh(\mathcal{D}), \mathcal{O}_\mathcal{D})$
が平坦であるための必要十分条件は、すべての $n \geq 0$ に対して
$a_n : (\Sh(\mathcal{C}_n), \mathcal{O}_n) \to
(\Sh(\mathcal{D}), \mathcal{O}_\mathcal{D})$
が平坦であることである。
\end{lemma}

\begin{proof}
$g_n : (\Sh(\mathcal{C}_n), \mathcal{O}_n) \to
(\Sh(\mathcal{C}_{total}), \mathcal{O})$ は平坦なので、$a$ が平坦ならば、
$a_n = a \circ g_n$ は合成として平坦である。逆に、すべての $n$ に対して
$a_n$ が平坦であると仮定する。$\mathcal{O}$ が
$a^{-1}\mathcal{O}_\mathcal{D}$-加群の層として平坦であることを確認すればよい。
$\mathcal{F} \to \mathcal{G}$ を $a^{-1}\mathcal{O}_\mathcal{D}$-加群の
単射とする。写像
$$
\mathcal{F} \otimes_{a^{-1}\mathcal{O}_\mathcal{D}} \mathcal{O}
\to
\mathcal{G} \otimes_{a^{-1}\mathcal{O}_\mathcal{D}} \mathcal{O}
$$
が単射であることを示さねばならない。これは $\mathcal{C}_n$ 上、すなわち
$g_n^{-1}$ を適用した後で確認できる。
$g_n^{-1}\mathcal{O} = \mathcal{O}_n$ なので $g_n^* = g_n^{-1}$ であり、
$$
g_n^{-1}\mathcal{F} \otimes_{g_n^{-1}a^{-1}\mathcal{O}_\mathcal{D}}
\mathcal{O}_n
\to
g_n^{-1}\mathcal{G} \otimes_{g_n^{-1}a^{-1}\mathcal{O}_\mathcal{D}}
\mathcal{O}_n
$$
を得る。これは
$g_n^{-1}a^{-1}\mathcal{O}_\mathcal{D} = a_n^{-1}\mathcal{O}_\mathcal{D}$
であり、$a_n$ が平坦であると仮定したので単射である。
\end{proof}

\begin{lemma}
\label{spaces-simplicial-lemma-simplicial-resolution-augmentation-modules}
% P10-E0631
上の記法を用いる。$\mathcal{O}_\mathcal{D}$-加群 $\mathcal{G}$ に対し、
$\mathcal{C}_{total}$ 上の $\mathcal{O}$-加群の層の完全複体
$$
\ldots \to
g_{2!}(a_2^*\mathcal{G}) \to
g_{1!}(a_1^*\mathcal{G}) \to
g_{0!}(a_0^*\mathcal{G}) \to
a^*\mathcal{G} \to 0
$$
が存在する。ここで $g_{n!}$ は補題
\ref{spaces-simplicial-lemma-restriction-module-to-components-site} のとおりである。
\end{lemma}

\begin{proof}
$a^*\mathcal{G}$ は、その $\mathcal{C}_m$ への制限が
$\mathcal{O}_m$-加群 $a_m^*\mathcal{G}$ であるような
$\mathcal{C}_{total}$ 上の $\mathcal{O}$-加群であることに注意する。
補題 \ref{spaces-simplicial-lemma-restriction-module-to-components-site} における加群上の
関手 $g_{n!}$ の記述から、$g_{n!}(a_n^*\mathcal{G})$ は、その
$\mathcal{C}_m$ への制限が $\mathcal{O}_m$-加群
$$
\bigoplus\nolimits_{\varphi : [n] \to [m]} f_\varphi^*a_n^*\mathcal{G} =
\bigoplus\nolimits_{\varphi : [n] \to [m]} a_m^*\mathcal{G}
$$
であるような $\mathcal{C}_{total}$ 上の $\mathcal{O}$-加群であることが分かる。
証明の残りは補題 \ref{spaces-simplicial-lemma-simplicial-resolution-augmentation} の証明と
まったく同じであり、$a_m^{-1}\mathcal{G}$ を $a_m^*\mathcal{G}$ で置き換える。
\end{proof}

\begin{lemma}
\label{spaces-simplicial-lemma-augmentation-cech-complex-modules}
% P10-E0632
上の記法を用いる。$\mathcal{C}_{total}$ 上の $\mathcal{O}$-加群
$\mathcal{F}$ に対し、$\mathcal{O}_\mathcal{D}$-加群の標準的な複体
$$
0 \to a_*\mathcal{F} \to a_{0, *}\mathcal{F}_0 \to a_{1, *}\mathcal{F}_1 \to
a_{2, *}\mathcal{F}_2 \to \ldots
$$
で、次数 $-1, 0$ で完全なものが存在する。$\mathcal{F}$ が入射的
$\mathcal{O}$-加群ならば、この複体はすべての次数で完全であり、任意の
$\mathcal{O}_\mathcal{D}$-加群 $\mathcal{G}$ に対して関手
$\Hom_{\mathcal{O}_\mathcal{D}}(\mathcal{G}, -)$ を適用しても完全なままである。
\end{lemma}

\begin{proof}
% P10-E0633
複体を構成するには、米田の補題により、$\mathcal{D}$ 上の任意の
$\mathcal{O}_\mathcal{D}$-加群 $\mathcal{G}$ に対して、複体
$$
0 \to \Hom_{\mathcal{O}_\mathcal{D}}(\mathcal{G}, a_*\mathcal{F}) \to
\Hom_{\mathcal{O}_\mathcal{D}}(\mathcal{G}, a_{0, *}\mathcal{F}_0) \to
\Hom_{\mathcal{O}_\mathcal{D}}(\mathcal{G}, a_{1, *}\mathcal{F}_1) \to \ldots
$$
を $\mathcal{G}$ に関して関手的に構成すれば十分である。そのためには補題
\ref{spaces-simplicial-lemma-simplicial-resolution-augmentation-modules} の完全複体に
$\Hom_\mathcal{O}(-, \mathcal{F})$ を適用し、引戻しと順像の随伴性を用いる。
次数 $-1, 0$ における完全性は、$\Hom_\mathcal{O}(-, \mathcal{F})$ が左完全なので
この構成から従う。$\mathcal{F}$ が入射的 $\mathcal{O}$-加群ならば、
$\Hom_\mathcal{O}(-, \mathcal{F})$ が完全なので複体は完全である。
\end{proof}

\begin{lemma}
\label{spaces-simplicial-lemma-augmentation-spectral-sequence-modules}
% P10-E0634
上の記法を用いる。任意の $K$ を $D^+(\mathcal{O})$ の対象とすると、
$\textit{Mod}(\mathcal{O}_\mathcal{D})$ におけるスペクトル系列
$(E_r, d_r)_{r \geq 0}$ で
$$
E_1^{p, q} = R^qa_{p, *} K_p
$$
をもち、$R^{p + q}a_*K$ に収束するものが存在する。このスペクトル系列は
$K$ に関して関手的である。
\end{lemma}

\begin{proof}
$K$ を表す入射的 $\mathcal{O}$-加群からなる下に有界な複体を
$\mathcal{I}^\bullet$ とする。項が
$$
A^{p, q} = a_{p, *}\mathcal{I}^q_p
$$
である二重複体を考える。ここで水平方向の射は補題
\ref{spaces-simplicial-lemma-augmentation-cech-complex-modules} から、垂直方向の射は複体
$\mathcal{I}^\bullet$ の微分から来る。この補題は、二重複体の各行が正の次数で
完全であり、次数 $0$ では $a_*\mathcal{I}^q$ になることを述べている。
ゆえに、この二重複体に付随する全複体は、『ホモロジー』補題
\ref{homology-lemma-double-complex-gives-resolution} により $Ra_*K$ を計算する。
一方、$\mathcal{C}_p$ への制限は完全なので（補題
\ref{spaces-simplicial-lemma-restriction-to-components-site}）、複体 $\mathcal{I}_p^\bullet$ は
$D(\mathcal{C}_p)$ において $K_p$ を表す。層 $\mathcal{I}_p^q$ は
$\mathcal{C}_p$ 上で完全非輪状である（補題
\ref{spaces-simplicial-lemma-restriction-injective-to-component-limp}）。したがって Leray の
非輪状性補題（『導来圏』補題 \ref{derived-lemma-leray-acyclicity}）および
『サイト上のコホモロジー』補題
\ref{sites-cohomology-lemma-limp-acyclic} により、各列のコホモロジーは層
$R^qa_{p, *}K_p$ である。『ホモロジー』補題
\ref{homology-lemma-first-quadrant-ss} を適用して結論を得る。
\end{proof}
\section{カルテジアン層と加群}
\label{spaces-simplicial-section-cartesian}

\noindent
定義は次のとおりである。

\begin{definition}
\label{spaces-simplicial-definition-cartesian-sheaf}
状況 \ref{spaces-simplicial-situation-simplicial-site} において、次を定義する。
\begin{enumerate}
\item $\mathcal{C}_{total}$ 上の集合またはアーベル群の層
$\mathcal{F}$ が {\it カルテジアン}であるとは、すべての
$\varphi : [m] \to [n]$ に対して写像
$\mathcal{F}(\varphi) : f_\varphi^{-1}\mathcal{F}_m \to \mathcal{F}_n$
が同型であることをいう。
\item $\mathcal{O}$ が $\mathcal{C}_{total}$ 上の環の層であるとき、
$\mathcal{O}$-加群の層 $\mathcal{F}$ が {\it カルテジアン}であるとは、
すべての $\varphi : [m] \to [n]$ に対して写像
$f_\varphi^*\mathcal{F}_m \to \mathcal{F}_n$ が同型であることをいう。
\item $D(\mathcal{C}_{total})$ の対象 $K$ が {\it カルテジアン}であるとは、
すべての $\varphi : [m] \to [n]$ に対して写像
$f_\varphi^{-1}K_m \to K_n$ が同型であることをいう。
\item $\mathcal{O}$ が $\mathcal{C}_{total}$ 上の環の層であるとき、
$D(\mathcal{O})$ の対象 $K$ が {\it カルテジアン}であるとは、すべての
$\varphi : [m] \to [n]$ に対して写像
$Lf_\varphi^*K_m \to K_n$ が同型であることをいう。
\end{enumerate}
\end{definition}

\noindent
もちろん、ファイバー圏のカルテジアン切断という一般概念があり、上記は
その単なる例である。
% P10-E0181
引戻しに関するこの性質は、退化写像についてだけ確認すればよい。

\begin{lemma}
\label{spaces-simplicial-lemma-check-cartesian-module}
状況 \ref{spaces-simplicial-situation-simplicial-site} において、次が成り立つ。
\begin{enumerate}
\item 集合またはアーベル群の層 $\mathcal{F}$ がカルテジアンであるための
必要十分条件は、写像
$(f_{\delta^n_j})^{-1}\mathcal{F}_{n - 1} \to \mathcal{F}_n$
が同型であることである。
\item $D(\mathcal{C}_{total})$ の対象 $K$ がカルテジアンであるための
必要十分条件は、写像
$(f_{\delta^n_j})^{-1}K_{n - 1} \to K_n$
が同型であることである。
\item $\mathcal{O}$ が $\mathcal{C}_{total}$ 上の環の層であるとき、
$\mathcal{O}$-加群の層 $\mathcal{F}$ がカルテジアンであるための必要十分条件は、
写像 $(f_{\delta^n_j})^*\mathcal{F}_{n - 1} \to \mathcal{F}_n$
が同型であることである。
\item $\mathcal{O}$ が $\mathcal{C}_{total}$ 上の環の層であるとき、
$D(\mathcal{O})$ の対象 $K$ がカルテジアンであるための必要十分条件は、写像
$L(f_{\delta^n_j})^*K_{n - 1} \to K_n$
が同型であることである。
% P10-E0182
\item ここにさらに追加する。
\end{enumerate}
\end{lemma}

\begin{proof}
いずれの場合にも鍵となるのは、引戻し関手の合成が引戻し関手になることである。
(4) については『サイト上のコホモロジー』補題
\ref{sites-cohomology-lemma-derived-pullback-composition} を参照されたい。
場合 (1) における議論を示し、他の場合の証明は省略する。
% P10-E0183
圏 $\Delta$ は射 $\delta^n_j$ と $\sigma^n_j$ により生成される。
『単体』補題 \ref{simplicial-lemma-face-degeneracy} を参照されたい。
したがって、写像
$(f_{\delta^n_j})^{-1}\mathcal{F}_{n - 1} \to \mathcal{F}_n$
および $(f_{\sigma^n_j})^{-1}\mathcal{F}_{n + 1} \to \mathcal{F}_n$ が
同型であることだけを確認すればよい。記法については『単体』補題
\ref{simplicial-lemma-characterize-simplicial-object} を参照されたい。
$\sigma^n_j \circ \delta_j^{n + 1} = \text{id}_{[n]}$ なので、合成
$$
\mathcal{F}_n =
(f_{\sigma^n_j})^{-1}
(f_{\delta_j^{n + 1}})^{-1}
\mathcal{F}_n \to
(f_{\sigma^n_j})^{-1}
\mathcal{F}_{n + 1} \to
\mathcal{F}_n
$$
は恒等写像である。したがって $\delta^{n + 1}_j$ に対する結果から
$\sigma^n_j$ に対する結果が従う。
\end{proof}

\begin{lemma}
\label{spaces-simplicial-lemma-augmentation-cartesian-module}
状況 \ref{spaces-simplicial-situation-simplicial-site} において、$a_0$ を注意
\ref{spaces-simplicial-remark-augmentation-site} のようなサイト $\mathcal{D}$ への増大とする。
\begin{enumerate}
\item $\mathcal{D}$ 上の集合またはアーベル群の層の引戻し
$a^{-1}\mathcal{G}$ はカルテジアンである。
\item $D(\mathcal{D})$ の対象 $K$ の引戻し
$a^{-1}K$ はカルテジアンである。
\end{enumerate}
$\mathcal{O}$ を $\mathcal{C}_{total}$ 上の環の層、
$\mathcal{O}_\mathcal{D}$ を $\mathcal{D}$ 上の環の層とし、
$a^\sharp : \mathcal{O}_\mathcal{D} \to a_*\mathcal{O}$ を節
\ref{spaces-simplicial-section-cohomology-augmentation-ringed-simplicial-sites} のような射とする。
\begin{enumerate}
\item[(3)] $\mathcal{O}_\mathcal{D}$-加群の層の引戻し
$a^*\mathcal{F}$ はカルテジアンである。
\item[(4)] $D(\mathcal{O}_\mathcal{D})$ の対象 $K$ の導来引戻し
$La^*K$ はカルテジアンである。
\end{enumerate}
\end{lemma}

\begin{proof}
これは、すべての $\varphi : [m] \to [n]$ に対する恒等式
$a_m \circ f_\varphi = a_n$ から直ちに従う。補題
\ref{spaces-simplicial-lemma-augmentation-site} および節
\ref{spaces-simplicial-section-cohomology-augmentation-ringed-simplicial-sites} の議論を
参照されたい。
\end{proof}

\begin{lemma}
\label{spaces-simplicial-lemma-characterize-cartesian}
状況 \ref{spaces-simplicial-situation-simplicial-site} において、カルテジアンな集合の層
（それぞれ、アーベル群の層）の圏は、対 $(\mathcal{F}, \alpha)$ の圏と同値で
ある。ここで $\mathcal{F}$ は $\mathcal{C}_0$ 上の集合の層
（それぞれ、アーベル群の層）であり、
$$
\alpha :
(f_{\delta_1^1})^{-1}\mathcal{F}
\longrightarrow (f_{\delta_0^1})^{-1}\mathcal{F}
$$
は $\mathcal{C}_1$ 上の集合の層（それぞれ、アーベル群の層）の同型であって、
$\mathcal{C}_2$ 上の層の写像として
$(f_{\delta^2_1})^{-1}\alpha =
(f_{\delta^2_0})^{-1}\alpha \circ (f_{\delta^2_2})^{-1}\alpha$
を満たす。
\end{lemma}

\begin{proof}
$d^n_j = f_{\delta^n_j} : \Sh(\mathcal{C}_n) \to
\Sh(\mathcal{C}_{n - 1})$ と略記する。補題の主張における $\alpha$ の条件が
意味をもつのは、トポスの射 $\Sh(\mathcal{C}_2) \to \Sh(\mathcal{C}_0)$
として
$$
d^1_1 \circ d^2_2 = d^1_1 \circ d^2_1, \quad
d^1_1 \circ d^2_0 = d^1_0 \circ d^2_2, \quad
d^1_0 \circ d^2_0 = d^1_0 \circ d^2_1
$$
が成り立つからである。『単体』注意
\ref{simplicial-remark-relations} を参照されたい。したがって、これらの写像を
次のように図示できる。
$$
\xymatrix@C=1em{
& (d^2_0)^{-1}(d^1_1)^{-1}\mathcal{F} \ar[r]_-{(d^2_0)^{-1}\alpha} &
(d^2_0)^{-1}(d^1_0)^{-1}\mathcal{F} \ar@{=}[rd] & \\
(d^2_2)^{-1}(d^1_0)^{-1}\mathcal{F} \ar@{=}[ru] & & &
(d^2_1)^{-1}(d^1_0)^{-1}\mathcal{F} \\
& (d^2_2)^{-1}(d^1_1)^{-1}\mathcal{F} \ar[lu]^{(d^2_2)^{-1}\alpha} \ar@{=}[r] &
(d^2_1)^{-1}(d^1_1)^{-1}\mathcal{F} \ar[ru]_{(d^2_1)^{-1}\alpha}
}
$$
そして条件はこの図式が可換であることを意味する。
$\mathcal{C}_{total}$ 上のカルテジアンな集合の層
$\mathcal{G}$（それぞれ、アーベル群の層）が与えられたとき、明らかに
$\mathcal{F} = \mathcal{G}_0$ とおき、$\alpha$ を合成
% P10-E0184
$$
(d_1^1)^{-1}\mathcal{G}_0 \to \mathcal{G}_1
\leftarrow (d_1^0)^{-1}\mathcal{G}_0
$$
に等しくとることができる。ここで $\mathcal{G}$ はカルテジアンなので、射は
可逆である。この関手が同値であることを示すため、擬逆を構成する。
擬逆の構成は『降下』節 \ref{descent-section-descent-modules} で論じられた
構成と類似しており、そこから記法
$\tau^n_i : [0] \to [n]$、$0 \mapsto i$ および
$\tau^n_{ij} : [1] \to [n]$、$0 \mapsto i$、$1 \mapsto j$ を借用する。
すなわち、補題のような対 $(\mathcal{F}, \alpha)$ が与えられたとき、
$\mathcal{G}_n = (f_{\tau^n_n})^{-1}\mathcal{F}$ とおく。
$\varphi : [n] \to [m]$ が与えられたとき、次を用いて
$\mathcal{G}(\varphi) : (f_\varphi)^{-1}\mathcal{G}_n \to \mathcal{G}_m$
を定義する。
$$
\xymatrix@C=1em{
(f_\varphi)^{-1}\mathcal{G}_n \ar@{=}[r] &
(f_\varphi)^{-1}(f_{\tau^n_n})^{-1}\mathcal{F} \ar@{=}[r] &
(f_{\tau^m_{\varphi(n)}})^{-1}\mathcal{F} \ar@{=}[r] &
(f_{\tau^m_{\varphi(n)m}})^{-1}(d^1_1)^{-1}\mathcal{F}
\ar[d]^{(f_{\tau^m_{\varphi(n)m}})^{-1}\alpha} \\
&
\mathcal{G}_m \ar@{=}[r] &
(f_{\tau^m_m})^{-1}\mathcal{F} \ar@{=}[r] &
(f_{\tau^m_{\varphi(n)m}})^{-1}(d^1_0)^{-1}\mathcal{F}
}
$$
上に表示した図式の可換性から写像が正しく合成され、したがって
$\mathcal{C}_{total}$ 上の層を定めることの確認は省略する。補題
\ref{spaces-simplicial-lemma-describe-sheaves-simplicial-site-site} を参照されたい。また、二つの
関手が互いに擬逆であることの確認も省略する。
\end{proof}

\begin{lemma}
\label{spaces-simplicial-lemma-characterize-cartesian-modules}
状況 \ref{spaces-simplicial-situation-simplicial-site} において、$\mathcal{O}$ を
$\mathcal{C}_{total}$ 上の環の層とする。カルテジアン
$\mathcal{O}$-加群の圏は、対 $(\mathcal{F}, \alpha)$ の圏と同値である。
ここで $\mathcal{F}$ は $\mathcal{O}_0$-加群であり、
$$
\alpha :
(f_{\delta_1^1})^*\mathcal{F}
\longrightarrow (f_{\delta_0^1})^*\mathcal{F}
$$
は $\mathcal{O}_1$-加群の同型であって、$\mathcal{O}_2$-加群の写像として
$(f_{\delta^2_1})^*\alpha =
(f_{\delta^2_0})^*\alpha \circ (f_{\delta^2_2})^*\alpha$
を満たす。
\end{lemma}

\begin{proof}
証明は補題 \ref{spaces-simplicial-lemma-characterize-cartesian} の証明において、アーベル群の層の
引戻しを加群の引戻しで置き換えたものと同一である。
\end{proof}

\begin{lemma}
\label{spaces-simplicial-lemma-Serre-subcat-cartesian-modules}
状況 \ref{spaces-simplicial-situation-simplicial-site} において、次が成り立つ。
\begin{enumerate}
\item カルテジアンなアーベル層の充満部分圏は
$\textit{Ab}(\mathcal{C}_{total})$ の弱 Serre 部分圏をなす。
カルテジアンなアーベル層の系の余極限はカルテジアンである。
\item $\mathcal{O}$ を $\mathcal{C}_{total}$ 上の環の層で、射
$$
f_{\delta^n_j} : (\Sh(\mathcal{C}_n), \mathcal{O}_n)
\to (\Sh(\mathcal{C}_{n - 1}), \mathcal{O}_{n - 1})
$$
が平坦であるものとする。カルテジアン $\mathcal{O}$-加群の充満部分圏は
$\textit{Mod}(\mathcal{O})$ の弱 Serre 部分圏をなす。
カルテジアン $\mathcal{O}$-加群の系の余極限はカルテジアンである。
\end{enumerate}
\end{lemma}

\begin{proof}
(1) で弱 Serre 部分圏が得られることを見るため、『ホモロジー』補題
\ref{homology-lemma-characterize-weak-serre-subcategory} に挙げられた条件を
確認する。まず、$\varphi : \mathcal{F} \to \mathcal{G}$ がカルテジアンな
アーベル層の間の写像ならば、制限関手
$\Sh(\mathcal{C}_{total}) \to \Sh(\mathcal{C}_n)$ と関手
$f_\varphi^{-1}$ が完全なので、$\Ker(\varphi)$ と $\Coker(\varphi)$ も
カルテジアンである。同様に、
$$
0 \to \mathcal{F} \to \mathcal{H} \to \mathcal{G} \to 0
$$
が $\mathcal{C}_{total}$ 上のアーベル層の短完全列で、
$\mathcal{F}$ と $\mathcal{G}$ がカルテジアンならば、五項補題から
$\mathcal{H}$ もカルテジアンである。余極限に関する性質については、引戻しが
左随伴なので余極限と可換することを用いる。加群の場合にも同様に論じるが、
平坦引戻しの完全性（『サイト上の加群』補題
\ref{sites-modules-lemma-flat-pullback-exact}）と、$f_{\delta^n_j}$ に対する条件を
確認すれば十分であるという事実（補題 \ref{spaces-simplicial-lemma-check-cartesian-module}）を用いる。
\end{proof}

\begin{remark}[警告]
\label{spaces-simplicial-remark-warning-cartesian-modules}
補題 \ref{spaces-simplicial-lemma-Serre-subcat-cartesian-modules} にもかかわらず、カルテジアン
$\mathcal{O}$-加群の圏がアーベル圏でありながら
$\textit{Mod}(\mathcal{O})$ の Serre 部分圏でないことが起こりうる。
実際、$f_{\delta_1^1}$ と $f_{\delta_0^1}$ が平坦であることだけが分かっていると
仮定しよう。すると補題 \ref{spaces-simplicial-lemma-characterize-cartesian-modules} から、
カルテジアン $\mathcal{O}$-加群の圏がアーベル圏であることは容易に従う。
しかし、たとえば $f_{\delta_0^2}$ が平坦でなければ、カルテジアン
$\mathcal{O}$-加群の圏からすべての $\mathcal{O}$-加群への包含関手が
完全である理由はない。
\end{remark}

\begin{lemma}
\label{spaces-simplicial-lemma-derived-cartesian-modules}
状況 \ref{spaces-simplicial-situation-simplicial-site} において、次が成り立つ。
\begin{enumerate}
\item $D(\mathcal{C}_{total})$ の対象 $K$ がカルテジアンであるための
必要十分条件は、すべての $q$ に対して $H^q(K)$ がカルテジアンな
アーベル層であることである。
\item $\mathcal{O}$ を $\mathcal{C}_{total}$ 上の環の層で、射
$f_{\delta^n_j} : (\Sh(\mathcal{C}_n), \mathcal{O}_n)
\to (\Sh(\mathcal{C}_{n - 1}), \mathcal{O}_{n - 1})$ が平坦であるものとする。
このとき $D(\mathcal{O})$ の対象 $K$ がカルテジアンであるための必要十分条件は、
すべての $q$ に対して $H^q(K)$ がカルテジアン $\mathcal{O}$-加群で
あることである。
\end{enumerate}
\end{lemma}

\begin{proof}
(1) は、引戻し関手 $(f_\varphi)^{-1}$ が完全なので成り立つ。(2) は補題
\ref{spaces-simplicial-lemma-check-cartesian-module} の特徴付けと、平坦性により
$L(f_{\delta^n_j})^* = (f_{\delta^n_j})^*$ であることから従う。
\end{proof}

\begin{lemma}
\label{spaces-simplicial-lemma-derived-cartesian-shriek}
状況 \ref{spaces-simplicial-situation-simplicial-site} において、次が成り立つ。
\begin{enumerate}
% P10-E0635
\item $D(\mathcal{C}_{total})$ の対象 $K$ がカルテジアンであるための
必要十分条件は、すべての $n$ に対して標準写像
$$
g_{n!}K_n \longrightarrow
g_{n!}\mathbf{Z} \otimes^\mathbf{L}_\mathbf{Z} K
$$
が同型であることである。
\item $\mathcal{O}$ を $\mathcal{C}_{total}$ 上の環の層で、すべての
$\varphi : [n] \to [m]$ に対して写像
$f_\varphi^{-1}\mathcal{O}_n \to \mathcal{O}_m$ が平坦であるものとする。
このとき $D(\mathcal{O})$ の対象 $K$ がカルテジアンであるための
必要十分条件は、すべての $n$ に対して標準写像
$$
g_{n!}K_n \longrightarrow
g_{n!}\mathcal{O}_n \otimes^\mathbf{L}_\mathcal{O} K
$$
が同型であることである。
\end{enumerate}
\end{lemma}

\begin{proof}
(1) の証明。$g_{n!}$ は完全なので、導来圏上の関手で $g_n^{-1}$ の
左随伴であるものを誘導する。この写像は、写像
$\mathbf{Z} \to g_n^{-1}g_{n!}\mathbf{Z}$ に対応する写像
$K_n \to (g_n^{-1}g_{n!}\mathbf{Z}) \otimes^\mathbf{L}_\mathbf{Z} K_n$
の随伴である。前者はさらに
$\text{id} : g_{n!}\mathbf{Z} \to g_{n!}\mathbf{Z}$ に随伴する。
補題 \ref{spaces-simplicial-lemma-restriction-to-components-site} で与えられた $g_{n!}$ の記述を
用いると、この写像の $\mathcal{C}_m$ への制限は
$$
\bigoplus\nolimits_{\varphi : [n] \to [m]} f_\varphi^{-1}K_n
\longrightarrow
\bigoplus\nolimits_{\varphi : [n] \to [m]} K_m
$$
であることが分かる。したがって主張は明らかである。

\medskip\noindent
(2) の証明。$g_{n!}$ は完全なので（補題
\ref{spaces-simplicial-lemma-exactness-g-shriek-modules}）、導来圏上の関手で $g_n^*$
（これも完全）の左随伴であるものを誘導する。この写像は、写像
$\mathcal{O}_n \to g_n^*g_{n!}\mathcal{O}_n$ に対応する写像
$K_n \to (g_n^*g_{n!}\mathcal{O}_n) \otimes^\mathbf{L}_{\mathcal{O}_n} K_n$
の随伴である。前者はさらに
$\text{id} : g_{n!}\mathcal{O}_n \to g_{n!}\mathcal{O}_n$ に随伴する。
補題 \ref{spaces-simplicial-lemma-restriction-module-to-components-site} で与えられた
$g_{n!}$ の記述を用いると、この写像の $\mathcal{C}_m$ への制限は
$$
\bigoplus\nolimits_{\varphi : [n] \to [m]} f_\varphi^*K_n
\longrightarrow
\bigoplus\nolimits_{\varphi : [n] \to [m]}
f_\varphi^*\mathcal{O}_n \otimes_{\mathcal{O}_m} K_m =
\bigoplus\nolimits_{\varphi : [n] \to [m]} K_m
$$
であることが分かる。したがって主張は明らかである。
\end{proof}

\begin{lemma}
\label{spaces-simplicial-lemma-quasi-coherent-sheaf}
状況 \ref{spaces-simplicial-situation-simplicial-site} において、$\mathcal{O}$ を
$\mathcal{C}_{total}$ 上の環の層とし、$\mathcal{F}$ を
$\mathcal{O}$-加群の層とする。このとき $\mathcal{F}$ が
『サイト上の加群』定義 \ref{sites-modules-definition-site-local} の意味で
準連接であるための必要十分条件は、$\mathcal{F}$ がカルテジアンであり、
すべての $n$ に対して $\mathcal{F}_n$ が準連接
$\mathcal{O}_n$-加群であることである。
\end{lemma}

\begin{proof}
$\mathcal{F}$ が準連接であると仮定する。準連接加群の引戻しは準連接なので
（『サイト上の加群』補題 \ref{sites-modules-lemma-local-pullback}）、すべての
$n$ に対して $\mathcal{F}_n$ は準連接 $\mathcal{O}_n$-加群である。
$\mathcal{F}$ がカルテジアンであることを示すため、ある $n$ に対して
$\mathcal{C}_n$ の対象 $U$ をとる。$U$ を $\mathcal{C}_{total}$ の対象と
みなそう。$\mathcal{F}$ は準連接なので、被覆 $\{U_i \to U\}$ と、各
$i$ に対する表示
$$
\bigoplus\nolimits_{j \in J_i} \mathcal{O}_{\mathcal{C}_{total}/U_i} \to
\bigoplus\nolimits_{k \in K_i} \mathcal{O}_{\mathcal{C}_{total}/U_i} \to
\mathcal{F}|_{\mathcal{C}_{total}/U_i} \to 0
$$
が存在する。サイト $\mathcal{C}_{total}$ の構成により、
$\{U_i \to U\}$ は $\mathcal{C}_n$ の被覆であることに注意する。
次に、ある $m$ に対して $\mathcal{C}_m$ の対象 $V$ をとり、
$\varphi : [n] \to [m]$ 上にある $\mathcal{C}_{total}$ の射
$V \to U$ をとる。ファイバー積 $V_i = V \times_U U_i$ が存在し、
$\mathcal{C}_m$ における誘導被覆 $\{V_i \to V\}$ を得る。
上の表示をサイト $\mathcal{C}_n/U_i$ と $\mathcal{C}_m/V_i$ に制限すると、
表示
% P10-E0185
$$
\bigoplus\nolimits_{j \in J_i} \mathcal{O}_{\mathcal{C}_m/U_i} \to
\bigoplus\nolimits_{k \in K_i} \mathcal{O}_{\mathcal{C}_m/U_i} \to
\mathcal{F}_n|_{\mathcal{C}_n/U_i} \to 0
$$
および
$$
\bigoplus\nolimits_{j \in J_i} \mathcal{O}_{\mathcal{C}_m/V_i} \to
\bigoplus\nolimits_{k \in K_i} \mathcal{O}_{\mathcal{C}_m/V_i} \to
\mathcal{F}_m|_{\mathcal{C}_m/V_i} \to 0
$$
を得る。これらの表示は写像
$\mathcal{F}(\varphi) : f_\varphi^*\mathcal{F}_n \to \mathcal{F}_m$
と両立する（この写像は、$\varphi$ 上にある $\mathcal{C}_{total}$ の射に沿う
$\mathcal{F}$ の制限写像を用いて定義される）。したがって
$\mathcal{F}(\varphi)|_{\mathcal{C}_m/V_i}$ は同型である。
$\{V_i \to V\}$ は被覆なので、
$\mathcal{F}(\varphi)|_{\mathcal{C}_m/V}$ は同型である。
$V$ と $U$ は任意であったから、$\mathcal{F}$ はカルテジアンである。
（場合 A では『サイト』補題
\ref{sites-lemma-morphism-of-sites-covering} を用いる。）

\medskip\noindent
逆に、すべての $n$ に対して $\mathcal{F}_n$ が準連接であり、
$\mathcal{F}$ がカルテジアンであると仮定する。すると任意の $n$ と
$\mathcal{C}_n$ の対象 $U$ に対して、$\mathcal{C}_n$ の被覆
$\{U_i \to U\}$ と、各 $i$ に対する表示
% P10-E0186
$$
\bigoplus\nolimits_{j \in J_i} \mathcal{O}_{\mathcal{C}_m/U_i} \to
\bigoplus\nolimits_{k \in K_i} \mathcal{O}_{\mathcal{C}_m/U_i} \to
\mathcal{F}_n|_{\mathcal{C}_n/U_i} \to 0
$$
を選べる。$\mathcal{C}_{total}/U_i$ へ引き戻すと、
$\mathcal{C}_{total}/U_i$ 上の加群の複体
$$
\bigoplus\nolimits_{j \in J_i} \mathcal{O}_{\mathcal{C}_{total}/U_i} \to
\bigoplus\nolimits_{k \in K_i} \mathcal{O}_{\mathcal{C}_{total}/U_i} \to
\mathcal{F}|_{\mathcal{C}_{total}/U_i} \to 0
$$
を得る。このとき $\mathcal{F}$ がカルテジアンであるという性質から、この複体は
完全である。詳細は省略する。
\end{proof}






\section{導来圏の単体的系}
\label{spaces-simplicial-section-glueing}

\noindent
本節では、アーベル層の導来圏という設定で
\cite[Proposition 3.2.9]{BBD} の特別な場合を証明する。加群の場合は節
\ref{spaces-simplicial-section-glueing-modules} で論じる。

\begin{definition}
\label{spaces-simplicial-definition-cartesian-derived}
状況 \ref{spaces-simplicial-situation-simplicial-site} において、
{\it 導来圏の単体的系}とは、次のデータ
\begin{enumerate}
\item すべての $n$ に対する $D(\mathcal{C}_n)$ の対象 $K_n$、
\item すべての $\varphi : [m] \to [n]$ に対する
$D(\mathcal{C}_n)$ における写像
$K_\varphi : f_\varphi^{-1}K_m \to K_n$
\end{enumerate}
で、$\Delta$ の任意の射 $\varphi : [m] \to [n]$ と
$\psi : [l] \to [m]$ に対して条件
$$
K_{\varphi \circ \psi} = K_\varphi \circ f_\varphi^{-1}K_\psi :
f_{\varphi \circ \psi}^{-1}K_l = f_\varphi^{-1} f_\psi^{-1}K_l
\longrightarrow
K_n
$$
を満たすものからなる。すべての $\varphi$ に対して写像 $K_\varphi$ が
同型であるとき、この単体的系を {\it カルテジアン}という。
導来圏の二つの単体的系が与えられたとき、
{\it 導来圏の単体的系の射}には明らかな意味がある。
\end{definition}

\noindent
この概念には意図的に途方もなく長い名前を付けた。目的は、ある仮定のもとで、
導来圏の単体的系が $D(\mathcal{C}_{total})$ の対象から生じることを示すことである。

\begin{lemma}
\label{spaces-simplicial-lemma-cartesian-objects-derived}
状況 \ref{spaces-simplicial-situation-simplicial-site} において、
$K \in D(\mathcal{C}_{total})$ が対象ならば、$(K_n, K(\varphi))$ は
導来圏の単体的系である。$K$ がカルテジアンならば、この系も
カルテジアンである。
\end{lemma}

\begin{proof}
明らかである。
\end{proof}

\begin{lemma}
\label{spaces-simplicial-lemma-construct-cartesian-objects-derived}
状況 \ref{spaces-simplicial-situation-simplicial-site} において、
$K_0 \in D(\mathcal{C}_0)$ と、余サイクル条件を満たす同型
$$
\alpha :
f_{\delta_1^1}^{-1}K_0
\longrightarrow
f_{\delta_0^1}^{-1}K_0
$$
が与えられていると仮定する。
$\tau^n_i : [0] \to [n]$、$0 \mapsto i$ とおき、
$K_n = f_{\tau^n_n}^{-1}K_0$ とおく。このとき $K_n$ は導来圏の
カルテジアン単体的系をなす。
\end{lemma}

\begin{proof}
補題 \ref{spaces-simplicial-lemma-characterize-cartesian} とその証明を比較されたい
（余サイクル条件の詳しい形もそこに書かれている）。構成は『降下』節
\ref{descent-section-descent-modules} で論じられた構成と類似しており、
そこから記法 $\tau^n_i : [0] \to [n]$、$0 \mapsto i$ および
$\tau^n_{ij} : [1] \to [n]$、$0 \mapsto i$、$1 \mapsto j$ を借用する。
$\varphi : [n] \to [m]$ が与えられたとき、次を用いて
$K_\varphi : f_\varphi^{-1}K_n \to K_m$ を定義する。
$$
\xymatrix{
f_\varphi^{-1}K_n \ar@{=}[r] &
f_\varphi^{-1} f_{\tau^n_n}^{-1}K_0 \ar@{=}[r] &
f_{\tau^m_{\varphi(n)}}^{-1}K_0 \ar@{=}[r] &
f_{\tau^m_{\varphi(n)m}}^{-1}f_{\delta^1_1}^{-1}K_0
\ar[d]_{f_{\tau^m_{\varphi(n)m}}^{-1}\alpha} \\
&
K_m \ar@{=}[r] &
f_{\tau^m_m}^{-1}K_0 \ar@{=}[r] &
f_{\tau^m_{\varphi(n)m}}^{-1}f_{\delta^1_0}^{-1}K_0
}
$$
余サイクル条件から写像が（それぞれの導来圏で）正しく合成され、したがって
導来圏の単体的系を与えることの確認は省略する。
\end{proof}

\begin{lemma}
\label{spaces-simplicial-lemma-abelian-postnikov}
状況 \ref{spaces-simplicial-situation-simplicial-site} において、$K$ を
$D(\mathcal{C}_{total})$ の対象とする。$D(\mathcal{C}_{total})$ の対象として
$$
X_n = (g_{n!}\mathbf{Z})
\otimes^\mathbf{L}_\mathbf{Z} K
\quad\text{および}\quad
Y_n =
(g_{n!}\mathbf{Z} \to \ldots \to g_{0!}\mathbf{Z})[-n]
\otimes^\mathbf{L}_\mathbf{Z} K
$$
とおく。ここで写像は補題 \ref{spaces-simplicial-lemma-simplicial-resolution-Z-site} のとおりである。
明らかな標準写像 $Y_n \to X_n$ および
$Y_0 \to Y_1[1] \to Y_2[2] \to \ldots$ とともに、次が成り立つ。
\begin{enumerate}
\item 完全三角形 $Y_n \to X_n \to Y_{n - 1} \to Y_n[1]$ は、
$\ldots \to X_2 \to X_1 \to X_0$ に対する Postnikov 系
（『導来圏』定義 \ref{derived-definition-postnikov-system}）を定める。
\item $D(\mathcal{C}_{total})$ において
$K = \text{hocolim} Y_n[n]$ である。
\end{enumerate}
\end{lemma}

\begin{proof}
まず $K = \mathbf{Z}$ ならば、これは $\textit{Ab}(\mathcal{C}_{total})$ の複体
$$
\ldots \to
g_{2!}\mathbf{Z} \to
g_{1!}\mathbf{Z} \to
g_{0!}\mathbf{Z}
$$
に適用した『導来圏』例 \ref{derived-example-key-postnikov} の構成と、補題
\ref{spaces-simplicial-lemma-simplicial-resolution-Z-site} によりこの複体が
$D(\mathcal{C}_{total})$ において $K = \mathbf{Z}$ を表すという事実を
組み合わせたものである。一般の場合はこれと、完全関手
$- \otimes^\mathbf{L}_\mathbf{Z} K$ が Postnikov 系を Postnikov 系へ送り、
$- \otimes^\mathbf{L}_\mathbf{Z} K$ がホモトピー余極限と可換するという
事実から従う。
\end{proof}

\begin{lemma}
\label{spaces-simplicial-lemma-nullity-cartesian-objects-derived}
状況 \ref{spaces-simplicial-situation-simplicial-site} において、
% P10-E0187
$K, K' \in D(\mathcal{C}_{total})$ とする。次を仮定する。
\begin{enumerate}
\item $K$ はカルテジアンである。
\item $i > 0$ に対して $\Hom(K_i[i], K'_i) = 0$ である。
\item $i \geq 0$ に対して $\Hom(K_i[i + 1], K'_i) = 0$ である。
\end{enumerate}
このとき、零写像 $K_0 \to K'_0$ を誘導する任意の写像 $K \to K'$ は零である。
\end{lemma}

\begin{proof}
補題 \ref{spaces-simplicial-lemma-abelian-postnikov} において $K$ に付随する対象 $X_n$ と
Postnikov 系 $Y_n$ を考える。$K = \text{hocolim} Y_n[n]$ なので、写像
$K \to K'$ は両立する射の族 $Y_n[n] \to K'$ を誘導する。
(1) と補題 \ref{spaces-simplicial-lemma-derived-cartesian-shriek} により
$X_n = g_{n!}K_n$ である。$Y_0 = X_0$ なので、$K_0 \to K'_0$ が零であることから
$Y_0 \to K'$ は零であることが分かる。ある $n \geq 0$ に対して写像
$Y_n[n] \to K'$ が零であることを示したと仮定する。完全三角形
$$
Y_n[n] \to Y_{n + 1}[n + 1] \to X_{n + 1}[n + 1] \to Y_n[n + 1]
$$
から完全列
$$
\Hom(X_{n + 1}[n + 1], K') \to
\Hom(Y_{n + 1}[n + 1], K') \to
\Hom(Y_n[n], K')
$$
を得る。$X_{n + 1}[n + 1] = g_{n + 1!}K_{n + 1}[n + 1]$ なので、最初の群は
$$
\Hom(K_{n + 1}[n + 1], K'_{n + 1})
$$
に等しく、仮定 (2) により零である。帰納法により、すべての写像
$Y_n[n] \to K'$ が零であることが分かる。ホモトピー余極限を定義する完全三角形
$$
\bigoplus Y_n[n] \to
\bigoplus Y_n[n] \to
K \to
(\bigoplus Y_n[n])[1]
$$
を考える。上と同様に論じると、すべての $n \geq 0$ に対して
$$
\Hom((\bigoplus Y_n[n])[1], K') = \prod \Hom(Y_n[n + 1], K')
$$
が零であることを示せば十分である。これを見るには、再び上と同様に論じて、
すべての $n \geq 0$ に対して
$$
\Hom(K_n[n + 1], K'_n)  = 0
$$
を示せば十分であり、これは条件 (3) から従う。
\end{proof}

\begin{lemma}
\label{spaces-simplicial-lemma-hom-cartesian-objects-derived}
状況 \ref{spaces-simplicial-situation-simplicial-site} において、
% P10-E0188
$K, K' \in D(\mathcal{C}_{total})$ とする。次を仮定する。
\begin{enumerate}
\item $K$ はカルテジアンである。
\item $i > 1$ に対して $\Hom(K_i[i - 1], K'_i) = 0$ である。
\end{enumerate}
このとき、$K$ と $K'$ に付随する単体的系の間の任意の写像
$\{K_n \to K'_n\}$ は、$D(\mathcal{C}_{total})$ における写像
$K \to K'$ から生じる。
\end{lemma}

\begin{proof}
$\{K_n \to K'_n\}_{n \geq 0}$ を導来圏の単体的系の射とする。
% P10-E0636
補題 \ref{spaces-simplicial-lemma-abelian-postnikov} において $K$ に付随する対象 $X_n$ と
Postnikov 系 $Y_n$ を考える。(1) と補題
\ref{spaces-simplicial-lemma-derived-cartesian-shriek} により $X_n = g_{n!}K_n$ である。
とくに、写像 $K_0 \to K'_0$ は射 $X_0 \to K'$ を誘導する。
$\{K_n \to K'_n\}$ は系の射なので、計算（省略する）により、合成
$$
X_1 \to X_0 \to K'
$$
は零である。$Y_0 = X_0$ であり、$Y_1$ は完全三角形
$$
Y_1 \to X_1 \to Y_0 \to Y_1[1]
$$
に入るので、射 $Y_1[1] \to K'$ で、その
$X_0 = Y_0 \to Y_1[1]$ との合成が上で与えた射 $X_0 \to K'$ になるものが
存在する。$n \geq 1$ に対する写像 $Y_n[n] \to K'$ が与えられたと仮定する。
完全三角形
$$
X_{n + 1}[n] \to Y_n[n] \to Y_{n + 1}[n + 1] \to X_{n + 1}[n + 1]
$$
から完全列
$$
\Hom(Y_{n + 1}[n + 1], K') \to
\Hom(Y_n[n], K') \to
\Hom(X_{n + 1}[n], K')
$$
を得る。$X_{n + 1}[n] = g_{n + 1!}K_{n + 1}[n]$ なので、最後の群は
$$
\Hom(K_{n + 1}[n], K'_{n + 1})
$$
に等しく、仮定 (2) により零である。帰納法により、遷移写像と両立し、
$Y_0$ 上では与えられた写像に帰着する写像の系 $Y_n[n] \to K'$ を得る。
これは写像
$$
\gamma :
K = \text{hocolim} Y_n[n]
\longrightarrow
K'
$$
を与える。この写像は、いずれにせよ図式
$$
\xymatrix{
X_0 \ar[rd] \ar[r] &
K \ar[d]^\gamma \\
& K'
}
$$
が可換であるという性質をもつ。$\mathcal{C}_0$ に制限すると、写像
$\gamma_0 : K_0 \to K'_0$ は単体的系の射の最初の写像
$K_0 \to K'_0$ と同じであることが分かる。$K$ はカルテジアンなので、これから
$\{\gamma_n\}$ が出発点の単体的系の写像であることが容易に従う。
\end{proof}

\begin{lemma}
\label{spaces-simplicial-lemma-cartesian-object-derived-from-simplicial}
状況 \ref{spaces-simplicial-situation-simplicial-site} において、
$(K_n, K_\varphi)$ を導来圏の単体的系とする。次を仮定する。
\begin{enumerate}
\item $(K_n, K_\varphi)$ はカルテジアンである。
\item $i \geq 0$ かつ $t > 0$ に対して
$\Hom(K_i[t], K_i) = 0$ である。
\end{enumerate}
このとき、付随する単体的系が $(K_n, K_\varphi)$ と同型であるような
$D(\mathcal{C}_{total})$ のカルテジアン対象 $K$ が存在する。
\end{lemma}

\begin{proof}
$D(\mathcal{C}_{total})$ において $X_n = g_{n!}K_n$ とおく。
各 $n \geq 1$ に対して
$$
\begin{aligned}
&\Hom(X_n, X_{n - 1})
\\
&=
\Hom(K_n, g_n^{-1}g_{n - 1!}K_{n - 1})
\\
&=
\bigoplus\nolimits_{\varphi : [n - 1] \to [n]}
\Hom(K_n, f_\varphi^{-1}K_{n - 1})
\end{aligned}
$$
である。したがって、$\varphi$ が $\delta^n_0, \ldots, \delta^n_n$ を走るときの
写像 $K_\varphi^{-1} : K_n \to f_\varphi^{-1}K_{n - 1}$ の交代和に対応する
写像 $X_n \to X_{n - 1}$ を得る。仮定 (1) により $K_\varphi$ は可逆なので、
この操作ができる。補題 \ref{spaces-simplicial-lemma-simplicial-resolution-Z-site} の証明における
写像の定義との類似に注意されたい。$D(\mathcal{C}_{total})$ における複体
$$
\ldots \to X_2 \to X_1 \to X_0
$$
を得る。合成が零であることを示す計算は省略する。『導来圏』補題
\ref{derived-lemma-existence-postnikov-system} により、
$$
\Hom(X_i[i - j - 2], X_j) = 0\text{、ただし }i > j + 2
$$
ならば、この複体を Postnikov 系に拡張できる。この群は
$$
\Hom(K_i[i - j - 2], g_i^{-1}g_{j!}K_j)
$$
に等しい。再び $(K_n, K_\varphi)$ がカルテジアンであることを用いると、
$g_i^{-1}g_{j!}K_j$ は $K_i$ の有限個のコピーの直和と同型であることが分かる。
したがって、この群は仮定 (2) により消滅する。Postnikov 系が
$Y_0 = X_0$ と、$n \geq 1$ に対する完全三角形
$Y_n \to X_n \to Y_{n - 1} \to Y_n[1]$ により与えられるものとする。
$$
K = \text{hocolim} Y_n[n]
$$
とおく。証明を終えるには、すべての $m$ に対して $g_m^{-1}K$ が
$K_m$ と、写像 $K_\varphi$ と両立するように同型であることを示さねばならない。
次に注意する。
$$
g_m^{-1} K = \text{hocolim} g_m^{-1}Y_n[n]
$$
また $g_m^{-1}Y_n[n]$ は $g_m^{-1}X_n$ に対する Postnikov 系である。
同型
$$
g_m^{-1}X_n =
\bigoplus\nolimits_{\varphi : [n] \to [m]} f_\varphi^{-1}K_n
\xrightarrow{\bigoplus K_\varphi}
\bigoplus\nolimits_{\varphi : [n] \to [m]} K_m
$$
を考える。これらの写像は $D(\mathcal{C}_m)$ における複体の同型
$$
\xymatrix{
\ldots \ar[r] &
g_m^{-1}X_2 \ar[r] \ar[d] &
g_m^{-1}X_1 \ar[r] \ar[d] &
g_m^{-1}X_0 \ar[d] \\
\ldots \ar[r] &
\bigoplus\nolimits_{\varphi : [2] \to [m]} K_m \ar[r] &
\bigoplus\nolimits_{\varphi : [1] \to [m]} K_m \ar[r] &
\bigoplus\nolimits_{\varphi : [0] \to [m]} K_m
}
$$
を定める。ここで下段の射は補題
\ref{spaces-simplicial-lemma-simplicial-resolution-Z-site} の証明のとおりである。
各正方形は複体 $\ldots \to X_2 \to X_1 \to X_0$ の射の選び方により
可換である。計算は省略する。下段の複体には
$$
Y'_{m, n} =
\left(\bigoplus\nolimits_{\varphi : [n] \to [m]} \mathbf{Z} \to
\ldots \to
\bigoplus\nolimits_{\varphi : [0] \to [m]} \mathbf{Z}\right)[-n]
\otimes^\mathbf{L}_\mathbf{Z} K_m
$$
により与えられる Postnikov 塔があり、
$\text{hocolim} Y'_{m, n} = K_m$ である
（補題 \ref{spaces-simplicial-lemma-abelian-postnikov} の証明および『導来圏』例
\ref{derived-example-key-postnikov} と比較されたい）。『導来圏』補題
\ref{derived-lemma-existence-postnikov-system} の第二部を適用すると、
$$
\Hom(g_m^{-1}X_i[i - j - 1], \bigoplus\nolimits_{\varphi : [j] \to [m]} K_m)
= 0\text{、ただし }i > j + 1
$$
である限り、大きな図式の垂直写像は Postnikov 系の同型へ拡張される。
% P10-E0189
これは $i > j + 1$ に対して
$\Hom(K_m[i - j - 1], K_m) = 0$ ならば成り立ち、その条件は仮定 (2) から従う。
$D(\mathcal{C}_m)$ における Postnikov 系の同型として
$\gamma_{m, n} : g_m^{-1}Y_n \to Y'_{m, n}$ により与えられる同型を選ぶ。
ホモトピー余極限の一意性により、$\gamma_{m, n}$ と両立する同型
$$
g_m^{-1} K = \text{hocolim} g_m^{-1}Y_n[n]
\xrightarrow{\gamma_m}
\text{hocolim} Y'_{m, n} = K_m
$$
を見つけることができる。

\medskip\noindent
さらに、すべての $\varphi : [m] \to [n]$ に対して、写像 $\gamma_m$ が
可換図式
$$
\xymatrix{
f_\varphi^{-1}g_m^{-1}K \ar[d]_{f_\varphi^{-1}\gamma_m} \ar[r]_{K(\varphi)} &
g_n^{-1}K \ar[d]^{\gamma_n} \\
f_\varphi^{-1}K_m \ar[r]^{K_\varphi} &
K_n
}
$$
に入ることを証明しなければならない。図式
$$
\xymatrix@C=0em{
f_\varphi^{-1}(\bigoplus_{\psi : [0] \to [m]} f_\psi^{-1}K_0)
\ar@{=}[r] \ar[d]_{f_\varphi^{-1}(\bigoplus K_\psi)} &
f_\varphi^{-1}g_m^{-1}X_0 \ar[d] \ar[r]_{X_0(\varphi)} &
g_n^{-1}X_0 \ar[d] &
\bigoplus_{\chi : [0] \to [n]} f_\chi^{-1}K_0
\ar@{=}[l] \ar[d]^{\bigoplus K_\chi} \\
f_\varphi^{-1}(\bigoplus_{\psi : [0] \to [m]} K_m) \ar@{=}[d] &
f_\varphi^{-1}g_m^{-1}K \ar[d]_{f_\varphi^{-1}\gamma_m} \ar[r]_{K(\varphi)} &
g_n^{-1}K \ar[d]^{\gamma_n} &
\bigoplus_{\chi : [0] \to [n]} K_n \ar@{=}[d] \\
f_\varphi^{-1}Y'_{0, m} \ar[r] &
f_\varphi^{-1}K_m \ar[r]^{K_\varphi} &
K_n &
Y'_{0, n} \ar[l]
}
$$
を考える。上中央の正方形は $X_0 \to K$ が単体的対象の射なので可換である。
$\gamma_m$（それぞれ $\gamma_n$）は $\gamma_{0, m}$（それぞれ
$\gamma_{0, n}$）と両立し、後者は図式中の射 $\bigoplus K_\psi$ と
$\bigoplus K_\chi$ なので、左の長方形（それぞれ右の長方形）は可換である。
% P10-E0190
$(K_n, K_\varphi)$ は単体的系であり、写像 $X_0(\varphi)$ は明らかな同一視
$f_\varphi^{-1}f_\psi^{-1}K_0 = f_{\varphi \circ \psi}^{-1}K_0$
により与えられるので、図式の外側の長方形は可換である。
射 $\bigoplus_\psi K_m \to Y'_{0, m} \to K_m$ は任意の直和因子上で
同型を誘導することに注意する（上での Postnikov 系 $Y'_{i, j}$ の明示的構成に
よる）。
% P10-E0191
したがって、左上隅の一つの直和因子をとると、$\gamma_m$ が同型なので、これは
$f_\varphi^{-1}g_m^{-1}K$ へ同型に写る。上の主張をこの直和因子だけに着目して
書き下すと、下中央の正方形も可換であることが分かる。これで証明が完了する。
\end{proof}






\section{導来圏の単体的系：加群}
\label{spaces-simplicial-section-glueing-modules}

\noindent
本節では、$\mathcal{O}$-加群の導来圏という設定で
\cite[Proposition 3.2.9]{BBD} の特別な場合を証明する。アーベル層の場合は
（少し）容易であり、節 \ref{spaces-simplicial-section-glueing} で論じた。

\begin{definition}
\label{spaces-simplicial-definition-cartesian-derived-modules}
状況 \ref{spaces-simplicial-situation-simplicial-site} において、$\mathcal{O}$ を
$\mathcal{C}_{total}$ 上の環の層とする。
{\it 加群の導来圏の単体的系}とは、次のデータ
\begin{enumerate}
\item すべての $n$ に対する $D(\mathcal{O}_n)$ の対象 $K_n$、
\item すべての $\varphi : [m] \to [n]$ に対する
$D(\mathcal{O}_n)$ における写像
$K_\varphi : Lf_\varphi^*K_m \to K_n$
\end{enumerate}
で、$\Delta$ の任意の射 $\varphi : [m] \to [n]$ と
$\psi : [l] \to [m]$ に対して条件
$$
K_{\varphi \circ \psi} = K_\varphi \circ Lf_\varphi^*K_\psi :
Lf_{\varphi \circ \psi}^*K_l = Lf_\varphi^* Lf_\psi^*K_l
\longrightarrow
K_n
$$
を満たすものからなる。すべての $\varphi$ に対して写像 $K_\varphi$ が
同型であるとき、この単体的系を {\it カルテジアン}という。
導来圏の二つの単体的系が与えられたとき、
{\it 加群の導来圏の単体的系の射}には明らかな意味がある。
\end{definition}

\noindent
この概念には意図的に途方もなく長い名前を付けた。目的は、ある仮定のもとで、
加群の導来圏の単体的系が $D(\mathcal{O})$ の対象から生じることを示すことである。

\begin{lemma}
\label{spaces-simplicial-lemma-cartesian-objects-derived-modules}
状況 \ref{spaces-simplicial-situation-simplicial-site} において、$\mathcal{O}$ を
$\mathcal{C}_{total}$ 上の環の層とする。
$K \in D(\mathcal{O})$ が対象ならば、$(K_n, K(\varphi))$ は
加群の導来圏の単体的系である。$K$ がカルテジアンならば、この系も
カルテジアンである。
\end{lemma}

\begin{proof}
定義から直ちに従う。
\end{proof}

\begin{lemma}
\label{spaces-simplicial-lemma-construct-cartesian-objects-derived-modules}
状況 \ref{spaces-simplicial-situation-simplicial-site} において、$\mathcal{O}$ を
$\mathcal{C}_{total}$ 上の環の層とする。
$K_0 \in D(\mathcal{O}_0)$ と、余サイクル条件を満たす同型
$$
\alpha :
L(f_{\delta_1^1})^*K_0
\longrightarrow
L(f_{\delta_0^1})^*K_0
$$
が与えられていると仮定する。$\tau^n_i : [0] \to [n]$、$0 \mapsto i$ とおき、
$K_n = Lf_{\tau^n_n}^*K_0$ とおく。対象 $K_n$ は加群の導来圏の
カルテジアン単体的系の成員をなす。
\end{lemma}

\begin{proof}
補題 \ref{spaces-simplicial-lemma-construct-cartesian-objects-derived}、
\ref{spaces-simplicial-lemma-characterize-cartesian} および後者の証明と比較されたい
（余サイクル条件の詳しい形もそこに書かれている）。構成は『降下』節
\ref{descent-section-descent-modules} で論じられた構成と類似しており、
そこから記法 $\tau^n_i : [0] \to [n]$、$0 \mapsto i$ および
$\tau^n_{ij} : [1] \to [n]$、$0 \mapsto i$、$1 \mapsto j$ を借用する。
$\varphi : [n] \to [m]$ が与えられたとき、次を用いて
$K_\varphi : L(f_\varphi)^*K_n \to K_m$ を定義する。
$$
\xymatrix@C=-0.6em{
L(f_\varphi)^*K_n \ar@{=}[r] &
L(f_\varphi)^* L(f_{\tau^n_n})^*K_0 \ar@{=}[r] &
L(f_{\tau^m_{\varphi(n)}})^*K_0 \ar@{=}[r] &
L(f_{\tau^m_{\varphi(n)m}})^* L(f_{\delta^1_1})^*K_0
\ar[d]_{L(f_{\tau^m_{\varphi(n)m}})^*\alpha} \\
&
K_m \ar@{=}[r] &
L(f_{\tau^m_m})^*K_0 \ar@{=}[r] &
L(f_{\tau^m_{\varphi(n)m}})^* L(f_{\delta^1_0})^*K_0
}
$$
% P10-E0192
余サイクル条件から写像が（それぞれの導来圏で）正しく合成され、したがって
加群の導来圏の単体的系を与えることの確認は省略する。
\end{proof}

\begin{lemma}
\label{spaces-simplicial-lemma-modules-postnikov}
状況 \ref{spaces-simplicial-situation-simplicial-site} において、$\mathcal{O}$ を
$\mathcal{C}_{total}$ 上の環の層とする。
% P10-E0193
$K$ を $D(\mathcal{C}_{total})$ の対象とする。$D(\mathcal{O})$ の対象として
$$
X_n = (g_{n!}\mathcal{O}_n)
\otimes^\mathbf{L}_\mathcal{O} K
\quad\text{および}\quad
Y_n =
(g_{n!}\mathcal{O}_n \to \ldots \to g_{0!}\mathcal{O}_0)[-n]
\otimes^\mathbf{L}_\mathcal{O} K
$$
とおく。ここで写像は
% P10-E0194
補題 \ref{spaces-simplicial-lemma-simplicial-resolution-Z-site} のとおりである。
明らかな標準写像 $Y_n \to X_n$ および
$Y_0 \to Y_1[1] \to Y_2[2] \to \ldots$ とともに、次が成り立つ。
\begin{enumerate}
\item 完全三角形 $Y_n \to X_n \to Y_{n - 1} \to Y_n[1]$ は、
$\ldots \to X_2 \to X_1 \to X_0$ に対する Postnikov 系
（『導来圏』定義 \ref{derived-definition-postnikov-system}）を定める。
\item $D(\mathcal{O})$ において $K = \text{hocolim} Y_n[n]$ である。
\end{enumerate}
\end{lemma}

\begin{proof}
まず $K = \mathcal{O}$ ならば、これは $\textit{Ab}(\mathcal{C}_{total})$ の複体
$$
\ldots \to
g_{2!}\mathcal{O}_2 \to
g_{1!}\mathcal{O}_1 \to
g_{0!}\mathcal{O}_0
$$
に適用した『導来圏』例 \ref{derived-example-key-postnikov} の構成と、補題
\ref{spaces-simplicial-lemma-simplicial-resolution-ringed} によりこの複体が
$D(\mathcal{C}_{total})$ において $K = \mathcal{O}$ を表すという事実を
組み合わせたものである。一般の場合はこれと、完全関手
$- \otimes^\mathbf{L}_\mathcal{O} K$ が Postnikov 系を Postnikov 系へ送り、
$- \otimes^\mathbf{L}_\mathcal{O} K$ がホモトピー余極限と可換するという
事実から従う。
\end{proof}

\begin{lemma}
\label{spaces-simplicial-lemma-nullity-cartesian-modules-derived}
状況 \ref{spaces-simplicial-situation-simplicial-site} において、$\mathcal{O}$ を
$\mathcal{C}_{total}$ 上の環の層とする。
% P10-E0195
$K, K' \in D(\mathcal{O})$ とする。次を仮定する。
\begin{enumerate}
% P10-E0196
\item $\varphi : [m] \to [n]$ に対して
$f_\varphi^{-1}\mathcal{O}_n \to \mathcal{O}_m$ は平坦である。
\item $K$ はカルテジアンである。
\item $i > 0$ に対して $\Hom(K_i[i], K'_i) = 0$ である。
\item $i \geq 0$ に対して $\Hom(K_i[i + 1], K'_i) = 0$ である。
\end{enumerate}
このとき、零写像 $K_0 \to K'_0$ を誘導する任意の写像 $K \to K'$ は零である。
\end{lemma}

\begin{proof}
証明は補題 \ref{spaces-simplicial-lemma-nullity-cartesian-objects-derived} の証明と
まったく同じであるが、補題 \ref{spaces-simplicial-lemma-modules-postnikov} を
補題 \ref{spaces-simplicial-lemma-abelian-postnikov} の代わりに用いる。
\end{proof}

\begin{lemma}
\label{spaces-simplicial-lemma-hom-cartesian-modules-derived}
状況 \ref{spaces-simplicial-situation-simplicial-site} において、$\mathcal{O}$ を
$\mathcal{C}_{total}$ 上の環の層とする。
% P10-E0197
$K, K' \in D(\mathcal{O})$ とする。次を仮定する。
\begin{enumerate}
% P10-E0196
\item $\varphi : [m] \to [n]$ に対して
$f_\varphi^{-1}\mathcal{O}_n \to \mathcal{O}_m$ は平坦である。
\item $K$ はカルテジアンである。
\item $i > 1$ に対して $\Hom(K_i[i - 1], K'_i) = 0$ である。
\end{enumerate}
このとき、$K$ と $K'$ に付随する単体的系の間の任意の写像
$\{K_n \to K'_n\}$ は、$D(\mathcal{O})$ における写像 $K \to K'$ から生じる。
\end{lemma}

\begin{proof}
証明は補題 \ref{spaces-simplicial-lemma-hom-cartesian-objects-derived} の証明と
まったく同じであるが、補題 \ref{spaces-simplicial-lemma-modules-postnikov} を
補題 \ref{spaces-simplicial-lemma-abelian-postnikov} の代わりに用いる。
\end{proof}

\begin{lemma}
\label{spaces-simplicial-lemma-cartesian-module-derived-from-simplicial}
状況 \ref{spaces-simplicial-situation-simplicial-site} において、$\mathcal{O}$ を
$\mathcal{C}_{total}$ 上の環の層とする。$(K_n, K_\varphi)$ を
加群の導来圏の単体的系とする。次を仮定する。
\begin{enumerate}
% P10-E0196
\item $\varphi : [m] \to [n]$ に対して
$f_\varphi^{-1}\mathcal{O}_n \to \mathcal{O}_m$ は平坦である。
\item $(K_n, K_\varphi)$ はカルテジアンである。
\item $i \geq 0$ かつ $t > 0$ に対して
$\Hom(K_i[t], K_i) = 0$ である。
\end{enumerate}
このとき、付随する単体的系が $(K_n, K_\varphi)$ と同型であるような
$D(\mathcal{O})$ のカルテジアン対象 $K$ が存在する。
\end{lemma}

\begin{proof}
証明は補題 \ref{spaces-simplicial-lemma-cartesian-object-derived-from-simplicial} の証明に、
次の変更を加えたものとまったく同じである。
\begin{enumerate}
\item 全体を通じて $g_n^{-1}$ の代わりに $g_n^* = Lg_n^*$ を用いる。
\item 全体を通じて $f_\varphi^{-1}$ の代わりに
$f_\varphi^* = Lf_\varphi^*$ を用いる。
\item 補題 \ref{spaces-simplicial-lemma-simplicial-resolution-ringed} を、補題
\ref{spaces-simplicial-lemma-simplicial-resolution-Z-site} の代わりに参照する。
\item $Y'_{m, n}$ の構成では $\mathbf{Z}$ の代わりに $\mathcal{O}_m$ を用いる。
\item 補題 \ref{spaces-simplicial-lemma-modules-postnikov} の証明と比較し、補題
\ref{spaces-simplicial-lemma-abelian-postnikov} の証明とは比較しない。
\end{enumerate}
これで証明が終わる。
\end{proof}

\section{半表現可能対象に付随するサイト}
\label{spaces-simplicial-section-semi-representable}

\noindent
$\mathcal{C}$ をサイトとする。$\mathcal{C}$ の {\it 半表現可能対象}とは、
単に $\mathcal{C}$ の対象の族 $\{U_i\}_{i \in I}$ のことである。
半表現可能対象の {\it 射
$\{U_i\}_{i \in I} \to \{V_j\}_{j \in J}$} とは、写像
$\alpha : I \to J$ と、各 $i \in I$ に対する $\mathcal{C}$ の射
$f_i : U_i \to V_{\alpha(i)}$ からなる。
$\mathcal{C}$ の半表現可能対象の圏を $\text{SR}(\mathcal{C})$ と書く。
詳しくは『超被覆』定義 \ref{hypercovering-definition-SR} および
それを含む節を参照されたい。

\medskip\noindent
$\mathcal{C}$ の半表現可能対象 $K = \{U_i\}_{i \in I}$ に対し、
$$
\mathcal{C}/K = \coprod\nolimits_{i \in I} \mathcal{C}/U_i
$$
を、$U_i$ における $\mathcal{C}$ の局所化の非交和とする。
この圏には自然なサイト構造があり、被覆は各局所化
$\mathcal{C}/U_i$ から継承される。サイト $\mathcal{C}/K$ を
{\it $K$ における $\mathcal{C}$ の局所化}と呼ぶ。
$\mathcal{C}/K$ 上の層を与えることは、各 $\mathcal{C}/U_i$ 上の
層 $\mathcal{F}_i$ の族を与えることと同じである。すなわち、
$$
\Sh(\mathcal{C}/K) = \prod\nolimits_{i \in I} \Sh(\mathcal{C}/U_i)
$$
である。この記述は、何が起きているかを理解するうえで時に有用である。

\medskip\noindent
$\mathcal{C}$ をサイトとし、$K = \{U_i\}_{i \in I}$ を
$\text{SR}(\mathcal{C})$ の対象とする。連続かつ余連続な局所化関手
$j : \mathcal{C}/K \to \mathcal{C}$ があり、これは局所化関手
% P10-E0198
$j_i : \mathcal{C}/V_i \to \mathcal{C}$ の積である。
『サイト』節 \ref{sites-section-localize} とまったく同様に、関手
$j_!$, $j^{-1}$, $j_*$ を得る。積分解
$\Sh(\mathcal{C}/K) = \prod\nolimits_{i \in I} \Sh(\mathcal{C}/U_i)$
を用いると、
$$
\begin{matrix}
j_! & : &
(\mathcal{F}_i)_{i \in I} &
\longmapsto &
\coprod j_{i, !}\mathcal{F}_i \\
j^{-1} & : &
\mathcal{G} &
\longmapsto &
(j_i^{-1}\mathcal{G})_{i \in I} \\
j_* & : &
(\mathcal{F}_i)_{i \in I} &
\longmapsto &
\prod j_{i, *}\mathcal{F}_i
\end{matrix}
$$
となることは容易に確かめられる。

\medskip\noindent
$f : K \to L$ を $\text{SR}(\mathcal{C})$ の射とする。このとき、
各成分に『サイト』補題 \ref{sites-lemma-relocalize} の構成を適用して、
連続かつ余連続な関手
$$
v : \mathcal{C}/K \longrightarrow \mathcal{C}/L
$$
を得る。より具体的に、$f = (\alpha, f_i)$ であり、
$K = \{U_i\}_{i \in I}$、$L = \{V_j\}_{j \in J}$、
$\alpha : I \to J$、および $f_i : U_i \to V_{\alpha(i)}$ とする。
このとき関手 $v$ は、上述の補題の構成を介して成分
$\mathcal{C}/U_i$ を成分 $\mathcal{C}/V_{\alpha(i)}$ へ写す。
とくに、トポスの射
$$
f : \Sh(\mathcal{C}/K) \to \Sh(\mathcal{C}/L)
$$
を得る。積分解\linebreak
$\Sh(\mathcal{C}/K) = \prod\nolimits_{i \in I} \Sh(\mathcal{C}/U_i)$ および
\linebreak
$\Sh(\mathcal{C}/L) = \prod\nolimits_{j \in J} \Sh(\mathcal{C}/V_j)$
を用いると、
$$
\begin{matrix}
f_! & : &
(\mathcal{F}_i)_{i \in I} &
\longmapsto &
(\coprod\nolimits_{i \in I, \alpha(i) = j} f_{i, !}\mathcal{F}_i)_{j \in J} \\
f^{-1} & : &
(\mathcal{G}_j)_{j \in J} &
\longmapsto &
(f_i^{-1}\mathcal{G}_{\alpha(i)})_{i \in I} \\
f_* & : &
(\mathcal{F}_i)_{i \in I} &
\longmapsto &
(\prod\nolimits_{i \in I, \alpha(i) = j} f_{i, *}\mathcal{F}_i)_{j \in J}
\end{matrix}
$$
となる。ここで $f_i : \Sh(\mathcal{C}/U_i) \to \Sh(\mathcal{C}/V_{\alpha(i)})$
は、$f_i : U_i \to V_{\alpha(i)}$ に対応する局所化関手
$\mathcal{C}/U_i \to \mathcal{C}/V_{\alpha(i)}$ に付随する射である。

\begin{lemma}
\label{spaces-simplicial-lemma-has-P}
$\mathcal{C}$ をサイトとする。
\begin{enumerate}
\item $\text{SR}(\mathcal{C})$ の $K$ に対して、関手
$j : \mathcal{C}/K \to \mathcal{C}$ は連続かつ余連続であり、
『サイト』注意 \ref{sites-remark-cartesian-cocontinuous} の性質 P をもつ。
\item $\text{SR}(\mathcal{C})$ の $f : K \to L$ に対して、関手
$v : \mathcal{C}/K \to \mathcal{C}/L$（上記参照）は連続かつ余連続であり、
『サイト』注意 \ref{sites-remark-cartesian-cocontinuous} の性質 P をもつ。
\end{enumerate}
\end{lemma}

\begin{proof}
(2) を証明する。補題の直前の議論の記法を用いると、局所化関手
$\mathcal{C}/U_i \to \mathcal{C}/V_{\alpha(i)}$ は
『サイト』節 \ref{sites-section-localize} により連続かつ余連続であり、
『サイト』注意 \ref{sites-remark-localization-cartesian-cocontinuous} により
$P$ を満たす。したがって、$v$ が連続かつ余連続で性質 $P$ をもつことは
形式的に従う。詳細は省略する。(1) の証明も省略する。
\end{proof}

\begin{lemma}
\label{spaces-simplicial-lemma-push-pull-localization}
$\mathcal{C}$ をサイトとし、$K$ を $\text{SR}(\mathcal{C})$ の対象とする。
$\Sh(\mathcal{C})$ の $\mathcal{F}$ に対して、
$$
j_*j^{-1}\mathcal{F} = \SheafHom(F(K)^\#, \mathcal{F})
$$
が成り立つ。ここで $F$ は『超被覆』定義
\ref{hypercovering-definition-SR-F} のものとする。
\end{lemma}

\begin{proof}
$K = \{U_i\}_{i \in I}$ と書く。上で与えた関手 $j^{-1}$ と $j_*$ の
記述を用いると、
$$
j_*j^{-1}\mathcal{F} =
\prod\nolimits_{i \in I} j_{i, *}(\mathcal{F}|_{\mathcal{C}/U_i}) =
\prod\nolimits_{i \in I} \SheafHom(h_{U_i}^\#, \mathcal{F})
$$
を得る。第二の等式は『サイト』補題
\ref{sites-lemma-hom-sheaf-hU} による。
% P10-E0199
$F(K) = \coprod h_{U_i}$ が $\textit{PSh}(\mathcal{C}$ において成り立つので、
$\Sh(\mathcal{C})$ において $F(K)^\# = \coprod h_{U_i}^\#$ である。
さらに $\SheafHom(-, \mathcal{F})$ は余積を積へ写す
（『サイト』節 \ref{sites-section-glueing-sheaves} の構成から直ちに従う）ので、
結論を得る。
\end{proof}

\begin{lemma}
\label{spaces-simplicial-lemma-localize-compare}
$\mathcal{C}$ をサイトとする。
\begin{enumerate}
\item $\text{SR}(\mathcal{C})$ の $K$ に対し、関手 $j_!$ は同値
$\Sh(\mathcal{C}/K) \to \Sh(\mathcal{C})/F(K)^\#$ を与える。
ここで $F$ は『超被覆』定義 \ref{hypercovering-definition-SR-F} のものとする。
\item 関手 $j^{-1} : \Sh(\mathcal{C}) \to \Sh(\mathcal{C}/K)$ は、
(1) の同一視のもとで
$\mathcal{F} \mapsto (\mathcal{F} \times F(K)^\# \to F(K)^\#)$
に対応する。
\item $\text{SR}(\mathcal{C})$ の $f : K \to L$ に対し、関手
$f^{-1}$ は (1) の同一視のもとで関手
$\Sh(\mathcal{C})/F(L)^\# \to \Sh(\mathcal{C})/F(K)^\#$、
$(\mathcal{G} \to F(L)^\#) \mapsto
(\mathcal{G} \times_{F(L)^\#} F(K)^\# \to F(K)^\#)$
に対応する。
\end{enumerate}
\end{lemma}

\begin{proof}
$K = \{U_i\}_{i \in I}$ ならば、圏 $\Sh(\mathcal{C}/K)$ は圏
$\Sh(\mathcal{C}/U_i)$ の積に分解される。また、
$F(K)^\# = \coprod_{i \in I} h_{U_i}^\#$ である（層の圏における余積）。
したがって $\Sh(\mathcal{C})/F(K)^\#$ は圏
$\Sh(\mathcal{C})/h_{U_i}^\#$ の積である。ゆえに (1) と (2) は
各 $i$ に対する対応する主張から従う。『サイト』補題
\ref{sites-lemma-essential-image-j-shriek} および
\ref{sites-lemma-compute-j-shriek-restrict} を参照されたい。
同様に、$L = \{V_j\}_{j \in J}$ であり、$f$ が
$\alpha : I \to J$ と $f_i : U_i \to V_{\alpha(i)}$ により与えられるなら、
各再局所化射 $\mathcal{C}/U_i \to \mathcal{C}/V_{\alpha(i)}$ に
『サイト』補題 \ref{sites-lemma-relocalize-explicit} を適用して (3) を得る。
\end{proof}

\begin{lemma}
\label{spaces-simplicial-lemma-localize-injective}
$\mathcal{C}$ をサイトとする。$\text{SR}(\mathcal{C})$ の $K$ に対し、
関手 $j^{-1}$ は入射的アーベル層を入射的アーベル層へ写す。同様に、
関手 $j^{-1}$ はアーベル層の K-入射的複体を
アーベル層の K-入射的複体へ写す。
\end{lemma}

\begin{proof}
第一の主張は、『サイト上のコホモロジー』補題
\ref{sites-cohomology-lemma-cohomology-of-open} の
半表現可能対象への自然な一般化である。実際、この補題から、
$\Sh(\mathcal{C}/K)$ の積分解と上で与えた関手 $j^{-1}$ の記述によって従う。
第二の主張は、『サイト上のコホモロジー』補題
\ref{sites-cohomology-lemma-restrict-K-injective-to-open} の
自然な一般化であり、トポスの積分解によってこの補題から従う。

\medskip\noindent
別の証明：補題 \ref{spaces-simplicial-lemma-localize-compare} の (1) により
$j$ はトポスの局所化を誘導するので、サイトを拡大することで
『サイト上のコホモロジー』補題
\ref{sites-cohomology-lemma-cohomology-of-open} および
\ref{sites-cohomology-lemma-restrict-K-injective-to-open} からも直ちに従う。
入射的層の場合については、『サイト上のコホモロジー』補題
\ref{sites-cohomology-lemma-cohomology-on-sheaf-sets} の証明と比較されたい。
\end{proof}

\begin{remark}[対象上の場合の変種]
\label{spaces-simplicial-remark-semi-representable-over-object}
$\mathcal{C}$ をサイトとし、$X \in \Ob(\mathcal{C})$ とする。
{\it $X$ 上の半表現可能対象}の圏 $\text{SR}(\mathcal{C}, X)$ を、式
$\text{SR}(\mathcal{C}, X) = \text{SR}(\mathcal{C}/X)$ により定義する。
『超被覆』定義 \ref{hypercovering-definition-SR} を参照されたい。
したがって、上の議論をサイト $\mathcal{C}/X$ に適用できる。
簡潔に述べると、上の構成は次を与える。
\begin{enumerate}
\item $\text{SR}(\mathcal{C}, X)$ の $K$ に対するサイト $\mathcal{C}/K$、
\item $K = \{U_i/X\}$ のときの分解
$\Sh(\mathcal{C}/K) = \prod \Sh(\mathcal{C}/U_i)$、
\item 局所化関手 $j : \mathcal{C}/K \to \mathcal{C}/X$、
\item $\text{SR}(\mathcal{C}, X)$ の $f : K \to L$ に対する射
$f : \Sh(\mathcal{C}/K) \to \Sh(\mathcal{C}/L)$。
\end{enumerate}
この状況では、この節のすべての結果が、至る所で $\mathcal{C}$ を
$\mathcal{C}/X$ に置き換えることによって成り立つ。
\end{remark}

\begin{remark}[環付きの場合の変種]
\label{spaces-simplicial-remark-semi-representable-ringed}
$\mathcal{C}$ をサイトとし、$\mathcal{O}_\mathcal{C}$ を
$\mathcal{C}$ 上の環の層とする。この場合、$\mathcal{C}$ の任意の
半表現可能対象 $K$ に対して、サイト $\mathcal{C}/K$ は環の層
$\mathcal{O}_K = j^{-1}\mathcal{O}_\mathcal{C}$ をもつ環付きサイトである。
上の構成は次を与える。
\begin{enumerate}
\item $\text{SR}(\mathcal{C})$ の $K$ に対する環付きサイト
$(\mathcal{C}/K, \mathcal{O}_K)$、
\item $K = \{U_i\}$ のときの分解
$\textit{Mod}(\mathcal{O}_K) =
\prod \textit{Mod}(\mathcal{O}_{U_i})$、
\item 環付きトポスの局所化射
$j : (\Sh(\mathcal{C}/K), \mathcal{O}_K) \to
(\Sh(\mathcal{C}), \mathcal{O}_\mathcal{C})$、
\item $\text{SR}(\mathcal{C})$ の $f : K \to L$ に対する環付きトポスの射
\linebreak
$f : (\Sh(\mathcal{C}/K), \mathcal{O}_K) \to
(\Sh(\mathcal{C}/L), \mathcal{O}_L)$。
\end{enumerate}
上の結果の多くはこの設定でも成り立つ。たとえば、関手 $j^*$ は完全な左随伴
$$
j_! : \textit{Mod}(\mathcal{O}_K) \to \textit{Mod}(\mathcal{O}_\mathcal{C}),
$$
をもち、(2) の積分解のもとで $(\mathcal{F}_i)_{i \in I}$ を
$\bigoplus j_{i, !}\mathcal{F}_i$ へ写す。同様に、上のような
$f : K \to L$ が与えられると、関手 $f^*$ は完全な左随伴
$f_! : \textit{Mod}(\mathcal{O}_K) \to \textit{Mod}(\mathcal{O}_L)$ をもつ。
したがって関手 $j^*$ と $f^*$ は完全である。すなわち、$j$ と $f$ は
環付きトポスの平坦射である（これは等式
$\mathcal{O}_K = j^{-1}\mathcal{O}_\mathcal{C}$ および
$\mathcal{O}_K = f^{-1}\mathcal{O}_L$ からも従う）。
\end{remark}

\begin{remark}[対象上の環付きの場合の変種]
\label{spaces-simplicial-remark-semi-representable-ringed-over-object}
$\mathcal{C}$ をサイトとし、$\mathcal{O}_\mathcal{C}$ を
$\mathcal{C}$ 上の環の層とする。$X \in \Ob(\mathcal{C})$ とし、
% P10-E0200
$\mathcal{O}_X = \mathcal{O}_\mathcal{C}|_{\mathcal{C}/U}$ と書く。
注意 \ref{spaces-simplicial-remark-semi-representable-over-object} と
\ref{spaces-simplicial-remark-semi-representable-ringed} の構成を組み合わせて、次を得る。
\begin{enumerate}
\item $\text{SR}(\mathcal{C}, X)$ の $K$ に対する環付きサイト
$(\mathcal{C}/K, \mathcal{O}_K)$、
\item $K = \{U_i\}$ のときの分解
$\textit{Mod}(\mathcal{O}_K) =
\prod \textit{Mod}(\mathcal{O}_{U_i})$、
\item 環付きトポスの局所化射
$j : (\Sh(\mathcal{C}/K), \mathcal{O}_K) \to
(\Sh(\mathcal{C}/X), \mathcal{O}_X)$、
\item $\text{SR}(\mathcal{C}, X)$ の $f : K \to L$ に対する環付きトポスの射
\linebreak
$f : (\Sh(\mathcal{C}/K), \mathcal{O}_K) \to
(\Sh(\mathcal{C}/L), \mathcal{O}_L)$。
\end{enumerate}
もちろん、注意 \ref{spaces-simplicial-remark-semi-representable-ringed} で述べた結果は
すべてこの設定でも成り立つ。
\end{remark}

\section{単体的半表現可能対象に付随するサイト}
\label{spaces-simplicial-section-simplicial-semi-representable}
% P10-E0201

\noindent
$\mathcal{C}$ をサイトとし、$K$ を $\text{SR}(\mathcal{C})$ の
単体的対象とする。通常どおり $K_n = K([n])$ とおき、
$\varphi : [m] \to [n]$ に付随する射を
$K(\varphi) : K_n \to K_m$ と書く。節 \ref{spaces-simplicial-section-semi-representable} の
構成により、対象がサイトで射が余連続関手である圏の単体的対象
$n \mapsto \mathcal{C}/K_n$ を得る。言い換えると、節
\ref{spaces-simplicial-section-simplicial-sites} の場合 B のデータを得る。
補題 \ref{spaces-simplicial-lemma-has-P} により、これらの関手は性質 P を満たす。
したがって補題 \ref{spaces-simplicial-lemma-simplicial-cocontinuous-site} を適用して、
サイト $(\mathcal{C}/K)_{total}$ を得る。

\medskip\noindent
サイト $(\mathcal{C}/K)_{total}$ は次のように明示的に記述できる。
$K_n = \{U_{n, i}\}_{i \in I_n}$ と書く。
$\varphi : [m] \to [n]$ に対して、射 $K(\varphi) : K_n \to K_m$ は、
写像 $\alpha(\varphi) : I_n \to I_m$ と、$i \in I_n$ ごとの射
$f_{\varphi, i} : U_{n, i} \to U_{m, \alpha(\varphi)(i)}$ により与えられる。
このとき次が成り立つ。
\begin{enumerate}
\item $(\mathcal{C}/K)_{total}$ の対象は、ある $n$ とある $i \in I_n$ に対する
$\mathcal{C}/U_{n, i}$ の対象 $(U/U_{n, i})$ に対応する。
\item $U/U_{n, i}$ から $V/U_{m, j}$ への射は対 $(\varphi, f)$ であり、
ここで $\varphi : [m] \to [n]$、$j = \alpha(\varphi)(i)$、かつ
$f : U \to V$ は $\mathcal{C}$ の射であって、
$$
\vcenter{
\xymatrix{
U \ar[r]_f \ar[d] & V \ar[d] \\
U_{n, i} \ar[r]^-{f_{\varphi, i}} &
U_{m, j}
}
}
$$
が可換である。
\item 対象 $U/U_{n, i}$ の被覆は、$\mathcal{C}$ の被覆
$\{f_j : U_j \to U\}$ から出発し、
$\{(\text{id}, f_j) : U_j/U_{n, i} \to U/U_{n, i}\}$ を
$(\mathcal{C}/K)_{total}$ の被覆とすることにより構成される。
\end{enumerate}
単体的サイトについて構成した一般論はすべて
$(\mathcal{C}/K)_{total}$ に適用できる。明らかな忘却関手
$$
j_{total} : (\mathcal{C}/K)_{total} \longrightarrow \mathcal{C}
$$
は連続かつ余連続である。実際、付随するトポスの射は
（明らかな）増大から生じる。

\begin{lemma}
\label{spaces-simplicial-lemma-augmentation-simplicial-semi-representable}
$\mathcal{C}$ をサイトとし、$K$ を $\text{SR}(\mathcal{C})$ の
単体的対象とする。局所化関手 $j_0 : \mathcal{C}/K_0 \to \mathcal{C}$ は、
注意 \ref{spaces-simplicial-remark-augmentation-site} の場合 (B) のように、増大
$a_0 : \Sh(\mathcal{C}/K_0) \to \Sh(\mathcal{C})$ を定める。
補題 \ref{spaces-simplicial-lemma-augmentation-site} の対応するトポスの射
$$
a_n : \Sh(\mathcal{C}/K_n) \longrightarrow \Sh(\mathcal{C}),\quad
a : \Sh((\mathcal{C}/K)_{total}) \longrightarrow \Sh(\mathcal{C})
$$
は、連続かつ余連続な局所化関手\linebreak
$j_n : \mathcal{C}/K_n \to \mathcal{C}$ および
\linebreak
$j_{total} : (\mathcal{C}/K)_{total} \to \mathcal{C}$ に付随する
トポスの射に等しい。
\end{lemma}

\begin{proof}
定義をたどれば直ちに分かる。とくに、$a$ と $j_{total}$ の関係については、
補題 \ref{spaces-simplicial-lemma-augmentation-site} の証明中の脚注を参照されたい。
\end{proof}

\begin{lemma}
\label{spaces-simplicial-lemma-comparison}
% P10-E0202
補題 \ref{spaces-simplicial-lemma-augmentation-simplicial-semi-representable} と同じ仮定と記法のもとで、
次の性質が成り立つ。
\begin{enumerate}
\item $a^{-1} : \Sh(\mathcal{C}) \to \Sh((\mathcal{C}/K)_{total})$ の
左随伴である関手\linebreak
$a^{Sh}_! : \Sh((\mathcal{C}/K)_{total}) \to \Sh(\mathcal{C})$ が存在する。
\item
$a^{-1} : \textit{Ab}(\mathcal{C}) \to \textit{Ab}((\mathcal{C}/K)_{total})$
の左随伴である関手\linebreak
$a_! : \textit{Ab}((\mathcal{C}/K)_{total}) \to \textit{Ab}(\mathcal{C})$
が存在する。
\item 関手 $a^{-1}$ は、$\Sh(\mathcal{C})$ の $\mathcal{F}$ に、
次数 $n$ で $a_n^{-1}\mathcal{F}$ に等しい
$(\mathcal{C}/K)_{total}$ 上の層を対応させる。
\item 関手 $a_*$ は、$\textit{Ab}((\mathcal{C}/K)_{total})$ の
$\mathcal{G}$ に、二つの写像
$j_{0, *}\mathcal{G}_0 \to j_{1, *}\mathcal{G}_1$ の等化子を対応させる。
\end{enumerate}
\end{lemma}

\begin{proof}
(3) と (4) は単体的サイトの任意の増大について成り立つ。補題
\ref{spaces-simplicial-lemma-augmentation-site} を参照されたい。
(1) と (2) は $j_{total}$ が連続かつ余連続であることから従う。
関手 $a^{Sh}_!$ は『サイト』補題 \ref{sites-lemma-when-shriek} で構成され、
関手 $a_!$ は『サイト上の加群』補題
\ref{sites-modules-lemma-g-shriek-adjoint} で構成される。
\end{proof}

\begin{lemma}
\label{spaces-simplicial-lemma-sanity-check-simplicial-semi-representable}
$\mathcal{C}$ をサイトとし、$K$ を $\text{SR}(\mathcal{C})$ の
単体的対象とする。$U/U_{n, i}$ を $\mathcal{C}/K_n$ の対象とし、
$\mathcal{F} \in \textit{Ab}((\mathcal{C}/K)_{total})$ とする。このとき
$$
H^p(U, \mathcal{F}) = H^p(U, \mathcal{F}_{n, i})
$$
である。ここで
\begin{enumerate}
% P10-E0203
\item 左辺では $U$ を $\mathcal{C}_{total}$ の対象とみなす。
\item 右辺では $\mathcal{F}_{n, i}$ は、節
\ref{spaces-simplicial-section-semi-representable} の分解
$\Sh(\mathcal{C}/K_n) = \prod \Sh(\mathcal{C}/U_{n, i})$ における、
$\mathcal{C}/K_n$ 上の層 $\mathcal{F}_n$ の第 $i$ 成分である。
\end{enumerate}
\end{lemma}

\begin{proof}
補題 \ref{spaces-simplicial-lemma-sanity-check} と節 \ref{spaces-simplicial-section-semi-representable} の
積分解から直ちに従う。
\end{proof}

\begin{remark}[対象上の場合の変種]
\label{spaces-simplicial-remark-augmentation-over-object}
$\mathcal{C}$ をサイトとし、$X \in \Ob(\mathcal{C})$ とする。
$X$ 上の半表現可能対象の圏
$\text{SR}(\mathcal{C}, X) = \text{SR}(\mathcal{C}/X)$ があることを思い出そう。
注意 \ref{spaces-simplicial-remark-semi-representable-over-object} を参照されたい。
上の議論をサイト $\mathcal{C}/X$ に適用できる。簡潔に述べると、
上の構成は次を与える。
\begin{enumerate}
% P10-E0204
\item $\text{SR}(\mathcal{C}, X)$ の単体的対象 $K$ に対するサイト
$(\mathcal{C}/K)_{total}$、
\item 局所化関手 $j_{total} : (\mathcal{C}/K)_{total} \to \mathcal{C}/X$、
\item 局所化関手 $j_n : \mathcal{C}/K_n \to \mathcal{C}/X$、
\item トポスの射
$a : \Sh((\mathcal{C}/K)_{total}) \to \Sh(\mathcal{C}/X)$、
\item トポスの射
$a_n : \Sh(\mathcal{C}/K_n) \to \Sh(\mathcal{C}/X)$、
\item $a^{-1}$ の左随伴である関手
$a^{Sh}_! : \Sh((\mathcal{C}/K)_{total}) \to \Sh(\mathcal{C}/X)$、
\item $a^{-1}$ の左随伴である関手
$a_! : \textit{Ab}((\mathcal{C}/K)_{total}) \to \textit{Ab}(\mathcal{C}/X)$。
\end{enumerate}
この節のすべての結果はこの設定でも成り立つ。これを証明するには、
至る所でサイト $\mathcal{C}$ を $\mathcal{C}/X$ に置き換えればよい。
\end{remark}

\begin{remark}[環付きの場合の変種]
\label{spaces-simplicial-remark-augmentation-ringed}
$\mathcal{C}$ をサイトとし、$\mathcal{O}_\mathcal{C}$ を環の層とする。
$\mathcal{C}$ の単体的半表現可能対象 $K$ が与えられたとき、
$\mathcal{O} = a^{-1}\mathcal{O}_\mathcal{C}$ とおく。ここで $a$ は補題
\ref{spaces-simplicial-lemma-augmentation-simplicial-semi-representable} および
\ref{spaces-simplicial-lemma-comparison} のものとする。注意
\ref{spaces-simplicial-remark-semi-representable-ringed} のように環の層も追跡すると、
上の構成は次を与える。
\begin{enumerate}
% P10-E0204
\item $\text{SR}(\mathcal{C})$ の単体的対象 $K$ に対する環付きサイト
$((\mathcal{C}/K)_{total}, \mathcal{O})$、
\item 環付きトポスの射
$a : (\Sh((\mathcal{C}/K)_{total}), \mathcal{O}) \to
(\Sh(\mathcal{C}), \mathcal{O}_\mathcal{C})$、
\item 環付きトポスの射
$a_n : (\Sh(\mathcal{C}/K_n), \mathcal{O}_n) \to
(\Sh(\mathcal{C}), \mathcal{O}_\mathcal{C})$、
\item $a^*$ の左随伴である関手
$a_! : \textit{Mod}(\mathcal{O}) \to \textit{Mod}(\mathcal{O}_\mathcal{C})$。
\end{enumerate}
関手 $a_!$ が存在する（ただし一般には完全でない）のは、
$a^{-1}\mathcal{O}_\mathcal{C} = \mathcal{O}$ であり、さらに
『サイト上の加群』補題 \ref{sites-modules-lemma-g-shriek-adjoint} を用いる
補題 \ref{spaces-simplicial-lemma-comparison} の証明を、『サイト上の加群』補題
\ref{sites-modules-lemma-lower-shriek-modules} を用いるものに
置き換えられるからである。注意
\ref{spaces-simplicial-remark-semi-representable-ringed} で論じたように、
$a_n^*$ の左随伴である完全関手
$a_{n!} : \textit{Mod}(\mathcal{O}_n) \to
\textit{Mod}(\mathcal{O}_\mathcal{C})$ が存在する。
したがって射 $a$ と $a_n$ は平坦である。注意
\ref{spaces-simplicial-remark-semi-representable-ringed} により、
$\varphi : [m] \to [n]$ に対する環付きトポスの射
$f_\varphi : (\Sh(\mathcal{C}/K_n), \mathcal{O}_n) \to
(\Sh(\mathcal{C}/K_m), \mathcal{O}_m)$ は平坦であり、$f_\varphi^*$ の
左随伴である完全関手
$f_{\varphi !} : \textit{Mod}(\mathcal{O}_n) \to \textit{Mod}(\mathcal{O}_m)$
が存在する。さらに、平坦な環付きトポスの射
$g_n : (\Sh(\mathcal{C}/K_n), \mathcal{O}_n) \to
(\Sh((\mathcal{C}/K)_{total}), \mathcal{O})$ に対して、$g_n^*$ の
左随伴である関手
$g_{n!} : \textit{Mod}(\mathcal{O}_n) \to
\textit{Mod}(\mathcal{O})$ は完全である。補題
\ref{spaces-simplicial-lemma-exactness-g-shriek-modules} を参照されたい。
\end{remark}

\begin{remark}[対象上の環付きの場合の変種]
\label{spaces-simplicial-remark-augmentation-ringed-over-object}
$\mathcal{C}$ をサイトとし、$\mathcal{O}_\mathcal{C}$ を環の層とする。
$X \in \Ob(\mathcal{C})$ とし、
$\mathcal{O}_X = \mathcal{O}_\mathcal{C}|_{\mathcal{C}/X}$ と書く。
注意 \ref{spaces-simplicial-remark-augmentation-over-object} と
\ref{spaces-simplicial-remark-augmentation-ringed} の構成を組み合わせて、次を得る。
\begin{enumerate}
% P10-E0204
\item $\text{SR}(\mathcal{C}, X)$ の単体的対象 $K$ に対する環付きサイト
$((\mathcal{C}/K)_{total}, \mathcal{O})$、
\item 環付きトポスの射
$a : (\Sh((\mathcal{C}/K)_{total}), \mathcal{O}) \to
(\Sh(\mathcal{C}/X), \mathcal{O}_X)$、
\item 環付きトポスの射
$a_n : (\Sh(\mathcal{C}/K_n), \mathcal{O}_n) \to
(\Sh(\mathcal{C}/X), \mathcal{O}_X)$、
\item $a^*$ の左随伴である関手
$a_! : \textit{Mod}(\mathcal{O}) \to \textit{Mod}(\mathcal{O}_X)$。
\end{enumerate}
もちろん、注意 \ref{spaces-simplicial-remark-augmentation-ringed} で述べた結果は
すべてこの設定でも成り立つ。
\end{remark}

\section{超被覆に対するコホモロジー的下降}
\label{spaces-simplicial-section-cohomological-descent-hypercoverings}

\noindent
$\mathcal{C}$ をサイトとする。この節では $\mathcal{C}$ が等化子と
ファイバー積をもつと仮定する。$K$ を『超被覆』定義
\ref{hypercovering-definition-hypercovering-variant} で定義された
超被覆とする。節 \ref{spaces-simplicial-section-simplicial-semi-representable} の増大
$$
a : \Sh((\mathcal{C}/K)_{total}) \longrightarrow \Sh(\mathcal{C})
$$
を調べる。

\begin{lemma}
\label{spaces-simplicial-lemma-hypercovering-descent-sheaves}
$\mathcal{C}$ を等化子とファイバー積をもつサイトとし、$K$ を超被覆とする。
このとき
\begin{enumerate}
\item $a^{-1} : \Sh(\mathcal{C}) \to \Sh((\mathcal{C}/K)_{total})$ は
充満忠実であり、その本質的像は集合のカルテジアン層である。
\item $a^{-1} : \textit{Ab}(\mathcal{C}) \to
\textit{Ab}((\mathcal{C}/K)_{total})$ は充満忠実であり、
その本質的像はアーベル群のカルテジアン層である。
\end{enumerate}
いずれの場合も $a_*$ が擬逆関手を与える。
\end{lemma}

\begin{proof}
アーベル層の場合は、関手 $a^{-1}$ が積と可換するので、
集合の層の場合から直ちに従う。以下では集合の層を扱う。
補題 \ref{spaces-simplicial-lemma-augmentation-cartesian-module} により、
$\Sh(\mathcal{C})$ の $\mathcal{F}$ に対して
$a^{-1}\mathcal{F}$ はカルテジアンである。随伴写像
$\mathcal{F} \to a_*a^{-1}\mathcal{F}$ が
% P10-E0205
$\Sh(\mathcal{C})$ の $\mathcal{F}$ に対して同型であること、および
$(\mathcal{C}/K)_{total}$ 上のカルテジアン層 $\mathcal{G}$ に対して
随伴写像 $a^{-1}a_*\mathcal{G} \to \mathcal{G}$ が同型であることを
示せば十分である。

\medskip\noindent
$\mathcal{F}$ を $\mathcal{C}$ 上の層とする。補題
\ref{spaces-simplicial-lemma-comparison} により、$a_*a^{-1}\mathcal{F}$ は二つの写像\linebreak
$a_{0, *}a_0^{-1}\mathcal{F} \to a_{1, *}a_1^{-1}\mathcal{F}$ の
等化子であることを思い出そう。補題 \ref{spaces-simplicial-lemma-push-pull-localization} により、
$$
a_{0, *}a_0^{-1}\mathcal{F} = \SheafHom(F(K_0)^\#, \mathcal{F})
\quad\text{および}\quad
a_{1, *}a_1^{-1}\mathcal{F} = \SheafHom(F(K_1)^\#, \mathcal{F})
$$
である。一方、超被覆の定義により、
$$
\xymatrix{
F(K_1)^\# \ar@<1ex>[r] \ar@<-1ex>[r] &
F(K_0)^\# \ar[r] & \Sh(\mathcal{C})\text{ の終対象 }*
}
$$
は集合の層における余等化子図式である。したがって、
$\SheafHom(-, \mathcal{F})$ が余等化子を等化子へ写すことを示せば十分であり、
これは『サイト』節 \ref{sites-section-glueing-sheaves} の構成から直ちに従う。

\medskip\noindent
$\mathcal{G}$ を $(\mathcal{C}/K)_{total}$ 上のカルテジアン層とする。
ある $\mathcal{C}$ 上の層 $\mathcal{F}$ に対して
$\mathcal{G} = a^{-1}\mathcal{F}$ であることを示す。そうすれば前段落の結果により
$a^{-1}a_*\mathcal{G} = a^{-1}a_*a^{-1}\mathcal{F} =
a^{-1}\mathcal{F} = \mathcal{G}$ となるので、証明は完了する。
$n \geq 0$ に対して $\mathcal{K}_n = F(K_n)^\#$ とおく。このとき、
$K$ が単体的半表現可能対象であることから層の写像
$$
\xymatrix{
\mathcal{K}_2
\ar@<1ex>[r]
\ar@<0ex>[r]
\ar@<-1ex>[r]
&
\mathcal{K}_1
\ar@<0.5ex>[r]
\ar@<-0.5ex>[r]
&
\mathcal{K}_0
}
$$
を得る。$K$ が超被覆であることは、
$$
\mathcal{K}_1 \to \mathcal{K}_0 \times \mathcal{K}_0
\quad\text{および}\quad
\mathcal{K}_2 \to
\left(\text{cosk}_1(
\xymatrix{
\mathcal{K}_1
\ar@<0.5ex>[r]
\ar@<-0.5ex>[r]
&
\mathcal{K}_0 \ar[l]
})\right)_2
$$
が層の全射であることを意味する。補題
\ref{spaces-simplicial-lemma-characterize-cartesian} で与えた
$(\mathcal{C}/K)_{total}$ 上のカルテジアン層の記述と、補題
\ref{spaces-simplicial-lemma-localize-compare} における $\Sh(\mathcal{C}/K_n)$ の記述を用いると、
問題は完全に次の言葉で定式化できることが分かる\footnote{正確な定式化が
何であるかは本質的でないが、明記しておく。問題は、
$\Sh(\mathcal{C})/\mathcal{K}_0$ の対象
$\mathcal{G}_0/\mathcal{K}_0$ と、$\mathcal{K}_1$ 上の同型
$$
\alpha :
\mathcal{G}_0 \times_{\mathcal{K}_0, \mathcal{K}(\delta^1_1)} \mathcal{K}_1 \to
\mathcal{G}_0 \times_{\mathcal{K}_0, \mathcal{K}(\delta^1_0)} \mathcal{K}_1
$$
であって $\Sh(\mathcal{C})/\mathcal{K}_2$ における余サイクル条件を
満たすものが与えられたとき、$\Sh(\mathcal{C})$ の $\mathcal{F}$ と、
$\alpha$ と両立する $\mathcal{K}_0$ 上の同型
$\mathcal{F} \times \mathcal{K}_0 \to \mathcal{G}_0$ が存在することを
示すことである。}。
\begin{enumerate}
\item トポス $\Sh(\mathcal{C})$、および
\item 各項が $\mathcal{K}_n$ である $\Sh(\mathcal{C})$ の単体的対象
$\mathcal{K}$。
\end{enumerate}
そこで、『サイト』補題 \ref{sites-lemma-topos-good-site} のように
$\mathcal{C}$ を別のサイト $\mathcal{C}'$ に置き換えることで、
$\mathcal{C}$ はすべての有限極限をもち、$\mathcal{C}$ 上の位相は
劣標準的であり、$\mathcal{C}$ の射の族 $\{V_j \to V\}$ が被覆であることと
$\coprod h_{V_j} \to V$ が全射であることが同値であり、さらに
$\mathcal{K}_n = h_{U_n}$ が単体的層として成り立つような
$\mathcal{C}$ の単体的対象 $U$ が存在すると仮定してよい。
同値を逆向きにたどることにより、すべての $n$ に対して
$K_n = \{U_n\}$ と仮定してよい。

\medskip\noindent
$X$ を $\mathcal{C}$ の終対象とする。このとき
$\{U_0 \to X\}$ は被覆、$\{U_1 \to U_0 \times U_0\}$ は被覆、かつ
$\{U_2 \to (\text{cosk}_1 \text{sk}_1 U)_2\}$ は被覆である。
『単体的対象』定義 \ref{simplicial-definition-face-degeneracy} と同様に、
$\delta^n_i$ と $\sigma^n_j$ に対応する射を
$d^n_i : U_n \to U_{n - 1}$ および $s^n_j : U_n \to U_{n + 1}$ と書く。
記法の濫用により、$\mathcal{C}$ の射 $c : V \to W$ が与えられたとき、
トポスの射 $c : \Sh(\mathcal{C}/V) \to \Sh(\mathcal{C}/W)$ も
同じ文字で表す。いま $\mathcal{G}$ は、$\mathcal{C}/U_0$ 上の層
$\mathcal{G}_0$ と、補題 \ref{spaces-simplicial-lemma-characterize-cartesian} で定式化された
$\mathcal{C}/U_2$ 上の余サイクル条件を満たす同型
$\alpha : (d^1_1)^{-1}\mathcal{G}_0 \to (d^1_0)^{-1}\mathcal{G}_0$
により与えられる。$\{U_2 \to (\text{cosk}_1 \text{sk}_1 U)_2\}$ は
被覆なので、対応する層の引き戻し関手は忠実である
（細部は省略する）。したがって $U$ を $\text{cosk}_1 \text{sk}_1 U$ で
置き換えてよい。これは $U_2$ を $(\text{cosk}_1 \text{sk}_1 U)_2$ で
置き換え、$U_1$ と $U_0$ を変えないからである。このとき
$$
(d^2_0, d^2_1, d^2_2) : U_2 \to U_1 \times U_1 \times U_1
$$
は単射であり、
% P10-E0206
$T$-値点上のその像は『単体的対象』補題
\ref{simplicial-lemma-work-out} で記述される。とくに、次の可換図式に
収まる射 $c$ が存在する。
$$
\xymatrix{
U_1 \times_{(d^1_1, d^1_0), U_0 \times U_0, (d^1_1, d^1_0)} U_1
\ar[d] \ar[rr]_c & & U_2 \ar[d] \\
U_1 \times U_1
\ar[rr]^{(\text{pr}_1, \text{pr}_2, s^0_0 \circ d^1_1 \circ \text{pr}_1)} & &
U_1 \times U_1 \times U_1
}
$$
これは、他方の経路を通ると $U_2$ の点が定まるからである。
% P10-E0207
$U_2$ 上の $\alpha$ の余サイクル条件を引き戻すと、
$U_1 \times_{(d^1_1, d^1_0), U_0 \times U_0, (d^1_1, d^1_0)} U_1$
への二つの射影による $\alpha$ の引き戻しが一致し、かつ
$s^0_0 \circ d^1_1 \circ \text{pr}_1$ による $\alpha$ の引き戻しが
恒等写像であるという条件になる（すなわち、$s^0_0$ による
$\alpha$ の引き戻しは恒等写像である）。『サイト』補題
\ref{sites-lemma-glue-maps} により、これは $\alpha$ が
$\mathcal{C}/U_0 \times U_0$ 上の層の同型
$$
\alpha' : \text{pr}_1^{-1}\mathcal{G}_0 \to \text{pr}_2^{-1}\mathcal{G}_0
$$
から生じることを意味する。最後に、上で与えた $U_2$ の記述と
$U_1 \to U_0 \times U_0$ の全射性を用いれば容易に分かるように、
射 $U_2 \to U_0 \times U_0 \times U_0$ は付随する層上で全射である。
したがって $\alpha'$ は $U_0 \times U_0 \times U_0$ 上で
余サイクル条件を満たす。被覆 $\{U_0 \to X\}$ に
『サイト』補題 \ref{sites-lemma-mapping-property-glue} を適用すれば
証明が完了する。
\end{proof}

\begin{lemma}
\label{spaces-simplicial-lemma-hypercovering-cech-complex}
$\mathcal{C}$ を等化子とファイバー積をもつサイトとし、$K$ を超被覆とする。
補題 \ref{spaces-simplicial-lemma-augmentation-cech-complex} の $a^{-1}\mathcal{F}$ に付随する
{\v C}ech 複体
$$
a_{0, *}a_0^{-1}\mathcal{F} \to a_{1, *}a_1^{-1}\mathcal{F} \to
a_{2, *}a_2^{-1}\mathcal{F} \to \ldots
$$
は複体 $\SheafHom(s(\mathbf{Z}_{F(K)}^\#), \mathcal{F})$ に等しい。
ここで $s(\mathbf{Z}_{F(K)}^\#)$ は『超被覆』定義
\ref{hypercovering-definition-homology} のものとする。
\end{lemma}

\begin{proof}
補題 \ref{spaces-simplicial-lemma-push-pull-localization} により
$$
a_{n, *}a_n^{-1}\mathcal{F} = \SheafHom'(F(K_n)^\#, \mathcal{F})
$$
である。ここで $\SheafHom'$ は『サイト』節
\ref{sites-section-glueing-sheaves} のものとする。補題
\ref{spaces-simplicial-lemma-augmentation-cech-complex} の複体の境界写像は単体的構造から生じる。
したがって複体の等しさは、$\Sh(\mathcal{C})$ の $\mathcal{G}$ に対する
標準的同一視
$\SheafHom'(\mathcal{G}, \mathcal{F}) =
\SheafHom(\mathbf{Z}_\mathcal{G}, \mathcal{F})$ から従う。
\end{proof}

\begin{lemma}
\label{spaces-simplicial-lemma-hypercovering-descent-bounded-abelian}
$\mathcal{C}$ を等化子とファイバー積をもつサイトとし、$K$ を超被覆とする。
$E \in D(\mathcal{C})$ に対して、写像
$$
E \longrightarrow Ra_*a^{-1}E
$$
は同型である。
\end{lemma}

\begin{proof}
まず、$\mathcal{I}$ を $\mathcal{C}$ 上の入射的アーベル層とする。
補題 \ref{spaces-simplicial-lemma-augmentation-spectral-sequence} の層
$a^{-1}\mathcal{I}$ に対するスペクトル系列は、補題
\ref{spaces-simplicial-lemma-localize-injective} により
$(a^{-1}\mathcal{I})_p = a_p^{-1}\mathcal{I}$ が入射的なので退化する。
したがって複体
$$
a_{0, *}a_0^{-1}\mathcal{I} \to
a_{1, *}a_1^{-1}\mathcal{I} \to
a_{2, *}a_2^{-1}\mathcal{I} \to\ldots
$$
が $Ra_*a^{-1}\mathcal{I}$ を計算する。補題
\ref{spaces-simplicial-lemma-hypercovering-cech-complex} により、これは複体
$\SheafHom(s(\mathbf{Z}_{F(K)}^\#), \mathcal{I})$ に等しい。
$K$ は超被覆なので、『超被覆』補題
\ref{hypercovering-lemma-acyclic-hypercover-sheaves} を単体的前層
$F(K)$ に適用すると、$s(\mathbf{Z}_{F(K)}^\#)$ は次数 $> 0$ で
完全である。$\mathcal{I}$ は入射的なので、関手
$\SheafHom(-, \mathcal{I})$ は完全であり、
$\SheafHom(s(\mathbf{Z}_{F(K)}^\#), \mathcal{I})$ は正次数で
完全であると結論できる。したがって $p > 0$ に対して
$R^pa_*a^{-1}\mathcal{I} = 0$ である。一方、補題
\ref{spaces-simplicial-lemma-hypercovering-descent-sheaves} により
$\mathcal{I} = a_*a^{-1}\mathcal{I}$ である。

\medskip\noindent
有界下の場合。$E \in D^+(\mathcal{C})$ とする。$E$ を表す
入射的対象からなる有界下複体 $\mathcal{I}^\bullet$ を選ぶ。
第一段落の結果と Leray の非輪状性補題
（『導来圏』補題 \ref{derived-lemma-leray-acyclicity}）により、
$Ra_*a^{-1}\mathcal{I}^\bullet$ は複体
$a_*a^{-1}\mathcal{I}^\bullet = \mathcal{I}^\bullet$ によって計算される。
したがって、この場合に補題が成り立つ。

\medskip\noindent
非有界の場合。収束の問題を避けるため二重複体による明示的表示を
用いることを除けば上と同じ議論なので、読者にはここを飛ばすよう勧める。
$E \in D(\mathcal{C})$ とする。写像 $E \to Ra_*a^{-1}E$ が同型であることを
示すには、$\mathcal{C}$ のすべての対象 $U$ に対して
$$
R\Gamma(U, E) = R\Gamma(U, Ra_*a^{-1}E)
$$
であることを示せば十分である。両辺を計算し、写像
$E \to Ra_*a^{-1}E$ が同型を誘導することを示す。$E$ を表す
K-入射的複体 $\mathcal{I}^\bullet$ を選ぶ。また、
$(\mathcal{C}/K)_{total}$ 上のある K-入射的複体 $\mathcal{J}^\bullet$ に対する
擬同型 $a^{-1}\mathcal{I}^\bullet \to \mathcal{J}^\bullet$ を選ぶ。
次が成り立つ。
$$
R\Gamma(U, E) = R\Hom(\mathbf{Z}_U^\#, E)
$$
また、
$$
R\Gamma(U, Ra_*a^{-1}E) = R\Hom(\mathbf{Z}_U^\#, Ra_*a^{-1}E) =
R\Hom(a^{-1}\mathbf{Z}_U^\#, a^{-1}E)
$$
である。補題 \ref{spaces-simplicial-lemma-simplicial-resolution-augmentation} により、擬同型
$$
\Big(\ldots \to
g_{2!}(a_2^{-1}\mathbf{Z}_U^\#) \to
g_{1!}(a_1^{-1}\mathbf{Z}_U^\#) \to
g_{0!}(a_0^{-1}\mathbf{Z}_U^\#)\Big)
\longrightarrow
a^{-1}\mathbf{Z}_U^\#
$$
がある。したがって $R\Hom(a^{-1}\mathbf{Z}_U^\#, a^{-1}E)$ は
$$
R\Gamma((\mathcal{C}/K)_{total},
R\SheafHom(
\ldots \to
g_{2!}(a_2^{-1}\mathbf{Z}_U^\#) \to
g_{1!}(a_1^{-1}\mathbf{Z}_U^\#) \to
g_{0!}(a_0^{-1}\mathbf{Z}_U^\#),
\mathcal{J}^\bullet))
$$
に等しい。『サイト上のコホモロジー』節
\ref{sites-cohomology-section-internal-hom} の構成と
$\mathcal{J}^\bullet$ の K-入射性により、これは項が
$$
\prod\nolimits_{p + q = n}
\Hom(g_{p!}(a_p^{-1}\mathbf{Z}_U^\#), \mathcal{J}^q) =
\prod\nolimits_{p + q = n}
\Hom(a_p^{-1}\mathbf{Z}_U^\#, g_p^{-1}\mathcal{J}^q)
$$
であるアーベル群の複体によって表される。詳しくは
『サイト上のコホモロジー』補題
\ref{sites-cohomology-lemma-RHom-into-K-injective} および
\ref{sites-cohomology-lemma-section-RHom-over-U} を参照されたい。
したがって $R\Gamma(U, Ra_*a^{-1}E)$ は積全化複体
$\text{Tot}_\pi(B^{\bullet, \bullet})$ により計算される。ここで
$B^{p, q} = \Hom(a_p^{-1}\mathbf{Z}_U^\#, g_p^{-1}\mathcal{J}^q)$ である。
他方についても同様に論じる。まず、『超被覆』補題
\ref{hypercovering-lemma-acyclic-hypercover-sheaves} により、
$$
s(\mathbf{Z}_{F(K)}^\#) \longrightarrow \mathbf{Z}
$$
は $\mathcal{C}$ 上の複体の擬同型である。
$\mathbf{Z}_U^\#$ は $\mathbf{Z}$-加群の平坦層なので、
$$
s(\mathbf{Z}_{F(K)}^\#) \otimes_\mathbf{Z} \mathbf{Z}_U^\#
\longrightarrow
\mathbf{Z}_U^\#
$$
は擬同型である。したがって $R\Hom(\mathbf{Z}_U^\#, E)$ は
$$
R\Gamma(\mathcal{C}, R\SheafHom(
s(\mathbf{Z}_{F(K)}^\#) \otimes_\mathbf{Z} \mathbf{Z}_U^\#,
\mathcal{I}^\bullet))
$$
に等しい。$R\SheafHom$ の構成と $\mathcal{I}^\bullet$ の
K-入射性により、これは項が
$$
\prod\nolimits_{p + q = n}
\Hom(\mathbf{Z}^\#_{K_p} \otimes_\mathbf{Z} \mathbf{Z}_U^\#, \mathcal{I}^q)
=
\prod\nolimits_{p + q = n}
\Hom(a_p^{-1}\mathbf{Z}_U^\#, a_p^{-1}\mathcal{I}^q)
$$
であるアーベル群の複体によって表される。項の等しさは、
『サイト上の加群』注意 \ref{sites-modules-remark-j-shriek-tensor} により
$\mathbf{Z}^\#_{K_p} \otimes_\mathbf{Z} \mathbf{Z}_U^\# =
a_{p!}a_p^{-1}\mathbf{Z}_U^\#$ であることから従う。
したがって $R\Gamma(U, E)$ は積全化複体
$\text{Tot}_\pi(A^{\bullet, \bullet})$ により計算される。ここで
$A^{p, q} = \Hom(a_p^{-1}\mathbf{Z}_U^\#, a_p^{-1}\mathcal{I}^q)$ である。

\medskip\noindent
$\mathcal{I}^\bullet$ は K-入射的なので、補題
\ref{spaces-simplicial-lemma-localize-injective} により $a_p^{-1}\mathcal{I}^\bullet$ は
K-入射的である。$\mathcal{J}^\bullet$ は K-入射的なので、補題
\ref{spaces-simplicial-lemma-restriction-injective-to-component-site} により
$g_p^{-1}\mathcal{J}^\bullet$ は K-入射的である。両者は対象
$a_p^{-1}E$ を表す。したがって、すべての $p \geq 0$ に対して、
$g_p^{-1}$ を所与の写像
$a^{-1}\mathcal{I}^\bullet \to \mathcal{J}^\bullet$ に適用して誘導される
複体の写像
$$
A^{p, \bullet} = \Hom(a_p^{-1}\mathbf{Z}_U^\#, a_p^{-1}\mathcal{I}^\bullet)
\longrightarrow
\Hom(a_p^{-1}\mathbf{Z}_U^\#, g_p^{-1}\mathcal{J}^\bullet) = B^{p, \bullet}
$$
は擬同型である。実際、これらの複体はいずれも
% P10-E0208
$$
R\Hom(a_p^{-1}\mathbf{Z}_U^\#, a_p^{-1}E)
$$
を計算する。『代数学続論』補題
\ref{more-algebra-lemma-prod-qis-gives-qis} により、可換図式
$$
\xymatrix{
R\Gamma(U, E) \ar[r] \ar[d] &
\text{Tot}_\pi(A^{\bullet, \bullet}) \ar[d] \\
R\Gamma(U, Ra_*a^{-1}E) \ar[r] &
\text{Tot}_\pi(B^{\bullet, \bullet})
}
$$
の右の垂直射は擬同型である。上で水平射が擬同型であることを見たので、
左の垂直射も擬同型である。
\end{proof}

\begin{lemma}
\label{spaces-simplicial-lemma-compare-cohomology-hypercovering}
$\mathcal{C}$ を等化子とファイバー積をもつサイトとし、$K$ を超被覆とする。
このとき $E \in D(\mathcal{C})$ に対して標準同型
$$
R\Gamma(\mathcal{C}, E) =
R\Gamma((\mathcal{C}/K)_{total}, a^{-1}E)
$$
がある。
\end{lemma}

\begin{proof}
補題 \ref{spaces-simplicial-lemma-hypercovering-descent-bounded-abelian} から従う。実際、
『サイト上のコホモロジー』注意
\ref{sites-cohomology-remark-before-Leray} により
$R\Gamma((\mathcal{C}/K)_{total}, -) =
R\Gamma(\mathcal{C}, -) \circ Ra_*$ である。
\end{proof}

\begin{lemma}
\label{spaces-simplicial-lemma-hypercovering-equivalence-bounded}
$\mathcal{C}$ を等化子とファイバー積をもつサイトとし、$K$ を超被覆とする。
$\mathcal{A} \subset \textit{Ab}((\mathcal{C}/K)_{total})$ を
カルテジアン・アーベル層からなる弱 Serre 部分圏とする。このとき関手
$a^{-1}$ は同値
$$
D^+(\mathcal{C}) \longrightarrow D_\mathcal{A}^+((\mathcal{C}/K)_{total})
$$
を定め、その擬逆は $Ra_*$ である。
\end{lemma}

\begin{proof}
補題 \ref{spaces-simplicial-lemma-Serre-subcat-cartesian-modules} により、
$\mathcal{A}$ は弱 Serre 部分圏である。この同値は、これまでに得た結果の
形式的帰結である。補題 \ref{spaces-simplicial-lemma-hypercovering-descent-sheaves}、
\ref{spaces-simplicial-lemma-hypercovering-descent-bounded-abelian}、および
『サイト上のコホモロジー』補題
\ref{sites-cohomology-lemma-equivalence-bounded} を用いればよい。
\end{proof}

\noindent
読者には次の注意を飛ばすよう勧める。

\begin{remark}
\label{spaces-simplicial-remark-compare-cohomology-hypercovering-presheaf}
$\mathcal{C}$ をサイトとし、$\mathcal{G}$ を $\mathcal{C}$ 上の
集合の前層とする。$\mathcal{C}$ が等化子とファイバー積をもつとき、
『超被覆』定義 \ref{hypercovering-definition-hypercovering-variant} で
$\mathcal{G}$ の超被覆という概念を定義した。この設定でも、この節の
すべての結果に対応する有効な主張がある。これを見るため、
$\mathcal{G}$ における $\mathcal{C}$ の局所化
$\mathcal{C}/\mathcal{G}$ を、『サイト』補題
\ref{sites-lemma-localize-topos-site} とまったく同様に定義する
（そこでは層についてのみ述べられているが、トポス
$\Sh(\mathcal{C}/\mathcal{G})$ は、層 $\mathcal{G}^\#$ におけるトポス
$\Sh(\mathcal{C})$ の局所化に等しい）。サイト
$\mathcal{C}/\mathcal{G}$ がファイバー積と等化子をもち、
$\mathcal{C}$ における $\mathcal{G}$ の超被覆がサイト
$\mathcal{C}/\mathcal{G}$ の超被覆と同じであることは容易に示せる。
したがって、上の超被覆に関する補題でサイト $\mathcal{C}$ を
$\mathcal{C}/\mathcal{G}$ に置き換えると、$\mathcal{G}$ の超被覆に関する
対応する結果の証明を得る。例として、$\mathcal{G}$ の超被覆 $K$ に対して
$$
R\Gamma(\mathcal{C}/\mathcal{G}, E) =
R\Gamma((\mathcal{C}/K)_{total}, a^{-1}E)
$$
が $E \in D^+(\mathcal{C}/\mathcal{G})$ に対して成り立つ。ここで
$a : \Sh((\mathcal{C}/K)_{total}) \to \Sh(\mathcal{C}/\mathcal{G})$ は
標準的な増大である。これは補題
\ref{spaces-simplicial-lemma-compare-cohomology-hypercovering} である。
関手 $R\Gamma(\mathcal{G}, -) : D(\mathcal{C}) \to D(\textit{Ab})$ を、
『超被覆』節 \ref{hypercovering-section-hypercoverings-verdier} および
『サイト上のコホモロジー』節 \ref{sites-cohomology-section-limp} で論じた関手
$H^0(\mathcal{G}, -) = H^0(\mathcal{G}^\#, -)$ の導来関手として定義する。
『サイト上のコホモロジー』補題
\ref{sites-cohomology-lemma-cohomology-of-open} の局所化関手
% P10-E0209
$j : \mathcal{C}/\mathcal{G} \to \mathcal{C}$ に対する類似により、
$$
R\Gamma(\mathcal{G}, E) = R\Gamma(\mathcal{C}/\mathcal{G}, j^{-1}E)
$$
である。すべてを合わせると、$E \in D^+(\mathcal{C})$ に対して
$$
R\Gamma(\mathcal{G}, E) =
R\Gamma((\mathcal{C}/K)_{total}, a^{-1}j^{-1}E) =
R\Gamma((\mathcal{C}/K)_{total}, g^{-1}E)
$$
を得る。ここで $g : \Sh((\mathcal{C}/K)_{total}) \to \Sh(\mathcal{C})$ は
$a$ と $j$ の合成である。
\end{remark}

\section{超被覆に対するコホモロジー的下降：加群}
\label{spaces-simplicial-section-cohomological-descent-hypercoverings-modules}

\noindent
$\mathcal{C}$ をサイトとし、$\mathcal{O}_\mathcal{C}$ を環の層とする。
$\mathcal{C}$ が等化子とファイバー積をもつと仮定し、$K$ を
『超被覆』定義 \ref{hypercovering-definition-hypercovering-variant} で
定義された超被覆とする。注意 \ref{spaces-simplicial-remark-augmentation-ringed} の増大
$$
a :
(\Sh((\mathcal{C}/K)_{total}), \mathcal{O})
\longrightarrow
(\Sh(\mathcal{C}), \mathcal{O}_\mathcal{C})
$$
に対するコホモロジー的下降を調べる。

\begin{lemma}
\label{spaces-simplicial-lemma-hypercovering-descent-modules}
$\mathcal{C}$ を等化子とファイバー積をもつサイトとし、
$\mathcal{O}_\mathcal{C}$ を環の層、$K$ を超被覆とする。
上の記法のもとで
$$
a^* : \textit{Mod}(\mathcal{O}_\mathcal{C}) \to \textit{Mod}(\mathcal{O})
$$
は充満忠実であり、その本質的像はカルテジアン
$\mathcal{O}$-加群である。関手 $a_*$ が擬逆を与える。
\end{lemma}

\begin{proof}
$a^{-1}\mathcal{O}_\mathcal{C} = \mathcal{O}$ なので
$a^* = a^{-1}$ である。したがって補題
\ref{spaces-simplicial-lemma-hypercovering-descent-sheaves} から直ちに従う。
\end{proof}

\begin{lemma}
\label{spaces-simplicial-lemma-hypercovering-descent-bounded-modules}
$\mathcal{C}$ を等化子とファイバー積をもつサイトとし、
$\mathcal{O}_\mathcal{C}$ を環の層、$K$ を超被覆とする。
$E \in D(\mathcal{O}_\mathcal{C})$ に対して、写像
$$
E \longrightarrow Ra_*La^*E
$$
は同型である。
\end{lemma}

\begin{proof}
$a^{-1}\mathcal{O}_\mathcal{C} = \mathcal{O}$ なので
$La^* = a^* = a^{-1}$ である。さらに、$Ra_*$ はアーベル層上の
$Ra_*$ と一致する。『サイト上のコホモロジー』補題
\ref{sites-cohomology-lemma-modules-abelian-unbounded} を参照されたい。
したがって補題
\ref{spaces-simplicial-lemma-hypercovering-descent-bounded-abelian} から直ちに従う。
\end{proof}

\begin{lemma}
\label{spaces-simplicial-lemma-compare-cohomology-hypercovering-modules}
$\mathcal{C}$ を等化子とファイバー積をもつサイトとし、
$\mathcal{O}_\mathcal{C}$ を環の層、$K$ を超被覆とする。
このとき $E \in D(\mathcal{O}_\mathcal{C})$ に対して標準同型
$$
R\Gamma(\mathcal{C}, E) =
R\Gamma((\mathcal{C}/K)_{total}, La^*E)
$$
がある。
\end{lemma}

\begin{proof}
補題 \ref{spaces-simplicial-lemma-hypercovering-descent-bounded-modules} から従う。実際、
『サイト上のコホモロジー』注意
\ref{sites-cohomology-remark-before-Leray} または
『サイト上のコホモロジー』補題
\ref{sites-cohomology-lemma-Leray-unbounded} により
$R\Gamma((\mathcal{C}/K)_{total}, -) =
R\Gamma(\mathcal{C}, -) \circ Ra_*$ である。
\end{proof}

\begin{lemma}
\label{spaces-simplicial-lemma-hypercovering-equivalence-bounded-modules}
$\mathcal{C}$ を等化子とファイバー積をもつサイトとし、
$\mathcal{O}_\mathcal{C}$ を環の層、$K$ を超被覆とする。
$\mathcal{A} \subset \textit{Mod}(\mathcal{O})$ をカルテジアン
$\mathcal{O}$-加群からなる弱 Serre 部分圏とする。このとき関手
$La^*$ は同値
$$
D^+(\mathcal{O}_\mathcal{C}) \longrightarrow D_\mathcal{A}^+(\mathcal{O})
$$
を定め、その擬逆は $Ra_*$ である。
\end{lemma}

\begin{proof}
補題 \ref{spaces-simplicial-lemma-Serre-subcat-cartesian-modules} により
$\mathcal{A}$ は弱 Serre 部分圏である（必要な仮定が成り立つことは注意
\ref{spaces-simplicial-remark-augmentation-ringed} の議論による）。この同値は、
これまでに得た結果の形式的帰結である。補題
\ref{spaces-simplicial-lemma-hypercovering-descent-modules}、
\ref{spaces-simplicial-lemma-hypercovering-descent-bounded-modules}、および
『サイト上のコホモロジー』補題
\ref{sites-cohomology-lemma-equivalence-bounded} を用いればよい。
\end{proof}

\section{対象の超被覆に対するコホモロジー的下降}
\label{spaces-simplicial-section-cohomological-descent-hypercoverings-X}

\noindent
この節では $\mathcal{C}$ がファイバー積をもち、
$X \in \Ob(\mathcal{C})$ であると仮定する。$K$ を
『超被覆』定義 \ref{hypercovering-definition-hypercovering} で定義された
$X$ の超被覆とする。注意 \ref{spaces-simplicial-remark-augmentation-over-object} の増大
$$
a : \Sh((\mathcal{C}/K)_{total}) \longrightarrow \Sh(\mathcal{C}/X)
$$
を調べる。$\mathcal{C}/X$ は等化子とファイバー積をもつサイトであり、
$K$ はサイト $\mathcal{C}/X$ の超被覆であることに注意する\footnote{逆は
成り立たないことがある。すなわち、$K$ が
$\text{SR}(\mathcal{C}, X) = \text{SR}(\mathcal{C}/X)$ の単体的対象で、
『超被覆』定義 \ref{hypercovering-definition-hypercovering-variant} の意味で
サイト $\mathcal{C}/X$ の超被覆を定めても、$K$ が $X$ の超被覆を
定めるとは限らない。たとえば $K_0 = \{U_{0, i}\}_{i \in I_0}$ ならば、
後者の条件は $\{U_{0, i} \to X\}$ が $\mathcal{C}$ の被覆であることを
保証するが、前者の条件が要求するのは層の写像
$\coprod h_{U_{0, i}}^\# \to h_X^\#$ が全射であることだけである。}。
これは『超被覆』補題 \ref{hypercovering-lemma-hypercovering-F} による。
したがって、節 \ref{spaces-simplicial-section-cohomological-descent-hypercoverings} で
超被覆について証明したすべての結果には、この節の状況における
直接の類似がある。

\begin{lemma}
\label{spaces-simplicial-lemma-hypercovering-X-descent-sheaves}
$\mathcal{C}$ をファイバー積をもつサイトとし、
$X \in \Ob(\mathcal{C})$ とする。$K$ を $X$ の超被覆とする。このとき
\begin{enumerate}
\item $a^{-1} : \Sh(\mathcal{C}/X) \to \Sh((\mathcal{C}/K)_{total})$ は
充満忠実であり、その本質的像は集合のカルテジアン層である。
\item $a^{-1} : \textit{Ab}(\mathcal{C}/X) \to
\textit{Ab}((\mathcal{C}/K)_{total})$ は充満忠実であり、
その本質的像はアーベル群のカルテジアン層である。
\end{enumerate}
いずれの場合も $a_*$ が擬逆関手を与える。
\end{lemma}

\begin{proof}
注意 \ref{spaces-simplicial-remark-semi-representable-over-object}、
\ref{spaces-simplicial-remark-augmentation-over-object}、およびこの節の導入部の議論を介して、
補題 \ref{spaces-simplicial-lemma-hypercovering-descent-sheaves} から従う。
\end{proof}

\begin{lemma}
\label{spaces-simplicial-lemma-hypercovering-X-descent-bounded-abelian}
% P10-E0210
$\mathcal{C}$ をファイバー積をもつサイトとし、
$X \in \Ob(\mathcal{C})$ とする。$K$ を $X$ の超被覆とする。
$E \in D(\mathcal{C}/X)$ に対して、写像
$$
E \longrightarrow Ra_*a^{-1}E
$$
は同型である。
\end{lemma}

\begin{proof}
注意 \ref{spaces-simplicial-remark-semi-representable-over-object}、
\ref{spaces-simplicial-remark-augmentation-over-object}、およびこの節の導入部の議論を介して、
補題 \ref{spaces-simplicial-lemma-hypercovering-descent-bounded-abelian} から従う。
\end{proof}

\begin{lemma}
\label{spaces-simplicial-lemma-compare-cohomology-hypercovering-X}
$\mathcal{C}$ をファイバー積をもつサイトとし、
$X \in \Ob(\mathcal{C})$ とする。$K$ を $X$ の超被覆とする。
このとき $E \in D(\mathcal{C}/X)$ に対して標準同型
$$
R\Gamma(X, E) = R\Gamma((\mathcal{C}/K)_{total}, a^{-1}E)
$$
がある。
\end{lemma}

\begin{proof}
注意 \ref{spaces-simplicial-remark-semi-representable-over-object} および
\ref{spaces-simplicial-remark-augmentation-over-object} を介して、補題
\ref{spaces-simplicial-lemma-compare-cohomology-hypercovering} から従う。
\end{proof}

\begin{lemma}
\label{spaces-simplicial-lemma-hypercovering-X-equivalence-bounded}
$\mathcal{C}$ をファイバー積をもつサイトとし、
$X \in \Ob(\mathcal{C})$ とする。$K$ を $X$ の超被覆とする。
$\mathcal{A} \subset \textit{Ab}((\mathcal{C}/K)_{total})$ を
カルテジアン・アーベル層からなる弱 Serre 部分圏とする。このとき関手
$a^{-1}$ は同値
$$
D^+(\mathcal{C}/X) \longrightarrow D_\mathcal{A}^+((\mathcal{C}/K)_{total})
$$
を定め、その擬逆は $Ra_*$ である。
\end{lemma}

\begin{proof}
注意 \ref{spaces-simplicial-remark-semi-representable-over-object} および
\ref{spaces-simplicial-remark-augmentation-over-object} を介して、補題
\ref{spaces-simplicial-lemma-hypercovering-equivalence-bounded} から従う。
\end{proof}

\section{対象の超被覆に対するコホモロジー的下降：加群}
\label{spaces-simplicial-section-cohomological-descent-hypercoverings-X-modules}

\noindent
この節では $\mathcal{C}$ がファイバー積をもち、
$X \in \Ob(\mathcal{C})$ であると仮定する。$K$ を
『超被覆』定義 \ref{hypercovering-definition-hypercovering} で定義された
$X$ の超被覆とする。$\mathcal{O}_\mathcal{C}$ を
$\mathcal{C}$ 上の環の層とし、
$\mathcal{O}_X = \mathcal{O}_\mathcal{C}|_{\mathcal{C}/X}$ とおく。
注意 \ref{spaces-simplicial-remark-augmentation-ringed-over-object} の増大
$$
a :
(\Sh((\mathcal{C}/K)_{total}), \mathcal{O})
\longrightarrow
(\Sh(\mathcal{C}/X), \mathcal{O}_X)
$$
を調べる。$\mathcal{C}/X$ は等化子とファイバー積をもつサイトであり、
$K$ はサイト $\mathcal{C}/X$ の超被覆である。したがって、この節の結果は
節 \ref{spaces-simplicial-section-cohomological-descent-hypercoverings-modules} の
対応する結果から直ちに従う。

\begin{lemma}
\label{spaces-simplicial-lemma-hypercovering-X-descent-modules}
$\mathcal{C}$ をファイバー積をもつサイトとし、
$X \in \Ob(\mathcal{C})$ とする。$\mathcal{O}_\mathcal{C}$ を環の層、
$K$ を $X$ の超被覆とする。上の記法のもとで
$$
a^* : \textit{Mod}(\mathcal{O}_X) \to \textit{Mod}(\mathcal{O})
$$
は充満忠実であり、その本質的像はカルテジアン
$\mathcal{O}$-加群である。関手 $a_*$ が擬逆を与える。
\end{lemma}

\begin{proof}
注意 \ref{spaces-simplicial-remark-semi-representable-ringed-over-object}、
\ref{spaces-simplicial-remark-augmentation-ringed-over-object}、およびこの節の導入部の議論を介して、
補題 \ref{spaces-simplicial-lemma-hypercovering-descent-modules} から従う。
\end{proof}

\begin{lemma}
\label{spaces-simplicial-lemma-hypercovering-X-descent-bounded-modules}
$\mathcal{C}$ をファイバー積をもつサイトとし、
$X \in \Ob(\mathcal{C})$ とする。$\mathcal{O}_\mathcal{C}$ を環の層、
$K$ を $X$ の超被覆とする。$E \in D(\mathcal{O}_X)$ に対して、写像
$$
E \longrightarrow Ra_*La^*E
$$
は同型である。
\end{lemma}

\begin{proof}
注意 \ref{spaces-simplicial-remark-semi-representable-ringed-over-object}、
\ref{spaces-simplicial-remark-augmentation-ringed-over-object}、およびこの節の導入部の議論を介して、
補題 \ref{spaces-simplicial-lemma-hypercovering-descent-bounded-modules} から従う。
\end{proof}

\begin{lemma}
\label{spaces-simplicial-lemma-compare-cohomology-hypercovering-X-modules}
$\mathcal{C}$ をファイバー積をもつサイトとし、
$X \in \Ob(\mathcal{C})$ とする。$\mathcal{O}_\mathcal{C}$ を環の層、
$K$ を $X$ の超被覆とする。このとき標準同型
$$
R\Gamma(X, E) = R\Gamma((\mathcal{C}/K)_{total}, La^*E)
$$
が、
% P10-E0211
$E \in D(\mathcal{O}_\mathcal{C})$ に対してある。
\end{lemma}

\begin{proof}
注意 \ref{spaces-simplicial-remark-semi-representable-ringed-over-object}、
\ref{spaces-simplicial-remark-augmentation-ringed-over-object}、およびこの節の導入部の議論を介して、
補題 \ref{spaces-simplicial-lemma-compare-cohomology-hypercovering-modules} から従う。
\end{proof}

\begin{lemma}
\label{spaces-simplicial-lemma-hypercovering-X-equivalence-bounded-modules}
$\mathcal{C}$ をファイバー積をもつサイトとし、
$X \in \Ob(\mathcal{C})$ とする。$\mathcal{O}_\mathcal{C}$ を環の層、
$K$ を $X$ の超被覆とする。
$\mathcal{A} \subset \textit{Mod}(\mathcal{O})$ をカルテジアン
$\mathcal{O}$-加群からなる弱 Serre 部分圏とする。このとき関手
$La^*$ は同値
$$
D^+(\mathcal{O}_X) \longrightarrow D_\mathcal{A}^+(\mathcal{O})
$$
を定め、その擬逆は $Ra_*$ である。
\end{lemma}

\begin{proof}
注意 \ref{spaces-simplicial-remark-semi-representable-ringed-over-object}、
\ref{spaces-simplicial-remark-augmentation-ringed-over-object}、およびこの節の導入部の議論を介して、
補題 \ref{spaces-simplicial-lemma-hypercovering-equivalence-bounded-modules} から従う。
\end{proof}

\section{サイトの単体的対象による超被覆}
\label{spaces-simplicial-section-hypercovering}

\noindent
$\mathcal{C}$ をファイバー積をもつサイトとし、
$X \in \Ob(\mathcal{C})$ とする。この節では、超被覆がサイトの
単体的対象によって与えられる場合について、節
\ref{spaces-simplicial-section-cohomological-descent-hypercoverings-X} の結果を明確にする。
$U$ を $\mathcal{C}$ の単体的対象とする。通常どおり
$U_n = U([n])$ と書き、$\varphi : [m] \to [n]$ に対応する射
$f_\varphi = U(\varphi)$ を $f_\varphi : U_n \to U_m$ と書く。
増大
$$
a : U \to X
$$
が与えられていると仮定する。これから単体的サイト
$(\mathcal{C}/U)_{total}$ と増大射
$$
a : \Sh((\mathcal{C}/U)_{total}) \longrightarrow \Sh(\mathcal{C}/X)
$$
を得る。すなわち、$U$ から $\text{SR}(\mathcal{C}, X)$ の
% P10-E0212
単体的対象 $K$ を得る。その次数 $n$ の部分は
$K_n = \{U_n \to X\}$ であり、注意
\ref{spaces-simplicial-remark-augmentation-over-object} の構成を適用できる。
% P10-E0215
より具体的には、サイト $(\mathcal{C}/U)_{total}$ の対象は
$V/U_n$ により与えられ、射
$(\varphi, f) : V/U_n \to W/U_m$ は、$\Delta$ の射
$\varphi : [m] \to [n]$ と、図式
$$
\xymatrix{
V \ar[r]_f \ar[d] & W \ar[d] \\
U_n \ar[r]^{f_\varphi} & U_m
}
$$
が可換となる射 $f : V \to W$ によって与えられる。
トポスの射 $a$ は、余連続関手 $V/U_n \mapsto V/X$ により与えられる。
以上である！

\medskip\noindent
この節では、増大 $a : U \to X$ が次を満たすとき、それを
{\it $\mathcal{C}$ における $X$ の超被覆}と呼ぶ。
\begin{enumerate}
\item $\{U_0 \to X\}$ は $\mathcal{C}$ の被覆である。
\item $\{U_1 \to U_0 \times_X U_0\}$ は $\mathcal{C}$ の被覆である。
\item $n \geq 1$ に対して
$\{U_{n + 1} \to (\text{cosk}_n\text{sk}_n U)_{n + 1}\}$ は
$\mathcal{C}$ の被覆である。
\end{enumerate}
これは、上の $K$ が $X$ の超被覆であるという条件と同値である。
『超被覆』例 \ref{hypercovering-example-hypercovering-in-C} を参照されたい。

\begin{lemma}
\label{spaces-simplicial-lemma-hypercovering-X-simple-descent-sheaves}
% P10-E0213
$\mathcal{C}$ をファイバー積をもつサイトとし、
$X \in \Ob(\mathcal{C})$ とする。$a : U \to X$ を、上で定義した
$\mathcal{C}$ における $X$ の超被覆とする。このとき
\begin{enumerate}
\item $a^{-1} : \Sh(\mathcal{C}/X) \to \Sh((\mathcal{C}/U)_{total})$ は
充満忠実であり、その本質的像は集合のカルテジアン層である。
\item $a^{-1} : \textit{Ab}(\mathcal{C}/X) \to
\textit{Ab}((\mathcal{C}/U)_{total})$ は充満忠実であり、
その本質的像はアーベル群のカルテジアン層である。
\end{enumerate}
いずれの場合も $a_*$ が擬逆関手を与える。
\end{lemma}

\begin{proof}
補題 \ref{spaces-simplicial-lemma-hypercovering-X-descent-sheaves} の特別な場合である。
\end{proof}

\begin{lemma}
\label{spaces-simplicial-lemma-hypercovering-X-simple-descent-bounded-abelian}
% P10-E0213
$\mathcal{C}$ をファイバー積をもつサイトとし、
$X \in \Ob(\mathcal{C})$ とする。$a : U \to X$ を、上で定義した
$\mathcal{C}$ における $X$ の超被覆とする。
$E \in D(\mathcal{C}/X)$ に対して、写像
$$
E \longrightarrow Ra_*a^{-1}E
$$
は同型である。
\end{lemma}

\begin{proof}
補題 \ref{spaces-simplicial-lemma-hypercovering-X-descent-bounded-abelian} の特別な場合である。
\end{proof}

\begin{lemma}
\label{spaces-simplicial-lemma-compare-cohomology-hypercovering-X-simple}
$\mathcal{C}$ をファイバー積をもつサイトとし、
$X \in \Ob(\mathcal{C})$ とする。$a : U \to X$ を、上で定義した
$\mathcal{C}$ における $X$ の超被覆とする。このとき標準同型
$$
R\Gamma(X, E) = R\Gamma((\mathcal{C}/U)_{total}, a^{-1}E)
$$
が $E \in D(\mathcal{C}/X)$ に対してある。
\end{lemma}

\begin{proof}
補題 \ref{spaces-simplicial-lemma-compare-cohomology-hypercovering-X} の特別な場合である。
\end{proof}

\begin{lemma}
\label{spaces-simplicial-lemma-hypercovering-X-simple-equivalence-bounded}
% P10-E0213
$\mathcal{C}$ をファイバー積をもつサイトとし、
$X \in \Ob(\mathcal{C})$ とする。$a : U \to X$ を、上で定義した
$\mathcal{C}$ における $X$ の超被覆とする。
$\mathcal{A} \subset \textit{Ab}((\mathcal{C}/U)_{total})$ を
カルテジアン・アーベル層からなる弱 Serre 部分圏とする。このとき関手
$a^{-1}$ は同値
$$
D^+(\mathcal{C}/X) \longrightarrow D_\mathcal{A}^+((\mathcal{C}/U)_{total})
$$
を定め、その擬逆は $Ra_*$ である。
\end{lemma}

\begin{proof}
補題 \ref{spaces-simplicial-lemma-hypercovering-X-equivalence-bounded} の特別な場合である。
% P10-E0214
\end{proof}

\begin{lemma}
\label{spaces-simplicial-lemma-sr-when-fibre-products}
$U$ を、ファイバー積をもつサイト $\mathcal{C}$ の単体的対象とする。
\begin{enumerate}
\item $\mathcal{C}/U$ は、対象がサイトで射がサイトの射である圏の
単体的対象の構造をもつ。
\item (1) の構造に補題 \ref{spaces-simplicial-lemma-simplicial-site-site} の構成を
適用すると、上のサイト $(\mathcal{C}/U)_{total}$ が再現される。
\item $a : U \to X$ が増大ならば、
$a_0 : \mathcal{C}/U_0 \to \mathcal{C}/X$ は注意
\ref{spaces-simplicial-remark-augmentation-site} の (A) の意味で増大であり、上と同じ
トポスの射
$a : \Sh((\mathcal{C}/U)_{total}) \to \Sh(\mathcal{C}/X)$ を与える。
\end{enumerate}
\end{lemma}

\begin{proof}
$\mathcal{C}$ の対象の射 $V \to W$ が与えられると、局所化射
$j : \mathcal{C}/V \to \mathcal{C}/W$ は基底変換関手
$\mathcal{C}/W \to \mathcal{C}/V$ の左随伴である。基底変換関手は連続であり、
$j$ と同じトポスの射を誘導する。『サイト』補題
\ref{sites-lemma-relocalize-given-fibre-products} を参照されたい。
これで (1) が示された。

\medskip\noindent
(2) が成り立つのは、補題 \ref{spaces-simplicial-lemma-simplicial-site-site} で構成された圏の射
$V/U_n \to W/U_m$ が、$U_n$ 上の射
$V \to W \times_{U_m, f_\varphi} U_n$、すなわち射
$f_\varphi : U_n \to U_m$ 上の射 $f : V \to W$ と同じものであり、
これは上で定義した圏 $(\mathcal{C}/U)_{total}$ の射と同じだからである。
被覆族の集合が等しいことは定義から直ちに分かる。

\medskip\noindent
(3) の証明は省略する。
\end{proof}

\section{サイトの単体的対象による超被覆：加群}
\label{spaces-simplicial-section-hypercovering-modules}

\noindent
$\mathcal{C}$ をファイバー積をもつサイトとし、
$X \in \Ob(\mathcal{C})$ とする。
$\mathcal{O}_\mathcal{C}$ を $\mathcal{C}$ 上の環の層とする。
$U \to X$ を、節 \ref{spaces-simplicial-section-hypercovering} で定義した
$\mathcal{C}$ における $X$ の超被覆とする。この節では、$U$ を
% P10-E0216
$\mathcal{C}/X$ の単体的半表現可能対象で、その次数 $n$ の部分が
一元集合 $\{U_n/X\}$ であるものとみなし、注意
\ref{spaces-simplicial-remark-augmentation-ringed-over-object} の構成を適用して得られる環付き増大
$$
a :
(\Sh((\mathcal{C}/U)_{total}), \mathcal{O})
\longrightarrow
(\Sh(\mathcal{C}/X), \mathcal{O}_X)
$$
を調べる。したがって、この節のすべての結果は、節
\ref{spaces-simplicial-section-cohomological-descent-hypercoverings-X-modules} の
対応する結果から直ちに従う。

\begin{lemma}
\label{spaces-simplicial-lemma-hypercovering-X-simple-descent-modules}
$\mathcal{C}$ をファイバー積をもつサイトとし、
$X \in \Ob(\mathcal{C})$ とする。
$\mathcal{O}_\mathcal{C}$ を環の層とする。
$U$ を $\mathcal{C}$ における $X$ の超被覆とする。上の記法のもとで
$$
a^* : \textit{Mod}(\mathcal{O}_X) \to \textit{Mod}(\mathcal{O})
$$
は充満忠実であり、その本質的像はカルテジアン $\mathcal{O}$-加群である。
関手 $a_*$ が擬逆を与える。
\end{lemma}

\begin{proof}
補題 \ref{spaces-simplicial-lemma-hypercovering-X-descent-modules} の特別な場合である。
\end{proof}

\begin{lemma}
\label{spaces-simplicial-lemma-hypercovering-X-simple-descent-bounded-modules}
$\mathcal{C}$ をファイバー積をもつサイトとし、
$X \in \Ob(\mathcal{C})$ とする。
$\mathcal{O}_\mathcal{C}$ を環の層とする。
$U$ を $\mathcal{C}$ における $X$ の超被覆とする。
$E \in D(\mathcal{O}_X)$ に対して、射
$$
E \longrightarrow Ra_*La^*E
$$
は同型である。
\end{lemma}

\begin{proof}
補題 \ref{spaces-simplicial-lemma-hypercovering-X-descent-bounded-modules} の特別な場合である。
\end{proof}

\begin{lemma}
\label{spaces-simplicial-lemma-compare-cohomology-hypercovering-X-simple-modules}
$\mathcal{C}$ をファイバー積をもつサイトとし、
$X \in \Ob(\mathcal{C})$ とする。
$\mathcal{O}_\mathcal{C}$ を環の層とする。
$U$ を $\mathcal{C}$ における $X$ の超被覆とする。
このとき標準同型
$$
R\Gamma(X, E) = R\Gamma((\mathcal{C}/U)_{total}, La^*E)
$$
が
% P10-E0217
$E \in D(\mathcal{O}_\mathcal{C})$ に対してある。
\end{lemma}

\begin{proof}
補題 \ref{spaces-simplicial-lemma-compare-cohomology-hypercovering-X-modules} の特別な場合である。
\end{proof}

\begin{lemma}
\label{spaces-simplicial-lemma-hypercovering-X-simple-equivalence-bounded-modules}
$\mathcal{C}$ をファイバー積をもつサイトとし、
$X \in \Ob(\mathcal{C})$ とする。
$\mathcal{O}_\mathcal{C}$ を環の層とする。
$U$ を $\mathcal{C}$ における $X$ の超被覆とする。
$\mathcal{A} \subset \textit{Mod}(\mathcal{O})$ をカルテジアン
$\mathcal{O}$-加群からなる弱 Serre 部分圏とする。このとき関手 $La^*$ は同値
$$
D^+(\mathcal{O}_X) \longrightarrow D_\mathcal{A}^+(\mathcal{O})
$$
を定め、その擬逆は $Ra_*$ である。
\end{lemma}

\begin{proof}
補題 \ref{spaces-simplicial-lemma-hypercovering-X-equivalence-bounded-modules} の特別な場合である。
\end{proof}

\section{超被覆に対する非有界コホモロジー的下降}
\label{spaces-simplicial-section-unbounded-cohomological-descent}

\noindent
この節では非有界コホモロジー的下降を論じる。結果そのものは、
前節までの有界コホモロジー的下降に関する結果と、
『サイト上のコホモロジー』補題
\ref{sites-cohomology-lemma-equivalence-unbounded-one} および（または）
\ref{sites-cohomology-lemma-equivalence-unbounded-two} から直ちに従う。
実質的な作業は、記法を整え、適切な仮定を選ぶことにある。
ここでの議論は \cite{six-I} の議論に動機づけられているが、
細部はかなり異なる。

\medskip\noindent
$(\mathcal{C}, \mathcal{O}_\mathcal{C})$ を環付きサイトとする。
$\mathcal{C}$ の各対象 $U$ に対し、次の性質を満たす弱 Serre 部分圏
$\mathcal{A}_U \subset \textit{Mod}(\mathcal{O}_U)$ が与えられていると仮定する。
\begin{enumerate}
\item
\label{spaces-simplicial-item-restriction}
$\mathcal{C}$ の射 $U \to V$ が与えられると、制限関手
$\textit{Mod}(\mathcal{O}_V) \to \textit{Mod}(\mathcal{O}_U)$ は
$\mathcal{A}_V$ を $\mathcal{A}_U$ に送る。
\item
\label{spaces-simplicial-item-local}
$\mathcal{C}$ の被覆 $\{U_i \to U\}_{i \in I}$ が与えられたとき、
$\textit{Mod}(\mathcal{O}_U)$ の対象 $\mathcal{F}$ が
$\mathcal{A}_U$ に属するための必要十分条件は、すべての $i \in I$ に対し、
$\mathcal{F}$ の $\mathcal{C}/U_i$ への制限が $\mathcal{A}_{U_i}$ に属することである。
\item
\label{spaces-simplicial-item-bounded-dimension}
次を満たす部分集合 $\mathcal{B} \subset \Ob(\mathcal{C})$ が存在する。
\begin{enumerate}
\item $\mathcal{C}$ の各対象は、その構成要素が $\mathcal{B}$ に属する被覆をもつ。
\item 各 $V \in \mathcal{B}$ に対して、整数 $d_V$ と、$V$ の被覆の余終系
$\text{Cov}_V$ が存在し、次を満たす。
$$
H^p(V_i, \mathcal{F}) = 0 \text{、ただし }
\{V_i \to V\} \in \text{Cov}_V,\ p > d_V, \text{ かつ }
\mathcal{F} \in \Ob(\mathcal{A}_V)
$$
\end{enumerate}
\end{enumerate}
ここで、この条件を単に「大域的」対象についてではなく、
$\mathcal{A}_V$ に属する $\mathcal{F}$ について要求していることに注意する。
したがって、これは『サイト上のコホモロジー』情況
\ref{sites-cohomology-situation-olsson-laszlo} で課された条件より強い。
この状況では、すべての $U \in \Ob(\mathcal{C})$ に対して
$\mathcal{C}/U$ への制限が $\mathcal{A}_U$ に属する対象からなる弱 Serre 部分圏
$\mathcal{A} \subset \textit{Mod}(\mathcal{O}_\mathcal{C})$ がある。
さらに、導来圏 $D_\mathcal{A}(\mathcal{O}_\mathcal{C})$ と
$D_{\mathcal{A}_U}(\mathcal{O}_U)$ があり、制限関手はこれらを互いに送る。

\begin{example}
\label{spaces-simplicial-example-quasi-coherent-spaces-etale}
$S$ をスキームとし、$X$ を $S$ 上の代数空間とする。\linebreak
$\mathcal{C} = X_{spaces, \etale}$ を、$X$ 上エタールな代数空間の圏上の
エタール・サイトとする。『空間の性質』定義
\ref{spaces-properties-definition-spaces-etale-site} を参照されたい。
$\mathcal{O}_\mathcal{C}$ を構造層、すなわち規則
$U \mapsto \Gamma(U, \mathcal{O}_U)$ によって与えられる層とする。
$\mathcal{A}_U$ を準連接 $\mathcal{O}_U$-加群の圏とする。
$\mathcal{B} = \Ob(\mathcal{C})$ とし、$V \in \mathcal{B}$ に対して
$d_V = 0$ とおき、$\text{Cov}_V$ を、すべての $i$ に対して
$V_i$ がアフィンである被覆 $\{V_i \to V\}$ とする。
このとき仮定 (1)、(2)、(3) は満たされる。性質 (1)、(2) については
『空間の性質』補題
\ref{spaces-properties-lemma-pullback-quasi-coherent} および
\ref{spaces-properties-lemma-properties-quasi-coherent} を参照されたい。
(3) の消滅は、『スキームのコホモロジー』補題
\ref{coherent-lemma-quasi-coherent-affine-cohomology-zero} と、
『空間のコホモロジー』節
\ref{spaces-cohomology-section-higher-direct-image} の議論から従う。
\end{example}

\begin{example}
\label{spaces-simplicial-example-etale}
$S$ を次のいずれかの型のスキームとする。
\begin{enumerate}
\item 有限体のスペクトル。
\item 分離閉体のスペクトル。
\item 狭義 Hensel な Noether 局所環のスペクトル。
\item 剰余体が有限である Hensel Noether 局所環のスペクトル。
\item % P10-E0218
ここにさらに追加する。
\end{enumerate}
$\Lambda$ を、その位数が $S$ 上可逆である有限環とする。
$\mathcal{C} \subset (\Sch/S)_\etale$ を、$S$ 上局所有限型のスキームからなる
全充満部分圏とし、エタール位相を入れる。
$\mathcal{O}_\mathcal{C} = \underline{\Lambda}$ を定数層とする。
$\mathcal{A}_U = \textit{Mod}(\mathcal{O}_U)$ とおく。
換言すれば、$\Lambda$-加群のすべてのエタール層を考える。
$\mathcal{B} \subset \Ob(\mathcal{C})$ を準コンパクト対象の集合とする。
$V \in \mathcal{B}$ に対して
$$
d_V = 1 + 2\dim(S) +
\sup\nolimits_{v \in V}(\text{trdeg}_{\kappa(s)}(\kappa(v)) +
2 \dim \mathcal{O}_{V, v})
$$
とおき、$\text{Cov}_V$ を、すべての $i$ に対して $V_i$ が準コンパクトである
エタール被覆 $\{V_i \to V\}$ とする。界 $d_V$ の選択は、
コホモロジー次元に関する Gabber の定理に由来する。
この選択で条件 (3) が成り立つことを見るには、
\cite[Expos\'e VIII-A, Corollary 1.2 and Lemma 2.2]{Traveaux}
と、体のコホモロジー次元に関する初等的な議論を用いる。
このリストには $S$ が有限剰余体をもつことを許す場合が含まれるため、
式に $1$ を加えている。後でこの例に戻る（% P10-E0219
将来の参照を挿入する）。
\end{example}

\noindent
$(\mathcal{C}, \mathcal{O}_\mathcal{C})$ を環付きサイトとする。
条件 (\ref{spaces-simplicial-item-restriction}) を満たす弱 Serre 部分圏
$\mathcal{A}_U \subset \textit{Mod}(\mathcal{O}_U)$ が与えられていると仮定する。
このとき次が成り立つ。
\begin{enumerate}
\item 半表現可能対象 $K = \{U_i\}_{i \in I}$ が与えられると、
\linebreak
$\prod \mathcal{A}_{U_i} \subset
\prod \textit{Mod}(\mathcal{O}_{U_i}) = \textit{Mod}(\mathcal{O}_K)$
とおくことにより、弱 Serre 部分圏
$\mathcal{A}_K \subset \textit{Mod}(\mathcal{O}_K)$ を得る。
\item 半表現可能対象の射 $f : K \to L$ が与えられると、引戻し写像
% P10-E0220
$f^* : \textit{Mod}(\mathcal{O}_L) \to \textit{Mod}(\mathcal{O}_L)$
は $\mathcal{A}_L$ を $\mathcal{A}_K$ に送る。
\end{enumerate}
記法と説明については注意 \ref{spaces-simplicial-remark-semi-representable-ringed} を参照されたい。
特に、単体的半表現可能対象 $K$ が与えられたとき、注意
\ref{spaces-simplicial-remark-augmentation-ringed} における $\textit{Mod}(\mathcal{O})$ の対象
$\mathcal{F}$ の制限 $\mathcal{F}_n$ がすべての $n$ に対して
$\mathcal{A}_{K_n}$ に属する、ということの意味は曖昧でない。

\begin{lemma}
\label{spaces-simplicial-lemma-hypercovering-equivalence-modules}
$(\mathcal{C}, \mathcal{O}_\mathcal{C})$ を環付きサイトとする。
上の条件 (\ref{spaces-simplicial-item-restriction})、(\ref{spaces-simplicial-item-local})、および
(\ref{spaces-simplicial-item-bounded-dimension}) を満たす弱 Serre 部分圏
$\mathcal{A}_U \subset \textit{Mod}(\mathcal{O}_U)$ が与えられていると仮定する。
$\mathcal{C}$ は等化子とファイバー積をもつと仮定し、$K$ を超被覆とする。
$((\mathcal{C}/K)_{total}, \mathcal{O})$ を注意
\ref{spaces-simplicial-remark-augmentation-ringed} のものとする。
$\mathcal{A}_{total} \subset \textit{Mod}(\mathcal{O})$ を、
すべての $n$ に対して制限 $\mathcal{F}_n$ が
$\mathcal{A}_{K_n}$ に属する（意味は上で定義した）カルテジアン
$\mathcal{O}$-加群 $\mathcal{F}$ からなる弱 Serre 部分圏とする。
このとき関手 $La^*$ は同値
$$
D_\mathcal{A}(\mathcal{O}_\mathcal{C})
\longrightarrow
D_{\mathcal{A}_{total}}(\mathcal{O})
$$
を定め、その擬逆は $Ra_*$ である。
\end{lemma}

\begin{proof}
補題 \ref{spaces-simplicial-lemma-Serre-subcat-cartesian-modules} により、
カルテジアン $\mathcal{O}$-加群は弱 Serre 部分圏をなす
（必要な仮定は注意 \ref{spaces-simplicial-remark-augmentation-ringed} の議論により成り立つ）。
制限関手
% P10-E0221
$g_n^* : \textit{Mod}(\mathcal{O}) \to \textit{Mod}(\mathcal{O}_n)$
は完全なので、$\mathcal{A}_{total}$ は弱 Serre 部分圏である。

\medskip\noindent
$a^* : \mathcal{A} \to \mathcal{A}_{total}$ が圏の同値であり、逆が
% P10-E0222
$La_*$ によって与えられることを示そう。有界版
（補題 \ref{spaces-simplicial-lemma-hypercovering-equivalence-bounded-modules}）により、すでに
$La_*a^*\mathcal{F} = \mathcal{F}$ であることが分かっている。
$\mathcal{F}$ が $\mathcal{A}$ に属するとき、
$a^*\mathcal{F}$ が $\mathcal{A}_{total}$ に属することは明らかである。
逆に $\mathcal{G} \in \mathcal{A}_{total}$ と仮定する。
$\mathcal{G}$ はカルテジアンなので、補題
\ref{spaces-simplicial-lemma-hypercovering-descent-modules} により、ある
$\mathcal{O}_\mathcal{C}$-加群 $\mathcal{F}$ に対して
$\mathcal{G} = a^*\mathcal{F}$ である。$\mathcal{F}$ が $\mathcal{A}$ に属することを示したい。
$U \in \Ob(\mathcal{C})$ をとる。$\mathcal{F}$ の
$\mathcal{C}/U$ への制限が $\mathcal{A}_U$ に属することを示せばよい。
通常どおり $K_0 = \{U_{0, i}\}_{i \in I_0}$ と書く。
$K$ は超被覆なので、写像
$\coprod_{i \in I_0} h_{U_{0, i}} \to *$ は層化後に全射となる。
したがって、被覆 $\{U_j \to U\}_{j \in J}$ と写像
$\tau : J \to I_0$ があり、各 $j \in J$ に対して射
$\varphi_j : U_j \to U_{0, \tau(j)}$ がある。
$\mathcal{G}_0 = a_0^*\mathcal{F}$ なので、$\mathcal{F}$ の
$\mathcal{C}/U_j$ への制限は、$\mathcal{G}_0$ の第 $\tau(j)$ 成分を射
% P10-E0223
$\varphi_j : U_j \to U_{0, \tau(i)}$ を介して
$\mathcal{C}/U_j$ へ制限したものに等しい。したがって
(\ref{spaces-simplicial-item-restriction}) により
$\mathcal{F}|_{\mathcal{C}/U_j}$ は $\mathcal{A}_{U_j}$ に属し、さらに
(\ref{spaces-simplicial-item-local}) により
$\mathcal{F}|_{\mathcal{C}/U}$ は $\mathcal{A}_U$ に属する。

\medskip\noindent
特に補題の主張は意味をもつ。補題は『サイト上のコホモロジー』補題
\ref{sites-cohomology-lemma-equivalence-unbounded-one} から従う。
仮定 (1) は明らかである（注意 \ref{spaces-simplicial-remark-augmentation-ringed} を参照されたい）。
仮定 (2)、(3) は前段落で証明した。仮定 (4) は
(\ref{spaces-simplicial-item-bounded-dimension}) から直ちに従う。
仮定 (5) について、$\mathcal{B}_{total}$ を、サイト
$(\mathcal{C}/K)_{total}$ の対象 $U/U_{n, i}$ であって
$U \in \mathcal{B}$ となるものの集合とする。ここで $\mathcal{B}$ は
(\ref{spaces-simplicial-item-bounded-dimension}) のものである。ここでは節
\ref{spaces-simplicial-section-simplicial-semi-representable} で与えたサイト
$(\mathcal{C}/K)_{total}$ の記述を用いている。
さらに、$\text{Cov}_{U/U_{n, i}}$ を $\text{Cov}_U$ と等しく、
$d_{U/U_{n, i}}$ を $d_U$ と等しくおく。ここで
$\text{Cov}_U$ と $d_U$ は (\ref{spaces-simplicial-item-bounded-dimension}) で与えられたものである。
これらの選択により条件 (5) が成り立つと主張する。これは補題
\ref{spaces-simplicial-lemma-sanity-check-simplicial-semi-representable} と、
$\mathcal{F} \in \mathcal{A}_{total}$ ならば
$\mathcal{F}_n \in \mathcal{A}_{K_n}$、したがって
$\mathcal{F}_{n, i} \in \mathcal{A}_{U_{n, i}}$ であるという事実から直ちに従う。
（アーベル層のコホモロジーと加群層のコホモロジーの差を懸念する読者は、
『サイト上のコホモロジー』補題
\ref{sites-cohomology-lemma-cohomology-modules-abelian-agree} を参照されたい。）
\end{proof}

\section{複体の貼り合わせ}
\label{spaces-simplicial-section-glueing-complexes}

\noindent
この節は『コホモロジー』節
\ref{cohomology-section-glueing-complexes} の続きである。
目標は \cite[Theorem 3.2.4]{BBD} のわずかな一般化を証明することである。
節 \ref{spaces-simplicial-section-glueing} と節 \ref{spaces-simplicial-section-glueing-modules} の内容を用いる点で、
ここでの方法は少しだけ異なる。また、\cite{six-I} における証明にならって
非有界版も改めて証明する。

\medskip\noindent
読者への助言：まず補題 \ref{spaces-simplicial-lemma-bbd-glueing-cech} の主張と、
この補題の第二の証明を見ることを勧める。

\medskip\noindent
以下が、ここで関心をもつ状況である。

\begin{situation}
\label{spaces-simplicial-situation-locally-given}
$(\mathcal{C}, \mathcal{O}_\mathcal{C})$ を環付きサイトとする。次が与えられている。
\begin{enumerate}
\item 圏 $\mathcal{B}$ と関手 $u : \mathcal{B} \to \mathcal{C}$。
\item $U \in \Ob(\mathcal{B})$ に対する
$D(\mathcal{O}_{u(U)})$ の対象 $E_U$。
\item $\mathcal{B}$ の $a : V \to U$ に対する
$D(\mathcal{O}_{u(V)})$ における同型
$\rho_a : E_U|_{\mathcal{C}/u(V)} \to E_V$。
\end{enumerate}
これらは、$\mathcal{B}$ の可合成な射 $b : W \to V$ と
$a : V \to U$ があるたびに
$\rho_{a \circ b} = \rho_b \circ \rho_a|_{\mathcal{C}/u(W)}$ を満たす。
\end{situation}

\noindent
さらに仮定を置かなければ、このデータについて何も証明できない。
興味深い場合の一つは、$\mathcal{B}$ が全充満部分圏であり、
$\mathcal{C}$ の各対象が、構成要素を $\mathcal{B}$ の対象とする被覆をもつ場合である
（\cite{BBD} で扱われているのはこの場合である）。
ここでは、そうでない場合も許すことが重要である。主要な別の場合は、
サイトの射 $f : \mathcal{C} \to \mathcal{D}$ があり、
$\mathcal{B}$ が $\mathcal{D}$ の全充満部分圏であって、
$\mathcal{D}$ の各対象が、構成要素を $\mathcal{B}$ の対象とする被覆をもつ場合である。

\medskip\noindent
情況 \ref{spaces-simplicial-situation-locally-given} における {\it 解}とは、対
$(E, \rho_U)$ であって、$E$ は $D(\mathcal{O}_\mathcal{C})$ の対象、
$U \in \Ob(\mathcal{B})$ に対する
$\rho_U : E|_{\mathcal{C}/u(U)} \to E_U$ は同型であり、
$\mathcal{B}$ の $a : V \to U$ に対して
$\rho_a \circ \rho_U|_{\mathcal{C}/u(V)} = \rho_V$ を満たすものをいう。

\begin{lemma}
\label{spaces-simplicial-lemma-prepare-bbd-glueing}
情況 \ref{spaces-simplicial-situation-locally-given} において、% P10-E0224
$D(\mathcal{O}_{u(U)})$ における $E_U$ の負次数の自己 Ext は消えると仮定する。
$L$ を $\text{SR}(\mathcal{B})$ の単体的対象とする。
$K = u(L)$ を $\text{SR}(\mathcal{C})$ の単体的対象とし、
$((\mathcal{C}/K)_{total}, \mathcal{O})$ を注意
\ref{spaces-simplicial-remark-augmentation-ringed} のものとする。
$D(\mathcal{O})$ のカルテジアン対象 $E$ で、次の性質をもつものが存在する。
$L_n = \{U_{n, i}\}_{i \in I_n}$ と書くと、$E$ の
$D(\mathcal{O}_{\mathcal{C}/u(U_{n, i})})$ への制限は、両立する仕方で
$E_{U_{n, i}}$ である（詳細は証明を参照されたい）。
さらに、$E$ は一意な同型を除いて一意である。
\end{lemma}

\begin{proof}
$\Sh(\mathcal{C}/K_n) = \prod_{i \in I_n} \Sh(\mathcal{C}/u(U_{n, i}))$
であり、加群の圏についても同様であることを思い出そう。
この積分解は加群層の導来圏にも継承される。
さらに、この積分解は単体的半表現可能対象 $K$ の射と両立する。
節 \ref{spaces-simplicial-section-semi-representable} を参照されたい。
したがって、$D(\mathcal{O}_n)$ において
$E_n = \prod_{i \in I_n} E_{U_{n, i}}$（「形式的」積）とおける。
情況 \ref{spaces-simplicial-situation-locally-given} の写像 $\rho_a$ の（形式的）積をとると、
同型 $E_\varphi : f_\varphi^*E_n \to E_m$ を得る。
% P10-E0225
写像 $\rho_a$ の合成に関する仮定から直ちに、
$(E_n, E_\varphi)$ は定義
\ref{spaces-simplicial-definition-cartesian-derived-modules} の意味で
加群の導来圏の単体的系を定める。
補題で仮定した負次数 Ext の消滅により、
$n \geq 0$ と $t > 0$ に対して $\Hom(E_n[t], E_n) = 0$ である。
したがって補題 \ref{spaces-simplicial-lemma-cartesian-module-derived-from-simplicial} により
$E$ を得る。一意な同型を除く一意性は、補題
\ref{spaces-simplicial-lemma-nullity-cartesian-modules-derived} と
\ref{spaces-simplicial-lemma-hom-cartesian-modules-derived} から従う。
\end{proof}

\begin{lemma}[BBD 貼り合わせ補題]
\label{spaces-simplicial-lemma-bbd-glueing}
情況 \ref{spaces-simplicial-situation-locally-given} において、次を仮定する。
\begin{enumerate}
\item $\mathcal{C}$ は等化子とファイバー積をもつ。
\item サイトの射 $f : \mathcal{C} \to \mathcal{D}$ が、連続関手
$u : \mathcal{D} \to \mathcal{C}$ により与えられ、次を満たす。
\begin{enumerate}
\item $\mathcal{D}$ は等化子とファイバー積をもち、$u$ はそれらと可換である。
\item $\mathcal{B}$ は $\mathcal{D}$ の全充満部分圏であり、
$u : \mathcal{B} \to \mathcal{C}$ は $u$ の制限である。
\item $\mathcal{D}$ の各対象は、その構成要素を $\mathcal{B}$ の対象とする被覆をもつ。
\end{enumerate}
\item $\mathcal{B}$ のすべての $U$ に対し、
$D(\mathcal{O}_{u(U)})$ における $E_U$ の負次数の自己 Ext はすべて消える。
\item $i < t$ および $U \in \Ob(\mathcal{B})$ に対して
$H^i(E_U) = 0$ となる $t \in \mathbf{Z}$ が存在する。
\end{enumerate}
このとき、一意な同型を除いて一意な解が存在する。
\end{lemma}

\begin{proof}
『超被覆』補題 \ref{hypercovering-lemma-hypercovering-site} により、
サイト $\mathcal{D}$ の超被覆 $L$ で、
% P10-E0226
$L_n = \{U_{n, i}\}_{i \in I_n}$ かつ
$U_{i, n} \in \Ob(\mathcal{B})$ となるものが存在する。
$K = u(L)$ とおく。補題 \ref{spaces-simplicial-lemma-prepare-bbd-glueing} を適用すると、
サイト $(\mathcal{C}/K)_{total}$ 上の $D(\mathcal{O})$ の
カルテジアン対象 $E$ で、両立する仕方で
$\mathcal{C}/u(U_{n, i})$ 上 $E_{U_{n, i}}$ に制限されるものを得る。
$t$ に関する仮定から $E \in D^+(\mathcal{O})$ である。
『超被覆』補題 \ref{hypercovering-lemma-hypercovering-morphism-sites} により、
$K$ も超被覆である。補題
\ref{spaces-simplicial-lemma-hypercovering-equivalence-bounded-modules} により、
% P10-E0227
$D^+(\mathcal{O}_\mathcal{C})$ のある $F$ に対して
$E = a^*F$ であることが分かる。

\medskip\noindent
$F$ が解であることを証明するため、『超被覆』補題
\ref{hypercovering-lemma-hypercovering-site} の証明で与えられた
$L_0$ と $L_1$ の構成を用いる。
（これはやや洗練を欠くが、完全に直接的な回避法はないようである。）

\medskip\noindent
すなわち $I_0 = \Ob(\mathcal{B})$ であり、したがって
$L_0 = \{U\}_{U \in \Ob(\mathcal{B})}$ である。
% P10-E0227
同型 $a^*F \to E$ を $\mathcal{C}/K_0$ の成分
$\mathcal{C}/u(U)$ に制限すると、$E$ の選択により、
$U \in \Ob(\mathcal{B})$ に対する同型
$\rho_U : F|_{\mathcal{C}/u(U)} \to E_U$ が定まる。

\medskip\noindent
$\rho_U$ が情況 \ref{spaces-simplicial-situation-locally-given} の写像 $\rho_a$ との
両立性の要件を満たすことを証明するため、$I_1$ が集合
$$
\Omega =
\{(U, V, W, a, b) \mid U, V, W \in \mathcal{B}, a : U \to V, b : U \to W\}
$$
を含み、$\Omega$ の $i = (U, V, W, a, b)$ に対して
$U_{1, i} = U$ であることを用いる。さらに、二つの射
$K_1 \to K_0$ の成分写像 $f_{\delta^1_0, i}$ と
$f_{\delta^1_1, i}$ は射
$$
a : U \to V \quad\text{および}\quad b : U \to V
$$
% P10-E0228
である。したがって、補題 \ref{spaces-simplicial-lemma-prepare-bbd-glueing} で述べた両立性から
$$
\rho_a \circ \rho_V|_{\mathcal{C}/u(U)} = \rho_U
\quad\text{および}\quad
\rho_b \circ \rho_W|_{\mathcal{C}/u(U)} = \rho_U
$$
を得る。例えば $i = (U, V, U, a, \text{id}_U) \in \Omega$ ととれば、
望む両立性を得る。$F$ の一意性は前補題における $E$ の一意性から従う
（細部を一つ省略した）。
\end{proof}

\begin{lemma}[非有界 BBD 貼り合わせ補題]
\label{spaces-simplicial-lemma-bbd-unbounded-glueing}
情況 \ref{spaces-simplicial-situation-locally-given} において、次を仮定する。
\begin{enumerate}
\item $\mathcal{C}$ は等化子とファイバー積をもつ。
\item サイトの射 $f : \mathcal{C} \to \mathcal{D}$ が、連続関手
$u : \mathcal{D} \to \mathcal{C}$ により与えられ、次を満たす。
\begin{enumerate}
\item $\mathcal{D}$ は等化子とファイバー積をもち、$u$ はそれらと可換である。
\item $\mathcal{B}$ は $\mathcal{D}$ の全充満部分圏であり、
$u : \mathcal{B} \to \mathcal{C}$ は $u$ の制限である。
\item $\mathcal{D}$ の各対象は、その構成要素を $\mathcal{B}$ の対象とする被覆をもつ。
\end{enumerate}
\item $D(\mathcal{O}_{u(U)})$ における $E_U$ の負次数の自己 Ext はすべて消える。
\item すべての $U \in \Ob(\mathcal{C})$ に対し、条件
(\ref{spaces-simplicial-item-restriction})、(\ref{spaces-simplicial-item-local})、および
(\ref{spaces-simplicial-item-bounded-dimension}) を満たす弱 Serre 部分圏
$\mathcal{A}_U \subset \textit{Mod}(\mathcal{O}_U)$ が存在する。
\item $E_U \in D_{\mathcal{A}_U}(\mathcal{O}_U)$ である。
\end{enumerate}
このとき、一意な同型を除いて一意な解が存在する。
\end{lemma}

\begin{proof}
証明は補題 \ref{spaces-simplicial-lemma-bbd-glueing} の証明と {\bf まったく}同じである。
唯一の変更は、$E$ が $D_{\mathcal{A}_{total}}(\mathcal{O})$ の対象であり、
したがって補題 \ref{spaces-simplicial-lemma-hypercovering-equivalence-modules} を用いて
% P10-E0227
$E = a^*F$ となる $F$ を得ることである。これは補題
\ref{spaces-simplicial-lemma-hypercovering-equivalence-bounded-modules} を用いる代わりである。
\end{proof}

\noindent
以上の一般論の応用例を一つ挙げる。

\begin{lemma}
\label{spaces-simplicial-lemma-bbd-glueing-cech}
\begin{reference}
2021 年 9 月 9 日付の Martin Olsson の電子メール。
\end{reference}
$(\mathcal{C}, \mathcal{O}_\mathcal{C})$ を環付きサイトとする。
$\mathcal{C}$ はファイバー積をもつと仮定する。
$\{U_i \to X\}_{i \in I}$ を $\mathcal{C}$ の被覆とする。
$i \in I$ に対して $E_i$ を $D(\mathcal{O}_{U_i})$ の対象とし、
$i, j \in I$ に対して
$$
\rho_{ij} : E_i|_{\mathcal{C}/U_{ij}} \longrightarrow E_j|_{\mathcal{C}/U_{ij}}
$$
を $D(\mathcal{O}_{U_{ij}})$ における同型とする。ここで
$U_{ij} = U_i \times_X U_j$ である。次を仮定する。
\begin{enumerate}
\item すべての $i, j, k \in I$ に対し、$\rho_{ij}$ は
$U_i \times_X U_j \times_X U_k$ 上で余サイクル条件を満たす。
\item すべての $p < 0$ と $i \in I$ に対し、
$\SheafExt^p_{\mathcal{O}_{U_i}}(E_i, E_i) = 0$ である。
\item $p < t$ およびすべての $i \in I$ に対して
$H^p(E_i) = 0$ となる $t \in \mathbf{Z}$ が存在する。
\end{enumerate}
このとき一意な対 $(E, \rho_i)$ が存在する。ここで $E$ は
$D(\mathcal{O}_X)$ の対象であり、
$\rho_i : E|_{U_i} \to E_i$ は $D(\mathcal{O}_{U_i})$ における同型で、
$\rho_{ij}$ と両立する。
\end{lemma}

\begin{proof}[第一の証明]
この証明では、非常に一般的な補題 \ref{spaces-simplicial-lemma-bbd-glueing} から
この補題を導く。読者には代わりに第二の証明を見ることを強く勧める。% P10-E0229

\medskip\noindent
$\mathcal{C}$ を $\mathcal{C}/X$ で置き換えてよい。
したがって、$X$ が $\mathcal{C}$ の終対象であり、
$\mathcal{C}$ がすべての有限極限をもつと仮定してよく、そう仮定する。

\medskip\noindent
$\mathcal{B}$ を、ある $i(U) \in I$ と射
$a_U : U \to U_{i(U)}$ が存在するような $U \in \Ob(\mathcal{C})$
からなる $\mathcal{C}$ の全充満部分圏とする。
$D(\mathcal{O}_U)$ において、$a_U$ による
% P10-E0230
$E_i$ の引戻し（制限）を
$E_U = a_U^*E_{i(U)}$ と書く。
$\mathcal{B}$ の射 $a : U \to U'$ が与えられると、射
$(a_{U'} \circ a, a_U) : U \to U_{i(U')} \times_X U_{i(U)} = U_{i(U')i(U)}$
を得て、したがって同型
$$
\rho_a :
a^*E_{U'} = a^*a_{U'}^*E_{i(U')}
\xrightarrow{(a_{U'} \circ a, a_U)^*\rho_{i(U')i(U)}}
a_{U}^*E_{i(U)} = E_{U}
$$
を $D(\mathcal{O}_U)$ において得る。データ
$\mathcal{B}, E_U, \rho_a$ は情況 \ref{spaces-simplicial-situation-locally-given} のものとなる。
同型 $\rho_a$ が余サイクル条件を満たすのは、補題の主張の条件 (1) による
（詳細は省略する）。

\medskip\noindent
上の $\mathcal{B}$、$E_U$、$\rho_a$、$\mathcal{D} = \mathcal{C}$、および恒等射
$f : \mathcal{C} \to \mathcal{D}$ に対して補題
\ref{spaces-simplicial-lemma-bbd-glueing} を適用する。
補題 \ref{spaces-simplicial-lemma-bbd-glueing} の仮定 (1) と (2)(a) は上で確認した。
補題 \ref{spaces-simplicial-lemma-bbd-glueing} の仮定 (2)(b) は明らかである。
補題 \ref{spaces-simplicial-lemma-bbd-glueing} の仮定 (2)(c) は
$\{U_i \to X\}$ が被覆であることから成り立つ\footnote{実際には、写像
$\coprod_{i \in I} h_{U_i} \to h_X$ が層化後に全射となれば十分であり、
この場合にも同じ証明で補題が成り立つ。}。
補題 \ref{spaces-simplicial-lemma-bbd-glueing} の仮定 (3) が成り立つのは、
$E_i$ のすべての負次数 Ext 層の消滅を仮定したからである。
これは確かに、$U_i$ 上の任意の対象 $U$ に対して
$E_i|_U$ の負次数の自己 Ext が消えることを含意する。
補題 \ref{spaces-simplicial-lemma-bbd-glueing} の仮定 (4) が成り立つのは、
各 $E_i$ のコホモロジー層が $t$ より左で消えると仮定したからである。

\medskip\noindent
一意な解 $(E, \rho_U)$ を得る。$\rho_i = \rho_{U_i}$ とおけば補題が従う。
\end{proof}

\begin{proof}[第二の証明]
より直接的な証明を略述する。$K$ を、被覆
$\{U_i \to X\}_{i \in I}$ に付随する $X$ の {\v C}ech 超被覆とする。
『超被覆』例 \ref{hypercovering-example-cech} を参照されたい。
したがって例えば $K_0 = \{U_i \to X\}_{i \in I}$ であり、
$K_1 = \{U_i \times_X U_j \to X\}_{i, j \in I}$ である、等々。
$((\mathcal{C}/K)_{total}, \mathcal{O})$、$a$、$a_n$ を注意
\ref{spaces-simplicial-remark-augmentation-ringed-over-object} のものとする。
対象 $E_i$ は $D(\mathcal{O}_0) = \prod D(\mathcal{O}_{U_i})$ の対象
$M_0$ を定める。同様に、同型 $\rho_{ij}$ は、余サイクル条件を満たす
$D(\mathcal{O}_1)$ における同型
$$
\alpha :
L(f_{\delta_1^1})^*M_0
\longrightarrow
L(f_{\delta_0^1})^*M_0
$$
を定める。補題 \ref{spaces-simplicial-lemma-construct-cartesian-objects-derived-modules} により、
導来圏のカルテジアン単体的系 $(M_n)$ を得る。
仮定した負次数 Ext 層の消滅から、対象 $M_n$ の負次数の自己 Ext は消える。
したがって補題 \ref{spaces-simplicial-lemma-cartesian-module-derived-from-simplicial} により、
付随する単体的系が $(M_n)$ と同型であるような
$D(\mathcal{O})$ のカルテジアン対象 $M$ を得る。
$M$ のコホモロジー層は次数 $< t$ で消えるから、補題
\ref{spaces-simplicial-lemma-hypercovering-X-equivalence-bounded-modules} により、
$D(\mathcal{O}_X)$ のある $E$ に対して $M = La^*E$ である。
同型 $La^*E \to M$ を $\mathcal{C}/U_i$ に制限すると、同型 $\rho_i$ が得られる。
同型 $\rho_{ij}$ との両立性の確認は省略する。
\end{proof}

\section{位相空間における固有超被覆}
\label{spaces-simplicial-section-proper-hypercovering}

\noindent
Hausdorff かつ局所準コンパクトな位相空間と連続写像の圏
$\textit{LC}$ で考える。『サイト上のコホモロジー』節
\ref{sites-cohomology-section-cohomology-LC} を参照されたい。
$X$ を $\textit{LC}$ の対象、$U$ を $\textit{LC}$ の単体的対象とする。
増大
$$
a : U \to X
$$
が与えられていると仮定する。次が成り立つとき、$U$ を $X$ の
{\it 固有超被覆}という。
\begin{enumerate}
\item $U_0 \to X$ は固有全射である。
\item $U_1 \to U_0 \times_X U_0$ は固有全射である。
\item $n \geq 1$ に対し、
$U_{n + 1} \to (\text{cosk}_n\text{sk}_n U)_{n + 1}$ は固有全射である。
\end{enumerate}
圏 $\textit{LC}$ はすべての有限極限をもつので、上の定式化で用いた余骨格は存在する。
$$
\fbox{原理：固有超被覆を用いてコホモロジーを計算できる。}
$$
この原理の証明の背後にある重要な考えは、$\textit{LC}$ 上に通常の位相より強い位相で、
次を満たすものを見つけることである。
(a) 全射固有写像は被覆を定める。
(b) この強い位相に関する通常の層のコホモロジーは、
通常のコホモロジーと一致する。
性質 (a)、(b) は qc 位相について成り立つ。『サイト上のコホモロジー』節
\ref{sites-cohomology-section-cohomology-LC} を参照されたい。
(a)、(b) が得られれば、この章で先に行った仕事から原理が従う。

\begin{lemma}
\label{spaces-simplicial-lemma-compare-simplicial-objects}
$U$ を $\textit{LC}$ の単体的対象とし、$a : U \to X$ を増大とする。
可換図式
$$
\xymatrix{
\Sh((\textit{LC}_{qc}/U)_{total}) \ar[r]_-h \ar[d]_{a_{qc}} &
\Sh(U_{Zar}) \ar[d]^a \\
\Sh(\textit{LC}_{qc}/X) \ar[r]^-{h_{-1}} &
\Sh(X)
}
$$
がある。左の垂直射は節 \ref{spaces-simplicial-section-hypercovering} で定義され、
右の垂直射は補題 \ref{spaces-simplicial-lemma-augmentation} で定義される。
\end{lemma}

\begin{proof}
$\Sh(X) = \Sh(X_{Zar})$ と書く。
$(\textit{LC}_{qc}/U)_{total}$ と $U_{Zar}$ はともに、情況
\ref{spaces-simplicial-situation-simplicial-site} の場合 A に該当することに注意する。
$U_{Zar}$ については節 \ref{spaces-simplicial-section-simplicial-top} の構成から直ちに分かり、
$(\textit{LC}_{qc}/U)_{total}$ については補題
\ref{spaces-simplicial-lemma-sr-when-fibre-products} から従う。
次に関手
$U_{n, Zar} \to \textit{LC}_{qc}/U_n$, $U \mapsto U/U_n$
および
$X_{Zar} \to \textit{LC}_{qc}/X$, $U \mapsto U/X$
を考える。これらがサイトの射を定めることは、『サイト上のコホモロジー』節
\ref{sites-cohomology-section-cohomology-LC} で見た。
したがって注意 \ref{spaces-simplicial-remark-morphism-augmentation-simplicial-sites} のように
増大と両立する単体的サイトの射を得るので、補題
\ref{spaces-simplicial-lemma-morphism-augmentation-simplicial-sites} を適用して結論を得る。
\end{proof}

\begin{lemma}
\label{spaces-simplicial-lemma-descent-sheaves-for-proper-hypercovering}
$U$ を $\textit{LC}$ の単体的対象とし、$a : U \to X$ を増大とする。
$a : U \to X$ が $X$ の固有超被覆を与えるならば、
$$
a^{-1} : \Sh(X) \to \Sh(U_{Zar})
\quad\text{および}\quad
a^{-1} : \textit{Ab}(X) \to \textit{Ab}(U_{Zar})
$$
は充満忠実であり、本質的像はカルテジアン層で、擬逆は $a_*$ により与えられる。
ここで $a : \Sh(U_{Zar}) \to \Sh(X)$ は補題
\ref{spaces-simplicial-lemma-augmentation} のものである。
\end{lemma}

\begin{proof}
集合の層について主張を証明する。これは、すでに確立した結果から
ほとんど形式的に従う。補題 \ref{spaces-simplicial-lemma-compare-simplicial-objects} の図式を考える。
『サイト上のコホモロジー』補題
\ref{sites-cohomology-lemma-describe-pullback-pi} により、関手
$(h_{-1})^{-1}$ は充満忠実で、擬逆は $h_{-1, *}$ である。
$h$ の成分 $h_n$ についても同じことが成り立つ。
補題 \ref{spaces-simplicial-lemma-morphism-simplicial-sites} における関手
$h^{-1}$ と $h_*$ の記述から、$h^{-1}$ は充満忠実で、擬逆は $h_*$ であると分かる。
$U$ は $\textit{LC}_{qc}$ における $X$ の超被覆であることに注意する
（節 \ref{spaces-simplicial-section-hypercovering} で定義した意味である）。これは
『サイト上のコホモロジー』補題
\ref{sites-cohomology-lemma-proper-surjective-is-qc-covering} による。
補題 \ref{spaces-simplicial-lemma-hypercovering-X-simple-descent-sheaves} により、
$a_{qc}^{-1}$ は充満忠実で、擬逆は $a_{qc, *}$、本質的像は
$(\textit{LC}_{qc}/U)_{total}$ 上のカルテジアン層である。
形式的な議論（図式を追うこと）により、$a^{-1}$ が充満忠実であると分かる。

\medskip\noindent
最後に、$\mathcal{G}$ を $U_{Zar}$ 上のカルテジアン層とする。
このとき $h^{-1}\mathcal{G}$ は、% P10-E0231
$\textit{LC}_{qc}/U$ 上のカルテジアン層である。
したがって、$\textit{LC}_{qc}/X$ 上のある層 $\mathcal{H}$ に対して
$h^{-1}\mathcal{G} = a_{qc}^{-1}\mathcal{H}$ である。次を計算する。
\begin{align*}
(h_{-1})^{-1}(a_*\mathcal{G})
& =
(h_{-1})^{-1}
\text{等化子}(
\xymatrix{
a_{0, *}\mathcal{G}_0
\ar@<1ex>[r] \ar@<-1ex>[r] &
a_{1, *}\mathcal{G}_1
}
) \\
& =
\text{等化子}(
\xymatrix{
(h_{-1})^{-1}a_{0, *}\mathcal{G}_0
\ar@<1ex>[r] \ar@<-1ex>[r] &
(h_{-1})^{-1}a_{1, *}\mathcal{G}_1
}
) \\
& =
\text{等化子}(
\xymatrix{
a_{qc, 0, *}h_0^{-1}\mathcal{G}_0
\ar@<1ex>[r] \ar@<-1ex>[r] &
a_{qc, 1, *}h_1^{-1}\mathcal{G}_1
}
) \\
& =
\text{等化子}(
\xymatrix{
a_{qc, 0, *}a_{qc, 0}^{-1}\mathcal{H}
\ar@<1ex>[r] \ar@<-1ex>[r] &
a_{qc, 1, *}a_{qc, 1}^{-1}\mathcal{H}
}
) \\
& =
a_{qc, *}a_{qc}^{-1}\mathcal{H} \\
& =
\mathcal{H}
\end{align*}
ここで第一の等式は補題 \ref{spaces-simplicial-lemma-augmentation} から従う。
第二の等式は $(h_{-1})^{-1}$ が完全関手であることから従う。
第三の等式は『サイト上のコホモロジー』補題
\ref{sites-cohomology-lemma-push-pull-LC} から従う
（ここで $a_0 : U_0 \to X$ と $a_1: U_1 \to X$ が固有であることを用いる）。
第四の等式は $a_{qc}^{-1}\mathcal{H} = h^{-1}\mathcal{G}$ から、
第五の等式は補題 \ref{spaces-simplicial-lemma-augmentation-site} から従い、
第六の等式は上で確認した。
$a_{qc}^{-1}\mathcal{H} = h^{-1}\mathcal{G}$ なので、
$h^{-1}\mathcal{G} \cong h^{-1}a^{-1}a_*\mathcal{G}$ を得る。
% P10-E0232
$h^{-1}$ の充満忠実性により、これで証明が終わる。
\end{proof}

\begin{lemma}
\label{spaces-simplicial-lemma-cohomological-descent-for-proper-hypercovering}
$U$ を $\textit{LC}$ の単体的対象とし、$a : U \to X$ を増大とする。
$a : U \to X$ が $X$ の固有超被覆を与えるならば、
$K \in D^+(X)$ に対して
$$
K \to Ra_*(a^{-1}K)
$$
は同型である。ここで $a : \Sh(U_{Zar}) \to \Sh(X)$ は補題
\ref{spaces-simplicial-lemma-augmentation} のものである。
\end{lemma}

\begin{proof}
補題 \ref{spaces-simplicial-lemma-compare-simplicial-objects} の図式を考える。
『サイト上のコホモロジー』補題
\ref{sites-cohomology-lemma-cohomological-descent-LC} により、
$D^+(U_n)$ 上で $Rh_{n, *}h_n^{-1}$ は恒等関手であることに注意する。
したがって補題 \ref{spaces-simplicial-lemma-direct-image-morphism-simplicial-sites} により、
$D^+(U_{Zar})$ 上で $Rh_*h^{-1}$ は恒等関手である。次を得る。
\begin{align*}
Ra_*(a^{-1}K)
& =
Ra_*Rh_*h^{-1}a^{-1}K \\
& =
Rh_{-1, *}Ra_{qc, *}a_{qc}^{-1}(h_{-1})^{-1}K \\
& =
Rh_{-1, *}(h_{-1})^{-1}K \\
& =
K
\end{align*}
% P10-E0233
第一の等式は上の議論による。第二の等式は補題
\ref{spaces-simplicial-lemma-compare-simplicial-objects} の図式の可換性による。
第三の等式は補題
\ref{spaces-simplicial-lemma-hypercovering-X-simple-descent-bounded-abelian} による
（$U$ が $\textit{LC}_{qc}$ における $X$ の超被覆であることは、
『サイト上のコホモロジー』補題
\ref{sites-cohomology-lemma-proper-surjective-is-qc-covering} による）。
最後の等式は、すでに用いた『サイト上のコホモロジー』補題
\ref{sites-cohomology-lemma-cohomological-descent-LC} による。
\end{proof}

\begin{lemma}
\label{spaces-simplicial-lemma-compute-via-proper-hypercovering}
$U$ を $\textit{LC}$ の単体的対象とし、$a : U \to X$ を増大とする。
$U$ が $X$ の固有超被覆ならば、
$$
R\Gamma(X, K) = R\Gamma(U_{Zar}, a^{-1}K)
$$
が $K \in D^+(X)$ に対して成り立つ。ここで
$a : \Sh(U_{Zar}) \to \Sh(X)$ は補題 \ref{spaces-simplicial-lemma-augmentation} のものである。
\end{lemma}

\begin{proof}
これは補題 \ref{spaces-simplicial-lemma-cohomological-descent-for-proper-hypercovering} から従う。
実際、『サイト上のコホモロジー』注意
\ref{sites-cohomology-remark-before-Leray} により
$R\Gamma(U_{Zar}, -) = R\Gamma(X, -) \circ Ra_*$ である。
\end{proof}

\begin{lemma}
\label{spaces-simplicial-lemma-proper-hypercovering-equivalence-bounded}
$U$ を $\textit{LC}$ の単体的対象とし、$a : U \to X$ を増大とする。
$\mathcal{A} \subset \textit{Ab}(U_{Zar})$ を、カルテジアン・アーベル層からなる
弱 Serre 部分圏とする。$U$ が $X$ の固有超被覆ならば、関手 $a^{-1}$ は同値
$$
D^+(X) \longrightarrow D_\mathcal{A}^+(U_{Zar})
$$
を定め、その擬逆は $Ra_*$ である。ここで
$a : \Sh(U_{Zar}) \to \Sh(X)$ は補題 \ref{spaces-simplicial-lemma-augmentation} のものである。
\end{lemma}

\begin{proof}
補題 \ref{spaces-simplicial-lemma-Serre-subcat-cartesian-modules} により、
$\mathcal{A}$ は弱 Serre 部分圏であることに注意する。
この同値は、ここまでに得た結果の形式的な帰結である。補題
\ref{spaces-simplicial-lemma-descent-sheaves-for-proper-hypercovering} と
\ref{spaces-simplicial-lemma-cohomological-descent-for-proper-hypercovering}、および
『サイト上のコホモロジー』補題
\ref{sites-cohomology-lemma-equivalence-bounded} を用いればよい。
\end{proof}

\begin{lemma}
\label{spaces-simplicial-lemma-spectral-sequence-proper-hypercovering}
$U$ を $\textit{LC}$ の単体的対象とし、$a : U \to X$ を増大とする。
$\mathcal{F}$ を $X$ 上のアーベル層とし、$\mathcal{F}_n$ を $U_n$ への引戻しとする。
$U$ が $X$ の固有超被覆ならば、標準スペクトル系列
$$
E_1^{p, q} = H^q(U_p, \mathcal{F}_p)
$$
が存在し、$H^{p + q}(X, \mathcal{F})$ に収束する。
\end{lemma}

\begin{proof}
補題 \ref{spaces-simplicial-lemma-compute-via-proper-hypercovering} と
\ref{spaces-simplicial-lemma-simplicial-sheaf-cohomology} の直ちに分かる帰結である。
\end{proof}

\section{単体的スキーム}
\label{spaces-simplicial-section-simplicial}

\noindent
{\it 単体的スキーム}とは、スキームの圏における単体的対象である。
『単体的対象』定義 \ref{simplicial-definition-simplicial-object} を参照されたい。
単体的スキームは次のような形をしていることを思い出そう。
$$
\xymatrix{
X_2
\ar@<2ex>[r]
\ar@<0ex>[r]
\ar@<-2ex>[r]
&
X_1
\ar@<1ex>[r]
\ar@<-1ex>[r]
\ar@<1ex>[l]
\ar@<-1ex>[l]
&
X_0
\ar@<0ex>[l]
}
$$
ここには二つの射 $d^1_0, d^1_1 : X_1 \to X_0$ と、一つの射
$s^0_0 : X_0 \to X_1$ があり、以下同様である。
これらの射は、例えば
$d^1_0 \circ s^0_0 = \text{id}_{X_0} = d^1_1 \circ s^0_0$ のような
所定の関係式を満たす。『単体的対象』補題
\ref{simplicial-lemma-characterize-simplicial-object} を参照されたい。
$d^n_i : X_n \to X_{n - 1}$ を「第 $i$ 座標を忘れる射影」、
$s^n_j : X_n \to X_{n + 1}$ を「第 $j$ 座標を繰り返す対角写像」
と考えると有用である。

\medskip\noindent
単体的スキームの {\it 射} $h : X \to Y$ は、スキームの圏における
単体的対象の射と同じものである。『単体的対象』定義
\ref{simplicial-definition-simplicial-object} を参照されたい。
したがって $h$ は、意味をもつすべての場合に
$h_{n - 1} \circ d^n_j = d^n_j \circ h_n$ および
$h_{n + 1} \circ s^n_j = s^n_j \circ h_n$ を満たすスキームの射
$h_n : X_n \to Y_n$ からなる。

\medskip\noindent
単体的スキーム $X$ の {\it 増大}とは、
$a_0 \circ d^1_0 = a_0 \circ d^1_1$ を満たすスキームの射
$a_0 : X_0 \to S$ である。『単体的対象』節
\ref{simplicial-section-augmentation} を参照されたい。

\medskip\noindent
$X$ を単体的スキームとする。基礎となる単体的位相空間に節
\ref{spaces-simplicial-section-simplicial-top} の構成を適用すると、サイト $X_{Zar}$ を得る。
一方、各 $n$ に対して小 Zariski サイト $X_{n, Zar}$ があり
（『位相』定義 \ref{topologies-definition-big-small-Zariski}）、
各射 $\varphi : [m] \to [n]$ に対して、スキームの射
$X(\varphi) : X_n \to X_m$ に付随するサイトの射
$f_\varphi = X(\varphi)_{small} : X_{n, Zar} \to X_{m, Zar}$ がある
（『位相』補題 \ref{topologies-lemma-morphism-big-small}）。
これにより、サイトの圏の単体的対象 $\mathcal{C}$ を得る。
補題 \ref{spaces-simplicial-lemma-simplicial-site-site} で、付随するサイト
$\mathcal{C}_{total}$ を構成した。開埋め込みにその像を対応させると、
層を同一視する同値 $\mathcal{C}_{total} \to X_{Zar}$、すなわち
$\Sh(\mathcal{C}_{total}) = \Sh(X_{Zar})$ を得る。
$\mathcal{C}_{total}$ と $X_{Zar}$ の違いは、小 Zariski サイト
$S_{Zar}$ と $S$ の基礎位相空間との違いに似ている。
以下では、これらのサイトを暗黙に同一視する。

\medskip\noindent
$X_{Zar}$ を単体的スキーム $X$ に付随するサイトとする。
$X_{Zar}$ 上には環の層 $\mathcal{O}$ があり、その $X_n$ への制限は
構造層 $\mathcal{O}_{X_n}$ である。これは補題
\ref{spaces-simplicial-lemma-describe-sheaves-simplicial-site} または補題
\ref{spaces-simplicial-lemma-describe-sheaves-simplicial-site-site} から従う。
{\it $\mathcal{O}$ は単体的スキーム $X$ の構造層である}という。
この時点で、単体的（環付き）サイトについて展開したすべての内容が適用できる。
節 \ref{spaces-simplicial-section-simplicial-sites}、
\ref{spaces-simplicial-section-augmentation-simplicial-sites}、
\ref{spaces-simplicial-section-morphism-simplicial-sites}、
\ref{spaces-simplicial-section-simplicial-sites-modules}、
\ref{spaces-simplicial-section-cohomology-simplicial-sites}、
\ref{spaces-simplicial-section-cohomology-augmentation-simplicial-sites}、
\ref{spaces-simplicial-section-cohomology-simplicial-sites-modules}、
\ref{spaces-simplicial-section-cohomology-augmentation-ringed-simplicial-sites}、
\ref{spaces-simplicial-section-cartesian}、
\ref{spaces-simplicial-section-glueing}、および
\ref{spaces-simplicial-section-glueing-modules} を参照されたい。

\medskip\noindent
$X$ を構造層 $\mathcal{O}$ をもつ単体的スキームとする。
任意の環付きトポスの場合と同様に、$X_{Zar}$ 上の
{\it 準連接 $\mathcal{O}$-加群}という概念がある。
『サイト上の加群』定義 \ref{sites-modules-definition-site-local} を参照されたい。
しかし、$X_{Zar}$ 上の準連接 $\mathcal{O}$-加群とは、各制限
$\mathcal{F}_n$ が $X_n$ 上準連接であるカルテジアン
$\mathcal{O}$-加群 $\mathcal{F}$ にほかならない。
補題 \ref{spaces-simplicial-lemma-quasi-coherent-sheaf} を参照されたい。

\medskip\noindent
$h : X \to Y$ を単体的スキームの射とする。補題
\ref{spaces-simplicial-lemma-simplicial-space-site-functorial} または補題
\ref{spaces-simplicial-lemma-morphism-simplicial-sites}（の証明）により、サイトの射
$h_{Zar} : X_{Zar} \to Y_{Zar}$ を得る。
$h_{Zar}^{-1}$ と $h_{Zar, *}$ が成分を用いて簡明に記述されることを思い出そう。
補題 \ref{spaces-simplicial-lemma-describe-functoriality} または補題
\ref{spaces-simplicial-lemma-morphism-simplicial-sites} を参照されたい。
$\mathcal{O}_X$ と $\mathcal{O}_Y$ を、それぞれ $X$ と $Y$ の構造層とする。
% P10-E0234
$h_{Zar}^\sharp : h_{Zar, *}\mathcal{O}_X \to \mathcal{O}_Y$ を、
$Y_n$ 上で $h_n^\sharp : h_{n, *}\mathcal{O}_{X_n} \to \mathcal{O}_{Y_n}$
により与えられる $Y_{Zar}$ 上の環の層の写像として定義する。
環付きサイトの射
$$
h_{Zar} : (X_{Zar}, \mathcal{O}_X) \longrightarrow (Y_{Zar}, \mathcal{O}_Y)
$$
を得る。

\medskip\noindent
$X$ を構造層 $\mathcal{O}$ をもつ単体的スキームとする。
$S$ をスキームとし、$a_0 : X_0 \to S$ を $X$ の増大とする。
補題 \ref{spaces-simplicial-lemma-augmentation} または補題
\ref{spaces-simplicial-lemma-augmentation-site} により、対応するトポスの射
$a : \Sh(X_{Zar}) \to \Sh(S)$ を得る。
$a^{-1}\mathcal{G}$ は、成分が $a_n^{-1}\mathcal{G}$ である
$X_{Zar}$ 上の層であることに注意する。したがって写像
$a_n^\sharp : a_n^{-1}\mathcal{O}_S \to \mathcal{O}_{X_n}$ を用いて、
写像 $a^\sharp : a^{-1}\mathcal{O}_S \to \mathcal{O}$、または同値に随伴により
写像 $a^\sharp : \mathcal{O}_S \to a_*\mathcal{O}$ を定義できる
（通常どおり同じ名前を用いる）。これで節
\ref{spaces-simplicial-section-cohomology-augmentation-ringed-simplicial-sites} で論じた状況となる。
したがって環付きトポスの射
$$
a : (\Sh(X_{Zar}), \mathcal{O}) \longrightarrow (\Sh(S), \mathcal{O}_S)
$$
を得る。

\medskip\noindent
最後に次を注意する。構造層 $\mathcal{O}_X$、$\mathcal{O}_Y$ をもつ
単体的スキーム $X$、$Y$ の射 $h : X \to Y$、増大
$a_0 : X_0 \to X_{-1}$、$b_0 : Y_0 \to Y_{-1}$、および射
$h_{-1} : X_{-1} \to Y_{-1}$ が与えられ、図式
$$
\xymatrix{
X_0 \ar[r]_{h_0} \ar[d]_{a_0} & Y_0 \ar[d]^{b_0} \\
X_{-1} \ar[r]^{h_{-1}} & Y_{-1}
}
$$
が可換であると仮定する。このとき、上で明らかにした構成から、
次の環付きトポスの射の可換図式を得る。
$$
\xymatrix{
(\Sh(X_{Zar}), \mathcal{O}_X) \ar[r]_{h_{Zar}} \ar[d]_a &
(\Sh(Y_{Zar}), \mathcal{O}_Y) \ar[d]^b \\
(\Sh(X_{-1}), \mathcal{O}_{X_{-1}}) \ar[r]^{h_{-1}} &
(\Sh(Y_{-1}), \mathcal{O}_{Y_{-1}})
}
$$

\section{単体的スキームによる下降の記述}
\label{spaces-simplicial-section-simplicial-descent}

\noindent
カルテジアン射を次のように定義する。

\begin{definition}
\label{spaces-simplicial-definition-cartesian-morphism}
$a : Y \to X$ を単体的スキームの射とする。
$\Delta$ の各射 $\varphi : [n] \to [m]$ に対応する図式
$$
\xymatrix{
Y_m \ar[r]_a \ar[d]_{Y(\varphi)} & X_m \ar[d]^{X(\varphi)}\\
Y_n \ar[r]^{a} & X_n
}
$$
がスキームの圏におけるファイバー平方であるとき、$a$ は
{\it カルテジアン}である、または {\it $Y$ は $X$ 上カルテジアンである}という。
\end{definition}

\noindent
カルテジアン射は下降データと関係する。まず、固定した単体的スキーム上の
カルテジアン単体的スキームの圏を記述する一般的な補題を証明する。
この補題では、スキームの射 $f : Y \to X$ に付随する基底変換関手を
$f^* : \Sch/X \to \Sch/Y$ と書く。

\begin{lemma}
\label{spaces-simplicial-lemma-characterize-cartesian-schemes}
$X$ を単体的スキームとする。$X$ 上カルテジアンな単体的スキームの圏は、
対 $(V, \varphi)$ の圏と同値である。ここで $V$ は $X_0$ 上のスキームであり、
$$
\varphi :
V \times_{X_0, d^1_1} X_1
\longrightarrow
X_1 \times_{d^1_0, X_0} V
$$
は $X_1$ 上の同型であって、$(s_0^0)^*\varphi = \text{id}_V$ および
$$
(d^2_1)^*\varphi = (d^2_0)^*\varphi \circ (d^2_2)^*\varphi
$$
を $X_2$ 上のスキームの射として満たす。
\end{lemma}

\begin{proof}
表示された等式の主張が意味をもつのは、$X_2 \to X_0$ の射として
$d^1_1 \circ d^2_2 = d^1_1 \circ d^2_1$、
$d^1_1 \circ d^2_0 = d^1_0 \circ d^2_2$、および
$d^1_0 \circ d^2_0 = d^1_0 \circ d^2_1$ だからである。
『単体的対象』注意 \ref{simplicial-remark-relations} を参照されたい。% P10-E0235
したがって、これらの写像は次のように描ける。
$$
\xymatrix@C=0.5em{
&
X_2 \times_{d^1_1 \circ d^2_0, X_0} V
\ar[r]_-{(d^2_0)^*\varphi} &
X_2 \times_{d^1_0 \circ d^2_0, X_0} V
\ar@{=}[rd] & \\
X_2 \times_{d^1_0 \circ d^2_2, X_0} V
\ar@{=}[ru] & & &
X_2 \times_{d^1_0 \circ d^2_1, X_0} V \\
&
X_2 \times_{d^1_1 \circ d^2_2, X_0} V
\ar[lu]^{(d^2_2)^*\varphi} \ar@{=}[r] &
X_2 \times_{d^1_1 \circ d^2_1, X_0} V
\ar[ru]_{(d^2_1)^*\varphi}
}
$$
この条件は図式が可換であることを意味する。
$X$ 上カルテジアンな単体的スキーム $Y$ が与えられたとき、
$V = Y_0$ とおき、$\varphi$ をカルテジアン構造が与える同一視の合成
$$
V \times_{X_0, d^1_1} X_1 =
Y_0 \times_{X_0, d^1_1} X_1 = Y_1 =
X_1 \times_{X_0, d^1_0} Y_0 =
X_1 \times_{X_0, d^1_0} V
$$
とおけることは明らかである。この関手が同値であることを証明するため、
擬逆を構成する。その構成は『下降』節
\ref{descent-section-descent-modules} で論じた構成と同様であり、そこから記法
$\tau^n_i : [0] \to [n]$, $0 \mapsto i$ および
$\tau^n_{ij} : [1] \to [n]$, $0 \mapsto i$, $1 \mapsto j$ を借用する。
すなわち、補題のような対 $(V, \varphi)$ が与えられたとき、
$Y_n = X_n \times_{X(\tau^n_n), X_0} V$ とおく。
次に $\beta : [n] \to [m]$ が与えられると、% P10-E0236
$V(\beta) : Y_m \to Y_n$ を、写像 $\varphi$ と射影
$X_m \times_{X(\beta), X_n} Y_n \to Y_n$ との合成を
$X(\tau^m_{\beta(n)m})$ によって引き戻したものとして定義する。
これが意味をもつのは、
$$
X_m \times_{X(\tau^m_{\beta(n)m}), X_1} X_1 \times_{d^1_1, X_0} V
=
X_m \times_{X(\tau^m_m), X_0} V = Y_m
$$
および
$$
X_m \times_{X(\tau^m_{\beta(n)m}), X_1} X_1 \times_{d^1_0, X_0} V =
X_m \times_{X(\tau^m_{\beta(n)}), X_0} V =
X_m \times_{X(\beta), X_n} Y_n.
$$
だからである。上の表示図式の可換性から写像が正しく合成されることの確認は省略する。
二つの関手が互いに擬逆であることの確認も省略する。
\end{proof}

\begin{definition}
\label{spaces-simplicial-definition-fibre-products-simplicial-scheme}
$f : X \to S$ をスキームの射とする。$f$ に {\it 付随する単体的スキーム}とは、
$(X/S)_\bullet$ と書かれる関手
$\Delta^{opp} \to \Sch$, $[n] \mapsto X \times_S \ldots \times_S X$
で、『単体的対象』例
\ref{simplicial-example-fibre-products-simplicial-object} で記述されたものをいう。
\end{definition}

\noindent
したがって $(X/S)_n$ は、$S$ 上の $X$ の $(n + 1)$ 重ファイバー積である。
射 $d^1_0 : X \times_S X \to X$ は写像
$(x_0, x_1) \mapsto x_1$ であり、射 $d^1_1$ はもう一方の射影である。
射 $s^0_0$ は対角射 $X \to X \times_S X$ である。

\begin{lemma}
\label{spaces-simplicial-lemma-cartesian-over}
$f : X \to S$ をスキームの射とする。
$\pi : Y \to (X/S)_\bullet$ を単体的スキームのカルテジアン射とする。
$V = Y_0$ とおき、$X$ 上のスキームとみなす。
射 $d^1_0, d^1_1 : Y_1 \to Y_0$ と射
$\pi_1 : Y_1 \to X \times_S X$ は同型
$$
\xymatrix{
V \times_S X & &
Y_1 \ar[ll]_-{(d^1_1, \text{pr}_1 \circ \pi_1)}
\ar[rr]^-{(\text{pr}_0 \circ \pi_1, d^1_0)} & &
X \times_S V.
}
$$
を誘導する。% P10-E0237
得られる同型を $\varphi : V \times_S X \to X \times_S V$ と書く。
このとき対 $(V, \varphi)$ は $X \to S$ に関する下降データである。
\end{lemma}

\begin{proof}
これは補題 \ref{spaces-simplicial-lemma-characterize-cartesian-schemes} の（一部の）特別な場合である。
実際、その補題の表示等式は『下降』定義
\ref{descent-definition-descent-datum} の余サイクル条件と同値である。
\end{proof}

\begin{lemma}
\label{spaces-simplicial-lemma-cartesian-equivalent-descent-datum}
$f : X \to S$ をスキームの射とする。補題 \ref{spaces-simplicial-lemma-cartesian-over} の構成
$$
\begin{matrix}
\text{カルテジアン・スキームの圏} \\
\text{（} (X/S)_\bullet \text{ 上）}
\end{matrix}
\longrightarrow
\begin{matrix}
\text{下降データの圏} \\
\text{（}X/S\text{ に関する）}
\end{matrix}
$$
は圏の同値である。
\end{lemma}

\begin{proof}
左から右への関手は補題 \ref{spaces-simplicial-lemma-cartesian-over} で与えられる。
したがって、これは補題 \ref{spaces-simplicial-lemma-characterize-cartesian-schemes} の特別な場合である。
\end{proof}

\noindent
『下降』補題 \ref{descent-lemma-pullback} の引戻しは、次のように再解釈できる。
単体的スキームの射 $f : X' \to X$ と、単体的スキームのカルテジアン射
$Y \to X$ が与えられているとする。このとき、単体的スキームのファイバー積
（「引戻し」とみなす）
$$
f^*Y = Y \times_X X'
$$
は $X'$ 上カルテジアンな単体的スキームである。
スキームの射の可換図式
$$
\xymatrix{
X' \ar[r]_f \ar[d] & X \ar[d] \\
S' \ar[r] & S.
}
$$
が与えられているとする。これは単体的スキームの射
$$
f_\bullet : (X'/S')_\bullet \longrightarrow (X/S)_\bullet.
$$
を誘導する。射 $f_\bullet : (X'/S')_\bullet \to (X/S)_\bullet$ に沿う
「引戻し」$f_\bullet^*$ は、補題
\ref{spaces-simplicial-lemma-cartesian-equivalent-descent-datum} を介して、前述の
『下降』補題 \ref{descent-lemma-pullback} において下降データにより定義された
引戻しに対応すると主張する。

\section{単体的スキーム上の準連接加群}
\label{spaces-simplicial-section-modules-simplicial}

\begin{lemma}
\label{spaces-simplicial-lemma-pullback-cartesian-module}
$f : V \to U$ を単体的スキームの射とする。
$U_{Zar}$ 上の準連接加群 $\mathcal{F}$ が与えられると、引戻し
$f^*\mathcal{F}$ は $V_{Zar}$ 上の準連接加群である。
\end{lemma}

\begin{proof}
$\mathcal{F}$ はカルテジアンで、$\mathcal{F}_n$ は準連接であることを思い出そう。
補題 \ref{spaces-simplicial-lemma-quasi-coherent-sheaf} を参照されたい。
補題 \ref{spaces-simplicial-lemma-describe-functoriality} により、
$(f^*\mathcal{F})_n = f_n^*\mathcal{F}_n$ である
（細部をいくらか省略した）。したがって $(f^*\mathcal{F})_n$ は準連接である。
同じ事実と $\mathcal{F}$ のカルテジアン性から、
$f^*\mathcal{F}$ のカルテジアン性が従う。% P10-E0238
したがって、再び補題 \ref{spaces-simplicial-lemma-quasi-coherent-sheaf} により、
$\mathcal{F}$ は準連接である。
\end{proof}

\begin{lemma}
\label{spaces-simplicial-lemma-pushforward-cartesian-module}
$f : V \to U$ を単体的スキームのカルテジアン射とする。
射 $d^n_j : U_n \to U_{n - 1}$ は平坦であり、射
$V_n \to U_n$ は準コンパクトかつ準分離であると仮定する。
$V_{Zar}$ 上の準連接加群 $\mathcal{G}$ に対して、直像
$f_*\mathcal{G}$ は $U_{Zar}$ 上の準連接加群である。
\end{lemma}

\begin{proof}
$\mathcal{F} = f_* \mathcal{G}$ とおくと、補題
\ref{spaces-simplicial-lemma-describe-functoriality} により
$\mathcal{F}_n = f_{n , *}\mathcal{G}_n$ である。
写像 $\mathcal{F}(\varphi)$ は基底変換写像を用いて定義される。
『コホモロジー』節 \ref{cohomology-section-base-change-map} を参照されたい。
層 $\mathcal{F}_n$ が準連接であることは、『スキーム』補題
\ref{schemes-lemma-push-forward-quasi-coherent} と、補題
\ref{spaces-simplicial-lemma-quasi-coherent-sheaf} により $\mathcal{G}_n$ が準連接であることから従う。
% P10-E0239
退化写像 $d^n_j$ に沿う基底変換写像は、『スキームのコホモロジー』補題
\ref{coherent-lemma-flat-base-change-cohomology} と、補題
\ref{spaces-simplicial-lemma-quasi-coherent-sheaf} により $\mathcal{G}$ がカルテジアンであることから、
同型である。したがって補題 \ref{spaces-simplicial-lemma-check-cartesian-module} により、
$\mathcal{F}$ はカルテジアンである。ゆえに補題
\ref{spaces-simplicial-lemma-quasi-coherent-sheaf} により、$\mathcal{F}$ は準連接である。
\end{proof}

\begin{lemma}
\label{spaces-simplicial-lemma-adjoint-functors-cartesian-modules}
$f : V \to U$ を単体的スキームのカルテジアン射とする。
射 $d^n_j : U_n \to U_{n - 1}$ は平坦であり、射
$V_n \to U_n$ は準コンパクトかつ準分離であると仮定する。
このとき $f^*$ と $f_*$ は、$U_{Zar}$ と $V_{Zar}$ 上の準連接加群の圏の間で
随伴関手の対をなす。
\end{lemma}

\begin{proof}
補題 \ref{spaces-simplicial-lemma-pullback-cartesian-module} と
\ref{spaces-simplicial-lemma-pushforward-cartesian-module} で、主張が意味をもつことを確認した。
随伴性は、各 $f_n^*$ が $f_{n, *}$ の左随伴であることから直ちに従う。
\end{proof}

\begin{lemma}
\label{spaces-simplicial-lemma-cartesian-modules-with-section}
$f : X \to S$ を切断をもつスキームの射とする\footnote{実際には、
準連接加群の下降があるため、$f$ が $S$ 上 fpqc 局所的に切断をもつと
仮定すれば十分である。『下降』節
\ref{descent-section-fpqc-descent-quasi-coherent} を参照されたい。}。
$(X/S)_\bullet$ を $X \to S$ に付随する単体的スキームとする。
定義 \ref{spaces-simplicial-definition-fibre-products-simplicial-scheme} を参照されたい。
このとき引戻しは、準連接 $\mathcal{O}_S$-加群の圏と、
$((X/S)_\bullet)_{Zar}$ 上の準連接加群の圏との同値を定める。
\end{lemma}

\begin{proof}
$\sigma : S \to X$ を $f$ の切断とする。$(\mathcal{F}, \alpha)$ を補題
\ref{spaces-simplicial-lemma-characterize-cartesian-modules} のような対とする。
$\mathcal{G} = \sigma^*\mathcal{F}$ とおく。図式
$$
\xymatrix{
X \ar[r]_-{(\sigma \circ f, 1)} \ar[d]_f &
X \times_S X \ar[d]^{\text{pr}_0} \ar[r]_-{\text{pr}_1} & X \\
S \ar[r]^\sigma & X
}
$$
を考える。$\text{pr}_0 = d^1_1$ および $\text{pr}_1 = d^1_0$ であることに注意する。
したがって $(\sigma \circ f, 1)^*\alpha$ は同型
$$
f^*\mathcal{G} = (\sigma \circ f, 1)^*\text{pr}_0^*\mathcal{F}
\longrightarrow
(\sigma \circ f, 1)^*\text{pr}_1^*\mathcal{F} = \mathcal{F}
$$
を定める。この同型が $\alpha$ および標準同型
$\text{pr}_0^*f^*\mathcal{G} \to \text{pr}_1^*f^*\mathcal{G}$ と
両立することの確認は省略する。
\end{proof}

\section{群亜群と単体的スキーム}
\label{spaces-simplicial-section-groupoids-simplicial}

\noindent
スキームにおける群亜群が与えられると、単体的スキームを構成できる。
群亜群上の準連接層の圏は、付随する単体的スキーム上の
カルテジアンな準連接層の圏と同値になることが後に分かる。

\begin{lemma}
\label{spaces-simplicial-lemma-groupoid-simplicial}
$(U, R, s, t, c, e, i)$ を $S$ 上の群亜群スキームとする。
次の性質をもつ $S$ 上の単体的スキーム $X$ が存在する。
\begin{enumerate}
\item $X_0 = U$, $X_1 = R$, $X_2 = R \times_{s, U, t} R$,
\item $s_0^0 = e : X_0 \to X_1$,
\item $d^1_0 = s : X_1 \to X_0$, $d^1_1 = t : X_1 \to X_0$,
\item $s_0^1 = (e \circ t, 1) : X_1 \to X_2$,
% P10-E0240
$s_1^1 = (1, e \circ t) : X_1 \to X_2$,
\item $d^2_0 = \text{pr}_1 : X_2 \to X_1$,
$d^2_1 = c : X_2 \to X_1$,
$d^2_2 = \text{pr}_0$, および
\item $X = \text{cosk}_2 \text{sk}_2 X$.
\end{enumerate}
すべての $n$ に対して
$X_n = R \times_{s, U, t} \ldots \times_{s, U, t} R$ であり、
右辺は $n$ 個の因子をもつ。写像 $d^n_j : X_n \to X_{n - 1}$ は、
点の関手上で
$$
(r_1, \ldots, r_n) \longmapsto (r_1, \ldots, c(r_j, r_{j + 1}), \ldots, r_n)
$$
と与えられる（$1 \leq j \leq n - 1$）。一方、
$d^n_0(r_1, \ldots, r_n) = (r_2, \ldots, r_n)$ および
$d^n_n(r_1, \ldots, r_n) = (r_1, \ldots, r_{n - 1})$ である。
\end{lemma}

\begin{proof}
(1)、(2)、(3)、(4)、(5) で指定された規則が $S$ 上の
$2$-切断単体的スキーム $U'$ を定めることだけを確認すればよい。
実際、そのとき (6) により $X = \text{cosk}_2 U'$ とおける。
『単体的対象』補題 \ref{simplicial-lemma-existence-cosk} を参照されたい。
点の関手による方法を用いると、確認すべきことは、
$(\text{対象}, \text{射}, s, t, c, e, i)$ が群亜群ならば
$$
\xymatrix{
\text{射} \times_{s, \text{対象}, t} \text{射}
\ar@<8ex>[d]^{\text{pr}_0}
\ar@<0ex>[d]_c
\ar@<-8ex>[d]_{\text{pr}_1}
\\
\text{射}
\ar@<4ex>[d]^t
\ar@<-4ex>[d]_s
\ar@<4ex>[u]^{1, e}
\ar@<-4ex>[u]_{e, 1}
\\
\text{対象}
\ar@<0ex>[u]_e
}
$$
が $2$-切断単体集合であることだけである。細部は省略する。

\medskip\noindent
最後に、$n > 2$ に対する $X_n$ の記述は、$X_0$, $X_1$, $X_2$ の記述と、
『単体的対象』注意 \ref{simplicial-remark-inductive-coskeleton} および
補題 \ref{simplicial-lemma-work-out} から帰納法により従う。別の方法として、
上に表示した $2$-切断単体集合に $\text{cosk}_2$ を施すと、
その第 $n$ 項が
$\text{射} \times_{s, \text{対象}, t} \ldots \times_{s, \text{対象}, t}
\text{射}$（$n$ 個の因子）に等しく、% P10-E0241
退化写像が補題で与えられたとおりである単体集合が得られることを示してもよい。
細部をいくらか省略する。
\end{proof}

\begin{lemma}
\label{spaces-simplicial-lemma-quasi-coherent-groupoid-simplicial}
$S$ をスキームとする。$(U, R, s, t, c)$ を $S$ 上の群亜群スキームとする。
$X$ を補題 \ref{spaces-simplicial-lemma-groupoid-simplicial} で構成された $S$ 上の単体的スキームとする。
このとき $(U, R, s, t, c)$ 上の準連接加群の圏は、
$X_{Zar}$ 上の準連接加群の圏と同値である。
\end{lemma}

\begin{proof}
これは補題 \ref{spaces-simplicial-lemma-quasi-coherent-sheaf}、
\ref{spaces-simplicial-lemma-characterize-cartesian-modules}、および
『群亜群』定義 \ref{groupoids-definition-groupoid-module} から明らかである。
\end{proof}

\noindent
次の補題では、定義 \ref{spaces-simplicial-definition-cartesian-morphism} で定めた
単体的スキームのカルテジアン射 $V \to U$ という概念を用いる。

\begin{lemma}
\label{spaces-simplicial-lemma-quasi-coherent-groupoid-R-cartesian}
$(U, R, s, t, c)$ をスキーム $S$ 上の群亜群スキームとする。
$X$ を補題 \ref{spaces-simplicial-lemma-groupoid-simplicial} で構成された $S$ 上の単体的スキームとする。
$(R/U)_\bullet$ を $s : R \to U$ に付随する単体的スキームとする。
定義 \ref{spaces-simplicial-definition-fibre-products-simplicial-scheme} を参照されたい。
低次数の射が
$$
\xymatrix{
R \times_{s, U, s} R \times_{s, U, s} R
\ar@<3ex>[r]_-{\text{pr}_{12}}
\ar@<0ex>[r]_-{\text{pr}_{02}}
\ar@<-3ex>[r]_-{\text{pr}_{01}}
\ar[dd]_{(r_0, r_1, r_2) \mapsto (r_0 \circ r_1^{-1}, r_1 \circ r_2^{-1})} &
R \times_{s, U, s} R
\ar@<1ex>[r]_-{\text{pr}_1} \ar@<-2ex>[r]_-{\text{pr}_0}
\ar[dd]_{(r_0, r_1) \mapsto r_0 \circ r_1^{-1}} &
R \ar[dd]^t
\\
\\
R \times_{s, U, t} R
\ar@<3ex>[r]_{\text{pr}_1}
\ar@<0ex>[r]_c
\ar@<-3ex>[r]_{\text{pr}_0} &
R \ar@<1ex>[r]_s \ar@<-2ex>[r]_t &
U
}
$$
で与えられる、単体的スキームのカルテジアン射
$t_\bullet : (R/U)_\bullet \to X$ が存在する。
\end{lemma}

\begin{proof}
% P10-E0242
任意の $n$ に対し、$(R/U)_\bullet \to X_n$ を規則
$$
(r_0, \ldots, r_n)
\longrightarrow
(r_0 \circ r_1^{-1}, \ldots, r_{n - 1} \circ r_n^{-1})
$$
により定義する。% P10-E0243
退化写像との両立性は、補題 \ref{spaces-simplicial-lemma-groupoid-simplicial} における
退化写像の記述から明らかである。
写像が射 $s^n_j$ を保つことの確認は省略する。
『群亜群』補題 \ref{groupoids-lemma-diagram-pull}
（$s$ と $t$ の役割を逆にする）から、右側の二つの正方形はカルテジアンである。
まったく同じ仕方で、他のすべての正方形もカルテジアンであることが分かる。
したがって、この射はカルテジアンである。
\end{proof}

\section{下降データが与える同値関係}
\label{spaces-simplicial-section-equivalence-relation}

\noindent
節 \ref{spaces-simplicial-section-simplicial-descent} では、$X \to S$ に関する下降データを、
$(X/S)_\bullet$ 上のカルテジアンな単体的スキームを用いて
どのように定式化できるかを見た。ここでは、これを次のように同値関係と結びつける。

\begin{lemma}
\label{spaces-simplicial-lemma-equivalence-relation}
$f : X \to S$ をスキームの射とする。
$\pi : Y \to (X/S)_\bullet$ を単体的スキームのカルテジアン射とする。
定義 \ref{spaces-simplicial-definition-cartesian-morphism} および
\ref{spaces-simplicial-definition-fibre-products-simplicial-scheme} を参照されたい。
このとき射
$$
j = (d^1_1, d^1_0) : Y_1 \to Y_0 \times_S Y_0
$$
は $S$ 上の $Y_0$ に同値関係を定める。
『群亜群』定義 \ref{groupoids-definition-equivalence-relation} を参照されたい。
\end{lemma}

\begin{proof}
$j$ はモノ射であることに注意する。実際、合成
$Y_1 \to Y_0 \times_S Y_0 \to Y_0 \times_S X$ は、
$\pi$ がカルテジアンであるため同型である。

\medskip\noindent
射
$$
(d^2_2, d^2_0) : Y_2 \to Y_1 \times_{d^1_0, Y_0, d^1_1} Y_1.
$$
を考える。これが意味をもつのは $d_0 \circ d_2 = d_1 \circ d_0$
だからである。『単体的対象』注意 \ref{simplicial-remark-relations} を参照されたい。
また、これは $(X/S)_2$ 上の射である。
$Y \to (X/S)_\bullet$ はカルテジアンなので、この射は同型である。
例えば右辺は、$\pi$ がカルテジアンであることから
$
Y_0 \times_{\pi_0, X, \text{pr}_1} (X \times_S X \times_S X) =
X \times_S Y_0 \times_S X
$
と同型であることに注意する。細部は省略する。

\medskip\noindent
『群亜群』定義 \ref{groupoids-definition-equivalence-relation} と同様に、
$t = \text{pr}_0 \circ j = d^1_1$ および
$s = \text{pr}_1 \circ j = d^1_0$ と書く。
上の同型と射 $d^2_1 : Y_2 \to Y_1$ を合わせると、% P10-E0244
$Y_0 \times_S Y_0$ 上の合成射
$$
c : Y_1 \times_{s, Y_0, t} Y_1 \longrightarrow Y_1
$$
が得られる。これから直ちに、任意のスキーム $T/S$ に対して関係
$Y_1(T) \subset Y_0(T) \times Y_0(T)$ が推移的であることが従う。

\medskip\noindent
反射性は、射 $j$ の対角射 $\Delta : X \to X \times_S X$ への制限が
同型であることから従う（ここでも $\pi$ のカルテジアン性を用いる）。

\medskip\noindent
対称性を見るため、射
$$
(d^2_2, d^2_1) : Y_2 \to Y_1 \times_{d^1_1, Y_0, d^1_1} Y_1.
$$
を考える。これが意味をもつのは $d_1 \circ d_2 = d_1 \circ d_1$
だからである。『単体的対象』注意 \ref{simplicial-remark-relations} を参照されたい。
$Y \to (X/S)_\bullet$ がカルテジアンなので、この射は同型である。
例えば右辺は、$\pi$ がカルテジアンであることから
$
Y_0 \times_{\pi_0, X, \text{pr}_0} (X \times_S X \times_S X) =
Y_0 \times_S X \times_S X
$
と同型であることに注意する。細部は省略する。

\medskip\noindent
$T/S$ をスキームとする。$a, b \in Y_0(T)$ に対して $a \sim b$ であるとは、
$(a, b) \in Y_1(T)$ であることと同義であるとする。
上の同型 $(d^2_2, d^2_1)$ から、$a \sim b$ かつ $a \sim c$ ならば
$b \sim c$ であることが従う。これと反射性を合わせると、
$\sim$ が同値関係であることが分かる。
\end{proof}

\section{一つの例}
\label{spaces-simplicial-section-example}

\noindent
この節では、アルティン環のスペクトルの非交和を、
スキームの準コンパクトな全射平坦射に沿って下降させられることを示す。

\begin{lemma}
\label{spaces-simplicial-lemma-equivalence-classes-points}
$X \to S$ をスキームの射とする。
$Y \to (X/S)_\bullet$ は単体的スキームのカルテジアン射であると仮定する。
点 $y \in Y_0$ に対し、$Y_0$ の部分集合
$$
T_y = \{y' \in Y_0 \mid \exists\ y_1 \in Y_1:
d^1_1(y_1) = y, d^1_0(y_1) = y'\}
$$
を定める。このとき $y \in T_y$ であり、
$T_y \cap T_{y'} \not = \emptyset \Rightarrow T_y = T_{y'}$ である。
\end{lemma}

\begin{proof}
補題 \ref{spaces-simplicial-lemma-equivalence-relation} と『群亜群』補題
\ref{groupoids-lemma-pre-equivalence-equivalence-relation-points} を合わせればよい。
\end{proof}

\begin{lemma}
\label{spaces-simplicial-lemma-quasi-compact}
$X \to S$ をスキームの射とする。
$Y \to (X/S)_\bullet$ は単体的スキームのカルテジアン射であると仮定する。
$y \in Y_0$ を点とする。$X \to S$ が準コンパクトならば、
$$
T_y = \{y' \in Y_0 \mid \exists\ y_1 \in Y_1:
d^1_1(y_1) = y, d^1_0(y_1) = y'\}
$$
は $Y_0$ の準コンパクトな部分集合である。
\end{lemma}

\begin{proof}
$F_y$ を $y$ における $d^1_1 : Y_1 \to Y_0$ のスキーム論的ファイバーとする。
すると $T_y$ は射
$$
\xymatrix{
F_y \ar[r] \ar[d] &
Y_1 \ar[r]^{d^1_0} \ar[d]^{d^1_1} &
Y_0 \\
y \ar[r] &
Y_0 &
}
$$
の像であることが分かる。$F_y$ は準コンパクトであることに注意する。
これで補題が証明された。
\end{proof}

\begin{lemma}
\label{spaces-simplicial-lemma-descent-disjoint-union-Artinian-along-fields}
$X \to S$ を準コンパクトな平坦全射とする。
$(V, \varphi)$ を $X \to S$ に関する下降データとする。
$V$ がアルティン環のスペクトルの非交和ならば、$(V, \varphi)$ は有効である。
\end{lemma}

\begin{proof}
$Y \to (X/S)_\bullet$ を、補題
\ref{spaces-simplicial-lemma-cartesian-equivalent-descent-datum} により
$(V, \varphi)$ に対応する単体的スキームのカルテジアン射とする。
$Y_0 = V$ であることに注意する。各 $A_i$ が局所アルティン環となるように
$V = \coprod_{i \in I} \Spec(A_i)$ と書く。
さらに、各 $i \in I$ に対し $v_i \in V$ を $\Spec(A_i)$ の唯一の閉点とする。
上の補題 \ref{spaces-simplicial-lemma-equivalence-classes-points} の記法により、
$v_i \in T_{v_j}$ であることと $i \sim j$ であることを同義とする。
補題 \ref{spaces-simplicial-lemma-equivalence-classes-points} と
\ref{spaces-simplicial-lemma-quasi-compact} により、これは有限な同値類をもつ同値関係である。
$\overline{I} = I/\sim$ とおく。このとき
$V = \coprod_{\overline{i} \in \overline{I}} V_{\overline{i}}$
と書け、ここで
$V_{\overline{i}} = \coprod_{i \in \overline{i}} \Spec(A_i)$ である。
構成から、$\varphi : V \times_S X \to X \times_S V$ は
開かつ閉な部分空間 $V_{\overline{i}} \times_S X$ を
開かつ閉な部分空間 $X \times_S V_{\overline{i}}$ に写す。
言い換えると、下降データ
$(V_{\overline{i}}, \varphi_{\overline{i}})$ が得られ、
$(V, \varphi)$ は下降データの圏におけるそれらの余積である。
各 $V_{\overline{i}}$ は局所アルティン環のスペクトルの有限和なので、
射 $V_{\overline{i}} \to X$ はアフィンである。
『射』補題 \ref{morphisms-lemma-Artinian-affine} を参照されたい。
$\{X \to S\}$ は fpqc 被覆だから、すべての下降データ
$(V_{\overline{i}}, \varphi_{\overline{i}})$ は
『下降』補題 \ref{descent-lemma-affine} により有効である。
\end{proof}

\noindent
念のため述べておくと、上の補題の適用範囲は非常に限られている！

\section{単体的代数空間}
\label{spaces-simplicial-section-simplicial-algebraic-spaces}

\noindent
$S$ をスキームとする。{\it 単体的代数空間}とは、
$S$ 上の代数空間の圏における単体的対象である。
『単体的対象』定義 \ref{simplicial-definition-simplicial-object} を参照されたい。
単体的代数空間は次のような形をしていることを思い出そう。
$$
\xymatrix{
X_2
\ar@<2ex>[r]
\ar@<0ex>[r]
\ar@<-2ex>[r]
&
X_1
\ar@<1ex>[r]
\ar@<-1ex>[r]
\ar@<1ex>[l]
\ar@<-1ex>[l]
&
X_0
\ar@<0ex>[l]
}
$$
ここには二つの射 $d^1_0, d^1_1 : X_1 \to X_0$ と
一つの射 $s^0_0 : X_0 \to X_1$ があり、以下同様である。
これらの射は
$d^1_0 \circ s^0_0 = \text{id}_{X_0} = d^1_1 \circ s^0_0$ のような
所定の関係式を満たす。『単体的対象』補題
\ref{simplicial-lemma-characterize-simplicial-object} を参照されたい。
$d^n_i : X_n \to X_{n - 1}$ は「第 $i$ 座標を忘れる射影」と、
$s^n_j : X_n \to X_{n + 1}$ は「第 $j$ 座標を繰り返す対角写像」と
考えると便利である。

\medskip\noindent
単体的代数空間の {\it 射} $h : X \to Y$ とは、
$S$ 上の代数空間の圏における単体的対象の射にほかならない。
『単体的対象』定義 \ref{simplicial-definition-simplicial-object} を参照されたい。
したがって $h$ は代数空間の射 $h_n : X_n \to Y_n$ からなり、
意味をもつすべての場合に
$h_{n - 1} \circ d^n_j = d^n_j \circ h_n$ および
$h_{n + 1} \circ s^n_j = s^n_j \circ h_n$ を満たす。

\medskip\noindent
単体的代数空間 $X$ の {\it 増大} $a : X \to X_{-1}$ は、
$a_0 \circ d^1_0 = a_0 \circ d^1_1$ を満たす代数空間の射
$a_0 : X_0 \to X_{-1}$ により与えられる。
『単体的対象』節 \ref{simplicial-section-augmentation} を参照されたい。
この状況では、$n \geq 0$ に対して誘導される射を常に
$a_n : X_n \to X_{-1}$ と書く。

\medskip\noindent
$X$ を単体的代数空間とする。各 $n$ に対してサイト
$X_{n, spaces, \etale}$（『代数空間の性質』定義
\ref{spaces-properties-definition-spaces-etale-site}）があり、
各射 $\varphi : [m] \to [n]$ に対してサイトの射
$$
f_\varphi = X(\varphi)_{spaces, \etale} :
X_{n, spaces, \etale} \to X_{m, spaces, \etale},
$$
がある。これは代数空間の射 $X(\varphi) : X_n \to X_m$ に付随する
（『代数空間の性質』補題
\ref{spaces-properties-lemma-functoriality-etale-site}）。
こうしてサイトの圏における単体的対象が得られる。
補題 \ref{spaces-simplicial-lemma-simplicial-site-site} では付随するサイトを構成した。
これを $X_{spaces, \etale}$ と書く。
% P10-E0245
サイト $X_{spaces, \etale}$ の対象とは、ある $n$ に対して
$X_n$ 上エタールな代数空間 $U$ である。射
$(\varphi, f) : U/X_n \to V/X_m$ は、$\Delta$ における射
$\varphi : [m] \to [n]$ と、図式
% P10-E0246
$$
\xymatrix{
U \ar[r]_f \ar[d] & V \ar[d] \\
X_n \ar[r]^{f_\varphi} & X_m
}
$$
を可換にする代数空間の射 $f : U \to V$ により与えられる。
対象が $U/X_n$ であって、$U$ がそれぞれアフィンまたはスキームである
充満部分圏
$$
X_{affine, \etale} \subset X_\etale \subset X_{spaces, \etale}
$$
を考える。これらの圏に自然な位相を入れると
（『代数空間の性質』補題 \ref{spaces-properties-lemma-alternative}、
定義 \ref{spaces-properties-definition-etale-site}、および
補題 \ref{spaces-properties-lemma-compare-etale-sites} を参照）、
これらの包含関手はトポスの同値
$$
\Sh(X_{affine, \etale}) = \Sh(X_\etale) = \Sh(X_{spaces, \etale})
$$
を定める。以下ではこれらのトポスを断りなく同一視する。
$X_\etale$ を $X$ の {\it 小エタール・サイト}、
そのトポスを $X$ の {\it 小エタール・トポス}と呼ぶ。

\medskip\noindent
$X_\etale$ を単体的代数空間 $X$ の小エタール・サイトとする。
$X_\etale$ 上には環の層 $\mathcal{O}$ があり、
その $X_n$ への制限は構造層 $\mathcal{O}_{X_n}$ である。
これは補題 \ref{spaces-simplicial-lemma-describe-sheaves-simplicial-site-site} から従う。
{\it $\mathcal{O}$ を単体的代数空間 $X$ の構造層}と呼ぶ。
この時点で、単体的（環付き）サイトについて展開したすべての内容が適用できる。
節 \ref{spaces-simplicial-section-simplicial-sites}、
\ref{spaces-simplicial-section-augmentation-simplicial-sites}、
\ref{spaces-simplicial-section-morphism-simplicial-sites}、
\ref{spaces-simplicial-section-simplicial-sites-modules}、
\ref{spaces-simplicial-section-cohomology-simplicial-sites}、
\ref{spaces-simplicial-section-cohomology-augmentation-simplicial-sites}、
\ref{spaces-simplicial-section-cohomology-simplicial-sites-modules}、
\ref{spaces-simplicial-section-cohomology-augmentation-ringed-simplicial-sites}、
\ref{spaces-simplicial-section-cartesian}、
\ref{spaces-simplicial-section-glueing}、および
\ref{spaces-simplicial-section-glueing-modules} を参照されたい。

\medskip\noindent
$X$ を構造層 $\mathcal{O}$ をもつ単体的代数空間とする。
任意の環付きトポスと同様に、$X_\etale$ 上の
{\it 準連接 $\mathcal{O}$-加群}という概念がある。
『サイト上の加群』定義 \ref{sites-modules-definition-site-local} を参照されたい。
しかし $X_\etale$ 上の準連接 $\mathcal{O}$-加群とは、
その制限 $\mathcal{F}_n$ が $X_n$ 上準連接であるカルテジアンな
$\mathcal{O}$-加群 $\mathcal{F}$ にほかならない。
補題 \ref{spaces-simplicial-lemma-quasi-coherent-sheaf} を参照されたい。

\medskip\noindent
$h : X \to Y$ を $S$ 上の単体的代数空間の射とする。
補題 \ref{spaces-simplicial-lemma-morphism-simplicial-sites} を、サイトの射
$(h_n)_{spaces, \etale} : X_{spaces, \etale} \to Y_{spaces, \etale}$
（『代数空間の性質』補題
\ref{spaces-properties-lemma-functoriality-etale-site}）に適用すると、
小エタール・トポスの射
$h_\etale : \Sh(X_\etale) \to \Sh(Y_\etale)$ が得られる。
$h_\etale^{-1}$ と $h_{\etale, *}$ は成分を用いて簡明に記述できることを
思い出そう。補題 \ref{spaces-simplicial-lemma-morphism-simplicial-sites} を参照されたい。
$\mathcal{O}_X$, resp.\ $\mathcal{O}_Y$ を、それぞれ $X$, resp.\ $Y$ の
構造層とする。% P10-E0247
環の層の写像 $h_\etale^\sharp : h_{\etale, *}\mathcal{O}_X \to \mathcal{O}_Y$
を、$Y_n$ 上で
$h_n^\sharp : h_{n, *}\mathcal{O}_{X_n} \to \mathcal{O}_{Y_n}$ により
与えられる $Y_\etale$ 上の写像として定義する。
こうして環付きトポスの射
$$
h_\etale :
(\Sh(X_\etale), \mathcal{O}_X)
\longrightarrow
(\Sh(Y_\etale), \mathcal{O}_Y)
$$
を得る。

\medskip\noindent
$X$ を構造層 $\mathcal{O}$ をもつ単体的代数空間とする。
$X_{-1}$ を $S$ 上の代数空間とし、$a_0 : X_0 \to X_{-1}$ を
$X$ の増大とする。サイトの射
$(a_0)_{spaces, \etale} :
X_{0, spaces, \etale} \to X_{-1, spaces, \etale}$ に
補題 \ref{spaces-simplicial-lemma-augmentation-site} を適用すると、対応するトポスの射
$a : \Sh(X_\etale) \to \Sh(X_{-1, \etale})$ が得られる。
$a^{-1}\mathcal{G}$ は、成分が $a_n^{-1}\mathcal{G}$ である
$X_\etale$ 上の層であることに注意する。したがって写像
$a_n^\sharp : a_n^{-1}\mathcal{O}_{X_{-1}} \to \mathcal{O}_{X_n}$ を用いて、
写像 $a^\sharp : a^{-1}\mathcal{O}_{X_{-1}} \to \mathcal{O}$、
または同値なこととして、随伴により写像
$a^\sharp : \mathcal{O}_{X_{-1}} \to a_*\mathcal{O}$ を定義できる
（通常どおり同じ名前を用いる）。これにより、節
\ref{spaces-simplicial-section-cohomology-augmentation-ringed-simplicial-sites} で
論じた状況に入る。したがって環付きトポスの射
% P10-E0248
$$
a :
(\Sh(X_\etale), \mathcal{O})
\longrightarrow
(\Sh(X_{-1}), \mathcal{O}_{X_{-1}})
$$
を得る。

\medskip\noindent
最後に次を注意しておく。構造層 $\mathcal{O}_X$, $\mathcal{O}_Y$ をもつ
単体的代数空間 $X$, $Y$ の射 $h : X \to Y$、増大
$a_0 : X_0 \to X_{-1}$, $b_0 : Y_0 \to Y_{-1}$、および射
$h_{-1} : X_{-1} \to Y_{-1}$ が与えられ、
$$
\xymatrix{
X_0 \ar[r]_{h_0} \ar[d]_{a_0} & Y_0 \ar[d]^{b_0} \\
X_{-1} \ar[r]^{h_{-1}} & Y_{-1}
}
$$
が可換であるとする。
このとき上で明らかにした構成から、環付きトポスの射の可換図式
$$
\xymatrix{
(\Sh(X_\etale), \mathcal{O}_X) \ar[r]_{h_\etale} \ar[d]_a &
(\Sh(Y_\etale), \mathcal{O}_Y) \ar[d]^b \\
(\Sh(X_{-1}), \mathcal{O}_{X_{-1}}) \ar[r]^{h_{-1}} &
(\Sh(Y_{-1}), \mathcal{O}_{Y_{-1}})
}
$$
が得られる。

\section{代数空間の fppf 超被覆}
\label{spaces-simplicial-section-fppf-hypercovering}

\noindent
この節は、代数空間と fppf 超被覆の場合における節
\ref{spaces-simplicial-section-proper-hypercovering} の類似である。% P10-E0249
望む読者は、全体を通じて「代数空間」を「スキーム」に置き換えても、
同様に正しい結果を得られる。その利点は、『代数空間のコホモロジーについてさらに』節
\ref{spaces-more-cohomology-section-fppf-etale} への参照を、
『エタール・コホモロジー』節
\ref{etale-cohomology-section-fppf-etale} への参照で置き換えられることである。

\medskip\noindent
基礎スキーム $S$ を固定する。$X$ を $S$ 上の代数空間、
$U$ を $S$ 上の単体的代数空間とする。増大
$$
a : U \to X
$$
があると仮定する。節 \ref{spaces-simplicial-section-simplicial-algebraic-spaces} を参照されたい。
次が成り立つとき、$U$ を $X$ の {\it fppf 超被覆}と呼ぶ。
\begin{enumerate}
\item $U_0 \to X$ は平坦、局所有限表示、かつ全射である。
\item $U_1 \to U_0 \times_X U_0$ は平坦、局所有限表示、かつ全射である。
\item $n \geq 1$ に対して
$U_{n + 1} \to (\text{cosk}_n\text{sk}_n U)_{n + 1}$ は
平坦、局所有限表示、かつ全射である。
\end{enumerate}
$S$ 上の代数空間の圏はすべての有限極限をもつので、
上の定式化で用いた余骨格は存在する。
$$
\fbox{原理：fppf 超被覆を用いてエタール・コホモロジーを計算できる。}
$$
この原理の証明の背後にある重要な考えは、圏 $\textit{Spaces}/S$ 上の
fppf 位相とエタール位相を比較することである。すなわち fppf 位相は
エタール位相より強く、次が成り立つ。
(a) 平坦、局所有限表示、全射な写像は fppf 被覆を定める。
(b) 小エタール・サイトから引き戻された層の fppf コホモロジーは、
『代数空間のコホモロジーについてさらに』節
\ref{spaces-more-cohomology-section-fppf-etale} で見たように
エタール・コホモロジーと一致する。

\begin{lemma}
\label{spaces-simplicial-lemma-compare-simplicial-objects-fppf-etale}
$S$ をスキームとする。$X$ を $S$ 上の代数空間とし、
$U$ を $S$ 上の単体的代数空間とする。$a : U \to X$ を増大とする。
可換図式
$$
\xymatrix{
\Sh((\textit{Spaces}/U)_{fppf, total}) \ar[r]_-h \ar[d]_{a_{fppf}} &
\Sh(U_\etale) \ar[d]^a \\
\Sh((\textit{Spaces}/X)_{fppf}) \ar[r]^-{h_{-1}} &
\Sh(X_\etale)
}
$$
がある。左の垂直射は節 \ref{spaces-simplicial-section-hypercovering} で定義され、
右の垂直射は節 \ref{spaces-simplicial-section-simplicial-algebraic-spaces} で定義される。
\end{lemma}

\begin{proof}
記法 $(\textit{Spaces}/U)_{fppf, total}$ は、
サイト $(\textit{Spaces}/S)_{fppf}$ とこのサイトの単体的対象 $U$ に
節 \ref{spaces-simplicial-section-hypercovering} の構成を用いていることを表す
\footnote{エタール位相を用いることもでき、その場合は
$(\textit{Spaces}/U)_{\etale, total}$ と書く。}。
右辺のトポスにはサイト $X_{spaces, \etale}$ と
$U_{spaces, \etale}$ を用いる。このことが許されるのは、節
\ref{spaces-simplicial-section-simplicial-algebraic-spaces} の議論による。

\medskip\noindent
$(\textit{Spaces}/U)_{fppf, total}$ と $U_{spaces, \etale}$ はともに
情況 \ref{spaces-simplicial-situation-simplicial-site} の場合 A に該当することに注意する。
$U_\etale$ については節 \ref{spaces-simplicial-section-simplicial-algebraic-spaces} の
構成から直ちに分かり、$(\textit{Spaces}/U)_{fppf, total}$ については
補題 \ref{spaces-simplicial-lemma-sr-when-fibre-products} から従う。
次に関手\linebreak
$U_{n, spaces, \etale} \to (\textit{Spaces}/U_n)_{fppf}$,
$U \mapsto U/U_n$ および
\linebreak
$X_{spaces, \etale} \to (\textit{Spaces}/X)_{fppf}$,
$U \mapsto U/X$ を考える。これらがサイトの射を定めることは、
『代数空間のコホモロジーについてさらに』節
\ref{spaces-more-cohomology-section-fppf-etale} で見た。
% P10-E0250
そこではそれぞれ $a_{U_n} = \epsilon_{U_n} \circ \pi_{u_n}$ および
$a_X = \epsilon_X \circ \pi_X$ と書かれていた。
したがって注意 \ref{spaces-simplicial-remark-morphism-augmentation-simplicial-sites} のように
増大と両立する単体的サイトの射を得るので、補題
\ref{spaces-simplicial-lemma-morphism-augmentation-simplicial-sites} を適用して結論を得る。
\end{proof}

\begin{lemma}
\label{spaces-simplicial-lemma-descent-sheaves-for-fppf-hypercovering}
$S$ をスキームとする。$X$ を $S$ 上の代数空間とし、
$U$ を $S$ 上の単体的代数空間とする。$a : U \to X$ を増大とする。
$a : U \to X$ が $X$ の fppf 超被覆ならば、
$$
a^{-1} : \Sh(X_\etale) \to \Sh(U_\etale)
\quad\text{および}\quad
a^{-1} : \textit{Ab}(X_\etale) \to \textit{Ab}(U_\etale)
$$
は充満忠実であり、本質的像はカルテジアン層で、擬逆は $a_*$ により与えられる。
ここで $a : \Sh(U_\etale) \to \Sh(X_\etale)$ は節
\ref{spaces-simplicial-section-simplicial-algebraic-spaces} のものである。
\end{lemma}

\begin{proof}
集合の層について主張を証明する。これは、すでに確立した結果から
ほとんど形式的に従う。補題
\ref{spaces-simplicial-lemma-compare-simplicial-objects-fppf-etale} の図式を考える。
この補題の証明で、$h_{-1}$ は
『代数空間のコホモロジーについてさらに』節
\ref{spaces-more-cohomology-section-fppf-etale} の射 $a_X$ であることを見た。
したがって『代数空間のコホモロジーについてさらに』補題
\ref{spaces-more-cohomology-lemma-comparison-fppf-etale} により、
$(h_{-1})^{-1}$ は充満忠実で、擬逆は $h_{-1, *}$ である。
$h$ の成分 $h_n$ についても同じことが成り立つ。
補題 \ref{spaces-simplicial-lemma-morphism-simplicial-sites} における関手 $h^{-1}$ と
$h_*$ の記述から、$h^{-1}$ は充満忠実で、擬逆は $h_*$ であると分かる。
$U$ は節 \ref{spaces-simplicial-section-hypercovering} で定義した意味で
$(\textit{Spaces}/S)_{fppf}$ における $X$ の超被覆であることに注意する。
補題 \ref{spaces-simplicial-lemma-hypercovering-X-simple-descent-sheaves} により、
$a_{fppf}^{-1}$ は充満忠実で、擬逆は $a_{fppf, *}$、
本質的像は $(\textit{Spaces}/U)_{fppf, total}$ 上のカルテジアン層である。
形式的な議論（図式を追うこと）により、$a^{-1}$ が充満忠実であると分かる。

\medskip\noindent
最後に、$\mathcal{G}$ を $U_\etale$ 上のカルテジアン層とする。
このとき $h^{-1}\mathcal{G}$ は次のサイト上のカルテジアン層である。
\linebreak
$(\textit{Spaces}/U)_{fppf, total}$。
したがって、$(\textit{Spaces}/X)_{fppf}$ 上のある層
$\mathcal{H}$ に対して\linebreak
$h^{-1}\mathcal{G} = a_{fppf}^{-1}\mathcal{H}$ である。
とくに\linebreak
$h_0^{-1}\mathcal{G}_0 = (a_{0, big, fppf})^{-1}\mathcal{H}$ である。
$h_0 = a_{U_0}$ であり、$U_0 \to X$ は平坦、局所有限表示、かつ全射なので、
『代数空間のコホモロジーについてさらに』補題
\ref{spaces-more-cohomology-lemma-descent-sheaf-fppf-etale} から、
% P10-E0251
$X_\etale$ 上の層 $\mathcal{F}$ と同型\linebreak
$\mathcal{H} = (h_{-1})^{-1}\mathcal{F}$ が存在する。
$a_{fppf}^{-1}\mathcal{H} = h^{-1}\mathcal{G}$ だから、\linebreak
$h^{-1}\mathcal{G} \cong h^{-1}a^{-1}\mathcal{F}$ が従う。
% P10-E0252
$h^{-1}$ の充満忠実性により、
$a^{-1}\mathcal{F} \cong \mathcal{G}$ と結論する。

\medskip\noindent
同型 $\theta : a^{-1}\mathcal{F} \to \mathcal{G}$ を固定する。
証明を終えるには $\mathcal{G} = a^{-1}a_*\mathcal{G}$ を示さなければならない
（擬逆が $a_*$ により与えられることを示すためであり、
それ以外はすべて上で証明した）。
$a^{-1}$ は充満忠実だから、『圏』補題
\ref{categories-lemma-adjoint-fully-faithful} により
$\text{id} \cong a_*a^{-1}$ である。
したがって $\mathcal{F} \cong a_*a^{-1}\mathcal{F}$ と
$a_*\theta : a_*a^{-1}\mathcal{F} \to a_*\mathcal{G}$ を合わせると、
同型 $\mathcal{F} \to a_*\mathcal{G}$ が得られる。
$a$ により引き戻し、$\theta^{-1}$ を前合成すると、望む同型が得られる。
\end{proof}

\begin{lemma}
\label{spaces-simplicial-lemma-cohomological-descent-for-fppf-hypercovering}
$S$ をスキームとする。$X$ を $S$ 上の代数空間とし、
$U$ を $S$ 上の単体的代数空間とする。$a : U \to X$ を増大とする。
$a : U \to X$ が $X$ の fppf 超被覆ならば、$K \in D^+(X_\etale)$ に対して
$$
K \to Ra_*(a^{-1}K)
$$
は同型である。ここで $a : \Sh(U_\etale) \to \Sh(X_\etale)$ は節
\ref{spaces-simplicial-section-simplicial-algebraic-spaces} のものである。
\end{lemma}

\begin{proof}
補題 \ref{spaces-simplicial-lemma-compare-simplicial-objects-fppf-etale} の図式を考える。
『代数空間のコホモロジーについてさらに』補題
\ref{spaces-more-cohomology-lemma-cohomological-descent-etale-fppf} により、
$D^+(U_{n, \etale})$ 上で $Rh_{n, *}h_n^{-1}$ は恒等関手であることに注意する。
したがって補題 \ref{spaces-simplicial-lemma-direct-image-morphism-simplicial-sites} により、
$D^+(U_\etale)$ 上で $Rh_*h^{-1}$ は恒等関手である。次を得る。
\begin{align*}
Ra_*(a^{-1}K)
& =
Ra_*Rh_*h^{-1}a^{-1}K \\
& =
Rh_{-1, *}Ra_{fppf, *}a_{fppf}^{-1}(h_{-1})^{-1}K \\
& =
Rh_{-1, *}(h_{-1})^{-1}K \\
& =
K
\end{align*}
% P10-E0253
第一の等式は上の議論による。第二の等式は補題
% P10-E0254
\ref{spaces-simplicial-lemma-compare-simplicial-objects} の図式の可換性による。
第三の等式は、$U$ が $(\textit{Spaces}/S)_{fppf}$ における
$X$ の超被覆なので、補題
\ref{spaces-simplicial-lemma-hypercovering-X-simple-descent-bounded-abelian} による。
最後の等式は、すでに用いた『代数空間のコホモロジーについてさらに』補題
\ref{spaces-more-cohomology-lemma-cohomological-descent-etale-fppf} による。
\end{proof}

\begin{lemma}
\label{spaces-simplicial-lemma-compute-via-fppf-hypercovering}
$S$ をスキームとする。$X$ を $S$ 上の代数空間とし、
$U$ を $S$ 上の単体的代数空間とする。$a : U \to X$ を増大とする。
$a : U \to X$ が $X$ の fppf 超被覆ならば、
$$
R\Gamma(X_\etale, K) = R\Gamma(U_\etale, a^{-1}K)
$$
が $K \in D^+(X_\etale)$ に対して成り立つ。
ここで $a : \Sh(U_\etale) \to \Sh(X_\etale)$ は節
\ref{spaces-simplicial-section-simplicial-algebraic-spaces} のものである。
\end{lemma}

\begin{proof}
これは補題 \ref{spaces-simplicial-lemma-cohomological-descent-for-fppf-hypercovering} から従う。
実際、『サイト上のコホモロジー』注意
\ref{sites-cohomology-remark-before-Leray} により
$R\Gamma(U_\etale, -) = R\Gamma(X_\etale, -) \circ Ra_*$ である。
\end{proof}

\begin{lemma}
\label{spaces-simplicial-lemma-fppf-hypercovering-equivalence-bounded}
$S$ をスキームとする。$X$ を $S$ 上の代数空間とし、
$U$ を $S$ 上の単体的代数空間とする。$a : U \to X$ を増大とする。
$\mathcal{A} \subset \textit{Ab}(U_\etale)$ を、
カルテジアン・アーベル層からなる弱 Serre 部分圏とする。
$U$ が $X$ の fppf 超被覆ならば、関手 $a^{-1}$ は同値
$$
D^+(X_\etale) \longrightarrow D_\mathcal{A}^+(U_\etale)
$$
を定め、その擬逆は $Ra_*$ である。
ここで $a : \Sh(U_\etale) \to \Sh(X_\etale)$ は節
\ref{spaces-simplicial-section-simplicial-algebraic-spaces} のものである。
\end{lemma}

\begin{proof}
補題 \ref{spaces-simplicial-lemma-Serre-subcat-cartesian-modules} により、
$\mathcal{A}$ は弱 Serre 部分圏であることに注意する。
この同値は、ここまでに得た結果の形式的な帰結である。補題
\ref{spaces-simplicial-lemma-descent-sheaves-for-fppf-hypercovering} と
\ref{spaces-simplicial-lemma-cohomological-descent-for-fppf-hypercovering}、および
『サイト上のコホモロジー』補題
\ref{sites-cohomology-lemma-equivalence-bounded} を用いればよい。
\end{proof}

\begin{lemma}
\label{spaces-simplicial-lemma-spectral-sequence-fppf-hypercovering}
$S$ をスキームとする。$X$ を $S$ 上の代数空間とし、
$U$ を $S$ 上の単体的代数空間とする。$a : U \to X$ を増大とする。
$\mathcal{F}$ を $X_\etale$ 上のアーベル層とし、
$\mathcal{F}_n$ を $U_{n, \etale}$ への引戻しとする。
$U$ が $X$ の fppf 超被覆ならば、標準スペクトル系列
$$
E_1^{p, q} = H^q_\etale(U_p, \mathcal{F}_p)
$$
が存在し、$H^{p + q}_\etale(X, \mathcal{F})$ に収束する。
\end{lemma}

\begin{proof}
補題 \ref{spaces-simplicial-lemma-compute-via-fppf-hypercovering} と
\ref{spaces-simplicial-lemma-simplicial-sheaf-cohomology-site} の直ちに分かる帰結である。
\end{proof}

\section{代数空間の fppf 超被覆：加群}
\label{spaces-simplicial-section-fppf-hypercovering-modules}

\noindent
節 \ref{spaces-simplicial-section-fppf-hypercovering} で始めた
fppf 超被覆に対する（コホモロジー的）下降の議論を続けるが、
本節では加群の層の場合に何が起こるかを論じる。
主として準連接加群を扱い、それらについては
非有界コホモロジー的下降も行えることが分かる。

\begin{lemma}
\label{spaces-simplicial-lemma-compare-simplicial-objects-fppf-etale-modules}
$S$ をスキームとする。$X$ を $S$ 上の代数空間とする。
$U$ を $S$ 上の単体的代数空間とする。
$a : U \to X$ を増大とする。このとき環付きトポスの可換図式
$$
\xymatrix{
(\Sh((\textit{Spaces}/U)_{fppf, total}), \mathcal{O}_{big, total})
\ar[r]_-h \ar[d]_{a_{fppf}} &
(\Sh(U_\etale), \mathcal{O}_U) \ar[d]^a \\
(\Sh((\textit{Spaces}/X)_{fppf}), \mathcal{O}_{big}) \ar[r]^-{h_{-1}} &
(\Sh(X_\etale), \mathcal{O}_X)
}
$$
が存在する。ここで左の垂直射は
節 \ref{spaces-simplicial-section-hypercovering-modules} で定義され、
右の垂直射は
節 \ref{spaces-simplicial-section-simplicial-algebraic-spaces} で定義される。
\end{lemma}

\begin{proof}
基礎にあるトポスの図式については、補題
\ref{spaces-simplicial-lemma-compare-simplicial-objects-fppf-etale}
の証明における議論を参照する。
層 $\mathcal{O}_U$ は、節
\ref{spaces-simplicial-section-simplicial-algebraic-spaces} で定義された
単体的代数空間 $U$ の構造層である。
層 $\mathcal{O}_X$ は代数空間 $X$ の通常の構造層である。
環の層 $\mathcal{O}_{big, total}$ と
$\mathcal{O}_{big}$ は、節
\ref{spaces-simplicial-section-hypercovering-modules} で説明した方法により
$(\textit{Spaces}/S)_{fppf}$ 上の構造層から得られる。
同節では $a_{fppf}$ も環付きトポスの射として構成される。
成分射 $h_n = a_{U_n}$ と $h_{-1} = a_X$ は、
『代数空間のコホモロジーについてさらに』節
\ref{spaces-more-cohomology-section-fppf-etale-modules}
により環付きトポスの射である。
最後に、$h$ の定義に用いた連続関手
$u : U_{spaces, \etale} \to (\textit{Spaces}/U)_{fppf, total}$
\footnote{これは補題
\ref{spaces-simplicial-lemma-compare-simplicial-objects-fppf-etale}
の証明で、補題
\ref{spaces-simplicial-lemma-morphism-augmentation-simplicial-sites}
を適用して行った。}
は $V/U_n \mapsto V/U_n$ により与えられるので、
$h_*\mathcal{O}_{big, total} = \mathcal{O}_U$ である。
これにより、注意
\ref{spaces-simplicial-remark-morphism-simplicial-sites-modules}
の意味で $h$ に環付き単体的サイトの射の構造を与える。
すると補題
\ref{spaces-simplicial-lemma-morphism-simplicial-sites-modules}
により、$h$ を環付きトポスの射として得る。
射 $h_n$ は実際に上で述べた射 $a_{U_n}$ と一致することに注意せよ。
この図式が環付きトポスの図式として可換であることの検証は省略する
（トポスの図式として可換であることはすでに分かっている）。
\end{proof}

\begin{lemma}
\label{spaces-simplicial-lemma-descent-qcoh-for-fppf-hypercovering}
$S$ をスキームとする。$X$ を $S$ 上の代数空間とする。
$U$ を $S$ 上の単体的代数空間とし、$a : U \to X$
を増大とする。$a : U \to X$ が $X$ の fppf 超被覆ならば、
% P10-E0255
$$
a^* : \QCoh(\mathcal{O}_X) \to \QCoh(\mathcal{O}_U)
$$
は充満忠実な同値であり、その擬逆は $a_*$ により与えられる。
ここで $a : \Sh(U_\etale) \to \Sh(X_\etale)$ は
節 \ref{spaces-simplicial-section-simplicial-algebraic-spaces} のものである。
\end{lemma}

\begin{proof}
補題 \ref{spaces-simplicial-lemma-compare-simplicial-objects-fppf-etale-modules}
の図式を考える。この補題の証明で、$h_{-1}$ は
『代数空間のコホモロジーについてさらに』節
\ref{spaces-more-cohomology-section-fppf-etale-modules}
の射 $a_X$ であることを見た。したがって
『代数空間のコホモロジーについてさらに』補題
\ref{spaces-more-cohomology-lemma-review-quasi-coherent}
より、
$$
(h_{-1})^* :
\QCoh(\mathcal{O}_X)
\longrightarrow
\QCoh(\mathcal{O}_{big})
$$
は同値であり、その擬逆は $h_{-1, *}$ である。
同じことが $h$ の成分 $h_n$ についても成り立つ。
$\QCoh(\mathcal{O}_U)$ と $\QCoh(\mathcal{O}_{big, total})$
は、成分が準連接であるカルテジアン加群からなることを思い出そう。
補題 \ref{spaces-simplicial-lemma-quasi-coherent-sheaf} を参照せよ。
補題 \ref{spaces-simplicial-lemma-morphism-simplicial-sites-modules} の関手
$h^*$ と $h_*$ は成分上で関手 $h_n^*$ と $h_{n, *}$ に一致するから、
$$
h^* :
\QCoh(\mathcal{O}_U)
\longrightarrow 
\QCoh(\mathcal{O}_{big, total})
$$
は同値であり、その擬逆は $h_*$ であると結論できる。
$U$ は、節 \ref{spaces-simplicial-section-hypercovering} で定義した
$(\textit{Spaces}/S)_{fppf}$ における $X$ の超被覆であることに注意する。
補題 \ref{spaces-simplicial-lemma-hypercovering-X-simple-descent-modules} により、
$a_{fppf}^*$ は充満忠実で、その擬逆は $a_{fppf, *}$ であり、
その本質的像はカルテジアンな
% P10-E0256
$\mathcal{O}_{fppf, total}$-加群の層である。
したがって、補題 \ref{spaces-simplicial-lemma-quasi-coherent-sheaf} による
$\QCoh(\mathcal{O}_{big})$ と
$\QCoh(\mathcal{O}_{big, total})$ の記述から、擬逆を
$a_{fppf, *}$ とする同値
$$
a_{fppf}^* :
\QCoh(\mathcal{O}_{big})
\longrightarrow
\QCoh(\mathcal{O}_{big, total})
$$
を得る。形式的な議論（図式を追うこと）により、
$a^*$ は $\QCoh(\mathcal{O}_X)$ 上で充満忠実であり、
その像は $\QCoh(\mathcal{O}_U)$ に含まれることが分かる。

\medskip\noindent
最後に、$\mathcal{G}$ が $\QCoh(\mathcal{O}_U)$ に属すると仮定する。
このとき $h^*\mathcal{G}$ は\linebreak
$\QCoh(\mathcal{O}_{big, total})$ に属する。したがって
$h^*\mathcal{G} = a_{fppf}^*\mathcal{H}$ であり、ここで\linebreak
$\mathcal{H} = a_{fppf, *}h^*\mathcal{G}$ は
$\QCoh(\mathcal{O}_{big})$ に属する（上を参照）。
さらに
$\mathcal{H} = (h_{-1})^*\mathcal{F}$ であり、ここで
$\mathcal{F} = h_{-1, *}\mathcal{H}$ は
$\QCoh(\mathcal{O}_X)$ に属する。
図式を一周すると
$h^*\mathcal{G} \cong h^*a^*\mathcal{F}$ を得る。
% P10-E0257
$h^*$ の充満忠実性により、
$a^*\mathcal{F} \cong \mathcal{G}$ と結論する。
さらに
$\mathcal{F} = h_{-1, *}a_{fppf, *}h^*\mathcal{G} =
a_*h_*h^*\mathcal{G} = a_*\mathcal{G}$ であるから、
擬逆が $a_*$ により与えられるという主張も得られる。
\end{proof}

\begin{lemma}
\label{spaces-simplicial-lemma-cohomological-descent-qcoh-for-fppf-hypercovering}
$S$ をスキームとする。$X$ を $S$ 上の代数空間とする。
$U$ を $S$ 上の単体的代数空間とし、$a : U \to X$
を増大とする。$a : U \to X$ が $X$ の fppf 超被覆ならば、
準連接 $\mathcal{O}_X$-加群 $\mathcal{F}$ に対して写像
$$
\mathcal{F} \to Ra_*(a^*\mathcal{F})
$$
は同型である。ここで $a : \Sh(U_\etale) \to \Sh(X_\etale)$ は
節 \ref{spaces-simplicial-section-simplicial-algebraic-spaces} のものである。
\end{lemma}

\begin{proof}
% P10-E0258
補題 \ref{spaces-simplicial-lemma-compare-simplicial-objects-fppf-etale}
の図式を考える。$\mathcal{F}_n = a_n^*\mathcal{F}$ を
$a^*\mathcal{F}$ の第 $n$ 成分とする。これは準連接
$\mathcal{O}_{U_n}$-加群である。
『代数空間のコホモロジーについてさらに』補題
\ref{spaces-more-cohomology-lemma-cohomological-descent-etale-fppf-modules}
により
$\mathcal{F}_n = Rh_{n, *}h_n^*\mathcal{F}_n$ である。
したがって補題
\ref{spaces-simplicial-lemma-direct-image-morphism-simplicial-sites-modules}
により
$a^*\mathcal{F} = Rh_*h^*a^*\mathcal{F}$ である。さらに
\begin{align*}
Ra_*(a^*\mathcal{F})
& =
Ra_*Rh_*h^*a^*\mathcal{F} \\
& =
Rh_{-1, *}Ra_{fppf, *}a_{fppf}^*(h_{-1})^*\mathcal{F} \\
& =
Rh_{-1, *}(h_{-1})^*\mathcal{F} \\
& =
\mathcal{F}
\end{align*}
となる。% P10-E0259
第一の等式は上の議論による。第二の等式は、% P10-E0260
補題 \ref{spaces-simplicial-lemma-compare-simplicial-objects} の図式の可換性による。
第三の等式は、補題
\ref{spaces-simplicial-lemma-hypercovering-X-simple-descent-bounded-modules}
による。実際、$U$ は $(\textit{Spaces}/S)_{fppf}$
における $X$ の超被覆であり、
$a_{fppf}$ が平坦であるため
$La_{fppf}^* = a_{fppf}^*$ であること
（すなわち
$a_{fppf}^{-1}\mathcal{O}_{big} = \mathcal{O}_{big, total}$；
注意 \ref{spaces-simplicial-remark-augmentation-ringed} を参照）。
最後の等式は、すでに用いた
『代数空間のコホモロジーについてさらに』補題
\ref{spaces-more-cohomology-lemma-cohomological-descent-etale-fppf-modules}
による。
\end{proof}

\begin{lemma}
\label{spaces-simplicial-lemma-coh-descent-qcoh-unbounded-for-fppf-hypercovering}
$S$ をスキームとする。$X$ を $S$ 上の代数空間とする。
$U$ を $S$ 上の単体的代数空間とし、$a : U \to X$
を増大とする。$a : U \to X$ が $X$ の fppf 超被覆であると仮定する。
このとき $\QCoh(\mathcal{O}_U)$ は
$\textit{Mod}(\mathcal{O}_U)$ の弱 Serre 部分圏であり、
$$
a^* : D_\QCoh(\mathcal{O}_X) \longrightarrow D_\QCoh(\mathcal{O}_U)
$$
は圏の同値で、その擬逆は $Ra_*$ により与えられる。
ここで\linebreak
$a : \Sh(U_\etale) \to \Sh(X_\etale)$ は
節 \ref{spaces-simplicial-section-simplicial-algebraic-spaces} のものである。
\end{lemma}

\begin{proof}
まず、写像 $a_n : U_n \to X$ と $d^n_i : U_n \to U_{n - 1}$ は、
『超被覆』注意 \ref{hypercovering-remark-P-covering}
により平坦、局所有限表示、かつ全射である。

\medskip\noindent
$\mathcal{O}_U$-加群 $\mathcal{F}$ が準連接であることと、
それがカルテジアンであり、かつすべての $n$ に対して
$\mathcal{F}_n$ が準連接であることとは同値であることを思い出そう。
補題 \ref{spaces-simplicial-lemma-quasi-coherent-sheaf} を参照せよ。
補題 \ref{spaces-simplicial-lemma-Serre-subcat-cartesian-modules} と
（上で示した写像 $d^n_i : U_n \to U_{n - 1}$ の平坦性）により、
% P10-E0261
カルテジアン加群は
$\textit{Mod}(\mathcal{O}_U)$ の弱 Serre 部分圏をなす。
他方、各 $n$ に対して
$\QCoh(\mathcal{O}_{U_n}) \subset \textit{Mod}(\mathcal{O}_{U_n})$
は弱 Serre 部分圏である
（『代数空間の性質』補題
\ref{spaces-properties-lemma-properties-quasi-coherent}）。
これらを合わせると、
$\QCoh(\mathcal{O}_U) \subset \textit{Mod}(\mathcal{O}_U)$
は弱 Serre 部分圏であることが分かる。

\medskip\noindent
証明を終えるため、『サイト上のコホモロジー』補題
\ref{sites-cohomology-lemma-equivalence-unbounded-one}
の条件 (1) -- (5) を順に確かめる。

\medskip\noindent
(1) について。% P10-E0262
これは $a_n$ が平坦（上で見た）であることから、
補題 \ref{spaces-simplicial-lemma-flat-augmentation-modules}
により $a$ が平坦であるため成り立つ。

\medskip\noindent
(2) について。これは補題
\ref{spaces-simplicial-lemma-descent-qcoh-for-fppf-hypercovering}
の内容である。

\medskip\noindent
(3) について。これは補題
\ref{spaces-simplicial-lemma-cohomological-descent-qcoh-for-fppf-hypercovering}
の内容である。

\medskip\noindent
(4) について。小エタール・トポス
$\Sh(U_\etale)$ の定義には、サイト $U_\etale$ と
$U_{spaces, \etale}$ のいずれを用いてもよいことを思い出そう。
節 \ref{spaces-simplicial-section-simplicial-algebraic-spaces} を参照せよ。
『サイト上のコホモロジー』情況
\ref{sites-cohomology-situation-olsson-laszlo}
の仮定は三つ組
$(U_{spaces, \etale}, \mathcal{O}_U, \QCoh(\mathcal{O}_U))$
について成り立ち、同じ議論により三つ組
$(U_\etale, \mathcal{O}_U, \QCoh(\mathcal{O}_U))$
についても成り立つ。実際、
$$
\mathcal{B} \subset \Ob(U_\etale) \subset \Ob(U_{spaces, \etale})
$$
をアフィン対象の集合とする。
$V/U_n \in \mathcal{B}$ に対して $d_{V/U_n} = 0$ とし、
$\text{Cov}_{V/U_n}$ を $V_i$ がアフィンであるエタール被覆
$\{V_i \to V\}$ の集合とする。このとき
$\mathcal{F} \in \QCoh(\mathcal{O}_U)$ と
任意の $V/U_n \in \mathcal{B}$ に対して
$$
H^p(V/U_n, \mathcal{F}) = H^p(V, \mathcal{F}_n)
$$
が補題 \ref{spaces-simplicial-lemma-sanity-check-modules} により成り立つので、
所望の消滅を得る。右辺は、$U_n$ 上の準連接層
$\mathcal{F}_n$ の、$U_{n, \etale}$ のアフィン対象 $V$
上におけるコホモロジーである。
これは $p > 0$ に対して、
『代数空間のコホモロジー』節
\ref{spaces-cohomology-section-higher-direct-image} の議論と
『スキームのコホモロジー』補題
\ref{coherent-lemma-quasi-coherent-affine-cohomology-zero}
により消滅する。

\medskip\noindent
(5) について。$\mathcal{B} \subset \Ob(X_{spaces, \etale})$
をアフィン対象の集合とし、上で挙げた参照を用いれば従う。
\end{proof}

\begin{lemma}
\label{spaces-simplicial-lemma-compute-via-fppf-hypercovering-modules}
$S$ をスキームとする。$X$ を $S$ 上の代数空間とする。
$U$ を $S$ 上の単体的代数空間とし、$a : U \to X$
を増大とする。$a : U \to X$ が $X$ の fppf 超被覆ならば、
$$
R\Gamma(X_\etale, K) = R\Gamma(U_\etale, a^*K)
$$
が $K \in D_\QCoh(\mathcal{O}_X)$ に対して成り立つ。
ここで $a : \Sh(U_\etale) \to \Sh(X_\etale)$ は
節 \ref{spaces-simplicial-section-simplicial-algebraic-spaces} のものである。
\end{lemma}

\begin{proof}
これは補題
\ref{spaces-simplicial-lemma-coh-descent-qcoh-unbounded-for-fppf-hypercovering}
から従う。実際、『サイト上のコホモロジー』注意
\ref{sites-cohomology-remark-before-Leray} により
$R\Gamma(U_\etale, -) = R\Gamma(X_\etale, -) \circ Ra_*$ である。
\end{proof}

\begin{lemma}
\label{spaces-simplicial-lemma-spectral-sequence-fppf-hypercovering-modules}
$S$ をスキームとする。$X$ を $S$ 上の代数空間とする。
$U$ を $S$ 上の単体的代数空間とし、$a : U \to X$
を増大とする。% P10-E0263
$\mathcal{F}$ を準連接 $\mathcal{O}_X$-加群とし、
$\mathcal{F}_n$ を $U_{n, \etale}$ への引戻しとする。
$U$ が $X$ の fppf 超被覆ならば、標準スペクトル系列
$$
E_1^{p, q} = H^q_\etale(U_p, \mathcal{F}_p)
$$
が存在し、$H^{p + q}_\etale(X, \mathcal{F})$ に収束する。
\end{lemma}

\begin{proof}
補題 \ref{spaces-simplicial-lemma-compute-via-fppf-hypercovering-modules} と
\ref{spaces-simplicial-lemma-simplicial-module-cohomology-site} の直ちに分かる帰結である。
\end{proof}

\section{複体の fppf 下降}
\label{spaces-simplicial-section-fppf-descent-derived}

\noindent
本節では、代数空間の fppf 被覆と準連接加群の導来圏について、
これまでに示した結果のいくつかをまとめる。

\begin{lemma}
\label{spaces-simplicial-lemma-fppf-neg-ext-zero-hom}
$X$ をスキーム $S$ 上の代数空間とする。
$K, E \in D_\QCoh(\mathcal{O}_X)$ とする。
$a : U \to X$ を fppf 超被覆とする。
すべての $n \geq 0$ に対して
$$
\Ext_{\mathcal{O}_{U_n}}^i(La_n^*K, La_n^*E) = 0
\text{、ただし } i < 0
$$
であると仮定する。このとき次が成り立つ。
\begin{enumerate}
\item $i < 0$ に対して $\Ext_{\mathcal{O}_X}^i(K, E) = 0$ であり、
\item 完全列
$$
0
\to
\Hom_{\mathcal{O}_X}(K, E)
\to
\Hom_{\mathcal{O}_{U_0}}(La_0^*K, La_0^*E)
\to
\Hom_{\mathcal{O}_{U_1}}(La_1^*K, La_1^*E)
$$
が存在する。
\end{enumerate}
\end{lemma}

\begin{proof}
$K_n = La_n^*K$、$E_n = La_n^*E$ と書く。これらは
$La^*K$ と $La^*E$ に付随する加群の導来圏の単体的系
（定義 \ref{spaces-simplicial-definition-cartesian-derived-modules}）である
（補題 \ref{spaces-simplicial-lemma-cartesian-objects-derived-modules}）。
ここで $a : U_\etale \to X_\etale$ は
節 \ref{spaces-simplicial-section-simplicial-algebraic-spaces} のものである。
まず (2) を証明しよう。補題
\ref{spaces-simplicial-lemma-coh-descent-qcoh-unbounded-for-fppf-hypercovering}
により
$$
\Hom_{\mathcal{O}_X}(K, E) =
\Hom_{\mathcal{O}_U}(La^*K, La^*E)
$$
である。したがって、この列は
$$
0
\to
\Hom_{\mathcal{O}_U}(La^*K, La^*E)
\to
\Hom_{\mathcal{O}_{U_0}}(K_0, E_0)
\to
\Hom_{\mathcal{O}_{U_1}}(K_1, E_1)
$$
という形になる。% P10-E0264
第一の矢印は補題
\ref{spaces-simplicial-lemma-nullity-cartesian-modules-derived}
により単射である。この矢印の像は、補題
\ref{spaces-simplicial-lemma-hom-cartesian-modules-derived}
により第二の矢印の核である。これで (2) の証明が完了する。
(1) は、$i > 0$ に対して (2) を $K[i]$ と $E$
に適用すれば従う。
\end{proof}

\begin{lemma}
\label{spaces-simplicial-lemma-fppf-glue-neg-ext-zero}
$X$ をスキーム $S$ 上の代数空間とする。
$a : U \to X$ を fppf 超被覆とする。
$K_0 \in D_\QCoh(U_0)$ と、同型
$$
\alpha :
L(f_{\delta_1^1})^*K_0
\longrightarrow
L(f_{\delta_0^1})^*K_0
$$
が与えられ、これが $U_1$ 上で余サイクル条件を満たすと仮定する。
% P10-E0265
$\tau^n_i : [0] \to [n]$、$0 \mapsto i$ とおき、
$K_n = Lf_{\tau^n_n}^*K_0$ とおく。
$i < 0$ に対して
$\Ext^i_{\mathcal{O}_{U_n}}(K_n, K_n) = 0$ と仮定する。
このとき対象 $K \in D_\QCoh(\mathcal{O}_X)$ と、
$\alpha$ と両立する同型
% P10-E0266
$La_0^*K \to K$ が存在する。
\end{lemma}

\begin{proof}
% P10-E0267
補題 \ref{spaces-simplicial-lemma-construct-cartesian-objects-derived-modules}
により、対象 $K_n$ は加群の導来圏の単体的系の成分をなす。
すると、$(K_n, K_\varphi)$ が $K'$ から導かれる系となるような対象
% P10-E0268
$K' \in D_\QCoh(\mathcal{O}_{U_\etale})$
を得る。補題
\ref{spaces-simplicial-lemma-cartesian-module-derived-from-simplicial}
を参照せよ。最後に補題
\ref{spaces-simplicial-lemma-coh-descent-qcoh-unbounded-for-fppf-hypercovering}
を適用すると、ある $K \in D_\QCoh(\mathcal{O}_X)$ に対して
$K' = La^*K$ であることが分かり、所望の結論を得る。
\end{proof}

\section{代数空間の固有超被覆}
\label{spaces-simplicial-section-proper-hypercovering-spaces}

\noindent
本節は、代数空間の場合における
節 \ref{spaces-simplicial-section-proper-hypercovering} の類似である。
% P10-E0269
望む読者は、随所で「代数空間」を「スキーム」に置き換えても、
同様に正しい結果を得られる。その利点は、
『代数空間のコホモロジーについてさらに』節
\ref{spaces-more-cohomology-section-ph-etale}
への参照を、『エタール・コホモロジー』節
\ref{etale-cohomology-section-ph-etale}
への参照に置き換えられることである。

\medskip\noindent
基礎スキーム $S$ を固定する。
$X$ を $S$ 上の代数空間とし、$U$ を $S$ 上の単体的代数空間とする。
増大
$$
a : U \to X
$$
が与えられていると仮定する。
節 \ref{spaces-simplicial-section-simplicial-algebraic-spaces} を参照せよ。
次が成り立つとき、$U$ は $X$ の {\it 固有超被覆} であるという。
\begin{enumerate}
\item $U_0 \to X$ は固有かつ全射である。
\item $U_1 \to U_0 \times_X U_0$ は固有かつ全射である。
\item $n \geq 1$ に対して
$U_{n + 1} \to (\text{cosk}_n\text{sk}_n U)_{n + 1}$
は固有かつ全射である。
\end{enumerate}
$S$ 上の代数空間の圏はすべての有限極限をもつから、
上の定式化で用いた余骨格は存在する。
$$
\fbox{原理：固有超被覆を用いてエタール・コホモロジーを計算できる。}
$$
この原理の証明の背後にある鍵となる考えは、圏
$\textit{Spaces}/S$ 上の ph 位相とエタール位相を比較することである。
すなわち、ph 位相はエタール位相より強く、次が成り立つ。
(a) 固有全射は ph 被覆を定める。
(b) 小エタール・サイトから引き戻された層の ph コホモロジーは、
『代数空間のコホモロジーについてさらに』節
\ref{spaces-more-cohomology-section-ph-etale}
で見たように、エタール・コホモロジーと一致する。

\medskip\noindent
本節のすべての結果は、$U \to X$ が単に「ph 超被覆」である場合、
すなわち節 \ref{spaces-simplicial-section-hypercovering} で定義したサイト
$(\textit{Spaces}/S)_{ph}$ における $X$ の超被覆である場合にも一般化する。
これが必要になれば、ここで正確に定式化して証明することにする。

\begin{lemma}
\label{spaces-simplicial-lemma-compare-simplicial-objects-ph-etale}
$S$ をスキームとする。$X$ を $S$ 上の代数空間とする。
$U$ を $S$ 上の単体的代数空間とする。
$a : U \to X$ を増大とする。このとき可換図式
$$
\xymatrix{
\Sh((\textit{Spaces}/U)_{ph, total}) \ar[r]_-h \ar[d]_{a_{ph}} &
\Sh(U_\etale) \ar[d]^a \\
\Sh((\textit{Spaces}/X)_{ph}) \ar[r]^-{h_{-1}} &
\Sh(X_\etale)
}
$$
が存在する。ここで左の垂直射は
節 \ref{spaces-simplicial-section-hypercovering} で定義され、
右の垂直射は
節 \ref{spaces-simplicial-section-simplicial-algebraic-spaces} で定義される。
\end{lemma}

\begin{proof}
記法 $(\textit{Spaces}/U)_{ph, total}$ は、サイト
$(\textit{Spaces}/S)_{ph}$ と、このサイトの単体的対象 $U$
に対して節 \ref{spaces-simplicial-section-hypercovering} の構成を用いることを表す
\footnote{fppf 位相を用いて節
\ref{spaces-simplicial-section-fppf-hypercovering} で定義した
$(\textit{Spaces}/U)_{fppf, total}$ と区別するためである。}。
右辺のトポスにはサイト $X_{spaces, \etale}$ と
$U_{spaces, \etale}$ を用いる。% P10-E0270
これが許されることについては、
節 \ref{spaces-simplicial-section-simplicial-algebraic-spaces} の議論を参照せよ。

\medskip\noindent
$(\textit{Spaces}/U)_{ph, total}$ と
$U_{spaces, \etale}$ はいずれも
情況 \ref{spaces-simplicial-situation-simplicial-site} の場合 A に入ることに注意する。
$U_\etale$ については、これは節
\ref{spaces-simplicial-section-simplicial-algebraic-spaces} の構成から直ちに従い、
$(\textit{Spaces}/U)_{ph, total}$ については補題
\ref{spaces-simplicial-lemma-sr-when-fibre-products} から従う。
次に関手\linebreak
$U_{n, spaces, \etale} \to (\textit{Spaces}/U_n)_{ph}$、
$U \mapsto U/U_n$ と、
\linebreak
$X_{spaces, \etale} \to (\textit{Spaces}/X)_{ph}$、
$U \mapsto U/X$ を考える。
これらがサイトの射を定めることは、
『代数空間のコホモロジーについてさらに』節
\ref{spaces-more-cohomology-section-ph-etale}
で見た。そこではこれらを
% P10-E0271
$a_{U_n} = \epsilon_{U_n} \circ \pi_{u_n}$ および
$a_X = \epsilon_X \circ \pi_X$ と表した。
したがって注意
\ref{spaces-simplicial-remark-morphism-augmentation-simplicial-sites}
の意味で、増大と両立する単体的サイトの射を得る。
そこで補題
\ref{spaces-simplicial-lemma-morphism-augmentation-simplicial-sites}
を適用して結論を得る。
\end{proof}

\begin{lemma}
\label{spaces-simplicial-lemma-descent-sheaves-for-ph-hypercovering}
$S$ をスキームとする。$X$ を $S$ 上の代数空間とする。
$U$ を $S$ 上の単体的代数空間とし、$a : U \to X$
を増大とする。$a : U \to X$ が $X$ の固有超被覆ならば、
$$
a^{-1} : \Sh(X_\etale) \to \Sh(U_\etale)
\quad\text{および}\quad
a^{-1} : \textit{Ab}(X_\etale) \to \textit{Ab}(U_\etale)
$$
は充満忠実であり、その本質的像はカルテジアン層で、
擬逆は $a_*$ により与えられる。
ここで $a : \Sh(U_\etale) \to \Sh(X_\etale)$ は
節 \ref{spaces-simplicial-section-simplicial-algebraic-spaces} のものである。
\end{lemma}

\begin{proof}
集合の層について主張を証明する。これは、すでに確立した結果の
ほとんど形式的な帰結となる。補題
\ref{spaces-simplicial-lemma-compare-simplicial-objects-ph-etale}
の図式を考える。この補題の証明で、$h_{-1}$ は
『代数空間のコホモロジーについてさらに』節
\ref{spaces-more-cohomology-section-ph-etale}
の射 $a_X$ であることを見た。したがって
『代数空間のコホモロジーについてさらに』補題
\ref{spaces-more-cohomology-lemma-comparison-ph-etale}
より、$(h_{-1})^{-1}$ は充満忠実で、その擬逆は $h_{-1, *}$ である。
同じことが $h$ の成分 $h_n$ についても成り立つ。
補題 \ref{spaces-simplicial-lemma-morphism-simplicial-sites} による
関手 $h^{-1}$ と $h_*$ の記述から、
$h^{-1}$ は充満忠実で、その擬逆は $h_*$ であると結論できる。
$U$ は、節 \ref{spaces-simplicial-section-hypercovering} で定義した
$(\textit{Spaces}/S)_{ph}$ における $X$ の超被覆であることに注意する。
実際、『代数空間上の位相』補題
\ref{spaces-topologies-lemma-surjective-proper-ph}
により、全射な固有射は ph 被覆を与える。
補題 \ref{spaces-simplicial-lemma-hypercovering-X-simple-descent-sheaves}
により、$a_{ph}^{-1}$ は充満忠実で、その擬逆は $a_{ph, *}$、
本質的像は $(\textit{Spaces}/U)_{ph, total}$ 上の
カルテジアン層である。形式的な議論（図式を追うこと）により、
$a^{-1}$ は充満忠実であると分かる。

\medskip\noindent
最後に、$\mathcal{G}$ が $U_\etale$ 上のカルテジアン層であると仮定する。
このとき $h^{-1}\mathcal{G}$ は
$(\textit{Spaces}/U)_{ph, total}$ 上のカルテジアン層である。
したがって、$(\textit{Spaces}/X)_{ph}$ 上のある層
$\mathcal{H}$ に対して
$h^{-1}\mathcal{G} = a_{ph}^{-1}\mathcal{H}$ となる。
いささか衒学的な記法を用いて計算すると
\begin{align*}
(h_{-1})^{-1}(a_*\mathcal{G})
& =
(h_{-1})^{-1}
\text{等化子}(
\xymatrix@C=0.8em{
a_{0, small, *}\mathcal{G}_0
\ar@<1ex>[r] \ar@<-1ex>[r] &
a_{1, small, *}\mathcal{G}_1
}
) \\
& =
\text{等化子}(
\xymatrix@C=0.8em{
(h_{-1})^{-1}a_{0, small, *}\mathcal{G}_0
\ar@<1ex>[r] \ar@<-1ex>[r] &
(h_{-1})^{-1}a_{1, small, *}\mathcal{G}_1
}
) \\
& =
\text{等化子}(
\xymatrix@C=0.8em{
a_{0, big, ph, *}h_0^{-1}\mathcal{G}_0
\ar@<1ex>[r] \ar@<-1ex>[r] &
a_{1, big, ph, *}h_1^{-1}\mathcal{G}_1
}
) \\
& =
\text{等化子}(
\xymatrix@C=0.8em{
a_{0, big, ph, *}(a_{0, big, ph})^{-1}\mathcal{H}
\ar@<1ex>[r] \ar@<-1ex>[r] &
a_{1, big, ph, *}(a_{1, big, ph})^{-1}\mathcal{H}
}
) \\
& =
a_{ph, *}a_{ph}^{-1}\mathcal{H} \\
& =
\mathcal{H}
\end{align*}
となる。第一の等式は補題 \ref{spaces-simplicial-lemma-augmentation-site} により、
第二の等式は $(h_{-1})^{-1}$ が完全関手であることにより、
第三の等式は
『代数空間のコホモロジーについてさらに』補題
\ref{spaces-more-cohomology-lemma-proper-push-pull-ph-etale}
により従う
（ここで $a_0 : U_0 \to X$ と
% P10-E0272
$a_1: U_1 \to X$ が固有であることを用いる）。
第四の等式は $a_{ph}^{-1}\mathcal{H} = h^{-1}\mathcal{G}$ から、
第五の等式は補題 \ref{spaces-simplicial-lemma-augmentation-site} から、
第六の等式は上で見たことから従う。
$a_{ph}^{-1}\mathcal{H} = h^{-1}\mathcal{G}$ だから、
$h^{-1}\mathcal{G} \cong h^{-1}a^{-1}a_*\mathcal{G}$ を得る。
% P10-E0273
$h^{-1}$ の充満忠実性により、これで証明が終わる。
\end{proof}

\begin{lemma}
\label{spaces-simplicial-lemma-cohomological-descent-for-ph-hypercovering}
$S$ をスキームとする。$X$ を $S$ 上の代数空間とする。
$U$ を $S$ 上の単体的代数空間とし、$a : U \to X$
を増大とする。$a : U \to X$ が $X$ の固有超被覆ならば、
$K \in D^+(X_\etale)$ に対して
$$
K \to Ra_*(a^{-1}K)
$$
は同型である。ここで $a : \Sh(U_\etale) \to \Sh(X_\etale)$ は
節 \ref{spaces-simplicial-section-simplicial-algebraic-spaces} のものである。
\end{lemma}

\begin{proof}
補題 \ref{spaces-simplicial-lemma-compare-simplicial-objects-ph-etale}
の図式を考える。
『代数空間のコホモロジーについてさらに』補題
\ref{spaces-more-cohomology-lemma-cohomological-descent-etale-ph}
により、$D^+(U_{n, \etale})$ 上で
$Rh_{n, *}h_n^{-1}$ は恒等関手であることに注意する。
したがって補題
\ref{spaces-simplicial-lemma-direct-image-morphism-simplicial-sites}
により、$D^+(U_\etale)$ 上で
$Rh_*h^{-1}$ は恒等関手である。さらに
\begin{align*}
Ra_*(a^{-1}K)
& =
Ra_*Rh_*h^{-1}a^{-1}K \\
& =
Rh_{-1, *}Ra_{ph, *}a_{ph}^{-1}(h_{-1})^{-1}K \\
& =
Rh_{-1, *}(h_{-1})^{-1}K \\
& =
K
\end{align*}
となる。% P10-E0274
第一の等式は上の議論による。第二の等式は、% P10-E0275
補題 \ref{spaces-simplicial-lemma-compare-simplicial-objects}
の図式の可換性による。第三の等式は、補題
\ref{spaces-simplicial-lemma-hypercovering-X-simple-descent-bounded-abelian}
により従う。実際、『代数空間上の位相』補題
\ref{spaces-topologies-lemma-surjective-proper-ph}
により、$U$ は $(\textit{Spaces}/S)_{ph}$
における $X$ の超被覆である。
最後の等式は、すでに用いた
『代数空間のコホモロジーについてさらに』補題
\ref{spaces-more-cohomology-lemma-cohomological-descent-etale-ph}
による。
\end{proof}

\begin{lemma}
\label{spaces-simplicial-lemma-compute-via-ph-hypercovering}
$S$ をスキームとする。$X$ を $S$ 上の代数空間とする。
$U$ を $S$ 上の単体的代数空間とし、$a : U \to X$
を増大とする。$a : U \to X$ が $X$ の固有超被覆ならば、
$$
R\Gamma(X_\etale, K) = R\Gamma(U_\etale, a^{-1}K)
$$
が $K \in D^+(X_\etale)$ に対して成り立つ。
ここで $a : \Sh(U_\etale) \to \Sh(X_\etale)$ は
節 \ref{spaces-simplicial-section-simplicial-algebraic-spaces} のものである。
\end{lemma}

\begin{proof}
これは補題
\ref{spaces-simplicial-lemma-cohomological-descent-for-ph-hypercovering}
から従う。実際、『サイト上のコホモロジー』注意
\ref{sites-cohomology-remark-before-Leray} により
$R\Gamma(U_\etale, -) = R\Gamma(X_\etale, -) \circ Ra_*$ である。
\end{proof}

\begin{lemma}
\label{spaces-simplicial-lemma-ph-hypercovering-equivalence-bounded}
$S$ をスキームとする。$X$ を $S$ 上の代数空間とする。
$U$ を $S$ 上の単体的代数空間とし、$a : U \to X$
を増大とする。
$\mathcal{A} \subset \textit{Ab}(U_\etale)$ を、
カルテジアン・アーベル層からなる弱 Serre 部分圏とする。
$U$ が $X$ の固有超被覆ならば、関手 $a^{-1}$ は同値
$$
D^+(X_\etale) \longrightarrow D_\mathcal{A}^+(U_\etale)
$$
を定め、その擬逆は $Ra_*$ である。
ここで $a : \Sh(U_\etale) \to \Sh(X_\etale)$ は
節 \ref{spaces-simplicial-section-simplicial-algebraic-spaces} のものである。
\end{lemma}

\begin{proof}
補題 \ref{spaces-simplicial-lemma-Serre-subcat-cartesian-modules} により、
$\mathcal{A}$ は弱 Serre 部分圏であることに注意する。
この同値は、ここまでに得た結果の形式的な帰結である。
補題 \ref{spaces-simplicial-lemma-descent-sheaves-for-ph-hypercovering} と
\ref{spaces-simplicial-lemma-cohomological-descent-for-ph-hypercovering}、および
『サイト上のコホモロジー』補題
\ref{sites-cohomology-lemma-equivalence-bounded}
を用いればよい。
\end{proof}

\begin{lemma}
\label{spaces-simplicial-lemma-spectral-sequence-ph-hypercovering}
$S$ をスキームとする。$X$ を $S$ 上の代数空間とする。
$U$ を $S$ 上の単体的代数空間とし、$a : U \to X$
を増大とする。$\mathcal{F}$ を $X_\etale$ 上のアーベル層とし、
$\mathcal{F}_n$ を $U_{n, \etale}$ への引戻しとする。
$U$ が $X$ の ph 超被覆ならば、標準スペクトル系列
$$
E_1^{p, q} = H^q_\etale(U_p, \mathcal{F}_p)
$$
が存在し、$H^{p + q}_\etale(X, \mathcal{F})$ に収束する。
\end{lemma}

\begin{proof}
補題 \ref{spaces-simplicial-lemma-compute-via-ph-hypercovering} と
\ref{spaces-simplicial-lemma-simplicial-sheaf-cohomology-site} の直ちに分かる帰結である。
\end{proof}
